\documentclass[11pt]{article}
\usepackage[margin=1in]{geometry}
\usepackage{amsmath,amssymb,amsthm,mathtools}
\usepackage{bm,booktabs,microtype}
\usepackage[hidelinks,pdfauthor={OpenAI},pdftitle={Universality of critical quench autocorrelations in the Sherrington--Kirkpatrick model}]{hyperref}
\numberwithin{equation}{section}
\newtheorem{theorem}{Theorem}[section]
\newtheorem{proposition}[theorem]{Proposition}
\newtheorem{lemma}[theorem]{Lemma}
\newtheorem{corollary}[theorem]{Corollary}
\theoremstyle{definition}
\newtheorem{definition}[theorem]{Definition}
\theoremstyle{remark}
\newtheorem{remark}[theorem]{Remark}
\newcommand{\E}{\mathbb E}
\newcommand{\mathscr}[1]{\mathcal{#1}}
\newcommand{\Prob}{\mathbb P}
\newcommand{\R}{\mathbb R}
\newcommand{\N}{\mathbb N}
\newcommand{\ind}{\mathbf 1}
\newcommand{\cE}{\mathcal E}
\newcommand{\cI}{\mathcal I}
\newcommand{\cH}{\mathcal H}
\newcommand{\cL}{\mathcal L}
\newcommand{\cF}{\mathcal F}
\newcommand{\cG}{\mathcal G}
\newcommand{\cP}{\mathcal P}
\newcommand{\cQ}{\mathcal Q}
\newcommand{\dd}{\,\mathrm d}
\DeclareMathOperator{\Law}{Law}
\DeclareMathOperator{\Ent}{Ent}
\DeclareMathOperator{\Var}{Var}
\DeclareMathOperator{\Cov}{Cov}
\DeclareMathOperator{\Tr}{Tr}
\DeclareMathOperator{\diag}{diag}
\DeclareMathOperator{\supp}{supp}
\DeclareMathOperator{\TV}{TV}
\DeclareMathOperator{\osc}{osc}
\DeclareMathOperator{\Lip}{Lip}
\DeclareMathOperator{\dist}{dist}
\newcommand{\ip}[2]{\langle #1,#2\rangle}
\newcommand{\norm}[1]{\lVert #1\rVert}
\newcommand{\abs}[1]{\lvert #1\rvert}
\newcommand{\set}[1]{\{#1\}}
\newcommand{\eqdist}{\mathrel{\overset{\mathrm d}=}}
\newcommand{\weakto}{\Longrightarrow}
\newcommand{\PiG}{\Pi_g}
\newcommand{\Sg}{S^g}
\newcommand{\Tn}{T_n}
\allowdisplaybreaks[2]
\title{Universality of critical quench autocorrelations\\
in the Sherrington--Kirkpatrick model}
\author{OpenAI}
\date{October 5, 2026}
\begin{document}
\maketitle
\begin{abstract}
We prove joint functional convergence of the stationary and quench
autocorrelations of zero-field Sherrington--Kirkpatrick heat-bath dynamics
at inverse temperature $\beta=1$, with the same random limit for Gaussian
and Rademacher couplings. Each site has a rate-one clock, mean spin
autocorrelations are multiplied by $n^{1/3}$, and waiting times and lags
are measured in units $n^{2/3}$. The quench starts from independent fair
spins, and convergence is uniform on compact sets of positive waiting
times and lags. The quench limit is selected by these initial states and
relaxes to the stationary limiting autocorrelation as the waiting time
tends to infinity.
\end{abstract}
\tableofcontents
\section{Introduction}\label{sec:introduction}

A quench starts a spin system from a simple nonequilibrium law and then
lets it evolve at a prescribed temperature. At the critical temperature
of the Sherrington--Kirkpatrick model, the relevant relaxation time grows
with the number of spins. We determine the two-time autocorrelation after
a quench from independent symmetric spins. The limit retains the random
spectral edge of the interaction matrix and is governed by the same
diffusion as the stationary critical autocorrelation.

\subsection{The model and the result}

Let \(D\) be either the standard Gaussian law or the uniform law on
\(\{-1,1\}\). For \(n\ge2\), let \(J_{ij}\), \(i<j\), be independent with
law \(D\), and put \(J_{ji}=J_{ij}\), \(J_{ii}=0\). The critical zero-field
Gibbs measure is
\begin{equation}\label{eq:gibbs}
 \pi_{n,J}(\sigma)
 =Z_{n,J}^{-1}
   \exp\left\{n^{-1/2}\sum_{i<j}J_{ij}\sigma_i\sigma_j\right\},
 \qquad \sigma\in\{-1,1\}^n .
\end{equation}
Every site has an independent rate-one clock. At a ring at site \(i\), its
spin is replaced by \(a\in\{-1,1\}\) with conditional probability
\[
 \frac{\exp(a h_i(\sigma))}{2\cosh h_i(\sigma)},
 \qquad
 h_i(\sigma)=n^{-1/2}\sum_{j\ne i}J_{ij}\sigma_j .
\]
Write \(P_t^J\) for this heat-bath semigroup. Thus one unit of time contains
\(n\) updates on average. Let \(\nu_n\) be the uniform law on the cube,
independent of \(J\), and set \(T_n=n^{2/3}\). The stationary and quench
autocorrelations in critical units are, respectively,
\begin{align}
 A_{n,J}(t)
   &=n^{-2/3}\sum_{i=1}^n
      \ip{\sigma_i}{P_{tT_n}^J\sigma_i}_{\pi_{n,J}},
       \qquad t>0,                                      \label{eq:finite-A}\\
 B_{n,J}(s,t)
   &=n^{-2/3}\sum_{i=1}^n
      \E_{\nu_n}^J[\sigma_i(sT_n)\sigma_i((s+t)T_n)],
       \qquad s,t>0.                                    \label{eq:finite-B}
\end{align}
Thermal and clock randomness are averaged in these functions; the
interaction \(J\) remains random. The variable \(s\) is the waiting time
after the quench and \(t\) is the subsequent lag.

We recall the edge data of the stationary theorem
\cite[Proposition~2.3, Proposition~5.1, and Theorem~6.8]{FUA}.
Let \(g=(g_a)_{a\ge1}\) be the increasing upper GOE edge process, the
joint limit in law of the leading coordinates: for each fixed \(m\),
\[
 \bigl(n^{2/3}(2-\lambda_a)\bigr)_{a=1}^m
       \ \weakto\ (g_a)_{a=1}^m ,
\]
where \(\lambda_1\ge\lambda_2\ge\cdots\) are the eigenvalues of a GOE
matrix with off-diagonal variance \(1/n\)
\cite[Theorem~1.1]{RamirezRiderVirag2011}. For any \(b+g_1>0\), put
\(\ell_a=(g_a+b)^{-1}\) and
\begin{equation}\label{eq:edge-D-intro}
 D_g(b)=\lim_{L\to\infty}
 \left\{\sqrt L-\sum_{a\ge1}
 \left(\frac1{g_a+b}-\frac1{g_a+L}\right)\right\}.
\end{equation}
The limit is in probability over \(g\), with the almost-sure version
chosen along a deterministic sequence as in the cited stationary theorem.
If \(Z_a\) are independent \(N(0,\ell_a)\), their centered square sum
converges. The continuous densities of this sum and its tails define
\[
 \Pi_g=\Law\left((Z_a)_{a\ge1}\,\middle|\,
              \sum_{a\ge1}(Z_a^2-\ell_a)=D_g(b)\right).
\]
This measure is independent of the admissible \(b\). The cylinder
gradient form
\[
 \cE_g(F)=\int\sum_{a\ge1}|\partial_aF|^2\,\dd\Pi_g
\]
is closable. Let \(H_g\) be the nonnegative operator of its closure, and
let \(c_*>0\) be the deterministic microscopic coefficient in
\cite[Theorem~6.8]{FUA}. We write
\begin{equation}\label{eq:edge-semigroup}
 S_t^g=e^{-t c_*H_g},
 \qquad
 A_g(t)=\sum_{a\ge1}\ip{x_a}{S_t^g x_a}_{\Pi_g},\quad t>0.
\end{equation}
The stationary theorem gives a Borel version of \(A_g\) in
\(C_{\mathrm{loc}}((0,\infty))\). Section~\ref{sec:background} records the
finite normalization and the stationary edge, static-coordinate, and
relaxed-energy inputs.

For an open set \(V\subset\R^d\), \(C_{\mathrm{loc}}(V)\) denotes the
space of continuous functions with uniform convergence on every compact
subset. Fix the countable family \(\mathscr V\) of bounded smooth cylinder
tests constructed in Section~\ref{sec:gaussian-limits}. It is dense in
\(L^2(\Pi_g)\) for each edge realization under consideration.
We call a family \(h_s^g\in L^2(\Pi_g)\) measurable
when \((g,s)\mapsto\Pi_g[h_s^g\phi]\) is measurable for each test
\(\phi\in\mathscr V\). These pairings determine each density class.

\newpage
\begin{theorem}[Critical quench universality]\label{thm:main}
There is a measurable family \((h_s^g)_{s>0}\), with
\(h_s^g\in L^2(\Pi_g)\), such that for almost every \(g\) the following
properties hold.
\begin{enumerate}
\item For every \(s>0\), \(h_s^g\ge0\) and \(\Pi_g[h_s^g]=1\). For
      \(s,t>0\),
      \[
       h_{s+t}^g=S_t^g h_s^g.
      \]
      Moreover,
      \[
       h_s^g\Pi_g\ \weakto\ \delta_0
       \quad\text{in the product topology on }\R^\N
       \quad(s\downarrow0),
       \qquad
       \|h_s^g-1\|_{L^1(\Pi_g)}\longrightarrow0
       \quad(s\to\infty).
      \]
\item The series
      \begin{equation}\label{eq:limit-B}
       B_g(s,t)=\sum_{a\ge1}
          \int h_s^g(x)\,x_a\,(S_t^g x_a)(x)\,\dd\Pi_g(x)
      \end{equation}
      converges locally uniformly on \((0,\infty)^2\) and defines a
      continuous function.
\item For both choices of \(D\),
      \[
       \Law_D(A_{n,J},B_{n,J})
       \ \weakto\
       \Law_g(A_g,B_g)
      \]
      in
      \(C_{\mathrm{loc}}((0,\infty))
        \times C_{\mathrm{loc}}((0,\infty)^2)\).
      The two coupling laws have the same \(g\), the same family
      \(h^g\), and the same coefficient \(c_*\).
\item For every \(0<a<b<\infty\),
      \[
       \sup_{a\le t\le b}|B_g(s,t)-A_g(t)|
          \longrightarrow0
       \quad\text{in probability over }g
       \quad(s\to\infty).
      \]
\end{enumerate}
The positive-time density limits of the Gaussian dynamics with initial
law \(\nu_n\) select this family, as made precise in
Theorem~\ref{thm:entrance-selection}. The normalization is exactly
\eqref{eq:finite-A}--\eqref{eq:finite-B}.
\end{theorem}

The boundary \(\delta_0\) is interpreted in the product topology.
At small times each fixed edge coordinate can vanish while the
renormalized constraint is carried by coordinates whose indices tend
to infinity. The proof identifies which such entrance is selected by
the uniform starts through a stronger, quantitative condition on these
coordinates.

\subsection{History and significance}

Sherrington and Kirkpatrick introduced the infinite-range spin-glass
model in 1975 \cite{SK1975}. The subsequent dynamical theory studied
memory through two-time correlation and response functions. In particular,
Sompolinsky and Zippelius developed a dynamical mean-field treatment
for soft spins \cite{SompolinskyZippelius1982}, and Cugliandolo and
Kurchan analyzed nonequilibrium relaxation after a quench in a
soft-spin Langevin formulation \cite{CugliandoloKurchan1994}. Their
asymptotic regime takes the system-size limit before the long-time limit
and differs from the finite-size critical window considered here.
For critical Ising heat-bath dynamics, Billoire and Campbell
\cite[Equation~(7)]{BilloireCampbell2011} numerically tested the
stationary finite-size form
\[
 q_c(n,t)\sim n^{-1/3}F(t/n^{2/3}),
\]
with time measured in updates per spin. Theorem~\ref{thm:main} studies
the two positive times following a uniform quench and identifies the
law of the sample-dependent limiting function.

The spectral mechanism is especially transparent in spherical spin
glasses. Kosterlitz, Thouless, and Jones
\cite{KosterlitzThoulessJones1976} analyzed the infinite-range Gaussian
spherical spin glass using properties of large random matrices.
Cugliandolo and Dean \cite{CugliandoloDean1995}
studied relaxation and aging for the spherical model, while
Ben Arous, Dembo, and Guionnet \cite{BenArousDemboGuionnet2001}
proved limiting dynamical and aging results with a soft spherical
constraint. Fyodorov, Perret, and Schehr
\cite{FyodorovPerretSchehr2015} later connected finite-size randomness
on the \(n^{2/3}\) scale to the edge in zero-temperature spherical
dynamics. These works concern different spin spaces or temperature
regimes; they explain why the critical edge should remain visible in
a nonequilibrium limit.

The stationary critical Ising theorem \cite{FUA} identifies the
conditioned edge measure, its closed dynamical form, and the
microscopic coefficient \(c_*\). Its equilibrium comparison uses
the conditional eigenframe averaging of Comets \cite[Equations~(2.1)--(2.2)]{Comets1996},
which equates the Haar-averaged cube and spherical partition integrals
when both use normalized reference measures and the same quadratic
Hamiltonian, and the
critical conditioned-Gaussian description of Du and Huang
\cite{DuHuang2026FreeEnergy,DuHuang2026Overlap}. Using the cap and
spectral estimates of \cite{FUA,CM,CS},
Sections~\ref{sec:edge-estimates} and~\ref{sec:warming} extend cap
control down to a tight lower threshold for \(np^3\) and derive warming
uniformly over the initial law at each positive critical time. The
order-\(n\) entropy argument uses Lemma~\ref{warm:posterior-mean}, which
strengthens the square-density posterior-mean estimate of
\cite[Proposition~7.3]{SG} to square-root dependence on the relative
heat-bath energy. The stationary relaxed-energy theorem is applied within its
polynomially bounded test class through bounded resolvent optimizers in
Section~\ref{sec:gaussian-limits}.

\subsection{The argument}

In the Gaussian model, adding an independent GOE diagonal changes the spin
Hamiltonian only by a constant. Let \(u_a\) be orthonormal eigenvectors of
the completed matrix and put \(x_a=n^{-1/3}\langle u_a,\sigma\rangle\).
If \(w_s\) is the Gibbs-relative
density at time \(sT_n\), orthogonality rewrites the quench correlation as
\[
 B_{n,J}(s,t)=\sum_{a=1}^n
       \pi_{n,J}[w_s\,x_a\,P_{tT_n}^Jx_a].
\]
This formula separates the law reached after the quench from subsequent
relaxation. The source identifies the relaxation operator for sufficiently
controlled tests. Three further arguments are needed.

We first regularize the law reached after the quench. At scale \(p\), a
nested Gaussian observation gives noisy information about the spin
projections onto modes with \(2-\lambda_a\le p^2\); conditioning caps
those eigenvalues at \(2-p^2\). On the retained spectral events, their
number is comparable to \(k_p=np^3\).
Section~\ref{sec:edge-estimates} bounds the variance left unresolved for
\(0\le f\le1\) by \(Cp^{-2}\) times the heat-bath energy plus
\(Ce^{-ck_p^a}\), down to a fixed threshold for \(k_p\) at each prescribed
disorder confidence. In the small-entropy part of
Section~\ref{sec:warming}, a contradiction argument assumes that a bounded
level test \(f\) has less energy than the claimed profile bound. At a later
scale \(R\), this makes the normalized tilt by \(f^2\) close to the
normalized observation tilt by \(F_R^2\), where
\(F_R=\E[f\mid h_R]\). Following its entropy through earlier observations,
together with the source path to a product posterior at higher entropy
levels, yields entropy and kernel bounds at every fixed positive critical
time, uniformly over the initial law.

For bounded cylinder sources, the finite Markov resolvent contraction
bounds the optimizers used in the stationary relaxed-energy theorem.
Section~\ref{sec:gaussian-limits} derives strong semigroup convergence
from that theorem and combines it with the warming bounds on densities
and transition kernels. This produces positive-time families on represented
Gaussian subsequences, evolving by \(S_t^g\), which may initially retain
randomness beyond \(g\).
To select one, choose a measurable admissible \(b\) with \(D_g(b)>0\),
put \(r_i=g_i+b\), and define
\[
 M_g(x)=\sum_i\frac{x_i^2}{2r_i}.
\]
We prove that \(0<M_g(x)<\infty\) for \(\Pi_g\)-almost every \(x\) and
therefore may define, on that full-measure set, the probability weights
\[
 p_i(x)=\frac{x_i^2}{2r_iM_g(x)}.
\]
For each represented Gaussian limit, the averaged uniform-start
estimate yields a deterministic sequence \(s_k\downarrow0\) along
which these weights leave every fixed finite set in probability under
\(h_{s_k}\Pi_g\), almost surely in the represented data.
Section~\ref{sec:entrance-selection} considers semigroup families whose
densities are locally uniformly bounded in \(L^\infty(\Pi_g)\) and satisfy
such sequential escape. Angular ratios and a common rescaling give linear
mode equations driven by independent Brownian motions with a common
diffusion factor. The square constraint gives a scalar Volterra equation
whose high-mode responses approach a common kernel. A stopped coupling
proves fixed-\(g\) uniqueness and the full \(s\downarrow0\) product-topology
boundary. The proof then constructs the common family measurably from \(g\).

Finally, we compare smooth functions of the stationary and quench traces
under entry replacement \cite{Chatterjee2006}. An order-\(k\) derivative
with respect to the normalized interaction \(J_{ij}/\sqrt n\) is
accompanied by \(n^{-k/2}\), while summing over entries contributes
order \(n^2\). The two coupling laws match through order three, so the
averaged fourth derivative must gain a vanishing factor.
Sections~\ref{sec:history-cutoffs}--\ref{sec:replacement} first impose
static overlap-tail estimates for every interaction matrix in a
cutoff's support. Under the associated normalized weighted path laws,
these estimates yield history-overlap bounds with failure probability
\(O(n^{-A})\) for any prescribed \(A>0\), after the cutoff parameters
are chosen. Forward and reverse conditional estimates for
endpoint-normalized likelihoods on a binary tree of transition intervals
then control the derivatives over the full critical time span. The
estimates also include derivatives of the cutoff weights. They
transfer the joint trace laws and the probability of retaining the
cutoffs to Rademacher disorder; smooth grid tests then transfer
tightness in the locally uniform topology.
Section~\ref{sec:completion} assembles the result.

\section{The stationary edge input}\label{sec:background}

We fix the finite normalization and record the stationary edge,
static-comparison, and relaxed-energy inputs used below. Source cap
estimates are stated at their points of use in
Section~\ref{sec:edge-estimates}; the following sections develop the
extensions, entrance construction, and disorder comparison.

\subsection{The finite form and spectral coordinates}

Write \(\mu=\pi_{n,J}\). For a function \(f\) on the cube let
\[
 D_i f(\sigma)=\frac{f(\sigma^{i,+})-f(\sigma^{i,-})}{2},
 \qquad
 t_i(\sigma)=\tanh h_i(\sigma),\qquad
 v_i(\sigma)=1-t_i(\sigma)^2,
\]
where \(\sigma^{i,\pm}\) is obtained by setting spin \(i\) to \(\pm1\).
The nonnegative heat-bath operator and its symmetric form are
\begin{equation}\label{eq:finite-form}
 H_n=\sum_{i=1}^n(\sigma_i-t_i)D_i,\qquad
 \cE(f,k)=\mu\left[\sum_{i=1}^n v_iD_ifD_ik\right],
 \qquad P_t^J=e^{-tH_n}.
\end{equation}
Indeed the summand of \(H_n\) is \(f-\mu[f\mid\sigma_{-i}]\).
Conditional two-point covariance gives the displayed form. Thus
\eqref{eq:finite-form} uses the rate-one clock at each site specified
in the theorem. We write \(\cE(f)=\cE(f,f)\) and, for \(z\ge0\),
\(\cI(z)=\cE(z,\log z)\), with its extended nonnegative value at zeros.

Until the disorder comparison, \(J\) is Gaussian. Put
\(W_{ij}=J_{ij}/\sqrt n\) off the diagonal and add independent
\(W_{ii}\sim N(0,2/n)\). The diagonal contributes a constant to
\(\sigma^{\mathsf T}W\sigma/2\), so neither \(\mu\) nor \(H_n\) changes.
After independent sign choices for the eigenvectors,
\[
 W=U\diag(\lambda_1,\ldots,\lambda_n)U^{\mathsf T},\qquad
 \lambda_1\ge\cdots\ge\lambda_n,
\]
has \(U\) Haar orthogonal conditionally on the eigenvalues. Denote its
columns by \(u_a\), and put
\begin{equation}\label{eq:finite-edge}
 d_a=2-\lambda_a,\qquad g_{a,n}=T_n d_a,\qquad
 x_a(\sigma)=n^{-1/3}\langle u_a,\sigma\rangle .
\end{equation}
We use the same \(x_a\) for the coordinate functions on \(\R^\N\).
For two spins, orthogonality gives the exact identity
\begin{equation}\label{eq:replica-identity}
 n^{-2/3}\sigma\cdot\sigma'
       =\sum_{a=1}^n x_a(\sigma)x_a(\sigma').
\end{equation}
In particular, if
\(w_s=d(\nu_nP_{sT_n}^J)/d\mu\), then
\begin{equation}\label{eq:finite-spectral-traces}
 A_{n,J}(t)=\sum_{a=1}^n\mu[x_aP_{tT_n}^Jx_a],
 \qquad
 B_{n,J}(s,t)=\sum_{a=1}^n\mu[w_sx_aP_{tT_n}^Jx_a].
\end{equation}

\subsection{Conditioned measure and relaxed energy}

The following statements are the parts of the stationary manuscript
that we use directly. They also specify the meaning of the conditional
measure in the introduction. A smooth compact cylinder is a smooth
function of finitely many coordinates, compactly supported in those
coordinates.

\begin{proposition}[Stationary inputs]\label{prop:stationary-inputs}
The following hold for the edge law and the Gaussian finite models.
\begin{enumerate}
\item Let \(b+g_1>0\), \(\ell_a=(g_a+b)^{-1}\). Almost surely
      \(\sum_a\ell_a=\infty\) and \(\sum_a\ell_a^2<\infty\).
      If \(Z_a\) are independent \(N(0,\ell_a)\), the densities \(q\)
      of \(\sum_a(Z_a^2-\ell_a)\) and \(q_m\) of
      \(\sum_{a>m}(Z_a^2-\ell_a)\) are continuous, bounded, and strictly
      positive. The first \(m\) coordinates of \(\Pi_g\) have density
      \begin{equation}\label{eq:edge-cylinder-density}
       \frac{q_m\!\left(D_g(b)-\sum_{a\le m}(x_a^2-\ell_a)\right)}
            {q(D_g(b))}
       \prod_{a\le m}
          \frac{e^{-x_a^2/(2\ell_a)}}{\sqrt{2\pi\ell_a}}.
      \end{equation}
      These consistent densities are measurable in \(g\) and
      independent of \(b\). For each fixed \(0\le j<\infty\),
      \(\Pi_g[|x_a|^j]\le C_{j,g,b}\ell_a^{j/2}\), with a finite
      constant uniform in \(a\).
      For sufficiently large fixed \(m\), the tail beyond \(m\)
      has a density bounded by a finite constant relative to
      \(\bigotimes_{a>m}N(0,\ell_a)\). Conditional on that tail, the
      first \(m\) coordinates have the Gaussian angular law on the
      sphere of squared radius
      \[
       D_g(b)+\sum_{a\le m}\ell_a-\sum_{a>m}(x_a^2-\ell_a).
      \]
      In particular the renormalized constraint holds \(\Pi_g\)-almost
      surely. These are the statements of
      \cite[Proposition~2.3]{FUA}.
\item With \(b_n=1+(-g_{1,n})_+\) and \(b=1+(-g_1)_+\), the spherical
      and cube edge marginals converge jointly with the edge points to
      \(\Pi_g\) on the represented subsequences described below.
      Cube integrals converge for every continuous cylinder of
      polynomial growth. Spherical coordinates obey the uniform bounds
      \[
       \mu^{\mathrm{sp}}[|x_a|^j]\le C_j(g_{a,n}+b_n)^{-j/2},
       \qquad
       \lim_{m\to\infty}\sup_n\sum_{a>m}(g_{a,n}+b_n)^{-2}=0
      \]
      after discarding finitely many \(n\) on such a representation.
      Here \(\mu^{\mathrm{sp}}\) is surface measure on
      \(\{|\sigma|^2=n\}\), tilted by
      \(\exp(\sigma^{\mathsf T}W\sigma/2)\).
      The sources are \cite[Proposition~2.3 and Corollary~3.3]{FUA}.
\item The Euclidean gradient on smooth compact cylinders is closable
      from \(L^2(\Pi_g)\) to \(L^2(\Pi_g;\ell^2)\). The associated
      closed form has kernel equal to the constants. Every coordinate
      belongs to its domain and satisfies
      \(\nabla x_a=e_a\), hence \(\cE_g(x_a)=1\).
      This is \cite[Proposition~5.1]{FUA}. The usual contraction rule
      for gradients passes to the closure, so the form is Dirichlet
      and \(S_t^g\) is a conservative symmetric Markov semigroup.
\item On a quenched represented subsequence, say that \(f_n\) converges
      weakly to \(F\in L^2(\Pi_g)\) if its \(L^2(\mu)\) norms are bounded
      and
      \(\mu[f_n\phi(x)]\to\Pi_g[F\phi]\) for every smooth compact
      cylinder \(\phi\); call the convergence strong when in addition
      \(\mu[f_n^2]\to\Pi_g[F^2]\). Suppose, for fixed \(C,D<\infty\),
      that \(\|f_n\|_\infty\le Cn^D\), that \(T_n\cE(f_n)\) is bounded,
      and that \(f_n\) converges weakly to \(F\).
      Then \(F\) belongs to the closed form domain and
      \begin{equation}\label{eq:relaxed-liminf}
       \liminf_n T_n\cE(f_n)\ge c_*\cE_g(F).
      \end{equation}
      Conversely, for every smooth compact cylinder \(F\), there are
      functions \(f_{n,j}\), polynomially bounded for each fixed \(j\),
      such that \(f_{n,j}\) converges strongly to \(F\) as \(n\to\infty\)
      for each \(j\), and
      \begin{equation}\label{eq:relaxed-recovery}
       \lim_{j\to\infty}\lim_{n\to\infty}T_n\cE(f_{n,j})
          =c_*\cE_g(F).
      \end{equation}
      The source asserts a common represented event for countably many
      chosen sources and tests. Recovery is used at each fixed \(j\),
      followed by the displayed \(j\)-limit
      \cite[Theorem~6.8]{FUA}.
\end{enumerate}
\end{proposition}

In Item 1, the tail density is
\[
 \frac{q_m^{\mathrm{head}}\!\left(
     D_g(b)-\sum_{a>m}(x_a^2-\ell_a)\right)}{q(D_g(b))},
\]
where \(q_m^{\mathrm{head}}\) is the density of the centered head square
sum. A sufficiently large head makes this density bounded. This concrete
description will be used for estimates on infinitely many coordinates.
For a candidate pair
\((F,G)\in L^2(\Pi_g)\times L^2(\Pi_g;\ell^2)\), the angular weak
identity underlying Item 3 is
\begin{equation}\label{eq:angular-weak}
 \int \psi\,(Ax)\cdot G\,\dd\Pi_g
 =-\int F\left\{(Ax)\cdot\nabla\psi
       -\psi\sum_a(g_a+b)x_a(Ax)_a\right\}\dd\Pi_g ,
\end{equation}
for every skew-symmetric matrix \(A\) with finitely many nonzero entries
and every smooth compact cylinder \(\psi\) whose coordinate set contains
every index appearing in \(A\). These identities hold exactly when
\(F\) belongs to the closed form domain and \(G=\nabla F\)
\cite[Proposition~5.1]{FUA}. The term with \(b\) cancels by skew-symmetry.

The source also proves that \(g\mapsto A_g\), defined in
\eqref{eq:edge-semigroup}, is a Borel
\(C_{\mathrm{loc}}((0,\infty))\)-valued map and identifies \(A_g\) as
the stationary limit for both entry laws
\cite[Proposition~9.2]{FUA}. We will obtain the joint limit with
\(B_{n,J}\) from the same finite coordinates. In particular,
\eqref{eq:relaxed-liminf} is used with its polynomial bound; the
strong resolvent convergence needed here is derived in
Section~\ref{sec:gaussian-limits}.

\subsection{High-probability estimates and subsequences}

We make one convention for random constants. An estimate holds with
\emph{tight constants} if, for every \(\chi>0\), there are deterministic
constants in that estimate and events \(\mathcal G_{n,\chi}\) such that
\[
 \liminf_{n\to\infty}\Prob(\mathcal G_{n,\chi})\ge1-\chi
\]
and the estimate holds on \(\mathcal G_{n,\chi}\) simultaneously for
every scale and test function in its stated class. The constants may
depend on \(\chi\) and on fixed parameters specified before the
estimate, but not on \(n\), the scale, or the test. A tight lower
threshold for the dimension parameter \(k_p=np^3\) has the same meaning with a deterministic
threshold \(K_\chi\) on each event. When an estimate is required above
a vanishing time threshold, this means that for each \(\chi\) a
deterministic threshold \(\varepsilon_{n,\chi}\downarrow0\) can be
chosen on those events.

We use the represented-subsequence convention of
\cite[Remark~2.2]{FUA}. The edge points, the countably many cylinder
integrals and error quantities in use, and bounds that are tight in
the preceding sense are included in a joint law before applying the
subsequence and Skorokhod representation theorems. A further diagonal
subsequence handles an iterated limit such as
\(\lim_{m\to\infty}\limsup_n\). On the resulting probability space,
the relevant bounds are finite and the chosen coordinates converge
almost surely. Conditional laws of the remaining finite data recover
the original finite models. This construction is used only to show
that every subsequence has a further subsequence with the identified
limit, which implies the asserted convergence in law.

For a disorder-dependent test, we use the separate class-uniform
quenched convention of \cite[Remark~2.2]{FUA}: each invoked inequality
holds on one event for every test in its stated norm, energy, and
polynomial-growth class. The source obtains this uniformity by bounding
the error by the stated norms times a common nonnegative random error,
which permits measurable choices of the test. The common represented
event is then chosen for the countable collection of sources and tests
used in this proof. The countable core encodes limits and density classes
after the uniform bounds have been established.

\section{Gaussian estimates at the bottom edge scales}
\label{sec:edge-estimates}

The observation below reveals noisy information about leading spin
projections.  We control the variance of a bounded spin function left
unresolved by that observation by \(Cp^{-2}\) times its heat-bath energy,
with a small remainder, at every scale for which \(np^3\) exceeds a tight
threshold.  The exponent \(2\) is important: a fixed
relative covariance error at every scale would accumulate to a power loss.
We first extend the Gaussian cap estimates to these scales and then show
that the relative error is integrable with respect to \(\mathrm d\log p\).

\subsection{Cap observations and the estimates}

Throughout this section the disorder is Gaussian, with the independent
GOE diagonal and randomized eigenvector signs specified in Section~2.
Set \(r_n=n^{-1/3+10^{-4}}\), the terminal scale of the source estimates,
and write \(d=2I-W\).  Define, in physical coordinates,
\[
\begin{aligned}
 P_p&=U\diag(\ind_{\{d_a\le p^2\}})U^{\mathsf T},&
 \cH_p&=\operatorname{ran}P_p,\\
 B_p&=U\diag((p^2-d_a)_+)U^{\mathsf T},&
 W_p&=W-B_p,\\
 k_p&=np^3,&M_p&=\sqrt{np}.
\end{aligned}
\]
Thus \(p^2M_p^2=k_p\).  A field at scale \(p\) lies in \(\cH_p\), with
length unit \(p^2M_p\).  For such a field let
\[
 \phi_p(h)=\log\nu_n
 \exp\left(\tfrac12 X^{\mathsf T}W_pX+h^{\mathsf T}X\right),
 \qquad m_p(h)=\nabla\phi_p(h),\qquad
 K_p(h)=\left[P_p\nabla^2\phi_p(h)P_p\right]_{\cH_p}.
\]
Brackets with subscript \(\cH_p\) denote compression to that space.
The probability with this log partition function is \(\mu_{p,h}\).
Thus \(m_p\) and \(K_p\) are its mean and its active covariance.
Superscript \({\rm sp}\) denotes the same definitions with normalized
surface measure on \(\{X:\|X\|^2=n\}\) in place of \(\nu_n\).
In the tight-constant convention of Section~\ref{sec:background}, a tight
lower threshold for \(k_p\) means that, for each \(\chi>0\), a deterministic
threshold \(K_\chi\) can be chosen on an event of limiting lower
probability at least \(1-\chi\).

For \(X\sim\mu\), form the Gaussian observation
\begin{equation}\label{edge:observation}
 h_p=B_pX+\beta_p,\qquad
 \operatorname{Cov}(\beta_p,\beta_q)=B_{\min(p,q)}.
\end{equation}
One construction uses an independent scalar Brownian motion in spectral
coordinate \(a\), evaluated at \((p^2-d_a)_+\).  Bayes' formula gives
\(\Law(X\mid h_q,\ q\le p)=\Law(X\mid h_p)=\mu_{p,h_p}\) at fixed \(W\).
Write \(Q_\mu\) for the joint law of
the spin and all observations at fixed \(W\).  The observation filtration
excludes the spin.  The path \(h_p\) is continuous at \(p^2=d_a\); the
new coordinate then has zero precision.  The finitely many activation
points may be omitted in scale integrals.  For \(0\le f\le1\), put
\begin{equation}\label{edge:posterior-variance}
 F_p(h)=\mu_{p,h}[f],\qquad
 V_p(f)=Q_\mu[\Var(f\mid h_p)].
\end{equation}

For \(u>-p^2\) define
\[
 R_{p,u}=U\diag\left((\max\{d_a,p^2\}+u)^{-1}\right)U^{\mathsf T}.
\]
On the spectral events and scales used below, the spherical saddle
\(\alpha_p(h)>-p^2\) is the unique solution of
\begin{equation}\label{edge:saddle}
 n=\Tr R_{p,\alpha_p(h)}+\|R_{p,\alpha_p(h)}h\|^2.
\end{equation}
Existence, uniqueness, and bounds on this solution are proved below.
Since \(R_{p,u}\) is scalar on \(\cH_p\), the length of a field whose
saddle is \(u\) is
\begin{equation}\label{edge:saddle-length}
 \ell_p(u)=(p^2+u)\sqrt{n-\Tr R_{p,u}}
\end{equation}
whenever the expression under the square root is nonnegative.

We use one change of probability law, always conditional on a fixed
spectrum.  Under the ordinary law \(U\) is Haar, and \(Q_\mu\) is the
joint law of the spin and its physical observations at the resulting
\(W\).  With probability reference measures in the partition functions,
set \(L_\circ(U)=Z_{\rm cube}(W)/Z_{\rm sp}(W)\).
Haar averaging gives \(E_UL_\circ=1\).  Under \(\mathbf P\), first choose
\(U\) with Haar density \(L_\circ\), and then use \(Q_\mu\) given \(W\);
write \(\mathbf E\) for its expectation.  The spherical identity of
\cite[Lemma~2.5 and Equation~(2.8)]{FUA} gives
\begin{equation}\label{edge:haar-bias}
 \Law_{\mathbf P}(U^{\mathsf T}X)
   =\mu^{\rm sp}\text{ in spectral coordinates},\qquad
 E_U(L_\circ-1)^2=o(1)
\end{equation}
on its retained eigenvalue events.  Consequently the spectral observation
path \((U^{\mathsf T}h_p)_p\) under \(\mathbf P\) has the spherical
observation law, denoted by \(Q^{\rm sp}\).  For comparison, let
\(\mathsf R\) choose \(U\) with Haar law independently of a spectral
path with law \(Q^{\rm sp}\), and transport that path to physical
coordinates by \(U\).
For the filtration generated by \(U\) and observations through \(p\),
but not by \(X\), the exact likelihood from \cite[Equation~(2.9)]{FUA} is
\begin{equation}\label{edge:path-likelihood}
 \frac{\mathrm d\mathbf P}{\mathrm d\mathsf R}\bigg|_{\cF_p}
 =L_p(h_p),\qquad
 L_p(h)=\exp\{\phi_p(h)-\phi_p^{\rm sp}(h)\}.
\end{equation}
The notation \(L_p(h)\) suppresses its dependence on \(U\): in a Haar
average the spectral field \(U^{\mathsf T}h\) is held fixed, while
\(L_p(h)\) still varies with \(U\).  At every such field \(E_UL_p(h)=1\).
To check the likelihood,
the observation density is a common Gaussian factor times the posterior
partition function divided by the original one, and \(L_\circ\) cancels
the original partition functions.  In particular, for every nonnegative
spin or observation functional \(G\),
\begin{equation}\label{edge:bias-transfer}
 E_U[\ind_{\{L_\circ\ge1/2\}}Q_\mu G]\le2\mathbf E G.
\end{equation}

\begin{proposition}[Cap estimates with a tight lower threshold]
\label{prop:edge-estimates}
Fix \(0<C_1<\infty\), \(0\le R<\infty\), positive tolerances
\(\epsilon_{\rm sp},\epsilon_0,\epsilon_v,\epsilon_g\), and \(\zeta>0\).
For every \(\varepsilon>0\) there are finite positive constants
\(K,C,c,a,v_0\), a width \(s_0>0\), a radius \(R_{\rm ann}<\infty\),
and Gaussian disorder events
\(\mathcal C_n\) with
\(\liminf_n\Prob_{\rm GOE}(\mathcal C_n)\ge1-\varepsilon\), such that
the following hold on \(\mathcal C_n\), simultaneously at every scale
\begin{equation}\label{edge:scale-range}
 K\le k_p=np^3,\qquad p\le C_1r_n.
\end{equation}
\begin{enumerate}
\item The active dimension is between \(ck_p\) and \(Ck_p\),
 \(d_1\ge-\epsilon_{\rm sp}p^2\), and, with \(b_0=4/(3\pi)\),
 \begin{equation}\label{edge:trace-estimates}
 \left|\frac{n-\Tr R_{p,0}}{np}-b_0\right|\le\epsilon_{\rm sp},
 \qquad
 \left|\frac p n\Tr R_{p,0}^2-2b_0\right|\le\epsilon_{\rm sp}.
 \end{equation}
 On \(\|h\|\le Rp^2M_p\), the saddle satisfies
 \(-1+c_R\le\alpha_p(h)/p^2\le C_R\).  On \(|u|\le s_0p^2\),
 \(\ell_p(u)\asymp p^2M_p\) and \(\ell'_p(u)\asymp M_p\).

\item On the field ball \(\|h\|\le Rp^2M_p\),
 \begin{equation}\label{edge:static-estimates}
 \begin{split}
 K_p(h)&\preceq C_Rp^{-2}I_{\cH_p},\\
 |\phi_p(h)-\phi_p^{\rm sp}(h)|&\le\epsilon_v k_p,\\
 \|P_pm_p(h)-R_{p,\alpha_p(h)}h\|&\le\epsilon_gM_p .
 \end{split}
 \end{equation}
 The constant \(C_R\) in the first line may be chosen independently of
 the requested annular tolerance \(\epsilon_0\).
 On the full annulus \(|\alpha_p(h)|\le s_0p^2\), the field is nonzero;
 for \(e=h/\|h\|\),
 \begin{equation}\label{edge:coarse-curvature}
 K_p(h)\preceq(1+\epsilon_0)p^{-2}I_{\cH_p},
 \qquad p^2\langle e,K_p(h)e\rangle\le0.9 .
 \end{equation}
 The annulus lies in \(\|h\|\le R_{\rm ann}p^2M_p\), and the event also
 includes the static comparisons on this ball, with constants allowed
 to depend on \(R_{\rm ann}\).

\item For every \(v\ge v_0\), the uncapped law \(\mu\), every posterior
 \(\mu_{p,h}\) in the stated field ball, and their spherical counterparts
 satisfy
 \begin{equation}\label{edge:projection-tails}
 \Pr\{\|P_pX\|>vM_p\}\le C e^{-c k_pv^6}.
 \end{equation}
 The event is empty for \(v>p^{-1/2}\).  The spherical pair bound of
 Lemma~\ref{edge:spherical-tails} also holds for each number of replicas
 fixed in advance, with constants allowed to depend on that number.  Moreover,
 \begin{equation}\label{edge:ordinary-annulus}
 Q_\mu\{|\alpha_p(h_p)|>\zeta p^2\}\le C e^{-c k_p}.
 \end{equation}

\item For every scale in \eqref{edge:scale-range} with \(p\le r_n\),
 simultaneously for every \(0\le f\le1\),
 \begin{equation}\label{edge:variance-estimates}
 \begin{split}
 V_p(f)&\le Cp^{-2}\cE(f)+C e^{-c k_p^a},\\
 p^2Q_\mu\|\nabla_{\cH_p}F_p(h_p)\|^2
       &\le C V_p(f)+C e^{-c k_p^a}.
 \end{split}
 \end{equation}
 The event and constants in these inequalities are independent of \(f\).
\end{enumerate}
All constants in the saddle, length, and comparison bounds are
deterministic on \(\mathcal C_n\) and may depend on the fixed inputs and
the requested failure probability, but not on \(n\), \(p\), or the field
in its stated region.
\end{proposition}

Every radius, tolerance, and finite multiplicative enlargement of a scale
interval is fixed before the constants are chosen.  Increasing \(K\) by a
fixed factor supplies any such enlargement at the lower endpoint.
A finite number of further requests for cap estimates may be imposed by
intersecting their events, with a smaller failure probability assigned to
each.  This preserves the constants and exponent of any variance estimate
already obtained.

\paragraph{Reduction to positive covariance error.}
Fix \(W\) and \(0\le f\le1\).  In the time variable \(t=p^2\), let
\[
 \mathcal F_t^{\rm obs}=\sigma(h_q:0\le q\le\sqrt t),
 \qquad M_t^f=F_{\sqrt t}(h_{\sqrt t}).
\]
This filtration excludes the spin and includes any information in
\(h_0\); the initial precision \(B_0\) need not vanish.  The bounded
posterior martingale \(M^f\) is continuous.  On positive time intervals,
its predictable bracket has density
\(\|\nabla_{\cH_{\sqrt t}}F_{\sqrt t}(h_{\sqrt t})\|^2\).
The coefficient can change at an activation, but the martingale has no
jump there.  Total variance and the bracket identity therefore give
\begin{equation}\label{edge:variance-identity}
 -\frac{\mathrm dV_p(f)}{\mathrm d(p^2)}
       =Q_\mu\|\nabla_{\cH_p}F_p(h_p)\|^2
\end{equation}
almost everywhere in \(p>0\), and in integrated form across activation
points.  This is the localization variance calculation in
\cite[proof of Theorem~1, Equation~(14) and the display after Equation~(15)]{EldanKoehlerZeitouni2022},
with deterministic rate \(P_p\).

Let \(\mathcal N\) be a jointly measurable set of triples \((p,W,h)\)
with \(h\in\cH_p\), chosen independently of \(f\), and let
\(\eta_p=\eta_p(W)\ge0\) be a measurable scalar fixed at each
\((W,p)\).  Suppose that
\(p^2K_p(h)\preceq(1+\eta_p)I_{\cH_p}\) for \(h\in\cH_p\) when
\((p,W,h)\notin\mathcal N\).  Put
\begin{equation}\label{edge:common-failure-cost}
 Z_p(X)=\frac{\|P_pX\|^2}{M_p^2},\qquad
 \rho_p(W)=k_pQ_\mu\!\left[
       \ind_{\mathcal N}(p,W,h_p)\mu_{p,h_p}(Z_p)\right].
\end{equation}
Conditional covariance and Cauchy--Schwarz give, at the observed field,
\begin{equation}\label{edge:gradient-bounds}
 \begin{split}
 \|\nabla_{\cH_p}F_p(h_p)\|^2
       &\le\|K_p(h_p)\|\Var_{\mu_{p,h_p}}(f),\\
 p^2\|\nabla_{\cH_p}F_p(h_p)\|^2
       &\le k_p\mu_{p,h_p}(Z_p).
 \end{split}
\end{equation}
For the second bound, write
\(\nabla F_p=\mu_{p,h_p}[P_pX(f-F_p)]\) and use \(|f-F_p|\le1\).
Its right side is a common cost for every such \(f\).  Splitting the
expectation in \eqref{edge:variance-identity} at \(\mathcal N\) yields
\begin{equation}\label{edge:variance-differential}
 -\frac{\mathrm dV_p}{\mathrm d\log p}
       \le2(1+\eta_p)V_p+2\rho_p
\end{equation}
almost everywhere.  If \(\eta_u\le Ck_u^{-a}\) for a fixed \(a>0\) on
\([p,q]\), its backward integrating factor is at most
\[
 \exp\left(2\int_p^q(1+Ck_u^{-a})\,\mathrm d\log u\right)
       \le C(q/p)^2,\qquad
 \int_p^qk_u^{-a}\,\mathrm d\log u
       =\frac{k_p^{-a}-k_q^{-a}}{3a}\le\frac{K^{-a}}{3a}
\]
whenever \(k_p\ge K\).  In particular,
\begin{equation}\label{edge:variance-amplification}
 V_p\le C(q/p)^2V_q+
       C\int_p^q(u/p)^2\rho_u\,\mathrm d\log u .
\end{equation}
The remaining work for Item~4 is to obtain this covariance bound with
\(\eta_p=Ck_p^{-a}\), control the common cost \(\rho_p\) on dyadic
intervals, and use the source variance estimate at \(r_n\).  The
integrable relative error is what preserves the exponent \(2\).

\subsection{Spectral counts, saddles, and spherical tails}

For a specified spectral error, being small above a tight threshold means
\begin{equation}\label{edge:threshold-convention}
 \lim_{K\to\infty}\limsup_{n\to\infty}
 \Prob_{\rm GOE}\left\{
 \sup_{\substack{np^3\ge K\\p\le C_1r_n}}
       \text{the specified error at \(p\)}>\eta\right\}=0
 \quad(\eta>0).
\end{equation}
Only finitely many errors at a time are needed.  This is a double-limit
assertion, not convergence with \(n\) at a fixed \(K\).

\begin{lemma}[Spectral estimates above a tight threshold]
\label{edge:spectral}
In the convention \eqref{edge:threshold-convention}, the measures
\(\nu_p=k_p^{-1}\sum_i\delta_{d_i/p^2}\) converge on fixed bounded
intervals to \(\pi^{-1}\sqrt v\,\ind_{\{v>0\}}\mathrm dv\).  Their
distribution functions converge uniformly on compact intervals.
The same convention gives \(d_1/p^2\ge-o(1)\),
\(\dim\cH_p/k_p\to2/(3\pi)\), and the two limits in
\eqref{edge:trace-estimates}.  On spectral events of arbitrarily large
limiting probability there is also a fixed \(C\) such that
\begin{equation}\label{edge:counting-tail}
 \begin{aligned}
 \#\{i:d_i\le u\}&\le Cnu^{3/2}
       &&(u\ge cp^2,\ u\le4),\\
 \int_{[A,\infty)}(1+v)^{-2}\nu_p(\mathrm dv)&\le CA^{-1/2}
       &&(A\ge1),
 \end{aligned}
\end{equation}
for every fixed \(c>0\), after increasing the threshold.
Every spectral restriction below also includes eigenvalue-only events
\(\mathcal S_n^{\rm glob}\) of probability \(1-o(1)\) on which, uniformly
over their spectra,
\begin{equation}\label{edge:global-spectral}
 \frac1n\sum_i\delta_{\lambda_i}\ \weakto\
 \rho_{\rm sc}(\dd x):=\frac1{2\pi}\sqrt{4-x^2}\,
             \ind_{\{|x|\le2\}}\dd x,\qquad
 \lambda_1\longrightarrow2,\qquad \|W\|\le3 .
\end{equation}
\end{lemma}

\begin{proof}
We use the eigenvalue-only event of \cite[Proposition~3.1]{CS}.
For every deterministic \(L_n\to\infty\) with \(L_n\le\log n\), that
proposition gives positive denominators and, uniformly for
\(L_nn^{-2/3}\le u\le M\) with fixed \(M\),
\begin{equation}\label{edge:source-resolvents}
 \begin{split}
 m_1(u):=\frac1n\sum_i(d_i+u)^{-1}
  &=\frac{2+u-\sqrt{u(4+u)}}2+o(\sqrt u),\\
 m_2(u):=\frac1n\sum_i(d_i+u)^{-2}
  &\le C_Mu^{-1/2},\\
 \lambda_j&=2+o(L_nn^{-2/3})\quad(1\le j\le3),
 \qquad \|W\|\le3 .
 \end{split}
\end{equation}
The errors have deterministic uniform bounds on an event of probability
at least \(1-Ce^{-cL_n^{1/4}}\).

For \eqref{edge:global-spectral}, take \(\mathcal S_n^{\rm glob}\) to be
the separate source event with \(L_n=\sqrt{\log n}\).  Its probability
tends to one, and the edge and
norm assertions follow directly from \eqref{edge:source-resolvents}.
For every fixed \(u>1\), the first resolvent converges uniformly there
to
\[
 \frac{2+u-\sqrt{u(4+u)}}2
 =\int\frac{\rho_{\rm sc}(\dd x)}{2+u-x}.
\]
The norm bound makes the empirical spectral measures a compact family.
Every weak subsequential limit therefore has the displayed transform
for every \(u>1\), and uniqueness of the Stieltjes transform identifies
it with \(\rho_{\rm sc}\).  A contradiction by subsequences gives this
weak convergence uniformly over the source event.  Intersecting with
this event does not change the tight-threshold convention.

Let \(K_n\to\infty\) be arbitrary and set
\(L_n=\min\{K_n^{1/3},\sqrt{\log n}\}\), changing finitely many initial
values if necessary.  For \(np^3\ge K_n\),
\(L_nn^{-2/3}=o(p^2)\).  Thus, for every fixed \(v>0\),
\[
 m_1(vp^2)=1-\sqrt v\,p+O_v(p^2)+o_v(p)
\]
uniformly in the scale range, and \(d_1/p^2\ge-o(1)\).  For
\(u\ge cp^2\), the counted denominators satisfy \(0<d_i+u\le2u\).
Consequently
\[
 \#\{d_i\le u\}\le4u^2\sum_i(d_i+u)^{-2}\le Cnu^{3/2}.
\]
Summing on dyadic bands gives the second bound in
\eqref{edge:counting-tail}; the last band above \(u=4\) costs \(O(p)\)
by \(\|W\|\le3\), and obeys the same bound.

For fixed \(v,w>0\), the transform difference is exactly
\begin{equation}\label{edge:transform-difference}
 \int\left(\frac1{x+v}-\frac1{x+w}\right)\nu_p(\mathrm dx)
 =\frac{m_1(vp^2)-m_1(wp^2)}p
 \longrightarrow\sqrt w-\sqrt v.
\end{equation}
To identify the limit including the tail, put
\(\omega_p(\mathrm dx)=(1+x)^{-2}\nu_p(\mathrm dx)\).
These finite measures have bounded mass, the tail bound
\eqref{edge:counting-tail}, and negative support in an interval shrinking
to zero.  Every weak limit \(\omega\) is therefore supported on
\([0,\infty)\).  Taking \(w=1\) in \eqref{edge:transform-difference} gives
\[
 \int\frac{1+x}{x+v}\,\omega(\mathrm dx)=\frac1{1+\sqrt v}.
\]
Letting \(v\to1\) gives \(\omega([0,\infty))=1/2\), and hence
\[
 \int\frac{\omega(\mathrm dx)}{x+v}
       =\frac1{2(1+\sqrt v)^2}.
\]
Uniqueness of the Stieltjes transform identifies
\(\omega(\mathrm dx)=\pi^{-1}\sqrt x(1+x)^{-2}\mathrm dx\).  The
continuous distribution function of the corresponding limit of
\(\nu_p\) gives uniformity in endpoints on compact intervals.  Failure
of uniformity over \(p\) would provide a sequence of scales to which
this same compactness argument applies, a contradiction.

The clipped second trace has limit
\[
 \frac p n\Tr R_{p,0}^2
 =\int(\max\{x,1\})^{-2}\nu_p(\mathrm dx)
 \longrightarrow
 \frac1\pi\left(\int_0^1\sqrt x\,\mathrm dx+
                 \int_1^\infty x^{-3/2}\,\mathrm dx\right)
 =\frac8{3\pi}.
\]
For the first trace, subtract an ordinary resolvent before integrating:
\[
 \frac{n-\Tr R_{p,0}}{np}
 =\frac{1-m_1(p^2)}p+
 \int\left(\frac1{x+1}-\frac1{\max\{x,1\}}\right)\nu_p(\mathrm dx)
 \longrightarrow\frac4{3\pi}.
\]
The difference kernel is \(O(x^{-2})\).  Its limiting integral is
\(4/(3\pi)-1\), obtained by substituting \(x=t^2\) on each of
\((0,1)\) and \((1,\infty)\).

These assertions hold for every arbitrarily slowly divergent \(K_n\).
If \eqref{edge:threshold-convention} failed, one could choose
\(K_j\to\infty\) and \(n_j\to\infty\) with failure probability bounded
away from zero and embed \(K_{n_j}=K_j\) in a divergent deterministic
sequence.  The preceding uniform conclusion contradicts that subsequence.
\end{proof}

\begin{lemma}[Saddles and squared-norm densities]
\label{edge:norm-density}
On the spectral restrictions of Lemma~\ref{edge:spectral}, the saddle
and length conclusions in Proposition~\ref{prop:edge-estimates} hold.
For a saddle field in a fixed ball let
\[
 Y\sim N(m,\widehat R),\qquad \widehat R=R_{p,\alpha_p(h)},\quad m=\widehat Rh,\quad
 \sigma^2=2\Tr\widehat R^2+4m^{\mathsf T}\widehat Rm.
\]
Then \(\sigma^2\asymp n/p\), and the density \(f_Y\) of \(\|Y\|^2\)
satisfies \(f_Y(n)\asymp\sqrt{p/n}=p^2/\sqrt{k_p}\).
In standard-deviation units this density converges uniformly on
\(\mathbb R\) to the standard normal density as \(K\to\infty\).

For a fixed number \(m_0\) of replicas, suppose the spectral rows
\(Y_i\in\mathbb R^{m_0}\) are independent Gaussian vectors and at least
\(L>2\) row covariances obey
\(cu^{-1}I_{m_0}\preceq C_i\preceq Cu^{-1}I_{m_0}\).
Their means are arbitrary.  The joint density of
\(\mathcal Q=(\sum_iY_{i1}^2,\ldots,\sum_iY_{im_0}^2)\) obeys
\begin{equation}\label{edge:replica-norm-density}
 \sup_z f_{\mathcal Q}(z)\le C_{m_0}(u/\sqrt L)^{m_0}.
\end{equation}
This applies after rowwise Gaussian tilts which preserve the stated band
and independence of the rows.
\end{lemma}

\begin{proof}
Put \(t=u/p^2\) and
\(\mathfrak d_p(t)=(n-\Tr R_{p,tp^2})/(np)\).
For \(y=\|h\|/(p^2M_p)\), the saddle equation becomes
\begin{equation}\label{edge:scaled-saddle}
 \mathfrak d_p(t)=\frac{y^2}{(1+t)^2},\qquad t>-1,\qquad
 \mathfrak d'_p(t)=\frac p n\Tr R_{p,tp^2}^2.
\end{equation}
Lemma~\ref{edge:spectral} gives \(\mathfrak d_p(0)\) arbitrarily close
to \(b_0\), and
\[
 \mathfrak d'_p(t)\ge
       \frac{\dim\cH_p/k_p}{(1+t)^2}.
\]
It follows by integration that
\(\mathfrak d_p(-1+\kappa)<0\) for a sufficiently small fixed
\(\kappa>0\); no saddle can lie there or below.  On \(t\ge0\),
\(\mathfrak d_p(t)\ge b_0/2\), so \eqref{edge:scaled-saddle} bounds
\(t\) above when \(y\le R\).  The right side of \eqref{edge:saddle},
viewed as a function of its parameter, is strictly decreasing and
continuous, tends to infinity at \(-p^2\), and tends to zero at infinity.
This also proves existence and uniqueness.  Denominator comparison gives
\(\Tr\widehat R^2\asymp n/p\), \(\|m\|^2\le C_Rnp\), and
\(\|\widehat R\|\le C_Rp^{-2}\), hence \(\sigma^2\asymp n/p\).
For a sufficiently small fixed \(s_0\), both
\(\mathfrak d_p(t)\) and \(\mathfrak d'_p(t)\) are bounded above and
below on \(|t|\le s_0\).  Differentiating
\(\ell_p(tp^2)=p^2M_p(1+t)\sqrt{\mathfrak d_p(t)}\) proves the length
and derivative bounds.

We record the density argument used here and later.  For independent
\(Y_i=a_i+\sqrt{\upsilon_i}G_i\), suppose
\[
 \max_i\upsilon_i\le Cp^{-2},\qquad
 \sigma^2:=2\sum_i\upsilon_i^2+4\sum_i\upsilon_i a_i^2
       \asymp k_pp^{-4},
\]
and at least \(ck_p\) variances are at least \(cp^{-2}\).  The Gaussian
quadratic formula, as in \cite[Section~15.3, Equation~(15.5)]{CM}, gives
\begin{equation}\label{edge:norm-characteristic}
 \left|\E e^{it\sum_iY_i^2}\right|
 =\prod_i(1+4t^2\upsilon_i^2)^{-1/4}
   \exp\left(-\sum_i
       \frac{2t^2\upsilon_i a_i^2}{1+4t^2\upsilon_i^2}\right).
\end{equation}
Expansion of the centered log characteristic function through order two
gives
\begin{equation}\label{edge:norm-expansion}
 \log\E e^{it(\|Y\|^2-\E\|Y\|^2)/\sigma}
       =-t^2/2+O(k_p^{-1/2}|t|^3)
       \quad (|t|\le c\sqrt{k_p}).
\end{equation}
Indeed its remainder is at most
\[
 C|t/\sigma|^3\sum_i(\upsilon_i^3+\upsilon_i^2a_i^2)
       \le Ck_p^{-1/2}|t|^3.
\]
The band in \eqref{edge:norm-characteristic} bounds the modulus in these
units by \((1+ct^2/k_p)^{-c'k_p}\).  Its real logarithm is at most
\(-c_0t^2\) up to a small fixed multiple of \(\sqrt{k_p}\), and its
integral outside that interval is exponentially small in \(k_p\).
Fourier inversion and \eqref{edge:norm-expansion} give
\begin{equation}\label{edge:norm-local-limit}
 \sup_z\left|\sigma f_Y(\E\|Y\|^2+\sigma z)
       -(2\pi)^{-1/2}e^{-z^2/2}\right|\le Ck_p^{-1/2}.
\end{equation}
At the saddle \(\E\|Y\|^2=n\), proving the lower and upper density bounds.
At any mean, the same characteristic majorant gives the upper bound
\(Cp^2/\sqrt{k_p}\).

For replicas let \(T_t=\diag(t_1,\ldots,t_{m_0})\) and
\(H_i=C_i^{1/2}T_tC_i^{1/2}\).  The modulus of the row characteristic
function is at most \(\det(I+4H_i^2)^{-1/4}\); its extra factor from a
nonzero mean is at most one.  The singular values of \(H_i\) are at
least \(cu^{-1}\) times those of \(T_t\), including when the entries of
\(T_t\) have different signs.  Independence of the \(L\) rows bounds
the joint characteristic function by
\(\prod_j(1+ct_j^2/u^2)^{-L/4}\).  Its integral proves
\eqref{edge:replica-norm-density}.  This is the argument of
\cite[Section~15.3, Equation~(15.6)]{CM};
it explains the independence requirement on the spectral rows.
\end{proof}

\begin{lemma}[Spherical projection and pair tails]
\label{edge:spherical-tails}
Fix a field radius and a fixed number of replicas.  On the spectral
restrictions above, for every \(p\) with \(k_p\ge K\) and \(p\le C_1r_n\),
both the uncapped zero-field spherical law and the capped spherical law
at a field in that ball satisfy
\begin{equation}\label{edge:spherical-tail-bounds}
 \begin{aligned}
 \Pr\{\|P_pX\|>vM_p\}&\le C e^{-c k_pv^6}
       &&(v\ge v_0),\\
 \Pr\{|n^{-1}X^{\mathsf T}X'|>w\}&\le C e^{-c n w^3}
       &&(Cp\le w\le\rho_0).
 \end{aligned}
\end{equation}
Here \(X,X'\) are independent replicas at the same field,
\(\rho_0>0\) is fixed and small, and the second bound holds for every
pair of a fixed finite family.  The constants are uniform in \(p,n\).
\end{lemma}

\begin{proof}
We give the scale check in the polar argument of
\cite[Section~15.3, Equations~(15.5)--(15.8)]{CM}.
For an interaction \(A\), write \(R_{A,u}=(2+u-A)^{-1}\) and
\[
 q_{A,u}(h)=c_n+\tfrac n2(2+u)+\tfrac12\log\det R_{A,u}
                      +\tfrac12h^{\mathsf T}R_{A,u}h,\qquad
 c_n=\log\!\left(2^{n/2}\Gamma(n/2)n^{1-n/2}\right).
\]
Coarea in squared radius gives, exactly,
\begin{equation}\label{edge:frozen-density}
 \phi_A^{\rm sp}(h)=q_{A,u}(h)+\log f_{A,u,h}(n),\qquad
 \partial_uq_{A,u}(h)
   =\tfrac12\bigl(n-\Tr R_{A,u}-\|R_{A,u}h\|^2\bigr),
\end{equation}
where \(f_{A,u,h}\) is the density of
\(\|N(R_{A,u}h,R_{A,u})\|^2\).  In the capped case take \(A=W_p\)
and start at its saddle \(u=\alpha\).  For \(u\ge p^2\),
\[
 n-\Tr R_{p,u}
 \le Cn\sqrt u+Ck_pp^2/u^2\le Cn\sqrt u.
\]
The first term is \eqref{edge:source-resolvents}; the second compares
the clipped active denominators with \(d_i+u\), using
\(\dim\cH_p\le Ck_p\).  For \(u\) in a bounded range in \(p^2\) units,
the upper bound is \(C_Rnp\).  Integrating
\eqref{edge:frozen-density} yields
\begin{equation}\label{edge:frozen-normalizer-cost}
 q_{W_p,u}(h)-q_{W_p,\alpha}(h)\le Cnu^{3/2}+C_Rk_p,
\end{equation}
where the first constant is independent of the field radius.

For a rowwise quadratic test \(\mathcal T\) of \(m_0\) replicas and
\(\lambda>0\), let \(f_{\mathcal Q}^{u,\lambda}\) be the norm density
under the normalized Gaussian tilt \(e^{\lambda\mathcal T}\) of
\(m_0\) independent copies of \(N(R_{p,u}h,R_{p,u})\).  Exponential
Markov and polar coarea give
\begin{equation}\label{edge:polar-tail}
 \begin{split}
 (\mu_{p,h}^{\rm sp})^{\otimes m_0}\{\mathcal T\ge b\}
 \le{}&
 e^{m_0(q_{W_p,u}(h)-q_{W_p,\alpha}(h))-\lambda b}
 \E_u e^{\lambda\mathcal T}
 \frac{f_{\mathcal Q}^{u,\lambda}(n\mathbf1)}
      {f_{W_p,\alpha,h}(n)^{m_0}} .
 \end{split}
\end{equation}
This compares norm densities before and after the normalized tilt; it
does not divide by the probability of a shrinking shell.

For a pair take \(\mathcal T=\pm X^{\mathsf T}X'\), \(b=nw\),
\(u=\iota w^2\), and \(\lambda=c'u\), with \(c'>0\) a sufficiently
small absolute constant.  We will choose \(\iota>0\) small and then
the lower multiple of \(w/p\) large, so that \(u\gg p^2\).
The Gaussian formula gives
\begin{equation}\label{edge:quadratic-mgf}
 \log\E_u e^{\lambda\mathcal T}
 \le C\lambda^2\Tr R_{p,u}^2+C\lambda\|R_{p,u}h\|^2
 \le Cnu^{3/2}+C_Rk_p.
\end{equation}
Indeed \(\Tr R_{p,u}^2\le Cn/\sqrt u\) and
\(\|R_{p,u}h\|^2\le C_Rnp^5/u^2\); the pair quadratic has zero first
trace in the centered Gaussian formula.  For
\(\mathcal T=\|P_pX\|^2\), \(b=nw\), its extra first trace costs only
\(Ck_p\).  The choice of \(c'\) preserves positive precision and
covariance comparable to \(u^{-1}I_{m_0}\) on \(ck_p\) active rows.
Lemma~\ref{edge:norm-density} bounds the density ratio in
\eqref{edge:polar-tail} by \(C_{m_0}(u/p^2)^{m_0}\).

For uncapped zero field start at \(u=p^2\).  The resolvent gives
\(n-\Tr(d+p^2)^{-1}\asymp np\), a leading variance comparable to
\(p^{-2}\), and variance \(O(n/p)\) for the other squares.  Chebyshev's
inequality puts the residual squared length to be supplied by that
coordinate in \([cnp,Cnp]\) with probability \(1-C/k_p\).  On this
interval its Gaussian-square density is at least
\(c\sqrt{p/n}e^{-Ck_p}\).  Convolution gives the same lower density
scale, with this extra exponential factor.  The uncapped resolvent
also gives \eqref{edge:frozen-normalizer-cost}, and the tilted upper
density has the same band bound.

The logarithm of the resulting probability is at most
\[
 -c'\iota nw^3+C\iota^{3/2}nw^3+C_Rk_p
       +O_{m_0}(\log(u/p^2)).
\]
Choose \(\iota\) so that \(C\sqrt\iota<c'/4\), and then \(w/p\) larger
than a sufficiently large fixed constant.  Since
\(nw^3=k_p(w/p)^3\), the last two terms are absorbed by another
\(c'\iota nw^3/4\) when \(K\) is large.  This proves the pair bound.
For projection put \(w=v^2p\le1\); then \(nw^3=k_pv^6\).  For
\(v>p^{-1/2}\) the event is empty because \(\|X\|=\sqrt n\).
\end{proof}

\begin{corollary}[Growth of the limiting spectral rates]
\label{edge:limit-growth}
For GOE-almost every \(g\), and every \(b\) with \(b+g_1>0\), there are
finite positive constants \(c_{g,b},C_{g,b}\) such that
\begin{equation}\label{edge:limit-rate-bounds}
 c_{g,b}(1+i)^{2/3}\le g_i+b\le C_{g,b}(1+i)^{2/3}
 \qquad(i\ge1).
\end{equation}
They can be chosen measurably when \(b\) is a measurable admissible
function of \(g\).
\end{corollary}

\begin{proof}
Fix errors in Lemma~\ref{edge:spectral} for which the number of points
with \(d_i\le p^2\) is between \(c np^3\) and \(Cnp^3\), and apply it
at the countable scales \(np^3=2^jK\).  For any finite list, finite
edge-coordinate convergence transfers the inequalities, with slightly
enlarged constants, to counts of limiting points below
\((2^jK)^{2/3}\).  A fixed small buffer in the endpoints may be used;
the limiting process has no point at any of these countably many fixed
endpoints.  Increase the finite list and use continuity of probability
for decreasing events.  The uniform finite-\(n\) event shows that all
the limiting inequalities hold with probability arbitrarily close to one
when \(K\) is large.  Letting that failure probability tend to zero
shows that almost surely, after a finite threshold,
\[
 cL^{3/2}\le\#\{i:g_i\le L\}\le CL^{3/2}
\]
on dyadic \(L\), and hence on all large \(L\) by monotonicity.  Invert
these inequalities, add \(b\), and adjust the bounds for the finitely
many remaining positive numbers \(g_i+b\).  The least integer threshold
and rational bounds give measurable choices.
\end{proof}

\subsection{Comparison with the cube}

The spherical estimates control deterministic fields.  We next obtain
one disorder event on which the corresponding cube estimates hold for
every field and scale.  The additional work is confined to
\(k_p=O(\sqrt{\log n})\).

\begin{lemma}[Two-replica overlap cells]\label{edge:gram-cell}
Let \(\mathscr R_n=\{-1+2j/n:0\le j\le n\}\) and
\(I_r=[r-1/n,r+1/n]\).  The overlap of two independent uniform cube
spins has mass
\[
 p_n(r)=2^{-n}\binom{n}{n(1+r)/2},\qquad r\in\mathscr R_n,
\]
whereas the overlap of two independent uniform spherical spins has
density
\[
 f_n(s)=\frac{\Gamma(n/2)}{\sqrt\pi\,\Gamma((n-1)/2)}
             (1-s^2)^{(n-3)/2},\qquad -1<s<1.
\]
For a sufficiently small fixed \(\rho_0>0\), uniformly for large \(n\)
and \(r\in\mathscr R_n\) with \(|r|\le\rho_0\),
\begin{equation}\label{edge:pair-cell-mass}
 p_n(r)\le e^{C+Cnr^4}\int_{I_r}f_n(s)\,\mathrm ds.
\end{equation}
Let \(H_2\) be the uniform law of an ordered orthonormal pair
\((q,q')\).  If a positive measurable angular weight
\(a_n(s,q,q')\) has finite integrals and satisfies, for \(H_2\)-almost
every frame,
\[
 |\log a_n(s,q,q')-\log a_n(r,q,q')|\le K_0
 \quad\text{for almost every }s\in I_r,
\]
then
\begin{equation}\label{edge:pair-cell-weighted}
 p_n(r)\int a_n(r,q,q')\,\mathrm dH_2
 \le e^{C+Cnr^4+K_0}
       \int_{I_r}\!\int a_n(s,q,q')\,\mathrm dH_2\,f_n(s)\,\mathrm ds.
\end{equation}
The cells have disjoint interiors.  Summing over their midpoints in an
overlap interval enlarges the spherical interval by at most \(1/n\)
at each endpoint.
\end{lemma}

\begin{proof}
We derive this form for cells centered at lattice points from the two-replica calculations in
\cite[Section~15.4, Equation~(15.10)]{CM} and
\cite[proof of Lemma~3.2]{FUA}.  Set
\[
 I(r)=\tfrac12\bigl((1+r)\log(1+r)+(1-r)\log(1-r)\bigr),
 \qquad J(r)=-\tfrac12\log(1-r^2).
\]
Uniform Stirling bounds on \(|r|\le\rho_0\) put \(p_n(r)\) between
fixed multiples of \(n^{-1/2}e^{-nI(r)}\).  The same bounds for the
gamma ratio, and the bounded logarithmic oscillation of \(f_n\) on
\(I_r\), put its cell integral between fixed multiples of
\(n^{-1/2}e^{-nJ(r)}\).  Since \(J(r)-I(r)=O(r^4)\), this proves
\eqref{edge:pair-cell-mass}.  The assumed oscillation of the positive
angular weight gives \eqref{edge:pair-cell-weighted} by integration.
The final assertion follows directly from the cell endpoints.
\end{proof}

\begin{lemma}[A Haar second moment at the low scales]
\label{edge:replica-comparison}
On the spectral restrictions above, for every fixed spectral field in
a prescribed ball and every \(K\le k_p\), \(p\le C_1r_n\),
\begin{equation}\label{edge:likelihood-second-moment}
 E_U(L_p(h)-1)^2\le C_Rn^{-1/4}.
\end{equation}
The field is fixed inside the expectation; the constant is uniform over
the indicated fields and scales.
\end{lemma}

\begin{proof}
We adapt the two-replica calculation of
\cite[Lemma~3.2, Equations~(3.3)--(3.4)]{FUA}, whose stated observation
range starts at \(r_n\).  The estimate below the source range is proved
here.
Conditional on \(\rho=n^{-1}X^{\mathsf T}X'\), the Haar angular
integral for a cube pair is the spherical angular integral.  On the
ordered orthonormal sum-and-difference frame \((q,q')\), its exponent is
\begin{equation}\label{edge:pair-angular-exponent}
 \Psi_{p,h}(s;q,q')=
 \frac n2\bigl((1+s)q^{\mathsf T}W_pq
       +(1-s){q'}^{\mathsf T}W_pq'\bigr)
       +\sqrt{2n(1+s)}\,h^{\mathsf T}q.
\end{equation}
On a fixed \(|s|\le2\rho_0<1\), the spectral norm bound and the field
bound \(\|h\|\le R\sqrt n\,p^{5/2}\) give
\(|\partial_s\Psi_{p,h}|\le C_Rn\).  Thus the positive angular weight
\(e^{\Psi_{p,h}}\) has a fixed logarithmic oscillation on every
width-\(2/n\) cell, as required in \eqref{edge:pair-cell-weighted}.
Retain
\(|\rho|\le C'r_n\), choosing \(C'\) larger than the lower tail
constant in Lemma~\ref{edge:spherical-tails} times \(C_1\).
On dyadic bands between \(C'r_n\) and a small fixed \(\rho_0\),
the weighted cell comparison costs \(e^{C+Cn\rho^4}\), while
Lemma~\ref{edge:spherical-tails} gives \(Ce^{-cn\rho^3}\).
The spherical band endpoints move by at most \(1/n\), negligible
compared with \(r_n\).
Decreasing \(\rho_0\) makes the latter pay the former and leaves
\(Ce^{-cnr_n^3}\).  At overlaps bounded away from zero, the global
spectral law, edge limit, and norm bound in \eqref{edge:global-spectral}
supply the global inputs for the uncapped zero-field angular comparison
of \cite[Section~15.4]{CM}.  Indeed, for every fixed \(w>2\), they give
\[
 \frac1n\log\det(wI-W)\longrightarrow
       \int\log(w-x)\rho_{\rm sc}(\dd x);
\]
interlacing gives the same per-site limit for every orthogonal-hyperplane
compression.  The cited deficit-filling lower bound then gives the full
critical free-energy normalization by taking \(n\to\infty\) at fixed
\(w>2\), followed by \(w\downarrow2\).  For clarity, the normalized
zero-field contribution being bounded is
\[
 \mathcal M_n(\rho_0;\boldsymbol\lambda)
   =E_U\!\left[L_\circ^2\,
       \mu_W^{\otimes2}\{|n^{-1}X^{\mathsf T}X'|\ge\rho_0\}\right].
\]
Here \(\boldsymbol\lambda=(\lambda_1,\ldots,\lambda_n)\) is fixed,
\(W=U\diag(\boldsymbol\lambda)U^{\mathsf T}\), and \(\mu_W\) is the
zero-field cube Gibbs law for this \(W\).
The comparison in
\cite[Section~15.4, ``Zero field and overlaps bounded away from zero'']{CM},
together with the lower normalization in its Section~15.3, gives
\(\mathcal M_n(\rho_0;\boldsymbol\lambda)\le
\exp\{-\kappa_{\rho_0}n+o(n)\}\) for a fixed \(\kappa_{\rho_0}>0\),
along the retained spectral sequences with the preceding per-site
errors uniform.  This is a consequence of that proof; the norm bound
supplies continuity at the endpoints.
Changing to \(W_p,h\) changes weights and log normalizers by
\(O_R(np^2+np^{5/2})=o(n)\), so that region still contributes \(e^{-cn}\).

On retained cells of width \(2/n\), Stirling's formula for the binomial
overlap mass and the spherical density proportional to
\((1-\rho^2)^{(n-3)/2}\) gives relative error
\begin{equation}\label{edge:retained-stirling}
 O(nr_n^4+r_n+n^{-1}).
\end{equation}
This is the cell estimate in the cited replica proof, now with cutoff
\(C'r_n\) at every \(p\).  More precisely than the fixed bound above,
the oscillation of \eqref{edge:pair-angular-exponent} on one cell is at most
\[
 C\bigl(|q^{\mathsf T}W_pq-{q'}^{\mathsf T}W_pq'|+p^{5/2}\bigr),
\]
which is uniformly bounded.  The normalized spherical integral of
this expression on retained cells is \(O(r_n)\).  Indeed,
\[
 q^{\mathsf T}W_pq-{q'}^{\mathsf T}W_pq'
 =\frac{2X^{\mathsf T}W_pX'
       -\rho(X^{\mathsf T}W_pX+{X'}^{\mathsf T}W_pX')}
        {n(1-\rho^2)}.
\]
The diagonal terms cost \(O(r_n)\).  For the cross term, tilt a Gaussian
pair at the cap saddle by \(\pm cp^2X^{\mathsf T}W_pX'\), with \(c\)
small enough to preserve a fixed precision margin.  The Gaussian
logarithmic moment is \(O_R(k_p)\), since
\(\Tr R_{p,\alpha}^2=O(n/p)\) and
\(\|R_{p,\alpha}h\|^2=O_R(np)\).  The replica density bound
\eqref{edge:replica-norm-density} costs a fixed factor.  Hence
\[
 \Pr_{\rm sp}\{|X^{\mathsf T}W_pX'|>t\}
       \le C e^{C_Rk_p-cp^2t}.
\]
Integrating from a sufficiently large multiple of \(np\) gives
\(\E_{\rm sp}|X^{\mathsf T}W_pX'|\le C_Rnp\le C_Rnr_n\).
This full probability integral also bounds its unnormalized restriction
to retained cells.

The bounded oscillation permits the mean value theorem for exponentials.
After integration and summation over disjoint cells, the error is the
preceding integrated \(O(r_n)\); no supremum over frames is taken.
Slightly enlarged boundary cells are paid by the pair tail.  Thus
\[
 E_UL_p(h)^2
 \le1+C_R\{nr_n^4+r_n+n^{-1}+e^{-cnr_n^3}\}.
\]
Use \(E_UL_p(h)=1\),
\(nr_n^4=n^{-1/3+4\cdot10^{-4}}\), and
\(nr_n^3=n^{3\cdot10^{-4}}\) to obtain
\eqref{edge:likelihood-second-moment}.
\end{proof}

\begin{lemma}[Uniform cube cap estimates]\label{edge:cube-caps}
The static, curvature, and projection assertions
\eqref{edge:static-estimates}--\eqref{edge:projection-tails} hold on
ordinary Gaussian disorder events of arbitrarily large limiting
probability, with a tight threshold and every fixed field radius,
tolerance, and scale buffer prescribed in
Proposition~\ref{prop:edge-estimates}.  The projection bound also holds
for \(P_qX/M_q\) under a cap interpolated linearly between \(B_p\) and
\(B_q\), whenever \(q\ge p\), \(q/p\) is bounded, and the field lies
in a fixed \(q\)-field ball.
\end{lemma}

\begin{proof}
On dyadic \(p\) and \(v\), the spherical tail and
\eqref{edge:bias-transfer} give
\[
 E_U[\ind_{\{L_\circ\ge1/2\}}
       \mu\{\|P_pX\|>vM_p\}]
       \le C e^{-c k_pv^6}.
\]
Markov's inequality with half the exponent has total failure at most
\[
 \sum_{j,l\ge0}C e^{-cK\,8^j64^l}\le C e^{-c'K}.
\]
The discarded \(L_\circ<1/2\) event has probability \(o(1)\) by
\eqref{edge:haar-bias}.  Monotonicity of \(P_p\), and a fixed-factor
change in \(M_p\), extend the uncapped result to every \(p,v\).

The posterior tails follow deterministically.  For \(q\ge p\) of bounded
ratio, let \(B_\theta\) lie between \(B_p\) and \(B_q\).  It is positive,
supported on \(\cH_q\), and has norm at most \(Cq^2\).  Restricting the
uncapped normalizer to \(\|P_qX\|\le VM_q\) gives
\[
 \mu e^{-X^{\mathsf T}B_\theta X/2}\ge e^{-Ck_q}
\]
for a fixed large \(V\).  The zero-field capped law is even, so its
normalizer for a linear field is at least one.  On the shell
\(\|P_qX\|\asymp vM_q\), a field of norm at most \(Rq^2M_q\) costs
at most \(Rk_qv\).  This and the denominator cost are absorbed by
the uncapped \(e^{-c k_qv^6}\) tail for large \(v\).  This proves the
posterior tail, including the interpolated assertion.  The spherical
argument is identical.  The first two projection moments under all
these posteriors are consequently bounded by \(C_RM_q\) and \(C_RM_q^2\).

For nearby scales and fields, integrating first derivatives of the log
partition function gives
\begin{equation}\label{edge:potential-interpolation}
 |\Delta\phi|\le C_Rk_q\left(
       |\Delta\log p|+\frac{\|\Delta h\|}{q^2M_q}\right),
\end{equation}
and likewise for \(\phi^{\rm sp}\).  Indeed,
\(\|B_q-B_p\|\le Cq^2|\log(q/p)|\) on \(\cH_q\); its derivative is
controlled by the second projection moment and the field derivative
by the first.  Fields in the smaller space are compared in \(\cH_q\).

Cover \(K\le k_p\le C_0\sqrt{\log n}\), with fixed \(C_0\) large enough
to overlap the source cutoff.  Net scales from above and their spectral
field balls at resolution
\(\eta/(C_R(1+C_0\sqrt{\log n}))\) in the units of
\eqref{edge:potential-interpolation}.  Conditional on the spectrum,
these are deterministic nets.  Since \(\dim\cH_p\le Ck_p\), their
total size is
\[
 \exp\{O_{\eta,R}(\sqrt{\log n}\log\log n)\}=n^{o(1)}.
\]
At a net point \eqref{edge:likelihood-second-moment} and Chebyshev give
\(\Pr_U(|\log L_p(h)|>\eta/2)\le C_\eta n^{-1/4}\).
The union bound and \eqref{edge:potential-interpolation} yield
\begin{equation}\label{edge:absolute-value-comparison}
 \sup_{\substack{K\le k_p\le C_0\sqrt{\log n}\\
                  \|h\|\le Rp^2M_p}}
 |\phi_p(h)-\phi_p^{\rm sp}(h)|\le\eta
\end{equation}
with failure \(n^{-1/4+o(1)}\), in addition to the projection event.
For \(k_p\ge c\sqrt{\log n}\), the cap and interpolation results of
\cite[Proposition~2.6, Equation~(2.11), and Lemma~2.7]{FUA} apply as
stated, since \(C_1r_n\) is eventually below their fixed upper cutoff.
Together these prove the value comparison at every scale.

For the mean, use this comparison on a larger ball.  For unit
\(v\in\cH_p\), freeze \(\alpha=\alpha_p(h)\).  The polar formula and
Lemma~\ref{edge:norm-density} give
\[
 \log\mu_{p,h}^{\rm sp}
       e^{\pm t v^{\mathsf T}(X-R_{p,\alpha}h)}
       \le\tfrac12t^2v^{\mathsf T}R_{p,\alpha}v+C.
\]
The cube bound has the additional \(2\epsilon_vk_p\).
Jensen with \(t=\theta p^2M_p\), in the direction of the projected
mean discrepancy, gives
\[
 \frac{\|P_pm_p(h)-R_{p,\alpha}h\|}{M_p}
       \le C_R\theta+\frac{2\epsilon_v+C/k_p}{\theta}.
\]
Choose \(\theta\) from the required mean tolerance, then impose a smaller
value tolerance and increase \(K\).

It remains to prove the Hessian assertion on the low range.  Put
\(\widehat R=R_{p,\alpha_p(h)}\), \(m=\widehat Rh\), and
\(\sigma^2=2\Tr\widehat R^2+4m^{\mathsf T}\widehat Rm\).  For a unit \(v\in\cH_p\)
and \(|\theta|\le1\), the shift \(h+\theta pv\) fits a fixed larger
ball because \(p/(p^2M_p)=k_p^{-1/2}\).  Freezing \(\alpha\), the
exact polar formula is
\[
 \log\mu_{p,h}^{\rm sp} e^{\theta p v^{\mathsf T}(X-m)}
 =\tfrac12\theta^2p^2v^{\mathsf T}\widehat Rv+
       \log\frac{f_{\alpha,h+\theta pv}(n)}{f_{\alpha,h}(n)}.
\]
The shifted mean \(m_\theta=m+\theta p\widehat Rv\) satisfies
\[
 \frac{\sigma_\theta^2}{\sigma^2}=1+O_R(k_p^{-1/2}),\qquad
 \frac{n-\E_\theta\|Y\|^2}{\sigma_\theta}
       =-\frac{2\theta p\langle m,\widehat Rv\rangle}{\sigma}
           +O_R(k_p^{-1/2}).
\]
The leading term is uniformly bounded.  The density estimate
\eqref{edge:norm-local-limit}, at this argument and at zero, gives
\begin{equation}\label{edge:frozen-mgf-curvature}
 \begin{split}
 \log\mu_{p,h}^{\rm sp} e^{\theta p v^{\mathsf T}(X-m)}
   &=\tfrac12\theta^2v^{\mathsf T}\Gamma_{p,h}v+O_R(k_p^{-1/2}),\\
 \Gamma_{p,h}
   &=p^2\left(\widehat R|_{\cH_p}
       -\frac4{\sigma^2}(\widehat Rm)\otimes(\widehat Rm)\right).
 \end{split}
\end{equation}
This is uniform in \(h,v,\theta\).
Cauchy--Schwarz and \(\sigma^2\ge4m^{\mathsf T}\widehat Rm\) imply
\(0\preceq\Gamma_{p,h}\preceq C_RI\).
Equation~\eqref{edge:absolute-value-comparison} changes the log moment
by at most \(2\eta\), so the cube error is at most
\(\delta=C_RK^{-1/2}+2\eta\).

To justify differentiating this approximation, suppose a real \(Z\) obeys
\[
 \sup_{|\theta|\le1}
 \left|\log\E e^{\theta Z}-a\theta^2/2\right|\le\delta,\qquad
 0\le a\le C,\quad 0<\delta\le1/16.
\]
The values at \(\pm1\) bound \(\E e^{|Z|}\), and hence bound the
fourth derivative of \(\psi(\theta)=\log\E e^{\theta Z}\) on
\([-1/2,1/2]\).  Taylor's formula gives
\[
 \left|\psi''(0)
       -\frac{\psi(t)+\psi(-t)-2\psi(0)}{t^2}\right|\le C_Ct^2.
\]
The assumed error changes this quotient by at most \(2\delta/t^2\).
With \(t=\delta^{1/4}\), this proves
\(|\Var Z-a|\le C_C\sqrt\delta\).  Applying it to
\(Z=pv^{\mathsf T}(X-m)\), for every unit \(v\), yields
\begin{equation}\label{edge:hessian-approximation}
 \|p^2K_p(h)-\Gamma_{p,h}\|
       \le C_R\sqrt{C_RK^{-1/2}+2\eta}.
\end{equation}

On the annulus put \(t=\alpha/p^2\),
\(A=(n-\Tr\widehat R)/(np)\), \(B=(p/n)\Tr\widehat R^2\), and \(e=h/\|h\|\).
Since \(\widehat R|_{\cH_p}=p^{-2}(1+t)^{-1}I\) and \(m\) is parallel to \(e\),
\[
 \langle e,\Gamma_{p,h}e\rangle=\frac{B}{(1+t)B+2A}.
\]
The trace estimates and denominator comparison give
\(A=b_0+O(\epsilon_{\rm sp}+s_0)\) and
\(B=2b_0+O(\epsilon_{\rm sp}+s_0)\), so this is
\(1/2+O(\epsilon_{\rm sp}+s_0)\).  All other quadratic directions are
bounded above by \((1-s_0)^{-1}\).  Choose the spectral accuracy and
\(s_0\) small, then \(\eta\) small and \(K\) large in
\eqref{edge:hessian-approximation}.  This proves the annular curvature
and radial deficit.  On the entire ball, the saddle gap and the same
estimate give \(C_Rp^{-2}\), with \(C_R\) fixed before the annular
tolerance.  The source curvature assertion covers the higher range.
\end{proof}

\subsection{Location and covariance of the observation path}

The uniform cap event controls external fields.  We now locate the random
field produced by the Gibbs observation.  A shrinking annulus will give
an averaged covariance estimate; a fixed annulus gives the stronger
exponential tail needed in Proposition~\ref{prop:edge-estimates}.

We first state the additional eigenvalue-only restriction used for norm
conditioning.  Put
\[
 b_n=1+(-g_{1,n})_+,\qquad
 \ell_{a,n}=(g_{a,n}+b_n)^{-1},\qquad
 D_n=n^{1/3}-\sum_a\ell_{a,n},
\]
and let \(q_n\) be the density of
\(\sum_a\ell_{a,n}(\xi_a^2-1)\) for independent standard normal
\(\xi_a\).  For every \(\chi>0\), one may fix sufficiently large
integers \(m,J_\chi\) and positive constants \(c_\chi,C_\chi\) so that
an eigenvalue-only event \(\mathcal S_n^{\rm ann}(\chi)\) satisfies
\(\liminf_n\Prob(\mathcal S_n^{\rm ann}(\chi))\ge1-\chi\) and, on it,
\begin{equation}\label{edge:conditioned-spectral-event}
 \begin{gathered}
 b_n\le C_\chi,\qquad \ell_{a,n}\ge c_\chi\quad(1\le a\le m),\\
 g_{a,n}\ge C_\chi^{-1}a^{2/3}\quad(J_\chi<a\le n),\qquad
 q_n(D_n)\ge c_\chi.
 \end{gathered}
\end{equation}
These restrictions follow from the eigenvalue bounds in
\cite[Lemma~2.1]{FUA}, finite edge-coordinate convergence, and the
density argument in its Proposition~2.3.  We intersect this event with the local spectral
event of Lemma~\ref{edge:spectral} and \(\mathcal S_n^{\rm glob}\),
assigning a portion of the prescribed failure probability to it.  The
eigenvalue events of \cite[Lemma~2.5]{FUA} needed for
\eqref{edge:haar-bias} are retained separately.  All constants below
may depend on the prescribed confidence; no fixed density lower bound
on an event of probability \(1-o(1)\) is asserted.

\begin{lemma}[The observation annuli]\label{edge:path-annuli}
On the fixed spectral restrictions, there are \(a_{\rm an},c,C>0\)
such that for every dyadic interval \([q,2q]\) with fixed buffers
inside the enlarged scale range,
\begin{equation}\label{edge:shrinking-annulus}
 \mathbf P\left\{\exists p\in[q,2q]:
       |\alpha_p(h_p)|>Ck_p^{-1/10}p^2\right\}
       \le C e^{-c k_q^{a_{\rm an}}}.
\end{equation}
For every fixed \(\zeta>0\), after increasing the threshold,
\begin{equation}\label{edge:fixed-annulus-biased}
 \mathbf P\left\{\exists p\in[q,2q]:
       |\alpha_p(h_p)|>\zeta p^2\right\}
       \le C_\zeta e^{-c_\zeta k_q}.
\end{equation}
Consequently \eqref{edge:ordinary-annulus} holds simultaneously at
every scale on ordinary disorder events of arbitrarily large limiting
probability.
\end{lemma}

\begin{proof}
By \eqref{edge:haar-bias}, under \(\mathbf P\) the spectral spin is
spherical.  Its Gaussian representation is that of
\cite[Proposition~2.3]{FUA}.  The annulus calculation in
\cite[Lemma~3.1, Equation~(3.1)]{FUA} is stated for the source observation
range beginning at \(r_n\); the tight-threshold dyadic estimates below
extend that calculation to the scales used here.
Set \(D_i'=d_i+b_nn^{-2/3}=n^{-2/3}(g_{i,n}+b_n)\).
Before conditioning, let \(Y_i\) be independent \(N(0,(D_i')^{-1})\);
conditioning on \(\sum_iY_i^2=n\) gives the spherical law.
The definition of \(b_n\) gives \(\min_iD_i'\ge n^{-2/3}\), while
\eqref{edge:conditioned-spectral-event} supplies the fixed leading
variance lower bounds.  The density of
\(n^{-2/3}\sum_iY_i^2\) at \(n^{1/3}\) is exactly \(q_n(D_n)\),
and hence is bounded below on the chosen spectral set.  The same set
and the active count give
\begin{equation}\label{edge:leading-resolvent-sum}
 \sum_{d_i\le p^2}(D_i')^{-1}
   =n^{2/3}\sum_{d_i\le p^2}\ell_{i,n}
   \le Cn^{2/3}(1+k_p^{1/3})\le Cnp.
\end{equation}
Indeed the first \(J_\chi\) terms are bounded by one, and the remaining
terms are bounded by \(C_\chi i^{-2/3}\).

Let \(\vartheta=p^2+b_nn^{-2/3}\) and, before conditioning, write
\(h_{p,i}=B_{p,i}Y_i+B_{p,i}^{1/2}Z_i\).
The quadratic form obtained from
\(\|h_p\|^2-\vartheta^2\|Y\|^2\) has mean
\(-\vartheta^2\Tr R_{p,b_nn^{-2/3}}\).  On an active coordinate the
coefficient of the standardized \(Y_i^2\) term is
\[
 \frac{B_{p,i}^2-\vartheta^2}{D_i'}=D_i'-2\vartheta=O(p^2);
\]
the cross coefficient is \(O(p^3/\sqrt{D_i'})\).
Equation~\eqref{edge:leading-resolvent-sum} and the capped second trace
bound therefore show that the matrix of this quadratic form, in the
standardized \((Y,Z)\) variables, satisfies
\[
 \|\mathsf A_p\|_{\rm F}^2\le Cp^4k_p,\qquad
 \|\mathsf A_p\|\le Cp^2k_p^{1/3}.
\]
The Gaussian quadratic-form moment formula for its centered version bounds a deviation of
\(p^2k_p^{9/10}\) by
\[
 C\exp\left[-c\min\{k_p^{4/5},k_p^{17/30}\}\right].
\]
This remains true after norm conditioning, with a constant loss.
Indeed use the same exponential tilt, with parameter at most one
eighth of \(\|\mathsf A_p\|^{-1}\).  Its joint covariance and its
\(Y\)-marginal covariance stay comparable to the originals.  A fixed
number of the leading marginal variances bounds the density of
\(n^{-2/3}\|Y\|^2\) from above by a constant, through the characteristic
formula \eqref{edge:norm-characteristic}.  The original density at
\(n^{1/3}\) is bounded below on the chosen spectral set.  The polar
density ratio therefore costs a constant.

After conditioning, the deterministic target for \(\|h_p\|^2\) is
\(\ell_p(b_nn^{-2/3})^2\asymp p^2k_p\).  The preceding error is a
relative \(O(k_p^{-1/10})\).  It first places the inferred saddle in
a fixed annulus.  On that annulus
\(\mathrm d(\ell_p(u)^2)/\mathrm du\asymp k_p\).  Hence
\[
 |\alpha_p(h_p)-b_nn^{-2/3}|\le Cp^2k_p^{-1/10}.
\]
Since \(b_nn^{-2/3}/p^2=O(k_p^{-2/3})\), this proves the fixed-scale
version of \eqref{edge:shrinking-annulus}.

For a whole interval use a mesh of spacing \(k_q^{-3}\) in \(\log p\).
Its size is polynomial in \(k_q\).  Except with \(Ce^{-ck_q}\)
probability, the spherical projection at the upper endpoint is at most
\(CM_q\), by Lemma~\ref{edge:spherical-tails}.  The signal increments
on a mesh gap are then at most \(Cq^2M_qk_q^{-3}\).
The noise increments have covariance norm \(O(q^2k_q^{-3})\).
A constant-resolution net of the sphere in \(O(k_q)\) dimensions and
the Gaussian maximal inequality show that their norm on every gap is
less than \(q^2M_qk_q^{-1/10}\), except with \(Ce^{-ck_q}\)
probability.  The target length is Lipschitz in the same scaled units:
between activations differentiate the capped trace, using the active
count and second trace; the target is continuous at activations.
These observations interpolate the fixed-scale bound.  The polynomial
mesh size is absorbed by decreasing the positive exponent, proving
\eqref{edge:shrinking-annulus}.

For a fixed \(\zeta\), it suffices by inclusion to treat \(\zeta\)
below a fixed small annulus width.  The stronger exponent follows from a
different Gaussian reference, following the source-range fixed saddle
test in \cite[Section~19.2, Equation~(19.9) and its proof]{CM}.
Fix a small \(\xi>0\), to be chosen after \(\zeta\), and let
\[
 R_\xi=(2+\xi p^2-W)^{-1},\qquad
 Y\sim N(0,R_\xi),\qquad h_p=B_pY+\beta_p .
\]
Increase \(K\) so that \(d_1\ge-\xi p^2/2\).
The vector \(a_{p,\xi}=R_{p,\xi p^2}h_p\) is centered Gaussian with
covariance
\[
 C_{p,\xi}=R_\xi-R_{p,\xi p^2}.
\]
The resolvent and clipped traces give
\[
 \left|\Tr C_{p,\xi}-(n-\Tr R_{p,0})\right|
       \le C\sqrt\xi\,np+o_K(np),\qquad
 \|C_{p,\xi}\|\le C(\xi p^2)^{-1},\qquad
 \Tr C_{p,\xi}^2\le Cn/(\sqrt\xi\,p).
\]
For a fixed multiple of \(np\), exponential Markov with both signs
of a sufficiently small multiple of \(\xi p^2\) therefore gives
a deviation cost \(e^{-c_\zeta\xi k_p}\), once \(\xi\) is sufficiently
small.  The centered logarithmic moment costs \(O(\xi^{3/2}k_p)\),
whereas the Markov gain is a fixed multiple of \(\zeta\xi k_p\).

We spell out the norm-conditioning cost.  The deficit
\(n-\Tr R_\xi\) is between fixed multiples of \(\sqrt\xi\,np\),
after increasing \(K\) for this fixed \(\xi\).  A leading variance is
comparable to \((\xi p^2)^{-1}\); the other squared norms have variance
at most \(Cn/(\sqrt\xi\,p)\).  Chebyshev places the residual length
in a fixed interval comparable to \(\sqrt\xi\,np\) when
\(K\gg\xi^{-3/2}\).  On this interval the leading Gaussian-square
density is at least
\[
 C_\xi^{-1}\sqrt{p/n}\exp(-C\xi^{3/2}k_p).
\]
This is a lower bound for the original norm density by convolution.
Under the quadratic tilt used above, the marginal covariance of \(Y\)
retains \(ck_p\) independent spectral rows with covariance comparable
to \(p^{-2}\); the tilt keeps the joint covariance comparable.
Equation~\eqref{edge:replica-norm-density} bounds its norm density by
\(C_\xi\sqrt{p/n}\).  Thus conditioning costs at most
\(C_\xi e^{C\xi^{3/2}k_p}\).

Since \(R_{p,0}h_p=(1+\xi)R_{p,\xi p^2}h_p\), compare the saddle
length at the two boundaries \(\alpha=\pm\zeta p^2\).  The derivative
bounds in \eqref{edge:scaled-saddle}, and monotonicity beyond those
boundaries, put
\(\{|\alpha_p(h_p)|>\zeta p^2\}\) inside the relaxed squared-length
deviation just bounded.  Choose \(\xi\) so that its \(O(\sqrt\xi)\)
target error fits the \(\zeta\) tolerance and
\(C\xi^{3/2}<c_\zeta\xi/2\).  The conditioned deviation then has
probability \(C_\zeta e^{-c_\zeta'k_p}\).
The same \(k_q^{-3}\) mesh proves \eqref{edge:fixed-annulus-biased}.

Finally take \(q_j=2^j(K/n)^{1/3}\).  By
\eqref{edge:bias-transfer} and Markov, the ordinary whole-interval
probability is at most \(Ce^{-c k_{q_j}}\) outside disorder probability
\(Ce^{-c k_{q_j}}\), after decreasing \(c\).  Their sum is
\[
 \sum_{j\ge0}Ce^{-cK8^j}\le Ce^{-c'K}.
\]
Use \eqref{edge:haar-bias} for \(L_\circ<1/2\) and increase \(K\).
Every scale lies in one such interval and \(k_p\asymp k_{q_j}\)
there, proving \eqref{edge:ordinary-annulus} simultaneously in \(p\).
\end{proof}

We now fix the finite set of auxiliary radii, buffers, and tolerances
used in the covariance argument.  Fix once a sufficiently small
\(\gamma>0\), as specified by the compression argument below.  Then
choose finite \(C_*\) so that the range through \(C_*r_n\) contains
all parent intervals \([q/(2\gamma),q/\gamma]\), \(q\le r_n\), and
their prescribed enlargements, as well as the range through \(C_1r_n\).
Choose \(\eta\) sufficiently small compared with \(\gamma^2\), and an
auxiliary curvature tolerance \(\epsilon_{\rm aux}\) sufficiently small
compared with \(\eta\); it may be smaller than the prescribed
\(\epsilon_0\).  The extension construction below prescribes an
auxiliary width \(s_{\rm aux}\), static tolerances, radii, and a threshold
for these choices.  They are further cap requests under the
finite-intersection convention after Proposition~\ref{prop:edge-estimates}.
Let \(\mathcal G_n^{\rm aux}\) be the event, measurable from \(W\), on which
Lemma~\ref{edge:cube-caps} holds with these choices.  Let
\(\mathcal G_n^{\rm FUA}\) be the entire good-disorder event
\(\mathcal G_n\) defined in \cite[Section~3]{FUA} for its
Proposition~4.2, with that source's fixed tolerances and full scale and
field ranges.  This event has ordinary Gaussian-disorder probability
tending to one.  Set
\(\mathcal G=\mathcal G_n^{\rm aux}\cap\mathcal G_n^{\rm FUA}\).
Below we use the source covariance conclusion on the terminal interval
\([r_n,C_\gamma r_n]\), which is contained in its stated interval
\([r_n,p_*]\) for all sufficiently large \(n\).
All probabilities below are measured under \(\mathbf P\)
before restriction to \(\mathcal G\).  On each chosen spectral set,
the constants are uniform.

\begin{lemma}[Integrated score and covariance error]
\label{edge:score-covariance}
Let \(I=[q,2q]\), with a prescribed fixed multiplicative enlargement
inside the scale range.  At path fields write
\[
 a_p=P_p(m_p-m_p^{\rm sp}),\qquad D_p=p^2K_p-I_{\cH_p}.
\]
There are \(c,C,a_1>0\) such that
\begin{equation}\label{edge:score-bound}
 \mathbf P\left\{
   \int_I\|a_p\|^2\,\mathrm d(p^2)>Ck_q^{1/2}\right\}
       \le C e^{-c k_q^{a_1}},
\end{equation}
and
\begin{equation}\label{edge:frobenius-bound}
 \mathbf P\left(
 \mathcal G\cap
 \left\{\int_q^{2q}\frac{\|D_p\|_{\rm F}^2}{k_p}\,\mathrm d\log p
              >Ck_q^{-1/5}\right\}\right)
       \le C e^{-c k_q^{a_1}} .
\end{equation}
The same statements hold for every fixed enlargement of \(I\), with
adjusted constants.  They are probability statements; no expectation
of the untruncated score on their failures is asserted.
\end{lemma}

\begin{proof}
Put \(t=p^2\), and include \(U\) in the observation filtration, so that
every disorder event is measurable at time zero.  The exact likelihood
identity \eqref{edge:path-likelihood} yields
\begin{equation}\label{edge:likelihood-sde}
 \mathrm d\log L_{\sqrt t}
       =a_{\sqrt t}\cdot\mathrm d\mathcal B_t
          +\tfrac12\|a_{\sqrt t}\|^2\,\mathrm dt,
\end{equation}
where \(\mathcal B\) is the continuous cube innovation martingale with
\(\mathrm d\langle\mathcal B\rangle_t=P_{\sqrt t}\,\mathrm dt\);
the coefficients lie in the active space.  This is the exact filtering
calculation in \cite[proof of Lemma~3.5, after Equation~(3.7)]{FUA};
it follows as well by differentiating the posterior partition function,
including the interaction change \(-\mathrm dB_{\sqrt t}\).  There is no jump
in \(\log L\) at activation.

At a deterministic endpoint and for \(H>0\),
\eqref{edge:path-likelihood} gives
\[
 \mathbf P(\log L_p<-H)\le e^{-H}.
\]
For a fixed \(R_0\) containing the path annulus, put
\(\mathcal U_p=\{\widetilde h:\|\widetilde h\|\le R_0p^2M_p\}\).
Lemma~\ref{edge:replica-comparison} gives
\[
 \mathbf P(\log L_p>H,h_p\in\mathcal U_p)
 \le e^{-H}\int_{\mathcal U_p}Q^{\rm sp}(\mathrm d\widetilde h)\,
          E_U L_p(\widetilde h)^2
 \le (1+Cn^{-1/4})e^{-H}.
\]
The field in the inner expectation is deterministic.
Lemma~\ref{edge:path-annuli} excludes the complementary endpoint fields
with probability \(Ce^{-c k_q^{a_{\rm an}}}\).

On an enlarged interval write
\(\mathcal J_I=\int\|a_p\|^2\,\mathrm d(p^2)\) and
\(M=\int a_p\cdot\mathrm d\mathcal B_{p^2}\).
On the two endpoint events just given,
\(M+\mathcal J_I/2\le2H\).  If \(\mathcal J_I\ge8H\), then
\(M\le-\mathcal J_I/4\).
The positive exponential supermartingale
\(\exp(-M/4-\mathcal J_I/32)\) bounds this event by \(e^{-H/4}\).
Taking \(H=Ck_q^{1/2}\) proves \eqref{edge:score-bound}.

Between activations the two filtering equations give, with
\(\Delta_t=K_{\sqrt t}-K_{\sqrt t}^{\rm sp}\),
\[
 \mathrm da_{\sqrt t}
       =\Delta_t\,\mathrm d\mathcal B_t
          -K_{\sqrt t}^{\rm sp}a_{\sqrt t}\,\mathrm dt .
\]
Choose a smooth nonnegative cutoff \(\chi(t)\), equal to one on the
target interval, zero at the ends of a fixed enlargement, and with
\(|t\chi'(t)|\le C\).  The integrated Itô identity is
\begin{equation}\label{edge:score-energy}
 \begin{split}
 \int\chi t\|\Delta_t\|_{\rm F}^2\,\mathrm dt
 ={}&-\int(\chi+t\chi')\|a_{\sqrt t}\|^2\,\mathrm dt
       +2\int\chi t\langle a_{\sqrt t},
                         K_{\sqrt t}^{\rm sp}a_{\sqrt t}\rangle\,\mathrm dt\\
 &-\mathcal M
   -\sum_j\chi(t_j)t_j
       \|\Delta P_j(m_{\sqrt{t_j}}-m_{\sqrt{t_j}}^{\rm sp})\|^2,\\
 \mathcal M={}&2\int\chi t
       \langle a_{\sqrt t},\Delta_t\,\mathrm d\mathcal B_t\rangle .
 \end{split}
\end{equation}
Here \(\Delta P_j\) is the jump of \(P_{\sqrt t}\) at the activation
time \(t_j\).
The sum has the favorable sign.  On \(\mathcal G\) and the whole-path
annulus event, the coarse bounds give
\(\|tK_{\sqrt t}\|+\|tK_{\sqrt t}^{\rm sp}\|\le C\).
Thus the left side on the target interval is bounded by
\(C\mathcal J_I+C|\mathcal M|\).

For concentration, stop only the stochastic integral at the first
annulus failure, and set its integrand to zero on \(\mathcal G^c\).
The latter event is measurable at time zero.  The stopped bracket is
at most \(C\mathcal J_I\), and the stopped integral agrees with \(\mathcal M\)
on \(\mathcal G\) and the whole-path annulus event.  The exponential
bracket inequality, \eqref{edge:score-bound}, and
\eqref{edge:shrinking-annulus} consequently bound the target integral
of \(\|p^2(K_p-K_p^{\rm sp})\|_{\rm F}^2\,\mathrm d\log p\) by
\(Ck_q^{1/2}\), with the required exceptional probability intersected
with \(\mathcal G\).  The full identity \eqref{edge:score-energy} was
not stopped, so it has no terminal boundary term.

On the shrinking annulus, rotational invariance inside the flat cap
gives one spherical transverse covariance eigenvalue
\(\kappa_\perp\) and one radial eigenvalue.  The frozen linear moment
bound at shifts \(\pm pe\) gives
\[
 |e^{\mathsf T}(m_p^{\rm sp}-R_{p,\alpha}h_p)|\le C/p.
\]
Since \(\kappa_\perp=(e^{\mathsf T}m_p^{\rm sp})/\|h_p\|\) and
\(\|h_p\|\asymp p^2M_p\),
\[
 |p^2\kappa_\perp-1|
       \le C(k_p^{-1/10}+k_p^{-1/2}).
\]
The scaled radial covariance is bounded by the spherical version of
\eqref{edge:hessian-approximation}.  There are \(O(k_p)\) transverse
dimensions and one radial dimension, so
\[
 \frac{\|p^2K_p^{\rm sp}-I\|_{\rm F}^2}{k_p}
       \le C(k_p^{-1/5}+k_p^{-1}).
\]
Combine this with the preceding cube--sphere covariance estimate,
divide by \(k_q\asymp k_p\), and decrease \(a_1>0\) if necessary.
This proves \eqref{edge:frobenius-bound}.
\end{proof}

\subsection{The covariance iteration and exact-scale variance}

The Frobenius estimate says that most directions have a small covariance
error on average over a scale interval.  To control the largest positive
error, compare a child cap to a larger parent cap.  The radial deficit
puts every positive error in the tangent space of the parent field;
the unused Haar rotation then dilutes it in the smaller child space.

Use the fixed small \(\gamma>0\) above.  Let \(q\) be a child scale and \(p\) a parent
scale with \(x=q/p\in[\gamma,2\gamma]\), and put \(\Delta=B_p-B_q\).
Conditional on \(W\) and observations through \(q\), \(h_p\) has law
\begin{equation}\label{edge:transition}
 \mathsf T_{q,p,h_q}(\mathrm dy)
 =\frac{e^{\phi_p(y)}N(h_q,\Delta)(\mathrm dy)}
        {\int e^{\phi_p(z)}N(h_q,\Delta)(\mathrm dz)} .
\end{equation}
The Gaussian is on the positive-increment space.  Its covariance
compressed to \(\cH_q\) is
\(\Delta_q=(1-x^2)p^2I_{\cH_q}\).  Differentiation of the convolution
gives the exact identity from \cite[Section~4.1, Equation~(4.3)]{FUA},
\begin{equation}\label{edge:transition-covariance}
 K_q=\Delta_q^{-1}
       [\Cov_{\mathsf T_{q,p,h_q}}(h_p)]_{\cH_q}\Delta_q^{-1}
       -\Delta_q^{-1}.
\end{equation}

\begin{lemma}[An extension on the parent annulus]
\label{edge:annular-extension}
Choose \(\eta>0\) sufficiently small compared with \(\gamma^2\), then
the coarse curvature tolerance \(\epsilon_{\rm aux}\) sufficiently small
compared with \(\eta\).  One can choose \(s_{\rm aux}\), the static tolerances,
and the threshold \(K\) so that the following holds on \(\mathcal G\).
For every child field satisfying \(|\alpha_q(h_q)|\le s_{\rm aux}q^2\), put
\[
 \mathcal A_p=\{h:|\alpha_p(h)|\le s_{\rm aux}p^2/2\},\qquad
 \mathcal A_p^+=\{h:|\alpha_p(h)|<3s_{\rm aux}p^2/4\}.
\]
There is a semiconcave extension \(\bar\phi_p\) with upper Hessian bound
\((1+\eta)p^{-2}I\), agreeing with \(\phi_p\) on \(\mathcal A_p^+\).
Moreover,
\(\dist(\mathcal A_p,(\mathcal A_p^+)^c)\ge c_{\rm buf}p^2M_p\)
for a fixed \(c_{\rm buf}>0\).  Under both \(\mathsf T_{q,p,h_q}\)
and the transition obtained by replacing \(\phi_p\) with \(\bar\phi_p\),
the probability of \(\mathcal A_p^c\) is at most \(Ce^{-c_1k_p}\).
The same bound holds with any fixed polynomial
moment of \(\|h_p\|/(p^2M_p)\) inserted, after decreasing \(c_1>0\).
All constants are uniform in the indicated fields and scales.
\end{lemma}

\begin{proof}
We verify the deterministic extension proof of \cite[Lemma~4.3]{FUA},
following \cite[Proposition~16.1]{CM}, with the new lower threshold.
On \(\{|\alpha_p|\le3s_{\rm aux}p^2/4\}\) take the infimum, over \(z\), of
\[
 \phi_p(z)+\nabla\phi_p(z)\cdot(y-z)
       +\frac{1+\eta}{2p^2}\|y-z\|^2.
\]
Choose \(s_{\rm aux}\ll\eta,\gamma^2\).  The length bounds in
Lemma~\ref{edge:norm-density} give a fixed
\(\beta_{\rm buf}>0\) such that a segment of length at most
\(\beta_{\rm buf}p^2M_p\) starting in this smaller annulus remains in
the full annulus.  For such pairs the supporting inequality follows
by integrating the coarse Hessian bound.  For
\(\|y-z\|\ge\beta_{\rm buf}p^2M_p\), freeze the spherical saddle at
\(z\).  The polar linear-moment bound gives the spherical quadratic
support with Hessian \((p^2+\alpha_p(z))^{-1}I\), its saddle gradient,
and a bounded norm-density error.  The static value and mean estimates
add \(2\epsilon_vk_p+\epsilon_gM_p\|y-z\|\).  Relative to
\(\|y-z\|^2/(2p^2)\), the total residual is bounded by
\[
 Cs_{\rm aux}+\frac{4\epsilon_v+C/k_p}{\beta_{\rm buf}^2}
       +\frac{2\epsilon_g}{\beta_{\rm buf}} .
\]
Choose \(s_{\rm aux}\), then \(\beta_{\rm buf}\), then \(\epsilon_v,\epsilon_g\)
so that this is less than \(\eta\), and finally increase \(K\).
Every support lies above \(\phi_p\) on the smaller annulus.  Its
infimum is finite and semiconcave with the stated upper Hessian, and
it agrees there by choosing \(z=y\).  The length derivative gives
\(\dist(\mathcal A_p,(\mathcal A_p^+)^c)\ge c_{\rm buf}p^2M_p\).

For the transition tails freeze the child saddle
\(\alpha=\alpha_q(h_q)\), and use as reference the Gaussian transition
tilted by the parent frozen quadratic \(q_{W_p,\alpha}\).
Writing \(R_u=R_{u,\alpha}\), Gaussian convolution gives
\begin{equation}\label{edge:reference-transition}
 R_ph_p\sim N(R_qh_q,R_q-R_p),\qquad R_q^{-1}=R_p^{-1}-\Delta.
\end{equation}
The child saddle equation shows that
\(\E\|h_p\|^2=\ell_p(\alpha)^2\), exactly.  The field covariance has
norm at most \(C\gamma^{-2}p^2\).  Gaussian concentration of its norm
therefore gives
\[
 \Pr_{\rm ref}\{|\|h_p\|-\ell_p(\alpha)|>vp^2M_p\}
       \le C e^{-c\gamma^2k_pv^2},
 \qquad v\ge C/(\gamma\sqrt{k_p}).
\]
The reference radius lies strictly inside \(\mathcal A_p\):
\(|\alpha|\le s_{\rm aux}q^2\le4s_{\rm aux}\gamma^2p^2\).
The length derivative gives a fixed positive gap in units \(p^2M_p\)
to the boundary of \(\mathcal A_p\).

Choose \(z_0\in\cH_p\) with \(\|z_0\|=\ell_p(\alpha)\) and calibrate the
reference quadratic by
\[
 q^\sharp(y)=q_{W_p,\alpha}(y)+\log f_{W_p,\alpha,z_0}(n).
\]
The active covariance is scalar, so the added density depends only on
\(\|z_0\|\); the additive constant leaves the reference transition
unchanged.  On bounded scaled radii, the noncentral upper norm-density
bound and the static value comparison give
\(\phi_p(y)-q^\sharp(y)\le\epsilon k_p+O(1)\), with any prescribed fixed
\(\epsilon>0\).
Outside a sufficiently large bounded radius, the zero-field projection
tail gives
\[
 \phi_p(y)-\phi_p(0)
 \le Ck_p\left(1+
       \left(\frac{\|y\|}{p^2M_p}\right)^{6/5}\right).
\]
This follows by maximizing the exponent \(uv-cv^6\) after integrating
the tail against a linear field of scaled length \(u\).  Its growth is
strictly subquadratic.  For the extension choose the support centered
at the point of radius \(\ell_p(\alpha)\) in the direction of \(y\).
At distance \(vp^2M_p\), its excess over the reference is at most
\[
 \epsilon k_p(1+v)+C(\eta+s_{\rm aux})k_pv^2+O(1).
\]
This uses the static mean error for the linear term and the difference
of the two quadratic coefficients for the last term.  Choose
\(\eta+s_{\rm aux}\ll\gamma^2\), and then \(\epsilon\) small compared with
the squared fixed gap to \(\mathcal A_p^c\).  Outside \(\mathcal A_p\) these errors
are smaller than the reference Gaussian deviation cost.

The normalizing integrals relative to the reference are at least
\(e^{-\epsilon'k_p}\) for any needed fixed \(\epsilon'>0\).
To see this, integrate over a thin fixed shell about the reference radius,
contained in \(\mathcal A_p\).  Its reference probability is bounded below
when \(K\) is large.  On that shell the frozen density formula
\eqref{edge:frozen-density} bounds the off-saddle logarithmic density
loss by \(C\tau^2k_p+O(1)\) when the shell width is \(\tau\) in scaled
units: the first parameter derivative vanishes at the saddle and the
second is one quarter of the norm-square variance \(O(n/p)\).
The static lower value error adds only \(\epsilon_vk_p\).
Take \(\tau\) and \(\epsilon_v\) sufficiently small.  Both potentials
agree on this shell.  Dividing the preceding upper tail integrals by
this lower bound proves the stated disagreement tails.  Polynomial
moments are absorbed by the same quadratic reference tail and the
subquadratic bound at large radii.
\end{proof}

\begin{lemma}[Positive covariance error]\label{edge:positive-curvature}
For some \(a,c,C>0\), uniformly at every deterministic scale
\(K\le k_p\), \(p\le r_n\), apart from activation points,
\begin{equation}\label{edge:positive-curvature-bound}
 \mathbf P\left(
 \mathcal G\cap
 \{p^2K_p\npreceq(1+Ck_p^{-a})I_{\cH_p}\}\right)
       \le C e^{-c k_p^a}.
\end{equation}
The probability is taken before restricting to \(\mathcal G\).
\end{lemma}

\begin{proof}
Fix a deterministic nonactivation scale \(p\) on the retained spectral
set.  For \(h\in\cH_p\), write \(D_p(h)=p^2K_p(h)-I_{\cH_p}\).
Define a Borel predicate \(\mathcal V_p(W_p,h)\) to be false unless
\(h\in\cH_p\), the saddle \(\alpha_p(h)\) exists, and \(h\ne0\).
Otherwise, with \(e=h/\|h\|\), it requires
\begin{equation}\label{edge:local-parent-validity}
 |\alpha_p(h)|\le s_{\rm aux}p^2,\qquad
 p^2K_p(h)\preceq(1+\epsilon_{\rm aux})I_{\cH_p},\qquad
 p^2\langle e,K_p(h)e\rangle\le0.9.
\end{equation}
This condition contains no restriction to \(\mathcal G\).  At a
nonactivation scale, \(P_p\) is the spectral projection of \(W_p\)
at its flat eigenvalue \(2-p^2\), and
\(R_{p,u}=(2I-W_p+uI)^{-1}\).  Thus the predicate, the saddle, and the
covariance used here are determined by \((p,W_p,h)\).
Their defining formulas are jointly Borel away from activation scales.
In the rest of the proof these objects are evaluated at \(h=h_p\),
and that argument is suppressed.

On \(\mathcal V_p\), use blocks relative to
\(\cH_p=\operatorname{span}\{e\}\oplus(\cH_p\cap e^\perp)\), put
\(z=(D_p)_{\perp e}\), and define on \(\cH_p\cap e^\perp\)
\begin{equation}\label{edge:tangent-error}
 E_p=\left((D_p)_{\perp\perp}
             +\frac{zz^{\mathsf T}}{|(D_p)_{ee}|}\right)_+,
 \qquad \widehat s_p=\|(D_p)_+\|.
\end{equation}
Extend \(E_p\) by zero on the radial direction and on \(\cH_p^\perp\).
Off \(\mathcal V_p\), set both \(E_p\) and \(\widehat s_p\) to zero,
before any restriction to \(\mathcal G\).  Covariance positivity and
\eqref{edge:local-parent-validity} give
\(-I\preceq D_p\preceq\epsilon_{\rm aux}I\) and
\((D_p)_{ee}\le-0.1\) on \(\mathcal V_p\).
Maximizing the quadratic form of \(D_p\) over its radial coordinate,
and then taking the positive part on \(e^\perp\), gives
\begin{equation}\label{edge:tangent-bounds}
 D_p\preceq E_p,\qquad
 \|E_p\|\le C\widehat s_p,\qquad
 \frac{\Tr E_p}{k_p}
       \le C\frac{\|D_p\|_{\rm F}}{\sqrt{k_p}}+\frac C{k_p}
\end{equation}
on \(\mathcal V_p\), with \(D_p\preceq E_p\) understood on \(\cH_p\).
For the norm bound, the maximizing radial coordinate
of a unit transverse vector is bounded because the radial deficit is
fixed and \(D_p\) is bounded; compare the resulting vector with
\(D_p\preceq\widehat s_pI\).  For the trace, the rank-one correction
has bounded trace, and
\(\Tr(((D_p)_{\perp\perp})_+)\le
\sqrt{\dim\cH_p}\|D_p\|_{\rm F}\); the displayed trace bound also
uses \(\dim\cH_p\le Ck_p\) from the retained spectral set.
These are the algebraic bounds of \cite[Lemma~4.4, Equation~(4.4)]{FUA}.
In particular \(\widehat s_p\le\epsilon_{\rm aux}\) and
\(\|E_p\|\le C\epsilon_{\rm aux}\) everywhere, including on disorder failures.
Include \(C\epsilon_{\rm aux}\ll\gamma^2\) in the initial auxiliary choice.

Apply Brascamp--Lieb's covariance inequality \cite{BrascampLieb1976}
to the transition with \(\bar\phi_p\).  After whitening, its Hessian
has a fixed positive margin because
\(\Delta\preceq(1-x^2)p^2I\) and \(\eta\ll\gamma^2\).
If \(\bar K\) is the extended Hessian, the covariance is at most the
transition expectation of
\[
 \Delta^{1/2}
 (I-\Delta^{1/2}\bar K\Delta^{1/2})^{-1}\Delta^{1/2}.
\]
Mollify the semiconcave extension to apply the smooth inequality and
pass to the limit using the positive margin.
Lemma~\ref{edge:annular-extension} replaces the extended Hessian and
transition by \(K_p\) and the true transition on agreement.  Its
polynomially weighted tails control first and second moments on the
complement.  In field units the covariance error is
\(O(p^4M_p^2e^{-c k_p})\); after
\eqref{edge:transition-covariance} and multiplication by \(q^2\),
the remaining polynomial factor in \(k_p\) is absorbed by decreasing
the exponent.

On \(\mathcal G\), the agreement annulus lies in \(\mathcal V_p\), so
\eqref{edge:tangent-bounds} gives
\(K_p\preceq p^{-2}(I+E_p)\).  Put
\(A=I-\Delta/p^2\) and
\(S=p^{-2}\Delta^{1/2}E_p\Delta^{1/2}\).
Then \(A\succeq x^2I\) and
\(\|A^{-1/2}SA^{-1/2}\|\ll1\), so
\[
 (A-S)^{-1}\preceq A^{-1}+2A^{-1}SA^{-1}.
\]
The baseline \(A^{-1}\), substituted in
\eqref{edge:transition-covariance}, is exactly \(q^{-2}I_{\cH_q}\).
For the correction, the two outer factors \(\Delta A^{-1}\) are the
scalar \(\Delta_q/x^2\) on \(\cH_q\).  Thus, after multiplication by
\(q^2\), the first insertion is at most \(2x^{-2}P_qE_pP_q\).
We obtain on \(\mathcal G\), for a child in its annulus,
\begin{equation}\label{edge:parent-inequality}
 \widehat s_q
 \le Cx^{-2}\mathbf E[\|P_qE_pP_q\|\mid W,\cF_q]
       +C_\gamma e^{-c_1k_p}.
\end{equation}
Only \(\Delta_q\) is inverted outside the covariance in
\eqref{edge:transition-covariance}; directions with arbitrarily small
new parent precision cause no loss.

We next compress \(E_p\) before imposing \(\mathcal G\).  The assertion
is trivial when \(E_p=0\); otherwise \(e\) is defined by
\(\mathcal V_p\).  The retained spectral set supplies
\(ck_p\le\dim\cH_p\le Ck_p\) and
\(\dim\cH_q\le Cx^3k_p\); after increasing the threshold,
\(\dim(\cH_p\cap e^\perp)\) is also comparable to \(k_p\).
Condition
under \(\mathbf P\) on the spectrum, \(W_p\) in physical coordinates,
the physical field \(h_p\), and its spectral coordinates.
Equation~\eqref{edge:path-likelihood} makes the remaining rotation in
\(e^\perp\cap\cH_p\) exactly Haar.  The operator \(E_p\) is fixed
under this conditioning: its definition uses only these parent data.
For a uniformly rotating unit vector \(v\) in this space, the Gaussian
representation of the uniform sphere gives
\[
 \mathbf P\left(
 v^{\mathsf T}E_pv>
       \frac{C(\Tr E_p+t\|E_p\|)}{k_p}
       \,\middle|\,\text{parent data}\right)
       \le C(e^{-t}+e^{-c_0k_p}).
\]
Indeed the numerator is a positive Gaussian quadratic form, and the
Gaussian squared norm in the denominator is at least a fixed fraction
of the dimension except with \(e^{-c_0k_p}\) probability.
A fixed-resolution net chosen in the spectral child unit sphere before
the remaining Haar rotation has at most
\(\exp(C\dim\cH_q)\le\exp(Cx^3k_p)\) points.  Project each rotated
net vector onto \(e^\perp\).  Conditional on the parent data, it has
a fixed norm at most one and a uniformly rotating direction.  With
\(t=C_2x^3k_p\), a sufficiently large \(C_2\) pays for the net, and
a sufficiently small \(\gamma\) pays for the second exponential term.
The net estimate for a positive quadratic form gives
\begin{equation}\label{edge:compression}
 \|P_qE_pP_q\|
 \le C\left(\frac{\Tr E_p}{k_p}+x^3\|E_p\|\right)
\end{equation}
outside a set of unconditional \(\mathbf P\)-probability
\(Ce^{-c_\gamma k_p}\).  This is the compression argument of
\cite[Lemma~4.5, Equation~(4.6)]{FUA}; no restriction to \(\mathcal G\)
was made in the Haar law.

Average \eqref{edge:parent-inequality} over
\(I_q=[q/(2\gamma),q/\gamma]\), with normalized logarithmic measure.
On \(\mathcal G\) outside the failure event in \eqref{edge:frobenius-bound},
\eqref{edge:tangent-bounds} and Cauchy--Schwarz give
\[
 \operatorname{Av}_{p\in I_q}\frac{\Tr E_p}{k_p}
       \le C_\gamma k_q^{-b},\qquad b=1/10.
\]
Here \(\operatorname{Av}\) denotes that normalized logarithmic integral.
On its failure the same averaged trace is bounded, because \(E_p\) is
bounded everywhere and \(\dim\cH_p\le Ck_p\).
Its contribution on that failure, with \(\ind_{\mathcal G}\) inserted,
has expectation at most
\(Ce^{-c k_q^{a_1}}\).  The failures of
\eqref{edge:compression} are charged in the same way by the bounded
\(\|E_p\|\).  These estimates are applied to the whole parent-scale
integral before its conditional expectation.  For every nonnegative
error cost \(Y\),
\begin{equation}\label{edge:conditioning-good-event}
 \mathbf E[\ind_{\mathcal G}\mathbf E(Y\mid W,\cF_q)]
       =\mathbf E[\ind_{\mathcal G}Y],
\end{equation}
because \(\mathcal G\) is measurable from \(W\).

The norm term in \eqref{edge:compression} has coefficient
\(Cx^{-2}x^3\le C\gamma\).  The initial choice of \(\gamma\) makes this coefficient
\(\theta<1/2\).  We have proved, on \(\mathcal G\),
\begin{equation}\label{edge:curvature-recurrence}
 \begin{split}
 \widehat s_q&\le
 \theta\operatorname{Av}_{p\in I_q}
       \mathbf E[\widehat s_p\mid W,\cF_q]
       +C_\gamma k_q^{-b}+\operatorname{err}_q,\\
 \operatorname{err}_q&\ge0,\qquad
 \mathbf E[\ind_{\mathcal G}\operatorname{err}_q]
       \le C e^{-c k_q^\beta}
 \end{split}
\end{equation}
for some \(\beta>0\).  Off the child annulus the left side is zero.
The bounded definition of \(E_p\) is essential here: no cost polynomial
in \(n\) is charged on a failed parent interval.

Use a barrier \(B k_q^{-\alpha}\), with \(0<\alpha\le b\) and \(B\)
large.  Since \(k_p=k_qx^{-3}\),
\[
 \theta\operatorname{Av}_{p\in I_q}B k_p^{-\alpha}
       \le\theta(2\gamma)^{3\alpha}B k_q^{-\alpha}.
\]
It absorbs the deterministic error in
\eqref{edge:curvature-recurrence}.  Taking positive parts, applying
conditional Jensen, and using \eqref{edge:conditioning-good-event},
we obtain
\[
 A(q):=\mathbf E[\ind_{\mathcal G}
       (\widehat s_q-Bk_q^{-\alpha})_+]
 \le\theta\operatorname{Av}_{p\in I_q}A(p)+Ce^{-c k_q^\beta}.
\]
For each deterministic nonactivation scale
\(r_n\le p\le C_\gamma r_n\), use
\cite[Proposition~4.2, Equation~(4.2)]{FUA} within its stated range.
Under the biased law \(\mathbf P\), it gives
\(p^2K_p\preceq(1+C(p^{a_0}+k_p^{-a_0}))I\) outside a set whose
intersection with the full event \(\mathcal G_n^{\rm FUA}\) has
probability at most \(Ce^{-n^{\xi_0}}\).
Since \(\mathcal G\subseteq\mathcal G_n^{\rm FUA}\),
the same failure bound holds after intersection with \(\mathcal G\).
Here \(pk_p=np^4\to0\), so \(p^{a_0}\le k_p^{-a_0}\), and
\(k_p\asymp n^{3\cdot10^{-4}}\).  By taking \(\alpha\le a_0\) and a
sufficiently small positive terminal exponential power, the same bound
for \(A(p)\) is \(Ce^{-c k_p^{\beta_0}}\).

Every parent generation multiplies \(k\) by at least
\(\Lambda=(2\gamma)^{-3}>1\).  Iteration, including the terminal
bound, gives for a suitably decreased \(\beta_*>0\)
\[
 A(q)\le C\sum_{j\ge0}\theta^j
       e^{-c(\Lambda^jk_q)^{\beta_*}}
       \le C e^{-c'k_q^{\beta_*}}.
\]
Markov's inequality at the barrier costs only a factor \(k_q^\alpha\),
absorbed by the exponential.  On \(\mathcal G\),
\eqref{edge:local-parent-validity} holds throughout the full auxiliary
annulus, and
\eqref{edge:shrinking-annulus} removes the annulus truncation.
Decrease to a common positive exponent in the barrier and probability.
This proves \eqref{edge:positive-curvature-bound}.
\end{proof}

\begin{proof}[Proof of Proposition~\ref{prop:edge-estimates}]
Lemmas~\ref{edge:spectral} and \ref{edge:norm-density} prove the
spectral and saddle assertions.  Lemmas~\ref{edge:spherical-tails},
\ref{edge:cube-caps}, and \ref{edge:path-annuli} prove the static,
curvature, projection, and ordinary annulus assertions.  Assign a
fixed fraction of the requested failure probability to each of their
spectral and disorder restrictions, and take their finite intersection.
The length bounds give a deterministic \(R_{\rm ann}\) containing the
full annulus; include a further cube-cap request for the static comparisons
on that ball, retaining the previously chosen width \(s_0\).
We select one additional event, independent of \(f\), for the variance.

In the reduction above take \(\eta_p=Ck_p^{-a}\) and let
\(\mathcal N\) be the jointly measurable field set where
\(p^2K_p(h)\npreceq(1+Ck_p^{-a})I_{\cH_p}\).  Set its indicator to one
at activation scales, which have zero scale measure for each
\(W\).  At every deterministic nonactivation scale its evaluation at
\(h_p\) is the failure event in \eqref{edge:positive-curvature-bound}.
Use \(Z_p\) and the common cost \(\rho_p\) from
\eqref{edge:common-failure-cost}; both are independent of \(f\).

Under \(\mathbf P\), conditional Jensen and the spherical identity give
\[
 \mathbf E[\mu_{p,h_p}(Z_p)^2]
       \le\mathbf E[\mu_{p,h_p}(Z_p^2)]
       =\mathbf E[Z_p^2]\le C,
\]
by the uncapped spherical projection tail.  Cauchy--Schwarz and
\eqref{edge:positive-curvature-bound}, absorbing \(k_p\) in the
exponential, imply
\begin{equation}\label{edge:weighted-failure}
 E_U[L_\circ\ind_{\mathcal G}\rho_p]\le C e^{-c k_p^a}.
\end{equation}
The left side is an ordinary path expectation averaged with the Haar
bias.  Take \(q_j=2^j(K/n)^{1/3}\) up to \(r_n\).
Fubini, \eqref{edge:weighted-failure}, and \(L_\circ\ge1/2\) bound
the expected ordinary cost on \([q_j,2q_j]\), intersected with
\(\mathcal G\), by \(Ce^{-c k_{q_j}^a}\).  Markov and
\[
 \sum_{j\ge0}e^{-c(K8^j)^a}\le C e^{-c'K^a}
\]
give an event of arbitrarily large limiting disorder probability on
which, simultaneously in \(j\),
\begin{equation}\label{edge:quenched-error-intervals}
 \int_{q_j}^{\min\{2q_j,r_n\}}\rho_p\,\mathrm d\log p
       \le C e^{-c k_{q_j}^a}.
\end{equation}
Increase \(K\) and use \eqref{edge:haar-bias} for the single discarded
event.  This event is independent of \(f\).

The conditional-variance estimate of
\cite[Proposition~4.1, Equation~(4.1)]{FUA} holds on an ordinary event
of arbitrarily large limiting probability, simultaneously for all real
cube functions and scales in its stated interval \([r_n,p_*]\).  Its
remainder is \(Ce^{-n^{\xi_0}}\|f\|_\infty^2\).  We use it only at
\(r_n\), for \(0\le f\le1\):
\[
 V_{r_n}(f)\le Cr_n^{-2}\cE(f)+Ce^{-n^{\xi_0}}.
\]
Intersect this event with the preceding ones.  In
\eqref{edge:variance-amplification} with terminal scale \(r_n\), its
amplified energy is
\(Cp^{-2}\cE(f)\).  Its amplified remainder is bounded by
\(Ce^{-c k_p^a}\) after decreasing \(a>0\), since
\(k_p\le k_{r_n}=n^{3\cdot10^{-4}}\).
For the other errors use \eqref{edge:quenched-error-intervals}.
The first partial dyadic interval has \(k_{q_j}\asymp k_p\); the rest
increase geometrically.  Hence
\[
 \sum_{q_j\gtrsim p}C(q_j/p)^2e^{-c k_{q_j}^a}
       \le C e^{-c'k_p^a}.
\]
Applying \eqref{edge:variance-amplification} proves
the first inequality of \eqref{edge:variance-estimates}.  It extends
to activation points by continuity of \(V_p\) in \(B_p\).

For the second inequality impose the additional fixed annulus estimate
with \(\zeta_*=\min\{\zeta,s_0/2\}\).  On that annulus the coarse covariance and
\eqref{edge:gradient-bounds} contribute \(CV_p\).  For its complement,
Cauchy--Schwarz and conditional Jensen give
\[
\begin{split}
 k_pQ_\mu[\ind_{\{|\alpha_p(h_p)|>\zeta_* p^2\}}\mu_{p,h_p}(Z_p)]
 &\le k_pQ_\mu\{|\alpha_p(h_p)|>\zeta_* p^2\}^{1/2}
       \mu(Z_p^2)^{1/2}\\
 &\le C e^{-c k_p}.
\end{split}
\]
Here \(\mu(Z_p^2)\le C\) holds at every scale on the chosen uncapped
cube projection event, and the ordinary annulus estimate is also
simultaneous in scale.  Weakening the exponent if needed proves the
second inequality of \eqref{edge:variance-estimates}, on the same
event for every \(f\).
\end{proof}

\section{Warming at the critical time scale}\label{sec:warming}

The first step toward a quench limit is to control the finite law after
any positive multiple of \(\Tn=n^{2/3}\), uniformly in its initial
distribution.  The edge estimates of Proposition~\ref{prop:edge-estimates} give the
required control at the smallest observation scales.  We turn them into
dissipation estimates by measuring how much relative entropy survives
between two observations.  Larger entropy levels are treated by the
observation path to a product measure from \cite[Section~19]{CM}.

For a probability \(\rho\) on the cube, write
\[
 w_s^\rho=\frac{\dd(\rho P_{s\Tn})}{\dd\mu},\qquad s\ge0.
\]
Thus \(w_s=w_s^{\nu_n}\) for the uniform initial law used elsewhere in
the paper.  Write
\[
 p_t(\sigma,\tau)=\frac{P_t(\sigma,\tau)}{\mu(\tau)}
\]
for the heat kernel relative to \(\mu\).  We use
\(\log_+z=\max\{\log z,0\}\) for \(z>0\), with
\(z(\log_+z)^m=0\) at \(z=0\).
For a probability \(\pi\) and a nonnegative integrable function \(z\), write
\[
 \Ent_\pi(z)=\pi[z\log z]-\pi[z]\log\pi[z],
\]
with \(0\log0=0\) and \(\Ent_\pi(0)=0\).  For a probability
\(\nu\ll\pi\), let
\(D(\nu\Vert\pi)=\Ent_\pi(\dd\nu/\dd\pi)\), and set this relative
entropy to \(+\infty\) otherwise.

\begin{theorem}[Warming from arbitrary initial laws]\label{thm:warming}
Suppose that the couplings are Gaussian.  Fix \(\vartheta>0\), an integer
\(m\ge1\), and \(0<\xi<1/100\).  There are measurable disorder events
\(\mathcal G_n\), a constant \(C<\infty\), and deterministic
\(\varepsilon_n\downarrow0\) such that
\(\liminf_{n\to\infty}\Prob(\mathcal G_n)\ge1-\vartheta\), and the
following inequalities hold on \(\mathcal G_n\), for all sufficiently
large \(n\).  The constant is independent of the initial law \(\rho\),
which may be chosen after the disorder.
\begin{align}
 \Ent_\mu(w_s^\rho)&\le C(1+s^{-1.54}),
       &&s>0,\quad\hbox{for every probability }\rho ,
       \label{warm:entropy}\\
 \|P_{s\Tn}\|_{L^1(\mu)\to L^\infty(\mu)}
       &=\max_{\sigma,\tau}p_{s\Tn}(\sigma,\tau)
       \le C\exp(C s^{-1.54}),
       &&s\ge\varepsilon_n ,
       \label{warm:kernel}\\
 \left(1+\mu\!\left[
       w_s^\rho(\log_+w_s^\rho)^m\right]\right)^{1/m}
       &\le C\left(1+s^{-1/(2/3-\xi)}\right),
       &&s>0,\quad\hbox{for every probability }\rho .
       \label{warm:log-moment}
\end{align}
The events and constants may depend on the displayed fixed parameters.
In particular the order \(m\) is fixed before \(n\) tends to infinity.
\end{theorem}

The distinction between \eqref{warm:entropy} and
\eqref{warm:kernel} is useful.  Entropy is controlled at every positive
time; the stronger kernel estimate begins at a deterministic time tending
to zero in the rescaled units.  The log-moment estimate supplies a
different form of control at the very short times used later in the
replacement argument.

\subsection{A restricted dissipation profile}\label{warm:profile-section}

Recall \(\cI(z)=\cE(z,\log z)\), with the extended value at zeros.  On a
two-point edge, the inequality
\[
 (a^2-b^2)(\log a^2-\log b^2)\ge4(a-b)^2,\qquad a,b>0,
\]
follows by integrating \(1/u\) between \(a\) and \(b\) and applying
Cauchy--Schwarz.  Summing it in the reversible edge representation gives
\begin{equation}
 \cI(f^2)\ge4\cE(f).                                      \label{warm:modified-classical}
\end{equation}
The inequality extends to nonnegative \(f\) by decreasing positive
regularization.

The next profile concerns a bounded function whose square tilt retains
a definite fraction of the logarithm of its inverse mass.  All its
inequalities hold simultaneously in the function.

\begin{proposition}[Restricted dissipation profiles]\label{warm:profiles}
Fix \(C_{\mathrm{lev}}<\infty\) and \(0<\xi<1/100\).  With tight constants, there
are a threshold \(K_0<\infty\) and a number \(\delta>0\) such that every
\(0\le f\le1\) satisfying
\begin{equation}
 N=\mu[f^2]=e^{-L},\qquad
 K_0\le L\le C_{\mathrm{lev}} n,\qquad
 D\!\left(\frac{f^2\mu}{N}\middle\Vert\mu\right)\ge\frac L2
                                                               \label{warm:restricted-class}
\end{equation}
obeys
\begin{align}
 \Tn\,\frac{\cI(f^2)}{N}&\ge c_\xi L^{5/3-\xi},
                                                               \label{warm:modified-profile}\\
 \Tn\,\frac{\cE(f)}{N}&\ge c
 \begin{cases}
 L^{1.662},&L\le n^\delta,\\
 L^{1.162},&L>n^\delta .
 \end{cases}                                                   \label{warm:classical-profile}
\end{align}
Here a threshold, a positive exponent, and constants are tight in the
following precise sense.  For each \(\vartheta>0\) they can be fixed,
independently of \(f\), on disorder events of limiting lower probability
at least \(1-\vartheta\).  They may depend on \(C_{\mathrm{lev}},\xi,\vartheta\).
\end{proposition}

The modified profile yields entropy dissipation and the higher logarithmic
moments after the reduction by level sets.  The classical profile yields a
support profile and hence the heat kernel bound.  We prove the profiles
in three ranges: \(L\le n^\delta\) by the argument with two observations;
\(n^\delta<L\le\lambda n\) by the proportional branch of the source
path; and \(\lambda n\le L\le C_{\mathrm{lev}}n\) by the product endpoint
and the estimate for a reweighted mean below, for a small fixed \(\lambda\)
chosen later.

\begin{lemma}[Entropy between two observation scales]\label{warm:small-profile}
Fix \(C_{\mathrm{lev}}<\infty\).  For every fixed \(\kappa>0\), there are tight \(K_0\) and
\(\delta>0\) such that every \(f\) satisfying
\eqref{warm:restricted-class} with \(L\le n^\delta\) satisfies
\begin{equation}
 \Tn\,\frac{\cE(f)}{N}\ge c_\kappa L^{5/3-\kappa}.
                                                               \label{warm:small-profile-bound}
\end{equation}
\end{lemma}

\begin{proof}
Fix first one instance of Proposition~\ref{prop:edge-estimates}, with
\(\epsilon_{\rm sp}\le1\), \(K\ge1\), and a variance-remainder exponent
\(0<a\le1\), decreasing the exponent if necessary.  Its constants are fixed on a
disorder event of arbitrarily high limiting probability.  Choose
\(M>1\) sufficiently large that
\[
 Ma>1,\qquad \frac{2M}{3}>\frac53-\kappa .
\]
We retain this variance estimate when imposing further cap tolerances
below; intersection with finitely many additional events does not change
its already chosen exponent \(a\).  Choose \(\delta>0\) so small that
\(M\delta/3<10^{-4}\).  This leaves every fixed factor buffer about the
scales below within \(r_n=n^{-1/3+10^{-4}}\) for large \(n\).
Use the same fixed small \(\gamma>0\) as in the covariance iteration
of Section~\ref{sec:edge-estimates}.  Choose the bridge tolerance
\(\epsilon>0\) after \(M\), with
\(2(M-1)\epsilon/3<\kappa\).

For a small fixed \(c_0>0\), to be chosen using only the projection-tail
constants, define
\[
 k_{p_1}=np_1^3=c_0L,\qquad k_R=nR^3=L^M.
\]
Increasing \(K_0\) puts these scales above the threshold of
Proposition~\ref{prop:edge-estimates}.  Suppose, for a contradiction, that
\(\Tn\cE(f)/N<L^{5/3-\kappa}\).  If
\(F_R(h)=\mu_{R,h}[f]\) and
\(V_R=Q_\mu[\Var(f\mid h_R)]\), that proposition gives
\begin{equation}
 \frac{V_R}{N}
 \le C L^{5/3-\kappa-2M/3}
       +C\exp\{L-cL^{Ma}\}=o_{K_0\to\infty}(1).
                                                               \label{warm:small-variance}
\end{equation}
The estimate is uniform in \(n,L,f\) in the present range, because
\(R^2\Tn=L^{2M/3}\).

Consider the joint ordinary law \(Q_\mu\) of the spin and its observation
path through \(R\).  Change this law to \(Q_{\rm spin}\) or \(Q_{\rm obs}\)
by the respective densities
\[
 \Lambda^{\rm spin}=\frac{f^2}{N},
 \qquad
 \Lambda^{\rm obs}=\frac{F_R(h_R)^2}{A},
 \qquad
 A=Q_\mu[F_R^2]=N-V_R .
\]
For \(K_0\) large, \(A\ge N/2\), so both densities are at most \(2e^L\).
For any of these laws, \(Q|_p\) denotes restriction to
\(\sigma(h_q:0\le q\le p)\), always excluding
the spin.  Define, for \(p<R\),
\begin{equation}\label{warm:residual-observation-entropy}
 H(p)=D(Q_{\rm obs}\Vert Q_\mu)
       -D(Q_{\rm obs}|_p\Vert Q_\mu|_p).
\end{equation}
The first entropy is on the full joint law.  Since its likelihood is
measurable from \(h_R\), it equals the relative entropy on the terminal
observation sigma-field.  Both terms are finite by the likelihood bound,
and the conditional entropy chain rule identifies \(H(p)\) with the
expected entropy of the terminal observation tilt remaining after the
prefix through \(p\).
To compare the two tilts, note that \(Q_\mu[fF_R]=A\), and hence
\[
 Q_\mu\left[\left(\frac f{\sqrt N}
                  -\frac{F_R}{\sqrt A}\right)^2\right]
 =2\left(1-\sqrt{\frac AN}\right)\le\frac{2V_R}{N}.
\]
Cauchy--Schwarz bounds the \(L^1(Q_\mu)\) distance between the two
densities by \(2\sqrt{2V_R/N}\).  This bound persists after projection
to any part of the observation path.  If two projected densities
\(u,v\le2e^L\) have \(L^1\) distance \(d\), then
\begin{equation}
 \left|Q_\mu[u\log u]-Q_\mu[v\log v]\right|
 \le (1+L+\log2)d+C.                                      \label{warm:entropy-continuity}
\end{equation}
Indeed \(x\log x\) has derivative bounded by \(1+L+\log2\) on
\([1,2e^L]\), while its absolute value is bounded on \([0,1]\).
It follows from \eqref{warm:small-variance} that the two relative
entropies, both before and after any projection, differ by \(o(L)\),
uniformly as \(K_0\) increases.

We bound the information revealed at \(p_1\) under the spin tilt.  The
field \(h_{p_1}\) is a sufficient statistic for the path through that
scale.  Let \(Q_0\) be the output of the same Gaussian channel with
zero signal.  Convexity of relative entropy and evenness of \(\mu\)
give
\[
 D(Q_{\rm spin}|_{p_1}\Vert Q_0)
       \le\frac12\frac{\mu[f^2X^{\mathsf T}B_{p_1}X]}N,
 \qquad
 \log\frac{\dd(Q_\mu|_{p_1})}{\dd Q_0}(h)
       \ge-\frac12\mu[X^{\mathsf T}B_{p_1}X].
\]
These formulas are on the range of \(B_{p_1}\), so zero precision
directions cause no difficulty.  Subtraction bounds
\(D(Q_{\rm spin}|_{p_1}\Vert Q_\mu|_{p_1})\) by half the sum of the two
quadratic moments.  The spectral bound \(d_1\ge-\epsilon_{\rm sp}p_1^2\)
gives \(B_{p_1}\preceq(1+\epsilon_{\rm sp})p_1^2P_{p_1}
\preceq2p_1^2P_{p_1}\).  The projection
tail in Proposition~\ref{prop:edge-estimates}, integrated against a density at most
\(e^L\), gives
\[
 \frac{\mu[f^2\|P_{p_1}X\|^2]}{NM_{p_1}^2}
       +\frac{\mu[\|P_{p_1}X\|^2]}{M_{p_1}^2}
 \le C\left(1+(L/k_{p_1})^{1/3}\right).
\]
For example, integrate
\(\min\{1,e^LCe^{-c k_{p_1}v^6}\}\) in the variable \(v^2\).
The information is therefore at most
\[
 Ck_{p_1}\left(1+(L/k_{p_1})^{1/3}\right)
   \le C(c_0+c_0^{2/3})L .
\]
This is the proportional-cap argument in
\cite[Equation~(19.12)]{CM}.  Choose \(c_0\) so that this last bound is
less than \(L/8\).  The full spin-tilt entropy is at least \(L/2\).
The two continuity comparisons in \eqref{warm:entropy-continuity},
combined into one uniform \(o(L)\), therefore give
\begin{equation}\label{warm:residual-lower}
 \begin{split}
 H(p_1)&\ge D(Q_{\rm spin}\Vert Q_\mu)
       -D(Q_{\rm spin}|_{p_1}\Vert Q_\mu|_{p_1})-o(L)\\
       &\ge L/2-L/8-o(L)\ge cL.
 \end{split}
\end{equation}

We follow that remaining entropy up to a fixed fraction of \(R\).
For \(p<R\), let \(\mathsf T_{p,R}^h\) be the ordinary conditional
law of \(h_R\) given \(h_p=h\).  It has density proportional to
\[
 e^{\phi_R(y)}\,\mathcal N(h,B_R-B_p)(\dd y),
\]
where the Gaussian is restricted to the range of its covariance.
Let \(\widetilde{\mathsf T}_{p,R}^h\) be its tilt by \(F_R(y)^2\), and
put
\[
 A_p(h)=\mathsf T_{p,R}^h[F_R^2],\qquad
 l_p(h)=D(\widetilde{\mathsf T}_{p,R}^h\Vert\mathsf T_{p,R}^h).
\]
Observation sufficiency gives the prefix density
\(\dd(Q_{\rm obs}|_p)/\dd(Q_\mu|_p)=A_p(h_p)/A\).  The conditional
chain rule in \eqref{warm:residual-observation-entropy} now reads
\begin{equation}\label{warm:residual-transition-entropy}
 H(p)=\frac1A Q_\mu\!\left[
       \Ent_{\mathsf T_{p,R}^{h_p}}(F_R^2)\right]
       =Q_{\rm obs}[l_p(h_p)].
\end{equation}

Write \(\Delta=B_R-B_p\).  On the currently active subspace
\(\cH_p\), \(\Delta=(R^2-p^2)I\).  Differentiating its Gaussian
convolution therefore gives
\[
 \nabla_{\cH_p}\log A_p(h)=\delta m_h,\qquad
 \delta m_h=\widetilde{\mathsf T}_{p,R}^h[S]
                    -\mathsf T_{p,R}^h[S],\qquad
 S=\frac{P_p(h_R-h)}{R^2-p^2}.
\]
In the parameter \(p^2\), the likelihood martingale \(A_p(h_p)/A\)
has innovation coefficient \((A_p/A)\delta m_h\).  The Gaussian
likelihood formula for its relative entropy and the chain rule yield
\begin{equation}
 -\frac{\dd H}{\dd\log p}
      =p^2Q_{\rm obs}[|\delta m_h|^2] .                   \label{warm:bridge-loss}
\end{equation}
This identity holds almost everywhere and in integrated form across
activation points.  One direct justification applies It\^o's formula
to \((u+\epsilon)\log(u+\epsilon)\) of the bounded likelihood,
then lets \(\epsilon\downarrow0\); zero-rate directions are omitted.
This is the observation entropy identity
\cite[Lemma~19.2, Equation~(19.6)]{CM}, applied here to a terminal
observation tilt instead of a spin tilt.

The ordinary log moment generating function of \(S\), for
\(v\in\cH_p\), is exactly
\[
 \log\mathsf T_{p,R}^h[e^{\ip vS}]
   =\phi_p(h+v)-\phi_p(h)+\frac{|v|^2}{2(R^2-p^2)} .
\]
Suppose that the Hessian of \(\phi_p\) in this subspace is at most
\(C_{\rm curv}p^{-2}I\) on all shifts of length at most
\(p\sqrt{2l_p(h)/C_{\rm curv}}\) from \(h\).  The entropy variational
inequality, tested against \(\ip vS\), then gives
\begin{equation}
 p^2|\delta m_h|^2
 \le2\left(C_{\rm curv}+\frac{p^2}{R^2-p^2}\right)l_p(h).
                                                               \label{warm:bridge-mean}
\end{equation}
To check the finite shift radius, the centered log moment generating
function is at most
\(\frac12(C_{\rm curv}p^{-2}+(R^2-p^2)^{-1})|v|^2\).
Testing at length \(\sqrt{2l_p/(C_{\rm curv}p^{-2}+(R^2-p^2)^{-1})}\)
in the mean-shift direction proves the inequality by contradiction;
this length is no larger than the assumed radius.

We next account for observations where that curvature bound is not
available.  The entropy bound \(l_p\le-\log A_p\), together with
\(A\ge e^{-L}/2\), implies
\begin{equation}
 Q_{\rm obs}(l_p>L+\log2+z)\le e^{-z},\qquad z\ge0.
                                                               \label{warm:bridge-entropy-tail}
\end{equation}
Under the ordinary joint law,
\[
 S=P_pX+(R^2-p^2)^{-1/2}G_p,
\]
with an independent standard Gaussian \(G_p\) on \(\cH_p\).
For \(p\le\gamma R\), the projection tails of
Proposition~\ref{prop:edge-estimates} and Gaussian concentration give,
for each prescribed \(a_0>0\), a fixed \(V\) such that
\begin{equation}\label{warm:bridge-input-tails}
 \begin{split}
 Q_\mu\!\left[\frac{|S|^2}{M_p^2}
       \ind_{\{|S|>VM_p\}}\right]&\le Ce^{-a_0k_p},\\
 Q_\mu\{\|h_p\|>Vp^2M_p\}&\le Ce^{-a_0k_p}.
 \end{split}
\end{equation}
To see the common scale, write the observation noise using a standard
Gaussian \(G'_p\) on \(\cH_p\).  Since \(B_p\preceq2p^2P_p\),
\[
 \frac{\|h_p\|}{p^2M_p}
 \le2\frac{\|P_pX\|}{M_p}+\sqrt2\frac{\|G'_p\|}{\sqrt{k_p}},
 \qquad
 \frac{|S|}{M_p}
 \le\frac{\|P_pX\|}{M_p}
       +\frac p{\sqrt{R^2-p^2}}\frac{\|G_p\|}{\sqrt{k_p}}.
\]
The last coefficient is bounded on this range and
\(\dim\cH_p\le Ck_p\).  Integration of the same tails allows any
fixed polynomial moment as well.  For a prefix event \(E\), conditional Jensen and
the full and prefix likelihood bounds at most \(2e^L\) give
\begin{equation}
\begin{split}
 p^2Q_{\rm obs}[|\delta m_h|^2\ind_E]
 \le Ck_p\bigg(
 V^2Q_{\rm obs}(E)
 +e^L Q_\mu\!\left[
       \frac{|S|^2}{M_p^2}\ind_{\{|S|>VM_p\}}\right]\bigg).
\end{split}                                                     \label{warm:discarded-bridge}
\end{equation}
Indeed Jensen bounds \(|\delta m_h|^2\) by twice the sum of the
second moments under the two conditional laws.  For their parts
\(|S|\le VM_p\) use the prefix probability of \(E\); for the
remaining parts use respectively the full likelihood \(F_R^2/A\)
and the prefix likelihood \(A_p/A\).  Finally \(p^2M_p^2=k_p\).

Choose \(p_2\) by \(k_{p_2}=C_2k_{p_1}\), where \(C_2\) will be
fixed below.  Once it is fixed, \(p_2<\gamma R\) for \(K_0\)
sufficiently large.  Fix \(A_0>0\), choose a tail rate
\(a_0>c_0^{-1}\) with a fixed margin, and then choose \(V\) in
\eqref{warm:bridge-input-tails}.  On \([p_1,p_2]\) retain
\(\|h_p\|\le Vp^2M_p\) and \(l_p\le L+\log2+A_0k_p\).
The trial shifts in \eqref{warm:bridge-mean}, in units \(p^2M_p\),
then have length at most
\(\sqrt{2(c_0^{-1}+A_0+1)}\), after enlarging the ball-curvature
constant to be at least one.  Choose a fixed ball containing the retained
fields and these shifts before requesting its curvature constant from
Proposition~\ref{prop:edge-estimates}.  Equations
\eqref{warm:bridge-entropy-tail}--\eqref{warm:discarded-bridge}
show that discarded observations contribute at most \(Ck_pe^{-c k_p}\).
Thus \eqref{warm:bridge-loss}--\eqref{warm:bridge-mean} imply
\[
 \frac{\dd H}{\dd\log p}\ge-C_1H-Ck_pe^{-c k_p}
       \qquad(p_1\le p\le p_2).
\]
The logarithmic length of this interval is fixed.

For \([p_2,\gamma R]\), make a fresh auxiliary cap and extension
request with the same \(\gamma\).  Choose \(\eta\ll\gamma^2\), then
\(\epsilon_{\rm aux}\le\epsilon\) sufficiently small compared with
\(\eta\).  Reapply the construction of Lemma~\ref{edge:annular-extension}
with these choices, and write \(s_0\) for its auxiliary width.
Let \(\mathcal G_{\rm ext}=\mathcal G_n^{\rm aux,ext}
\cap\mathcal G_n^{\rm FUA}\) be the resulting event for the required
static tolerances, radii, threshold, and scale buffers.
Lemma~\ref{edge:cube-caps} gives \(\mathcal G_n^{\rm aux,ext}\) with
arbitrarily high limiting ordinary probability, and
\(\mathcal G_n^{\rm FUA}\) costs \(o(1)\).  Assign a further fraction
of the failure probability and intersect with the retained variance
event; its exponent \(a\) and constants are unchanged.

On the full annulus \(|\alpha_p(h)|\le s_0p^2\), this request gives
\(K_p(h)\preceq(1+\epsilon_{\rm aux})p^{-2}I
\preceq(1+\epsilon)p^{-2}I\).  The same annulus specifies admissible
child fields for the extension.  Choose \(0<\zeta<s_0\) and a fixed
\(b_{\rm rad}>0\) so that shifts of length at most
\(b_{\rm rad}p^2M_p\) from \(|\alpha_p(h)|\le\zeta p^2\) stay in the
full annulus.  After \(\zeta\) is fixed, impose its ordinary annulus
estimate from Lemma~\ref{edge:path-annuli} as one further fixed request
and intersection.  Thus the ordinary failure probability is
\(Ce^{-a_\zeta k_p}\), with a fixed \(a_\zeta>0\); only this probability
bound is used from the further request.  The trial shifts fit whenever
\(l_p\le b k_p\), where \(b=(1+\epsilon)b_{\rm rad}^2/2\).
Choose \(C_2\) so large that
\[
 \frac{L}{k_p}\le\frac1{C_2c_0}
       <\frac14\min\{b,a_\zeta\}\qquad(p\ge p_2).
\]
The likelihood bound and \eqref{warm:bridge-entropy-tail} give
exponentially small tilted probabilities for both discarded events.
Using \eqref{warm:discarded-bridge} with a further fixed \(S\)-tail
threshold proves
\[
 \frac{\dd H}{\dd\log p}
 \ge-\left(2(1+\epsilon)+\frac{2p^2}{R^2-p^2}\right)H
        -Ck_pe^{-c k_p}.
\]
The extra coefficient has bounded integral,
\[
 \int_{p_2}^{\gamma R}\frac{2p^2}{R^2-p^2}\,\dd\log p
       \le-\log(1-\gamma^2).
\]
In both intervals the integrating factor charges a fixed power of
\(p/p_1\), whereas \(k_p=np^3\) grows as its cube.  Consequently the
weighted errors are bounded by a polynomial in \(k_{p_1}\) times
\((e^{-c k_{p_1}}+e^{-c k_{p_2}})\).  They are \(o(L)\) after
increasing \(K_0\).  The first interval costs a fixed multiplicative
factor.  We obtain
\begin{equation}
 H(\gamma R)\ge cL(p_1/R)^{2+2\epsilon}.                  \label{warm:retained-small}
\end{equation}

It remains to bound this entropy by the energy of \(f\).  For a child
field satisfying \(|\alpha_{\gamma R}(h)|\le s_0(\gamma R)^2\), the
fresh request to Lemma~\ref{edge:annular-extension} supplies an extension
of the parent potential with upper Hessian \((1+\eta)R^{-2}I\).  Put
\[
 \mathcal A_R=\{|\alpha_R|\le s_0R^2/2\},\qquad
 \mathcal A_R^+=\{|\alpha_R|<3s_0R^2/4\}.
\]
The extension agrees with the true potential on \(\mathcal A_R^+\).
The distance from \(\mathcal A_R\) to \((\mathcal A_R^+)^c\) is at least
\(c_{\rm buf}R^2M_R\).  Both transition laws assign
\(\mathcal A_R^c\) probability \(Ce^{-c k_R}\), also with fixed
polynomial moments of \(\|h_R\|/(R^2M_R)\).
The precision increment satisfies
\(\Delta\preceq(1-\gamma^2)R^2I\).  On its range the extended
transition, denoted by \(\overline{\mathsf T}\), has negative log-density Hessian at least
\[
 \Delta^{-1}-(1+\eta)R^{-2}I\succeq c_\gamma R^{-2}I .
\]
The Bakry--\'Emery logarithmic Sobolev criterion
\cite[Corollary~5.7.2]{BakryGentilLedoux2014} gives
\[
\Ent_{\overline{\mathsf T}}(g^2)
\le C_\gamma R^2\overline{\mathsf T}[|\nabla g|^2].
\]
It applies to the semiconcave extension by mollification.

Choose a smooth cutoff \(\chi\) equal to one on \(\mathcal A_R\)
and supported in \(\mathcal A_R^+\).  The fixed buffer and the
radial saddle derivative in Proposition~\ref{prop:edge-estimates} give
\(|\nabla\chi|\le C/(R^2M_R)\).  Apply the preceding inequality to
\(\chi F_R\).  Its cutoff-gradient term is at most \(Ce^{-c k_R}\),
because \(R^2|\nabla\chi|^2\le C/k_R\) and the transition tail has
that exponential rate.  On agreement the ratio of the two transition
densities is constant; their normalization and their exponentially small
tails make that constant \(1+O(e^{-c k_R})\).  To compare the entropies,
use
\(\Ent_{\mathsf T}(g)=\mathsf T[g\log g]-\mathsf T[g]\log\mathsf T[g]\):
\(x\log x\) is bounded and H\"older continuous on \([0,1]\), and
\(0\le F_R,\chi\le1\).  The comparison, with a possibly smaller
positive exponent, proves
\begin{equation}
 \Ent_{\mathsf T_{\,\gamma R,R}^{h}}(F_R^2)
 \le C_\gamma R^2\mathsf T_{\,\gamma R,R}^{h}
                   [|\nabla_{\cH_R}F_R|^2]+Ce^{-c k_R}.
                                                               \label{warm:transition-lsi}
\end{equation}
Failure of the child condition is contained in the ordinary
\(\zeta\)-annulus failure, since \(\zeta<s_0\).  Its probability is
\(Ce^{-c k_R}\) after changing \(c\), since
\(k_{\gamma R}=\gamma^3k_R\).
On it the conditional entropy of \(F_R^2\le1\) is at most \(1/e\).
Average \eqref{warm:transition-lsi} under the ordinary child law and
use \eqref{warm:residual-transition-entropy}.  The tower property,
\(A\ge N/2\), and the two inequalities of
\eqref{edge:variance-estimates} at \(R\) give
\begin{equation}
 \begin{split}
 H(\gamma R)
 &\le\frac{C_\gamma R^2}{A}
       Q_\mu\!\left[\|\nabla_{\cH_R}F_R(h_R)\|^2\right]
       +\frac{Ce^{-c k_R}}A\\
 &\le C\frac{V_R}{N}+C\exp\{L-c k_R^a\}\\
 &\le C R^{-2}\frac{\cE(f)}N+C\exp\{L-c k_R^a\}.
 \end{split}                                                \label{warm:small-upper}
\end{equation}
Here \(0<a\le1\) allows the transition error to be included in the
displayed remainder, and \(k_R=L^M\) makes it \(C e^{L-cL^{Ma}}\).
The error is negligible compared with
\eqref{warm:retained-small}, even if the latter is a negative power
of \(L\).  Comparing the two bounds and using
\(\Tn p_1^2=(c_0L)^{2/3}\) gives
\[
 \Tn\,\frac{\cE(f)}N
 \ge c\,\Tn R^2L(p_1/R)^{2+2\epsilon}
 =c' L^{5/3-2(M-1)\epsilon/3}.
\]
The choice \(2(M-1)\epsilon/3<\kappa\) contradicts the assumed upper
bound for \(K_0\) large.  All ball radii, tail rates, annulus
tolerances and factor buffers have been fixed before this last
increase of \(K_0\) and before \(n\) tends to infinity.
\end{proof}

\subsection{The product endpoint and larger entropy levels}
\label{warm:large-section}

We record first two entropy bounds at a product posterior.  Let
\(\nu=f^2\mu/N\).  For each path used below, let \(Q_\mu,Q_\nu\)
denote the corresponding joint laws of the initial spin and its
observations, started from these two spin laws.  Expectations of
observation-only functions use their observation marginals, and
\(Q|_u\) again denotes the observation prefix, excluding the spin.
At observation time \(u\), write \(\mu_u,\nu_u\) for their spin
posteriors and \(N_u=\mu_u[f^2]\).  Use the posterior heat-bath form and
the weighted Euclidean norm
\[
 \cE_{\mu_u}(f)=\sum_i\mu_u[\Var_{\mu_u}(f\mid X_{-i})],
 \qquad \|v\|_A^2=v^{\mathsf T}Av.
\]
Bayes' formula and the observation
identities of \cite[Lemma~19.2, Equations~(19.6)--(19.8)]{CM} are
\begin{align}
 \frac{\dd(Q_\nu|_u)}{\dd(Q_\mu|_u)}
       &=\frac{N_u}{N},\qquad \nu_u=\frac{f^2\mu_u}{N_u},
                                                               \label{warm:posterior-likelihood}\\
 -\dot H(u)
       &=\frac12 Q_\nu\!\left[
             \|\nu_u[X]-\mu_u[X]\|_{\dot B_u}^{\,2}\right],
 \qquad H(u)=Q_\nu[D(\nu_u\Vert\mu_u)],
                                                               \label{warm:spin-loss}\\
 Q_\nu\!\left[\frac{\cE_{\mu_u}(f)}{N_u}\right]
       &=\frac1N Q_\mu[\cE_{\mu_u}(f)]
       \le\frac{\cE(f)}N .                                \label{warm:conditioned-energy}
\end{align}
The posterior \(\nu_u\) is defined where \(N_u>0\); the remaining
observation paths have \(Q_\nu\)-mass zero and contribute zero to
these averaged formulas.
The loss identity has the corresponding stopped-path form, with the
information in precision jumps added.  The last inequality is
conditional-variance contraction, site by site, under the ordinary
joint spin-observation law.  The form at every posterior uses the
same rate-one-per-site normalization as \(\cE\).

\begin{lemma}[Product-posterior entropy bounds]\label{warm:product}
Fix the at most two product endpoints used below.  Each has precision
\(t'I+W\) for a fixed \(t'>3\), chosen from the fixed path parameters
before \(f\).  On the common disorder event of
\cite[Equation~(19.10)]{CM} for these endpoints, including its
polynomial-moment field-tail control, the tests
\eqref{warm:restricted-class} satisfy
\begin{equation}
 H_{\rm end}\le
 C\sqrt{\log n+L}\,\frac{\cE(f)}N+O(n^{-3}),
 \qquad
 H_{\rm end}\le\frac{\cI(f^2)}N .                         \label{warm:product-bounds}
\end{equation}
The constants in the first inequality, including its remainder, may
depend on these endpoints, the fixed path parameters, and
\(C_{\mathrm{lev}}\), but not on \(n\) or \(f\).
The second inequality holds for every observation channel whose
posterior at the endpoint is a product measure; it does not require
a field-tail estimate.
\end{lemma}

\begin{proof}
The first inequality is the product calculation of
\cite[Equation~(19.18)]{CM}, applied at each chosen endpoint.
Its precision is \(t'I+W\) for the corresponding fixed \(t'>3\).
The posterior interaction is then a scalar diagonal, hence its spin
law is a product with fields \(h_i\).  Tensorizing the two-point
log-Sobolev bound gives
\[
 D(\nu_h\Vert\mu_h)
 \le C(1+\|h\|_\infty)\frac{\cE_{\mu_h}(f)}{N_h}.
\]
The field tail of \cite[Equation~(19.10)]{CM}, on the common event for
the chosen endpoints, restricts
\(\|h\|_\infty\le C\sqrt{\log n+L}\).  Its ordinary complement has
probability at most \(e^{-2L}n^{-10}\), with polynomial moments;
after the likelihood bound \(N_h/N\le e^L\) its entropy contribution
is \(O(n^{-3})\).  Equation \eqref{warm:conditioned-energy} gives the
displayed bound.

For the second inequality, entropy tensorization at a product posterior
and Jensen's inequality for the logarithm give, for positive \(z\),
\[
 \Ent_{\mu_h}(z)
 \le\sum_i\mu_h[\Ent_{\mu_h(\cdot\mid X_{-i})}(z)]
 \le\sum_i\mu_h[\Cov_{\mu_h(\cdot\mid X_{-i})}(z,\log z)]
 =\cI_{\mu_h}(z).
\]
In fact \(\Ent_\pi(z)\le\Cov_\pi(z,\log z)\) follows from
\(\pi[\log z]\le\log\pi[z]\).  The posterior entropy formula is
\[
 H_{\rm end}=\frac1N Q_\mu[\Ent_{\mu_h}(f^2)].
\]
It remains to compare the averaged modified form to its initial value.
For a fixed site and fixed \(x_{-i}\), let
\[
 a_h=\mu_h(X_{-i}=x_{-i},X_i=+1),\qquad
 b_h=\mu_h(X_{-i}=x_{-i},X_i=-1)
\]
be the two joint posterior masses.  Their contribution to
\(\cI_{\mu_h}(z)\) is
\[
 \frac{a_hb_h}{a_h+b_h}
       (z_+-z_-)(\log z_+-\log z_-).
\]
The second factor is nonnegative and independent of the observation.
The function \(c(a,b)=ab/(a+b)\), extended by zero on the coordinate
axes, is concave: in the positive quadrant its Hessian is
\(-2(b,-a)(b,-a)^{\mathsf T}/(a+b)^3\).
Since \(Q_\mu[a_h]\) and \(Q_\mu[b_h]\) are the corresponding
initial masses, Jensen's inequality proves
\(Q_\mu[\cI_{\mu_h}(z)]\le\cI_\mu(z)\).
Take \(z=f^2+\epsilon\) and let \(\epsilon\downarrow0\) to include
zeros.  This proves the second bound.
\end{proof}

At entropy of order \(n\), the observation path must be stopped before
its field leaves a region where the posterior mean is controlled.
The following strengthening of the square-density estimate in
\cite[Proposition~7.3]{SG} is what permits the necessary energy range.
We display the stability condition used at subcritical posterior
interactions; the finite word estimates and vector recipes of \cite{SG}
are used in their stated form.

Let \(J_*=\beta W_{\rm off}\), \(0<\beta<1\), and \(j=\beta^2\).
For \(y,h\in\R^n\) put
\[
\begin{gathered}
 m(y)=\tanh y,\qquad q(y)=\frac{\|m(y)\|^2}{n},\qquad
 V(y)=\diag(1-m_i(y)^2),\\
 F_h(y)=y+\bigl(j(1-q(y))I-J_*\bigr)m(y)-h .
\end{gathered}
\]
The stability condition of \cite[Equation~(17.2)]{CM}, with fixed
positive \(b,\eta,\epsilon_A,\rho_0\), is the implication
\begin{equation}
\begin{gathered}
 n^{-1/2}\|F_h(y)\|\le2\rho_0,\qquad 0\preceq A\preceq(1+\eta)I,\\
 A\ \hbox{diagonal},\qquad
 n^{-1/2}\|A-V(y)\|_{\rm F}\le\epsilon_A\\
\Longrightarrow\quad
 I+A^{1/2}\left(j(1-q(y))I-J_*
       -\frac{2j}{n}m(y)m(y)^{\mathsf T}\right)A^{1/2}\succeq bI .
\end{gathered}
                                                               \label{warm:stability}
\end{equation}
for every \(y,A\).  The ordinary word diagnostics below are the
simultaneous estimates of \cite[Lemma~3.1]{SG}: for a prescribed
finite length they bound the operator norm, the off-diagonal entrywise
\(\ell^4\) norm, and the Euclidean norm of the diagonal error from
the noncrossing prediction for products of \(J_*\) and bounded
diagonal matrices.  The ordinary vector recipes are those of
\cite[Definition~6.2]{SG}; their square-density argument uses no
inverse word.

\begin{lemma}[Square-root energy in a reweighted posterior mean]
\label{warm:posterior-mean}
Fix an operator-norm bound for \(J_*\) and the parameters in
\eqref{warm:stability}, with \(j(1+\eta)^2<1\).
For every sufficiently small fixed \(\epsilon>0\), impose the
ordinary word diagnostics of a finite length depending on \(\epsilon\).
For any field \(h\) satisfying \eqref{warm:stability}, taking
\(A=V(y)\) verifies the hypothesis of \cite[Lemma~5.7]{SG} with
residual threshold \(2\rho_0\sqrt n\) and lower bound \(b\).
Let \(y_*\) be its unique root of \(F_h\), and set
\(m_*=\tanh y_*\).  This root existence is the source conclusion;
the square-root energy dependence below is the new estimate.  Define
\[
 \mu_h(\sigma)\ \propto\
   \exp\{\tfrac12\sigma^{\mathsf T}J_*\sigma+h^{\mathsf T}\sigma\},
 \qquad
 \mathcal D_h(f):=\cE_{\mu_h}(f)
       =\sum_i\mu_h[\Var(f\mid\sigma_{-i})].
\]
For every real \(f\) with \(N_h=\mu_h[f^2]>0\) and
\(\nu_h=f^2\mu_h/N_h\),
\begin{equation}
 \frac{\|\nu_h[\sigma]-m_*\|}{\sqrt n}
 \le C_{\rm lead}\epsilon+
       \frac{C_\epsilon}{\sqrt n}
          \left(1+\sqrt{\frac{\mathcal D_h(f)}{N_h}}\right).
                                                               \label{warm:sqrt-posterior}
\end{equation}
The constants are uniform in \(n,h,f\) under these hypotheses.
\(C_{\rm lead}\) depends only on the fixed stability and norm bounds;
the finite diagnostic length and \(C_\epsilon\) may depend on
\(\epsilon\).  The statement is simultaneous on a fixed compact
subinterval of \((0,1)\) carrying a common strict margin, as in the
last paragraph of the proof of \cite[Proposition~17.1]{CM}.
\end{lemma}

\begin{proof}
We prove the stronger square-density residual estimate that replaces
the linear energy dependence in \cite[Lemma~6.5]{SG}.  Use the
zero-diagonal interaction \(J_*\), so the cube coordinates in this
proof are spins, not the scaled edge coordinates.  For a spin flip
\(\sigma^{(i)}\), put
\[
 d_i g(\sigma)=\frac{g(\sigma)-g(\sigma^{(i)})}{2\sigma_i}.
\]
The primary iterates of \cite[Section~6]{SG} are
\[
\begin{gathered}
 m^0=\sigma,\qquad h^1=h+J_*\sigma,\qquad
 m^l=\tanh h^l,\qquad b_l=1-\|m^l\|^2/n,\\
 h^{l+1}=h+J_*m^l-jb_lm^{l-1},\qquad
 R_l=m^{l-1}-m^l .
\end{gathered}
\]
For a fixed ordinary recipe admissible from level \(l\), let \(U\)
be a terminal vector in that recipe.  We claim that, whenever
\(\sup_\sigma(|H_0|+\|dH_0\|)\le C_0\),
\begin{equation}
 \left|\mu_h[f^2H_0R_l^{\mathsf T}U]\right|
 \le C\left(N_h+\sqrt{N_h\mathcal D_h(f)}\right).
                                                               \label{warm:sqrt-residual}
\end{equation}
The constant has the dependence on the fixed recipe, depth, coefficient
bounds and finite word diagnostics allowed in the source lemma, and
is independent of \(f\).  The residual claim itself does not require
root stability.

Here are the base case and the induction, including the point where
the improvement occurs.  The pointwise estimate
\cite[Lemma~6.4, Equation~(6.20)]{SG} gives
\[
 \|U\|+\sum_i|d_iU_i|\le C .
\]
It also gives the same bound for
\(\widetilde U=H_0U\): indeed
\[
 d_i(H_0U_i)=H_0\,d_iU_i+U_i(\sigma^{(i)})d_iH_0,
 \qquad
 \|(U_i(\sigma^{(i)}))_i\|
       \le\|U\|+2\|\diag\nabla U\|_2\le C,
\]
and Cauchy--Schwarz bounds the sum of the added diagonal derivatives.
Write \(C_U\) for a bound on
\(\|\widetilde U\|+\sum_i|d_i\widetilde U_i|\).

Let \(p_i^\pm=\mu_h(\sigma_i=\pm1\mid\sigma_{-i})\).
Then \(m_i^1=p_i^+-p_i^-\) and the conditional variance is
\(v_i^1=4p_i^+p_i^-\), in particular \(v_i^1\le4p_i^\pm\)
for each sign.  Since a function of one spin is affine with slope
\(d_if\), the form is
\[
 \mathcal D_h(f)=\mu_h\sum_i v_i^1(d_if)^2 .
\]
The exact one-site integration by parts
\cite[Equation~(6.31)]{SG} gives
\begin{align*}
 \mu_h[f^2R_1^{\mathsf T}\widetilde U]
 &=\mu_h\sum_i v_i^1d_i(f^2\widetilde U_i)\\
 &=\mu_h\sum_i v_i^1 f^2d_i\widetilde U_i
   +\mu_h\sum_i v_i^1(f_{i,+}+f_{i,-})d_if\,
                                  \widetilde U_i(\sigma^{(i)}).
\end{align*}
Here \(f_{i,\pm}=f(\sigma^{i,\pm})\), and
\(d_i(f^2)=(f_{i,+}+f_{i,-})d_if\), also for signed \(f\).
The first term is at most \(C_UN_h\) in absolute value.
For the second set
\[
 M_i(\sigma_{-i})=\max_\pm|\widetilde U_i(\sigma^{i,\pm})|.
\]
At each configuration,
\[
 \left(\sum_iM_i(\sigma_{-i})^2\right)^{1/2}
 \le\|\widetilde U(\sigma)\|
       +2\|\diag\nabla\widetilde U(\sigma)\|_2\le 2C_U.
\]
Weighted Cauchy--Schwarz bounds that second term by
\[
 \mathcal D_h(f)^{1/2}
 \left(\mu_h\sum_i v_i^1(f_{i,+}+f_{i,-})^2M_i^2\right)^{1/2}.
\]
The two separate bounds \(v_i^1\le4p_i^\pm\) imply
\[
 v_i^1(f_{i,+}+f_{i,-})^2
       \le8(p_i^+f_{i,+}^2+p_i^-f_{i,-}^2).
\]
Since \(M_i\) is independent of \(\sigma_i\), conditioning term by
term gives
\[
 \mu_h\sum_i v_i^1(f_{i,+}+f_{i,-})^2M_i^2
       \le8\mu_h\!\left[f^2\sum_iM_i^2\right]\le32C_U^2N_h .
\]
This proves \eqref{warm:sqrt-residual} for \(l=1\).

The finite induction in the source preserves this new right side.
For clarity, let \(U\) be admissible from level \(l+1\), put
\(a_l=j(b_l-b_{l-1})\), and define
\[
 U_i'=\Phi(h_i^l;h_i^{l+1},a_l)U_i,\qquad
 \Phi(z;v,a)=\int_0^1
 e^{a(1-u^2)/2}\frac{\cosh(u(z-v))}{\cosh z\cosh v}\,\dd u .
\]
Here \(|a_l|\le j\).  The integral formula bounds \(\Phi\) and every
needed fixed-order derivative uniformly in its two real field arguments:
the hyperbolic derivatives introduce bounded hyperbolic-tangent factors,
and the exponential factor and its derivatives are bounded for
\(|a|\le j\).  Together with the explicit level indices, these bounds
make \(U'\) admissible from level \(l\).  The scalar identity in
\cite[Section~6]{SG} then gives the residual expansion below.  Adjoin
the auxiliary field
\[
 y'=J_*U'-j\sum_q\theta_qm^{q-1}
              -j\sum_b\omega_bw^b,\qquad
 \theta_q=\frac1n\sum_i\partial_qU'_i,\quad
 \omega_b=\frac1n\sum_i\partial_{y^b}U'_i .
\]
Here \(w^b,y^b\) are the finitely many earlier auxiliary sources and
fields of the recipe; the partials differentiate their explicitly
displayed local arguments, holding parameters fixed, as in
\cite[Definition~6.2]{SG}.  Equation~(6.33) of that source is
\begin{align*}
 R_{l+1}^{\mathsf T}U
 ={}&R_l^{\mathsf T}y'+a_lR_l^{\mathsf T}U'
       -jb_{l-1}R_{l-1}^{\mathsf T}U'
       +j\sum_b\omega_bR_l^{\mathsf T}w^b\\
 &+j\sum_{q\ge l}\theta_qR_l^{\mathsf T}m^{q-1}
       -a_ln\theta_l ,
\end{align*}
with the \(R_{l-1}\) term absent when \(l=1\).  All first-line terms
are earlier residuals with ordinary admissible sources.  Their scalar
multipliers and their discrete-gradient norms are bounded by the
fixed-recipe bounds.  For \(q>l\) in the second line, use the source
\(m^{q-1}/\sqrt n\) and multiplier \(j\sqrt n\,\theta_q\);
\cite[Equation~(6.21)]{SG} gives the same boundedness for that
multiplier.  Finally the exact identity
\[
 R_l^{\mathsf T}m^{l-1}
       =n(b_l-b_{l-1})-R_l^{\mathsf T}m^l
\]
cancels the displayed \(-a_ln\theta_l\).  The remainder again uses
an earlier residual, the source \(m^l/\sqrt n\), and a bounded
multiplier.  A bounded multiplier \(\zeta\) with bounded \(\|d\zeta\|\)
can always be absorbed into \(H_0\), because
\[
 d_i(\zeta H_0)=\zeta(\sigma^{(i)})d_iH_0+H_0(\sigma)d_i\zeta .
\]
Every induction call therefore retains the original \(f,N_h,\mathcal D_h(f)\).
There are finitely many calls with a smaller residual index.  This
proves \eqref{warm:sqrt-residual} at every fixed depth.

We spell out its consequences for selection and the mean.  Put
\(S_l=R_l^{\mathsf T}m^l/\sqrt n\).
The primary derivative bounds in \cite[Lemma~7.1]{SG} give
\(|S_l|\le2\sqrt n\) and \(\|dS_l\|\le C_l\).  Apply
\eqref{warm:sqrt-residual} with \(U=m^l/\sqrt n\) and
\(H_0=S_l/\sqrt n\).  Division by \(N_h\sqrt n\) yields the
strengthened form of \cite[Equation~(7.1)]{SG}:
\begin{equation}
 \nu_h\!\left[\left(\frac{S_l}{\sqrt n}\right)^2\right]
 \le\frac{C_l}{\sqrt n}
       \left(1+\sqrt{\frac{\mathcal D_h(f)}{N_h}}\right).
                                                               \label{warm:sqrt-selection-moment}
\end{equation}
For a small fixed \(\rho>0\), choose a fixed depth \(d\) with
\(\lfloor d/2\rfloor\rho^2/4>2\), and then \(e_\rho>0\) with
\(2de_\rho<1\).  The telescoping identity
\[
 \frac{\|R_l\|^2}{n}=b_l-b_{l-1}-\frac{2S_l}{\sqrt n}
\]
shows that \(\max_{l\le d}|S_l|/\sqrt n\le e_\rho\) forces a
consecutive pair \(R_{k-1},R_k\) to have both norms at most
\(\rho\sqrt n/2\).  Otherwise the disjoint pairs would make the sum
of squared norms exceed \(2\), whereas telescoping bounds it by
\(1+2de_\rho<2\).

As in that source lemma, choose smooth \(C_k\in[0,1]\), equal to one
on this smaller pair region and supported where both norms are at
most \(\rho\sqrt n\).  The primary derivative bounds give
\(\|dC_k\|\le C_\rho/\sqrt n\).  Put
\[
 \widehat C_k=C_k\prod_{i=2}^{k-1}(1-C_i),\qquad
 D_{\rm sel}=1-\sum_{k=2}^d\widehat C_k=\prod_{k=2}^d(1-C_k).
\]
Positive deficit implies \(\max_{l\le d}|S_l|/\sqrt n>e_\rho\).
Chebyshev's inequality and \eqref{warm:sqrt-selection-moment} give
\begin{equation}
 \nu_h(D_{\rm sel}>0)
 \le\frac{C_\rho}{\sqrt n}
       \left(1+\sqrt{\frac{\mathcal D_h(f)}{N_h}}\right).
                                                               \label{warm:sqrt-deficit}
\end{equation}
The cutoffs are admissible bounded multipliers in
\eqref{warm:sqrt-residual}.

Choose \(\rho\le\epsilon\) with
\((\|J_*\|+4j)\rho<2\rho_0\).
On the support of \(C_k\), the recursion gives
\[
 F_h(h^k)=J_*R_k-jb_{k-1}(R_{k-1}+R_k)
                  +j(b_k-b_{k-1})m^k,
\]
so \(\|F_h(h^k)\|\le(\|J_*\|+4j)\rho\sqrt n\).
The root estimate of \cite[Lemma~5.7]{SG},
\[
 \|y-y_*\|\le
 \left(1+\frac{\|J_*\|+3j}{b}\right)\|F_h(y)\|
 \quad\hbox{when }\|F_h(y)\|<2\rho_0\sqrt n,
\]
and the Lipschitz property of \(\tanh\) therefore give
\(\|m^k-m_*\|\le C_{\rm lead}\rho\sqrt n\), with a constant
independent of the depth.  For a deterministic unit vector \(e\),
sum \eqref{warm:sqrt-residual} for \(R_1,\ldots,R_k\), with source
\(e\) and multiplier \(\widehat C_k\).  Since
\(\sigma-m^k=\sum_{l\le k}R_l\), this bounds its reweighted mean
against \(e\) by \(C_\rho(1+\sqrt{\mathcal D_h(f)/N_h})\).
The weighted root errors together cost at most
\(C_{\rm lead}\epsilon\sqrt n\), because the cutoffs sum to at most
one.  The remaining deficit costs at most
\(2\sqrt n\,\nu_h(D_{\rm sel}>0)\), controlled by
\eqref{warm:sqrt-deficit}.  Sum over the fixed number of cutoffs,
divide by \(\sqrt n\), and take the supremum over \(e\).
This proves \eqref{warm:sqrt-posterior}.
\end{proof}

\begin{lemma}[Retention at order-\(n\) entropy]\label{warm:large-retention}
Fix \(0<\lambda<C_{\mathrm{lev}}<\infty\), \(C_E<\infty\), and \(0<\omega<1\).
On Gaussian disorder events of probability tending to one, every
test \eqref{warm:restricted-class} with \(\lambda n\le L\le C_{\mathrm{lev}}n\)
and
\begin{equation}
 \frac{\cE(f)}N\le C_E n^{1-\omega}                       \label{warm:large-energy-assumption}
\end{equation}
has an observation path to a product posterior for which
\(H_{\rm end}\ge c_{\lambda,C_{\mathrm{lev}}}n\).  The path and disorder events
are chosen independently of \(f\).
\end{lemma}

\begin{proof}
We use the path and the field tests in
\cite[Section~19.6]{CM}; we verify that its entropy argument remains
valid under \eqref{warm:large-energy-assumption}.  On source disorder
events independent of the initial spin law, the calculation around
Equation~(19.20) there, using its Lemma~17.4, Equation~(17.12), and
Proposition~17.5, gives the following early field-test bound.
On a fixed interval \([t_*,t_b]\) of small
positive scalar precision, the diagnostic and stability failures
have integrated probability \(o(1)\), uniformly under every initial
spin law.  This follows there by conditioning on the spin: the
centered Gaussian-noise diagnostics have integrated failure \(o(1)\),
the additional spin shift has normalized norm \(O(t)\), and the
Gaussian low-band small-ball estimate is uniform in its center.
No factor \(e^L\) is used for this interval.  From \(t_b\) until
the interaction coefficient reaches a fixed globally controlled
range, the buffered stability test fails anywhere on the path with
probability at most \(e^{-a n}\) under the ordinary law \(Q_\mu\),
for a fixed \(a>0\).
This is the later field bound in \cite[Lemma~17.3, Section~19.4, and
Section~19.6 after Equation~(19.21)]{CM}.
The word diagnostics and stability constants on the retained
regions are those of Lemma~\ref{warm:posterior-mean}.

Here is the order of the fixed choices in that construction.  Choose
first a large scalar endpoint \(t_{\rm end}\) and a small interaction
coefficient below which covariance is globally bounded.  Choose
\(t_b>0\) for the small-positive-field estimate, and obtain the
buffer and rate \(a\) for the later test.  The norm bound and the
fixed total duration give constants \(C_{\rm jump}\) and
\(C_{\rm aft}\) for the jump loss divided by \(t_*n\) and the
loss after failure divided by \(n\), uniformly when
\(1-\beta'\le t_*/3\).  Choose a target drift budget \(\eta_1>0\)
so that
\[
 3\eta_1(1+C_{\rm aft}/a)<\lambda/4.
\]
Next choose \(0<t_*\ll t_b\) with
\(C_{\rm jump}t_*\le\eta_1\), and finally \(1-\beta'>0\)
sufficiently small relative to these choices, including
\(1-\beta'\le t_*/3\).  Start with precision
\[
 B_*=t_*I+(1-\beta')W,
\]
which is positive when \(1-\beta'\) is small enough.  Increase scalar
precision to \(t_{\rm end}I+(1-\beta')W\), and then traverse
\[
 B_v=(t_{\rm end}+4v)I+(1-\beta'+v)W,\qquad 0\le v\le\beta'.
\]
The latter rate \(4I+W\) is positive on the norm event.  The endpoint
is \( (t_{\rm end}+4\beta')I+W\), so its posterior is a product.
Before the globally controlled portion, the posterior interaction
coefficient is either \(\beta'<1\) or in a fixed compact subinterval
of \((0,1)\).  Reduce the available diagonal enlargement in
\eqref{warm:stability} to \(0<\eta<(\beta')^{-1}-1\).
Once these stability margins and \(C_{\rm lead}\) are fixed, choose
the mean accuracy \(\epsilon\) in Lemma~\ref{warm:posterior-mean}
so that its leading error, integrated over the bounded path and
converted to a quadratic mean bound, is less than \(\eta_1\).
Then fix its residual tolerance, selection depth, finite word
diagnostics and constant \(C_\epsilon\), and finally take \(n\)
large.

On a valid posterior, apply \eqref{warm:sqrt-posterior} to \(f\)
and to \(1\).  With
\(a_u=\|\nu_u[X]-\mu_u[X]\|/\sqrt n\), this gives
\[
 a_u\le2C_{\rm lead}\epsilon+
       \frac{C_\epsilon}{\sqrt n}
          \left(2+\sqrt{\frac{\cE_{\mu_u}(f)}{N_u}}\right).
\]
By \eqref{warm:conditioned-energy} and Cauchy--Schwarz, the
\(Q_\nu\)-mean of the square root, also with any validity indicator,
is at most \(\sqrt{\cE(f)/N}\).  Under
\eqref{warm:large-energy-assumption} the resulting error is
\[
 \frac{C_\epsilon}{\sqrt n}
       \left(2+\sqrt{\frac{\cE(f)}N}\right)
       =O_\epsilon(n^{-1/2}+n^{-\omega/2})=o(1).
\]
This bound is applied at deterministic observation times; inserting
\(\ind_{\{u<\tau\}}\) for a stopping time only decreases its
Cauchy--Schwarz bound.  Since \(a_u\le2\) always, \(a_u^2\le2a_u\).
The early invalid paths cost at most four times their integrated
probability, hence \(o(1)\).  Thus the normalized quadratic
mean-shift integral, on valid paths and before any later stopping,
is at most \(\eta_1+o(1)\) by the preceding choice of \(\epsilon\).

Let \(u_g\) begin the globally controlled portion and let \(\tau\)
be the first failure of the later buffered test in \([t_b,u_g]\),
with the immediate and no-failure conventions of the source.  The
initial jump leaks at most \(C_{\rm jump}t_*n\) entropy: compare to pure noise
as above and use \(\|B_*\|\le Ct_*\) and \(\|X\|^2=n\).
The stopped form of \eqref{warm:spin-loss}, with the preceding
integrated estimate, gives, with \(Q^\tau\) denoting the law of the
stopped observations, excluding the spin,
\begin{equation}
 D(Q_\nu^\tau\Vert Q_\mu^\tau)
       \le(C_{\rm jump}t_*+\eta_1+o(1))n .              \label{warm:stopped-entropy}
\end{equation}
For the event \(E\) of a failure before \(u_g\), the ordinary
probability is at most \(e^{-an}\).  Data processing for its indicator
then gives
\[
 D(Q_\nu^\tau\Vert Q_\mu^\tau)
 \ge Q_\nu(E)\log\frac1{Q_\mu(E)}-\log2,\qquad
 Q_\nu(E)\le\frac{C_{\rm jump}t_*+\eta_1+o(1)}a .
\]
These are exactly the deductions in
\cite[Equations~(19.21)--(19.22)]{CM}, now with the enlarged energy
range justified by Lemma~\ref{warm:posterior-mean}.
After a failure the remaining precision rate and duration are bounded,
and both spin means have norm at most \(\sqrt n\); the remaining
loss per failed path is therefore at most \(C_{\rm aft}n\).  The total loss
before \(u_g\) is bounded by
\[
 \left(C_{\rm jump}t_*+\eta_1+o(1)\right)
       (1+C_{\rm aft}/a)n .
\]
The fixed choices above make this less than \(\lambda n/4\).
The initial entropy is at least \(L/2\ge\lambda n/2\), hence
\(H(u_g)\ge\lambda n/4\).  In the remaining portion the spin
covariance is bounded for every field.  The linear entropy test
then gives \(-\dot H\le CH\); over its bounded duration it retains
a fixed fraction of \(H(u_g)\).  This proves the lemma.
\end{proof}

\begin{proof}[Proof of Proposition~\ref{warm:profiles}]
Fix \(C_{\mathrm{lev}},\xi\) and a desired disorder confidence.  Apply
Lemma~\ref{warm:small-profile} with
\(\kappa<\min\{\xi,\,5/3-1.662\}\).  After fixing its \(M\), choose
its \(\delta>0\) as in that proof.  For \(L\le n^\delta\),
\eqref{warm:small-profile-bound} and
\eqref{warm:modified-classical} imply both claimed inequalities.

For \(n^\delta<L\le\lambda n\), use the proportional branch of
\cite[Equation~(19.16) and Section~19.4]{CM}.  It gives, at the
product endpoint,
\begin{equation}
 H_{\rm end}\ge cL(L/n)^{(2+2\epsilon)/3}.                \label{warm:moderate-retention}
\end{equation}
Here \(\epsilon>0\) is any fixed sufficiently small tolerance,
and \(\lambda>0\) is chosen afterward.  We justify the exact test
class being used.  The starting dimension in that source is
\(\max\{k_{\min},c_0L\}\), with \(k_{\min}\asymp\sqrt{\log n}\).
It equals \(c_0L\) throughout the present range for large \(n\).
The proportional part of its initial jump estimate is the pure-noise
argument reproduced in Lemma~\ref{warm:small-profile}; it requires only
the entropy at least \(L/2\) and density at most \(e^L\).
The source's support condition is used in its other, minimum-dimension
branch.  Its subsequent propagation uses the likelihood and
conditional-entropy tails in \cite[Equation~(19.8)]{CM} and the
ordinary path tests.  These also require only the same density and
entropy bounds.  The continuation through its finite remaining
scales uses those likelihood tails and conditional Pinsker bounds,
then globally bounded covariance, and so preserves this test class.
Thus \eqref{warm:moderate-retention} applies to
\eqref{warm:restricted-class} without a support restriction.

Since \(L\ge n^\delta\gg\log n\), comparison with the first bound
of \eqref{warm:product-bounds} gives
\[
 \Tn\,\frac{\cE(f)}N
 \ge c n^{-2\epsilon/3}L^{7/6+2\epsilon/3}
 \ge c L^{7/6-(2\epsilon/3)(1/\delta-1)}.
\]
The additive \(O(n^{-3})\) is smaller than
\eqref{warm:moderate-retention}.  The second product bound similarly
gives
\[
 \Tn\,\frac{\cI(f^2)}N
 \ge c n^{-2\epsilon/3}L^{5/3+2\epsilon/3}
 \ge c L^{5/3-(2\epsilon/3)(1/\delta-1)}.
\]
Choose \(\epsilon\) after \(\delta\) so that
\[
 \frac{2\epsilon}{3}(1/\delta-1)
       <\min\{\xi,\,7/6-1.162\}.
\]
Then choose \(\lambda\) below the finitely many fixed rates and
shift buffers in that source argument.  This proves the profiles
in the moderate range.

Finally let \(\lambda n\le L\le C_{\mathrm{lev}}n\).  If the classical bound
were false, then
\(\cE(f)/N\le Cn^{1.162-2/3}\).  This meets
\eqref{warm:large-energy-assumption} for a fixed positive
\(\omega\), whereas the first product bound would give
\[
 H_{\rm end}\le Cn^{1/2+1.162-2/3}+O(n^{-3})=o(n),
\]
contrary to Lemma~\ref{warm:large-retention}.  If the modified bound
were false, \eqref{warm:modified-classical} would instead give
\(\cE(f)/N\le Cn^{1-\xi}\), so that lemma applies with
\(\omega=\xi\).  The second product bound would give
\(H_{\rm end}\le Cn^{1-\xi}=o(n)\), again a contradiction.

All the choices are uniform in \(f\).  More explicitly, the
variance exponent and \(M\) are fixed first, then the small
\(\delta\), then the sharp cap and moderate-entropy tolerances,
then \(\lambda\) and the fixed parameters of the order-\(n\) path.
The finite word length in that last path is fixed after its mean
accuracy.  The final threshold \(K_0\) and \(n\) are chosen last.
Intersecting this finite collection of disorder events preserves
arbitrarily high limiting probability.  This proves the proposition.
\end{proof}

\subsection{Passage from profiles to arbitrary initial laws}
\label{warm:layers-section}

On the norm event \(\|W\|\le3\), comparison of the largest and
smallest spin Hamiltonians gives
\begin{equation}
 \min_\sigma\mu(\sigma)\ge e^{-C_a n}                     \label{warm:atom}
\end{equation}
for an absolute \(C_a\).  In the remainder of the proof choose
\(C_{\mathrm{lev}}>2C_a+1\) in Proposition~\ref{warm:profiles}.  For any density
\(z\ge0\), \(\mu[z]=1\) implies \(z(\sigma)\le1/\mu(\sigma)\);
thus every nonempty positive logarithmic level of \(z\) lies below
\(C_a n\).

We first isolate the bounded test associated with one such level.
This same construction will be used for entropy and higher log moments.

\begin{lemma}[A bounded test from a logarithmic level]\label{warm:log-layer}
Fix \(b>0\) and \(r>0\).  Let \(z>0\) be a density and suppose that,
for some \(d\ge v/2\) and sufficiently large \(v\),
\[
 B=\mu[z\ind_{\{d\le\log z\le1.1d\}}]
       \ge b(v/d)^r>0 .
\]
Define
\begin{equation}
 u=\mathcal U_d(z)
   :=\min\{e^{1.1d},\,z\min(1,ze^{-0.8d})\},
 \qquad f^2=e^{-1.1d}u .                                 \label{warm:level-transform}
\end{equation}
Then \(0\le f\le1\), \(B\le\mu[u]\le1\), and, if
\(N=\mu[f^2]=e^{-L}\),
\[
 1.1d\le L\le1.1d+C_{b,r}\bigl(1+\log_+(d/v)\bigr),
 \qquad
 D\!\left(\frac{f^2\mu}{N}\middle\Vert\mu\right)\ge L/2 .
\]
For large \(v,n\), this test lies in the range
\(K_0\le L\le C_{\mathrm{lev}}n\) of Proposition~\ref{warm:profiles}.  In addition,
\begin{equation}
 \cI(\mathcal U_d(z))\le4\cI(z).                          \label{warm:level-dissipation}
\end{equation}
\end{lemma}

\begin{proof}
In logarithmic coordinates \(l=\log z\), the logarithm of
\(\mathcal U_d(e^l)\) is
\[
 2l-0.8d\quad(l\le0.8d),\qquad
 l\quad(0.8d\le l\le1.1d),\qquad
 1.1d\quad(l\ge1.1d).
\]
Consequently \(u\le z\), and \(u=z\) on the indicated level.
This proves \(B\le\mu[u]\le1\) and
\(L=1.1d-\log\mu[u]\), with the stated upper bound from the lower
bound on \(B\).

Let \(E=\{\log z\ge0.7d\}\).  Since \(u\le e^{-0.1d}z\) on its
complement,
\[
 \frac{\mu[u\ind_{E^c}]}{\mu[u]}\le\frac{e^{-0.1d}}B,
 \qquad \mu(E)\le e^{-0.7d}.
\]
Set \(e_d=e^{-0.1d}/B\).  Data processing for the indicator of
\(E\) gives
\[
 D\!\left(\frac{u\mu}{\mu[u]}\middle\Vert\mu\right)
 \ge (1-e_d)\log\frac1{\mu(E)}-\log2
 \ge0.7d(1-e_d)-\log2 .
\]
Uniformly for \(d\ge v/2\), the hypothesis on \(B\) makes
\(e_d\to0\) and \(1+\log_+(d/v)=o(d)\) as \(v\to\infty\).
The last lower bound therefore exceeds \(L/2\) for large \(v\).
The level is nonempty, so \(d\le C_an\) by \eqref{warm:atom};
the bound on \(L\), with \(C_{\mathrm{lev}}>2C_a+1\), puts it below \(C_{\mathrm{lev}}n\)
for large \(n\).  Its lower bound puts it above \(K_0\) for large
\(v\).

Finally \(\mathcal U_d\) is nondecreasing and \(2\)-Lipschitz as a
function of \(z\), while \(\log\mathcal U_d(e^l)\) is nondecreasing
and \(2\)-Lipschitz as a function of \(l\).  The product of their
increments on an edge is at most four times the product of the
increments of \(z\) and \(\log z\).  Summing proves
\eqref{warm:level-dissipation}.
\end{proof}

\begin{lemma}[Entropy dissipation for every density]\label{warm:entropy-profile}
With tight constants, every strictly positive density \(z\) with
\(\Ent_\mu(z)\ge K_1\) satisfies
\begin{equation}
 \Tn\cI(z)\ge c\,\Ent_\mu(z)^{1.662}.                    \label{warm:full-entropy-profile}
\end{equation}
\end{lemma}

\begin{proof}
Put \(H=\Ent_\mu(z)\) and \(d_j=(H/2)(1.1)^j\).
The part of \(\mu[z\log z]\) with \(\log z<H/2\) is at most \(H/2\).
Thus the levels \([d_j,1.1d_j]\) account for at least \(H/2\).
If their tilted masses \(B_j\) were all less than
\(c(H/d_j)^{1.1}\), their contribution would be at most
\[
 1.1\sum_{j\ge0}d_jB_j
 <1.1cH^{1.1}\sum_{j\ge0}d_j^{-0.1}\le CcH .
\]
Choose \(c\) small.  One level therefore has
\(B\ge c(H/d)^{1.1}\).  Apply Lemma~\ref{warm:log-layer} with \(v=H\)
and \(r=1.1\), followed by \eqref{warm:modified-profile} with any
fixed \(\xi<5/3-1.662\).  Homogeneity of \(\cI\) gives
\[
 \Tn\cI(z)\ge\tfrac14\Tn\cI(u)
       \ge c\,\mu[u]L^{1.662}
       \ge cB d^{1.662}
       \ge c'H^{1.1}d^{0.562}\ge c''H^{1.662}.
\]
The constants and the required lower threshold on \(H\) are uniform.
\end{proof}

For the kernel we need a classical support profile.  The next
deduction is the equal-height and dyadic-layer argument of
\cite[Lemma~19.1]{CM}, keeping the two powers from
\eqref{warm:classical-profile}.

\begin{lemma}[A support profile]\label{warm:support-profile}
With the constants and \(\delta\) of Proposition~\ref{warm:profiles}, every
nonnegative \(g\), supported on a set of \(\mu\)-mass \(e^{-L_{\rm supp}}\)
with \(L_{\rm supp}\ge K_0\), satisfies
\begin{equation}
 \Tn\cE(g)\ge c\mu[g^2]\,\Phi_n(L_{\rm supp}),\qquad
 \Phi_n(L)=
 \begin{cases}
 \min\{L^{1.654},\,n^{1.154\delta}\},&L\le n^\delta,\\
 L^{1.154},&L>n^\delta .
 \end{cases}                                                \label{warm:support-bound}
\end{equation}
The zero function satisfies the assertion by convention.
\end{lemma}

\begin{proof}
First let \(0\le g\le1\) equal one on a nonempty set of mass
\(a=e^{-L_a}\), and let its support have mass at most \(e^{-L_{\rm supp}}\).
Take
\[
 K=\left\lceil\log_2(1+L_a/L_{\rm supp})\right\rceil,\qquad
 g_i=K\bigl((g-i/K)_+\wedge K^{-1}\bigr),\quad 0\le i<K .
\]
Let \(S_i=\mu(g>i/K)\), \(T_i=\mu(g\ge i/K)\).
Some \(i\) satisfies \(T_{i+1}\ge S_i^2\); otherwise induction gives
\(a\le T_K<e^{-2^KL_{\rm supp}}<a\).  For this \(i\), with \(N_i=\mu[g_i^2]\),
\[
 a\le T_{i+1}\le N_i\le S_i,\qquad
 D\!\left(\frac{g_i^2\mu}{N_i}\middle\Vert\mu\right)
       \ge\log(1/S_i)\ge\tfrac12\log(1/N_i).
\]
Its logarithmic inverse mass \(L_i\) lies in \([L_{\rm supp},L_a]\),
and \(L_a\le C_an\) by \eqref{warm:atom}.  For \(L\ge K_0\)
large enough, the function
\[
 e^{-L}
 \begin{cases}L^{1.662},&L\le n^\delta,\\
               L^{1.162},&L>n^\delta
 \end{cases}
\]
is decreasing within each range and has a downward jump at their
boundary.  Since \(\cE(g_i)\le K^2\cE(g)\),
\eqref{warm:classical-profile} implies
\[
 \Tn\cE(g)\ge
 \frac{ca}{(1+\log L_a)^2}
 \begin{cases}L_a^{1.662},&L_a\le n^\delta,\\
               L_a^{1.162},&L_a>n^\delta .
 \end{cases}
\]
Absorbing the logarithmic denominator by lowering each exponent by
\(0.008\) gives
\begin{equation}
 \Tn\cE(g)\ge ca
 \begin{cases}L_a^{1.654},&L_a\le n^\delta,\\
               L_a^{1.154},&L_a>n^\delta .
 \end{cases}                                                \label{warm:set-capacity}
\end{equation}
Indeed \(L^{0.008}/(1+\log L)^2\) has a positive infimum for
\(L\ge K_0\).

For arbitrary nonnegative \(g\) with the stated support, use
\(\varphi_j=(g-2^j)_+\wedge2^j\), \(j\in\mathbb Z\), and let
\(a_{j+1}=\mu(g\ge2^{j+1})\).  Omit empty upper level sets.
The normalized layer \(2^{-j}\varphi_j\) is one on that set and
has the original support restriction.  Its \(L_a=\log(1/a_{j+1})\)
is at least \(L_{\rm supp}\).  The right side of
\eqref{warm:set-capacity}, divided by \(a\), is at least
\(\Phi_n(L_{\rm supp})\): before the boundary it is at least \(L_{\rm supp}^{1.654}\),
and after the boundary it is at least \(n^{1.154\delta}\), or at
least \(L_{\rm supp}^{1.154}\) when \(L_{\rm supp}>n^\delta\).
For each oriented edge the nonnegative increments of the
\(\varphi_j\)'s sum to the increment of \(g\), so the sum of their
squares is at most its square.  Therefore
\[
 \Tn\cE(g)\ge\sum_j\Tn\cE(\varphi_j)
       \ge c\Phi_n(L_{\rm supp})\sum_j4^ja_{j+1}.
\]
Pointwise,
\(\sum_j4^j\ind_{\{g\ge2^{j+1}\}}\ge g^2/12\).
Integration proves \eqref{warm:support-bound}.
\end{proof}

\begin{proof}[Proof of Theorem~\ref{thm:warming}]
Fix the parameters in the theorem and intersect the finitely many
events needed for the preceding results.  The finite heat-bath chain
has a strictly positive kernel at every positive time.  Hence
\(w_s^\rho>0\) for \(s>0\), for every initial law, and symmetry of
the generator gives
\[
 \frac{\dd}{\dd s}\Ent_\mu(w_s^\rho)
       =-\Tn\cI(w_s^\rho).
\]
When the entropy exceeds the threshold of
Lemma~\ref{warm:entropy-profile}, its derivative is at most
\(-c\Ent_\mu(w_s^\rho)^{1.662}\).  Integrating this differential
inequality from any earlier positive time, and then letting that
time decrease to zero, yields uniformly in the initial law
\[
 \Ent_\mu(w_s^\rho)
       \le C\left(1+s^{-1/0.662}\right)
       \le C'\left(1+s^{-1.54}\right).
\]
This proves \eqref{warm:entropy}.

For the kernel, put \(Q_s=\mu[(w_s^\rho)^2]=e^{L_s}\).
The atom bound gives \(Q_s\le e^{C_an}\), so its logarithm is
within the range where the support profile was proved.
Whenever \(L_s\) is large, set
\(g=(w_s^\rho-Q_s/4)_+\).  Its support has mass at most
\(4e^{-L_s}\) by \(\mu[w_s^\rho]=1\), and
\[
 \mu[g^2]\ge Q_s-2(Q_s/4)\mu[w_s^\rho]=Q_s/2,\qquad
 \cE(g)\le\cE(w_s^\rho).
\]
Its actual support logarithm is
\[
 L_{\rm supp}=\log\frac1{\mu\{w_s^\rho>Q_s/4\}}
       \ge L_s-\log4.
\]
The support is nonempty by the preceding square-mass bound, so
\eqref{warm:atom} also gives \(L_{\rm supp}\le C_an\).
Since \(\dot Q_s=-2\Tn\cE(w_s^\rho)\),
Lemma~\ref{warm:support-profile} and monotonicity of \(\Phi_n\) give
\[
 -\dot L_s\ge c\,\Phi_n(L_s-\log4)
\]
above a fixed threshold.  A shift of the argument by \(\log4\)
only changes constants there.  For any large \(K'\), the time
allowed by this differential inequality above \(K'\) is at most
\begin{equation}
 \int_{K'}^\infty\frac{C\,\dd L}{\Phi_n(L)}
       \le C(K')^{-0.654}+C n^{-0.154\delta}.
                                                               \label{warm:hitting-integral}
\end{equation}
For completeness, below
\(n^{\delta(1.154/1.654)}\) integrate \(L^{-1.654}\).
The following constant portion of \(\Phi_n\), up to \(n^\delta\),
costs at most \(n^\delta/n^{1.154\delta}\), and the integral of
\(L^{-1.154}\) above \(n^\delta\) has the same order.
Inverting \eqref{warm:hitting-integral} gives
\[
 L_s\le C(1+s^{-1/0.654})
       \qquad\hbox{for }s\ge Cn^{-0.154\delta},
\]
uniformly even for point initial laws.  Reversibility and two half
steps show that
\[
 p_{s\Tn}(\sigma,\tau)
 =\mu[p_{s\Tn/2}(\sigma,\cdot)p_{s\Tn/2}(\tau,\cdot)]
 \le\bigl(\mu[p_{s\Tn/2}(\sigma,\cdot)^2]\,
          \mu[p_{s\Tn/2}(\tau,\cdot)^2]\bigr)^{1/2}.
\]
The preceding bound, with \(1/0.654<1.54\), proves
\eqref{warm:kernel} after enlarging the vanishing threshold.

We finish by proving \eqref{warm:log-moment}.  For the fixed integer
\(m\ge1\), let \(\Psi_m\) be the convex function on \((0,\infty)\)
determined by
\[
 \Psi_m(1)=\Psi_m'(1)=0,\qquad
 \Psi_m''(z)=\frac{(1+|\log z|)^{m-1}}{z}.
\]
It extends continuously to zero: in its twice-integrated formula the
possible singularity is bounded by
\(\int_0^1(1+|\log z|)^{m-1}\dd z<\infty\).
For a density \(z\), put
\[
 V_m(z)=\mu[\Psi_m(z)],\qquad
 \mathsf M_m(z)=\left(1+\mu[z(\log_+z)^m]\right)^{1/m}.
\]
There are constants depending only on \(m\) such that
\begin{equation}
 c_m\mathsf M_m(z)^m-C_m
       \le V_m(z)\le C_m\mathsf M_m(z)^m.                 \label{warm:moment-comparison}
\end{equation}
Indeed for \(z\ge1\),
\(\Psi_m'(z)=((1+\log z)^m-1)/m\).  Integrating this on
\([1,z]\) gives upper and lower bounds of order
\(z(\log z)^m\) when \(\log z\) is large; the remaining range
contributes at most a constant times \(1+z\).  On \([0,1]\)
\(\Psi_m\) is bounded.

For a strictly positive density \(z\), define the corresponding
dissipation \(\mathcal D_m(z)=\cE(z,\Psi_m'(z))\).
We claim that, for \(\mathsf M_m(z)\) above a fixed threshold,
\begin{equation}
 \Tn\mathcal D_m(z)
       \ge c_{m,\xi}\mathsf M_m(z)^{m+2/3-\xi}.
                                                               \label{warm:moment-dissipation}
\end{equation}
Put \(V=\mathsf M_m(z)\) and use levels
\(d_j=(V/2)(1.1)^j\).  If \(V\) is large, the part of
\(\mu[z(\log_+z)^m]\) below \(V/2\) is at most \(2^{-m}V^m\).
Geometric summation as in Lemma~\ref{warm:entropy-profile} gives a level
\(d\ge V/2\) with
\[
 B=\mu[z\ind_{\{d\le\log z\le1.1d\}}]
       \ge c_m(V/d)^{m+0.1}.
\]
Indeed otherwise
\((1.1)^m\sum_jd_j^mB_j\le C_m c_m V^m\) would be too small.
The test \(u=\mathcal U_d(z)\) from Lemma~\ref{warm:log-layer} therefore
satisfies \eqref{warm:restricted-class}, and
\eqref{warm:modified-profile} gives
\begin{equation}
 \Tn\cI(u)\ge c\,\mu[u]L^{5/3-\xi}
       \ge cB d^{5/3-\xi}.                               \label{warm:moment-layer-profile}
\end{equation}

We check that the new dissipation controls this layer with its
additional weight \(d^{m-1}\).  On an edge write
\(a\ge b>0\), \(x=\log a\), \(y=\log b\).  Then
\[
 \Psi_m'(a)-\Psi_m'(b)
       =\int_y^x(1+|t|)^{m-1}\dd t.
\]
If \(x\ge0.4d\), the last integral is at least
\(c_m d^{m-1}(x-y)\).  To verify this uniformly in \(y\), use the
top half of \([0,x]\) when \(-x\le y\le x/2\), the entire interval
when \(y\ge x/2\), and the negative half nearest \(y\) when
\(y<-x\); in each case a fixed fraction of the interval has
\(|t|\) at least a fixed fraction of \(x\).  The two
Lipschitz bounds in the proof of Lemma~\ref{warm:log-layer} then apply.
If \(x<0.4d\), both endpoints are in the lower branch of
\(\mathcal U_d\).  On \([b,a]\) its derivative is at most
\(2e^{-0.4d}\), while its logarithm has slope two in log
coordinates.  Since \(d^{m-1}e^{-0.4d}\le C_m\), the same
conclusion follows using \((1+|t|)^{m-1}\ge1\).
Thus in both cases
\[
 d^{m-1}\bigl(\mathcal U_d(a)-\mathcal U_d(b)\bigr)
       \bigl(\log\mathcal U_d(a)-\log\mathcal U_d(b)\bigr)
 \le C_m(a-b)\bigl(\Psi_m'(a)-\Psi_m'(b)\bigr).
\]
Summation gives
\(\mathcal D_m(z)\ge c_m d^{m-1}\cI(u)\).
Combining this with \eqref{warm:moment-layer-profile} and the
lower bound on \(B\), and using \(2/3-\xi>0.1\), gives
\[
 \Tn\mathcal D_m(z)
 \ge cB d^{m+2/3-\xi}
 \ge c_m V^{m+0.1}d^{2/3-\xi-0.1}
 \ge c_m' V^{m+2/3-\xi}.
\]
This proves \eqref{warm:moment-dissipation}.

Along the finite semigroup at positive times,
\[
 \frac{\dd}{\dd s}V_m(w_s^\rho)
       =-\Tn\mathcal D_m(w_s^\rho).
\]
Equations \eqref{warm:moment-comparison}--\eqref{warm:moment-dissipation}
give, above a fixed threshold,
\[
 \frac{\dd}{\dd s}V_m(w_s^\rho)
       \le-c V_m(w_s^\rho)^{1+(2/3-\xi)/m}.
\]
Integrating from an earlier positive time and letting that time
decrease to zero yields
\(V_m(w_s^\rho)\le C(1+s^{-m/(2/3-\xi)})\).
The comparison \eqref{warm:moment-comparison} proves
\eqref{warm:log-moment}.  This argument allows initial zeros,
because \(\Psi_m\) is continuous at zero and the kernel is strictly
positive afterward.

The three conclusions \eqref{warm:entropy}--\eqref{warm:log-moment}
depend only on the off-diagonal couplings.  Let \(\mathcal G_n^{\rm aug}\)
be the final event in the augmented Gaussian space, and put
\[
 \mathcal G_n^{\rm off}
 =\left\{W_{\rm off}:
       \Prob(\mathcal G_n^{\rm aug}\mid W_{\rm off})>0\right\}.
\]
Then \(\Prob(\mathcal G_n^{\rm off})\ge
\Prob(\mathcal G_n^{\rm aug})\).  For fixed off-diagonal couplings the
three inequalities are deterministic, so their validity on an augmented
realization in \(\mathcal G_n^{\rm aug}\) descends to
\(\mathcal G_n^{\rm off}\).  Taking this as the
theorem event proves the stated disorder quantifiers.
\end{proof}

\begin{corollary}[A polynomial tail bound for the heat-kernel logarithm]
\label{warm:kernel-log-tail}
Fix \(\beta>0\), \(A>0\), \(0<C_\star<\infty\), and
\(\vartheta>0\).  On Gaussian disorder events of limiting lower
probability at least \(1-\vartheta\), with a fixed constant \(C\),
\begin{equation}
 \sup_{\sigma}\ \sup_{1\le d\le C_\star\Tn}
 \sum_\tau P_d(\sigma,\tau)
 \ind_{\left\{\log p_d(\sigma,\tau)>
       \frac12 n^\beta(\Tn/d)^{3/2}\right\}}
 \le Cn^{-A}                                                \label{warm:kernel-tail}
\end{equation}
for all sufficiently large \(n\).  For any initial law \(\rho\),
averaging the displayed row probabilities over \(\sigma\sim\rho\)
preserves the same bound.
\end{corollary}

\begin{proof}
Choose \(\xi>0\) so small that
\[
 \frac{\xi}{2/3-\xi}<\frac{\beta}{2},
\]
then choose a fixed integer \(m>2A/\beta\).  Apply
\eqref{warm:log-moment} to the point initial law at \(\sigma\)
and time \(s=d/\Tn\).  Set \(v=\Tn/d\).  Uniformly for
\(1\le d\le C_\star\Tn\), so \(C_\star^{-1}\le v\le n^{2/3}\),
\[
 \frac{1+v^{1/(2/3-\xi)}}{n^\beta v^{3/2}}
 \le C_{\xi,C_\star}
       n^{-\beta+\xi/(2/3-\xi)}
 \le C_{\xi,C_\star}n^{-\beta/2}.
\]
Markov's inequality under the transition law
\(P_d(\sigma,\tau)=\mu(\tau)p_d(\sigma,\tau)\) now bounds the
left side of \eqref{warm:kernel-tail} by
\[
 \left(
 \frac{2\{1+\mu[p_d(\sigma,\cdot)
                  (\log_+p_d(\sigma,\cdot))^m]\}^{1/m}}
      {n^\beta(\Tn/d)^{3/2}}\right)^m
 \le C_m n^{-m\beta/2}\le C_m n^{-A}.
\]
The choices of \(\xi,m\) are fixed in that order before \(n\).
The statement for \(\sigma\sim\rho\) follows by averaging these
row-specific probabilities.
\end{proof}

\section{Positive-time Gaussian limits and the uniform start}
\label{sec:gaussian-limits}

Two inputs give the positive-time edge evolution.  For bounded cylinder
sources, the finite Markov resolvent contraction bounds the optimizers
used in the stationary relaxed-energy theorem.  A variational argument
then gives the required strong semigroup convergence.  The warming
kernel bound controls the evolved densities, whose weak limits inherit
that evolution.  Warming also controls the total mass in a growing
block of edge modes.  A rotation argument distributes that mass among
individual modes and retains information specific to the uniform start.

Throughout this section we use Gaussian couplings and the completed
GOE matrix of Section~\ref{sec:background}, writing
\(\Lambda=\diag(\lambda_1,\ldots,\lambda_n)\).  Write
\(w_{n,s}=d(\nu_nP_{sT_n})/d\mu\) when the system size needs emphasis.
We use the represented-subsequence convention of that section and
denote the probability and expectation on a representation by
\(\widehat{\Prob}\) and \(\widehat{\E}\).

\begin{proposition}[Positive-time Gaussian subsequential limits]
\label{prop:gaussian-subsequence}
Every sequence of Gaussian system sizes tending to infinity has a further represented
subsequence on which the following assertions hold almost surely.  The
edge coordinates converge to \(g\), and there is a measurable random
family \(h_s\in L^2(\Pi_g)\), \(s>0\), with the following properties.
The family may at this stage depend on the represented data in addition
to \(g\).

There is a finite random \(C\ge1\), independent of \(s,t\), such that
\begin{equation}
 0\le h_s\le \mathcal K(s):=C\exp(C s^{-77/50}),\qquad
 \Pi_g[h_s]=1,\qquad
 h_{s+t}=\Sg_t h_s,\qquad s,t>0,
 \label{glim:density-family}
\end{equation}
where \(\Sg_t=\exp(-t c_*H_g)\).  The map \(s\mapsto h_s\) is continuous
in \(L^2(\Pi_g)\).  For each \(t>0\), \(\Sg_t\) has a nonnegative symmetric
kernel \(k_t^g\) relative to \(\Pi_g\) satisfying
\begin{equation}
 0\le k_t^g(x,y)\le \mathcal K(t)
 \quad\text{for \(\Pi_g\otimes\Pi_g\)-almost every \((x,y)\)}.
 \label{glim:limit-kernel}
\end{equation}
The constant \(77/50\) is the exponent \(1.54\) in
Theorem~\ref{thm:warming}.

On this subsequence,
\begin{equation}
 A_{n,J}\longrightarrow A_g
 \quad\text{in }C_{\mathrm{loc}}((0,\infty)),\qquad
 B_{n,J}\longrightarrow B_h
 \quad\text{in }C_{\mathrm{loc}}((0,\infty)^2),
 \label{glim:functional-subsequence}
\end{equation}
where
\begin{equation}
 B_h(s,t)=\sum_{a\ge1}\Pi_g[h_s x_a\,\Sg_t x_a].
 \label{glim:correlation-series}
\end{equation}
The series in \eqref{glim:correlation-series} converges locally uniformly
and is continuous.  Finally,
\begin{equation}
 \|h_s-1\|_{L^2(\Pi_g)}\longrightarrow0,\qquad
 B_h(s,\cdot)\longrightarrow A_g(\cdot)
 \quad\text{in }C_{\mathrm{loc}}((0,\infty))
 \quad(s\to\infty).
 \label{glim:relaxation}
\end{equation}
\end{proposition}

\begin{proof}
We first pass the semigroups on bounded sources, then construct the
densities, and finally justify the sum over all edge coordinates.

\paragraph{Varying \(L^2\) spaces.}
Choose a countable family \(\mathscr C\) of smooth compact cylinders
with this approximation property: every smooth compact cylinder
is approximated uniformly, together with its first derivatives, by
elements of \(\mathscr C\) supported in one fixed compact coordinate set.
Rational smooth approximations on rational boxes give such a family.
The stated approximation and the definition of the closed form make
\(\mathscr C\) dense in its form norm, and also dense in \(L^2(\Pi_g)\).
For linear tests put
\(\mathscr V=\operatorname{span}_{\mathbb Q}(\mathscr C\cup\{1\})\).
This is a countable family of bounded smooth cylinders, dense in
\(L^1(\Pi_g)\) and \(L^2(\Pi_g)\).  If \(\phi\in\mathscr V\) depends on the
first \(d\) coordinates, let
\(\phi_n=\phi(x_1,\ldots,x_d)\) on the finite cube.

The static cylinder and moment convergence of
\cite[Corollary~3.3]{FUA} gives
\begin{equation}
 \mu[\phi_n\psi_n]\longrightarrow\Pi_g[\phi\psi]
 \quad(\phi,\psi\in\mathscr V).
 \label{glim:static-pairings}
\end{equation}
On sequences bounded in \(L^2(\mu)\), the weak convergence of
Section~\ref{sec:background} can be tested on this countable family:
the pairings with every \(\phi\in\mathscr V\) converge to the
corresponding pairings with \(f\) if and only if the pairings with
every smooth compact cylinder do.  To pass from \(\mathscr V\) to an
arbitrary smooth compact cylinder, approximate it uniformly on a
common support by \(\mathscr C\), and use the uniform \(L^2\) bound
to control the pairing error in both spaces.  The converse follows
from the same approximation, with cutoffs tending to one for the
constant test.
A sequence bounded in \(L^2(\mu)\) has a weakly convergent subsequence:
diagonal extraction of its pairings defines a bounded linear functional
on the cylinder span by \eqref{glim:static-pairings}, and the Riesz
representation theorem supplies its \(L^2(\Pi_g)\) representative.
The same argument gives
\[
 \|f\|_{L^2(\Pi_g)}\le\liminf_n\|f_n\|_{L^2(\mu)}.
\]
As in Section~\ref{sec:background}, the convergence is strong when
the norms also converge.  For such a sequence,
\begin{equation}
 \lim_n\|f_n-\phi_n\|_{L^2(\mu)}^2
       =\|f-\phi\|_{L^2(\Pi_g)}^2
       \quad(\phi\in\mathscr C).
 \label{glim:common-approximation}
\end{equation}
Consequently strong convergence is equivalent to approximation, in both
spaces, by the same cylinders.  In particular, the pairing of an
\(L^2\)-bounded weakly convergent sequence with a strongly convergent
sequence converges.

\paragraph{Resolvents from bounded optimizers.}
Set
\[
 \mathfrak q_n=T_n\cE,\quad \mathfrak A_n=T_nH_n,\qquad
 \mathfrak q=c_*\Pi_g[|\nabla\cdot|^2],\quad \mathfrak A=c_*H_g.
\]
For \(\phi\in\mathscr C\) and rational \(z>0\), consider
\[
 u_n=(z+\mathfrak A_n)^{-1}\phi_n,\qquad
 u=(z+\mathfrak A)^{-1}\phi.
\]
The finite Markov property gives
\(\|u_n\|_\infty\le z^{-1}\|\phi\|_\infty\).
The resolvent identity also bounds \(\|u_n\|_2\) and
\(\mathfrak q_n(u_n)\) uniformly.  Thus these particular optimizers
satisfy the polynomial supremum bound as well as the norm and energy
bounds in the relaxed-energy theorem
\cite[Theorem~6.8]{FUA}.  That theorem applies to disorder-dependent
tests under the class-uniform quenched convention recorded in
Section~\ref{sec:background}.  This proof invokes it only for the
countably specified optimizer sequences, one for each
\(\phi\in\mathscr C\) and rational \(z>0\), on their common
represented event.  Its lower bound therefore gives
\begin{equation}
 \liminf_n\mathfrak q_n(u_n)\ge\mathfrak q(U)
 \quad\text{whenever \(u_n\) has weak limit \(U\)}.
 \label{glim:optimizer-liminf}
\end{equation}

Here is the variational passage, including the norm conclusion.  Let
\[
 v_n=\sup_f\{2\langle\phi_n,f\rangle_\mu-z\|f\|_2^2-\mathfrak q_n(f)\}
     =\langle\phi_n,u_n\rangle_\mu
     =z\|u_n\|_2^2+\mathfrak q_n(u_n),
\]
and define \(v\) by the analogous supremum in \(L^2(\Pi_g)\).  Take an
optimizer subsequence attaining the upper limit and then a weak limit
\(U\).  The norm lower bound and
\eqref{glim:optimizer-liminf} imply
\[
 \limsup_n v_n
 \le 2\langle\phi,U\rangle_{\Pi_g}
       -z\|U\|_2^2-\mathfrak q(U)
 \le v.
\]
For the reverse inequality, take \(F\in\mathscr C\).
The recovery part of \cite[Theorem~6.8]{FUA} supplies, for each fixed
recovery index \(j\), polynomially bounded finite functions strongly
converging to \(F\), whose energies converge to \(\mathfrak q(F)\) in
the order \(n\to\infty\), then \(j\to\infty\).  Testing the variational
supremum with these functions in exactly that order gives
\[
 \liminf_n v_n\ge
 2\langle\phi,F\rangle_{\Pi_g}-z\|F\|_2^2-\mathfrak q(F).
\]
The form density of \(\mathscr C\) allows the supremum over \(F\),
proving \(v_n\to v\).
Only the countably many sources in \(\mathscr C\) and rational \(z\)
are used on the common source event.
The limiting variational problem is strictly concave, so every weak
optimizer limit is its unique optimizer \(u\).  Moreover,
\[
 \limsup_n z\|u_n\|_2^2
 \le v-\mathfrak q(u)=z\|u\|_2^2.
\]
Together with norm lower semicontinuity this proves strong convergence
of \(u_n\) to \(u\).  This is the bounded-source version of the
variational argument in \cite[proof of Lemma~7.2]{FUA}.

The bound \(\|(z+\mathfrak A_n)^{-1}\|_{2\to2}\le z^{-1}\) and
\eqref{glim:common-approximation} now extend this resolvent convergence
to every strongly converging input.  No additional use of the lower
energy bound is needed for that extension.  To pass to semigroups, put
\(R_n=(1+\mathfrak A_n)^{-1}\), \(R=(1+\mathfrak A)^{-1}\), and, for
\(t>0\), define on \([0,1]\)
\[
 F_t(q)=
 \begin{cases}
 \exp[-t(q^{-1}-1)],&q>0,\\
 0,&q=0.
 \end{cases}
\]
This function is continuous.  Polynomial approximation of \(F_t\) and
convergence of every power of \(R_n\) on strongly converging inputs
prove
\begin{equation}
 P_{t_nT_n}f_n\longrightarrow \Sg_t f
 \quad\text{strongly whenever \(f_n\to f\) strongly and \(t_n\to t>0\)}.
 \label{glim:semigroup-strong}
\end{equation}
For the assertion with varying times, use
\[
 \sup_{\lambda\ge0}|e^{-t\lambda}-e^{-u\lambda}|
 \le\frac{|t-u|}{e\delta},\qquad t,u\ge\delta>0.
\]
This proof identifies the limit of \(P_{tT_n}\) as
\(\exp(-t c_*H_g)\) with the stated \(T_n=n^{2/3}\).

\paragraph{Densities and the limiting kernel.}
On the represented sequence, Theorem~\ref{thm:warming}, with fixed
admissible auxiliary parameters, gives a finite
random \(C\) and numbers \(\delta_n\downarrow0\) such that
\(\|P_{sT_n}\|_{1\to\infty}\le \mathcal K(s)\) for all \(s\ge\delta_n\).
Since \(w_{n,0}\) is a nonnegative density and the semigroup is
self-adjoint,
\begin{equation}
 0\le w_{n,s}\le \mathcal K(s),\qquad
 \mu[w_{n,s}]=1,\qquad \|w_{n,s}\|_2^2\le \mathcal K(s)
 \quad(s\ge\delta_n).
 \label{glim:finite-density}
\end{equation}
Include the pairings of these densities with
\(\mathscr V\) at each positive rational \(s\) in the
representation.  Every
limit functional \(L_s\) obeys
\[
 |L_s(\phi)|\le \mathcal K(s)\Pi_g[|\phi|],\qquad
 L_s(\phi)\ge-\mathcal K(s)\Pi_g[\phi^-],\qquad L_s(1)=1.
\]
Here \(\phi^-=\max\{-\phi,0\}\).  The bound follows from
\eqref{glim:finite-density} and static convergence also for the
continuous cylinders \(|\phi|\) and \(\phi^-\).
It extends \(L_s\) to \(L^1(\Pi_g)\).  The second inequality proves
positivity of the extension: if cylinders \(\phi_j\) converge in
\(L^1\) to \(f\ge0\), then
\(\Pi_g[\phi_j^-]\le\|\phi_j-f\|_1\to0\).
The extension is therefore represented by a density
\(0\le h_s\le \mathcal K(s)\) of mass one.  In particular
\(\|h_s\|_2^2\le \mathcal K(s)\).  This also justifies normalization if only
compact cylinders were initially retained: their cutoffs tending to
one have complements bounded by \(\mathcal K(s)\) times the corresponding
static probabilities.  Measurability can be seen concretely.  Take nested
finite cylinder partitions which generate the product Borel sets and
whose cell boundaries have zero \(\Pi_g\) measure; rational coordinate
boxes have this property by the cylinder densities of
\cite[Proposition~2.3]{FUA}.  The value \(L_s(\ind_A)\) for each cell is
a measurable limit of its values on fixed smooth approximations to
\(\ind_A\), which can in turn be approximated by \(\mathscr C\).
On each cell \(A\) put the ratio
\(L_s(\ind_A)/\Pi_g(A)\), with value zero if the denominator is zero.
These ratios are measurable, bounded by \(\mathcal K(s)\), and are the
conditional expectations of \(h_s\) on the nested partitions.  Their
bounded martingale limit gives a measurable representative of \(h_s\).

For positive rational \(s,t\), reversibility and
\eqref{glim:semigroup-strong} imply for every \(\phi\in\mathscr C\)
\[
 \Pi_g[h_{s+t}\phi]
 =\lim_n\mu[w_{n,s}P_{tT_n}\phi_n]
 =\Pi_g[h_s\Sg_t\phi].
\]
We used the weak--strong pairing rule; the \(w_{n,s}\) are weakly
convergent and bounded in \(L^2\).  Hence \(h_{s+t}=\Sg_t h_s\).
The closed cylinder-gradient form is a Dirichlet form: smooth normal
contractions obey its contraction inequality, and closure preserves it.
Constants belong to the form and have zero energy, by
\cite[Proposition~5.1]{FUA}.  Thus \(\Sg_t\) is Markov and preserves one.
For a real \(s>0\), choose a rational \(q<s\) and set
\(h_s=\Sg_{s-q}h_q\).  The rational relation makes this definition
independent of \(q\).  Strong continuity of the semigroup proves
\(L^2\) continuity in \(s>0\), as well as the relation for all \(s,t>0\).
Positivity and mass one persist.  Taking rationals \(q\uparrow s\) in
the density bound, using \(L^2\) convergence and an almost-everywhere
subsequence, proves \(h_s\le \mathcal K(s)\).  For deterministic dyadic rationals
\(q_j(s)\to s\), every cylinder integral at \(s\) is the limit of its
measurable values at \(q_j(s)\).  It is therefore jointly measurable
in the environment and \(s\).  The same partition construction
therefore gives jointly measurable representatives for the real-time
family.

The finite convergence against cylinders also holds along every
sequence \(s_n\to s>0\).  Indeed, for \(s\ge\delta>0\), the identity
\(w_{n,s}=P_{(s-\delta/2)T_n}w_{n,\delta/2}\) and spectral calculus give
\begin{equation}
 \|\partial_s w_{n,s}\|_2
 \le \frac{2}{e\delta}\sqrt{\mathcal K(\delta/2)}
 \quad(n\ \text{sufficiently large}).
 \label{glim:time-equicontinuity}
\end{equation}
Approximation of \(s\) by a rational time, followed by the weak
convergence there and the \(L^2\) continuity of \(h_s\), proves the
claim.

For completeness we obtain the kernel needed to sum coordinates.
For nonnegative smooth compact cylinders \(\phi,\psi\), the finite kernel
bound yields
\[
 0\le\mu[\phi_nP_{tT_n}\psi_n]
       \le \mathcal K(t)\mu[\phi_n]\mu[\psi_n].
\]
Passing to the limit using \eqref{glim:semigroup-strong}, and then
approximating indicators in \(L^2(\Pi_g)\), gives the same inequality
with \(\Pi_g[\phi\Sg_t\psi]\) in place of the left side.  Define on
rectangles the nonnegative functional
\[
 \ind_{A\times B}\longmapsto
       \Pi_g[\ind_A\Sg_t\ind_B].
\]
For a simple function on rectangles use common disjoint partitions
of the two factors.  The preceding rectangle inequality then bounds
the functional on every nonnegative such function by \(\mathcal K(t)\) times
its \(L^1(\Pi_g\otimes\Pi_g)\) norm.  It extends to a positive bounded
functional on that \(L^1\) space, so has a density \(k_t^g\) with
\(0\le k_t^g\le \mathcal K(t)\).  Self-adjointness makes the density symmetric.
Its marginals are \(\Pi_g\), since \(\Sg_t1=1\).  This proves
\eqref{glim:limit-kernel}; it does not require a joint pointwise version
of the kernels for all \(t\).

\paragraph{The tail of two equilibrium copies.}
Let \(x_a'=x_a(\sigma')\), and set
\[
 Q_{n,m}(\sigma,\sigma')=\sum_{a>m}x_a x_a',\qquad
 V_{n,m}=(\mu\otimes\mu)[Q_{n,m}^2].
\]
To see why this equilibrium quantity controls the quench trace, define
the relative transition kernel and the two-time density
\[
 p_{n,t}(\sigma,\sigma')
   =\frac{P_{tT_n}(\sigma,\sigma')}{\mu(\sigma')},\qquad
 d_{n,s,t}(\sigma,\sigma')
   =w_{n,s}(\sigma)p_{n,t}(\sigma,\sigma').
\]
The latter is the density of the two spins at times \(sT_n\) and
\((s+t)T_n\) relative to \(\mu\otimes\mu\).  On every compact subset
of \(s,t>0\), the warming bounds give a uniform bound on this density
for all sufficiently large \(n\) on the represented sequence.
Hence
\[
 \left|\sum_{a>m}\mu[w_{n,s}x_aP_{tT_n}x_a]\right|
 =\left|(\mu\otimes\mu)[d_{n,s,t}Q_{n,m}]\right|
 \le\|d_{n,s,t}\|_{L^2(\mu\otimes\mu)}V_{n,m}^{1/2}.
\]
The stationary trace has the same bound with \(d_{n,s,t}\) replaced
by \(p_{n,t}\).  Thus it suffices to control \(V_{n,m}\) uniformly
as \(m\) increases.
We claim that the representation may be further extracted so that
\begin{equation}
 \lim_{m\to\infty}\limsup_n V_{n,m}=0
 \quad\text{almost surely}.
 \label{glim:static-tail}
\end{equation}
Recall \(r_n=n^{-1/3+10^{-4}}\) from
Section~\ref{sec:edge-estimates}, and fix the spectrum in one of the
tight spectral sets of \cite[Proposition~2.3]{FUA}.  In this paragraph
take its canonical parameters
\[
 b_n^\circ=1+(-g_{1,n})_+,\qquad
 b^\circ=1+(-g_1)_+,\qquad
 \ell_{a,n}^\circ=(g_{a,n}+b_n^\circ)^{-1}.
\]
Let \(\mu^{\mathrm{sp}}\) be the spherical Gibbs law, and
\(L_\circ=Z_{\mathrm{cube}}/Z_{\mathrm{sp}}\), with both partition
functions relative to their normalized uniform measures.  Average
over the full Haar eigenframe at this fixed spectrum.  Multiplication
by \(L_\circ^2\) cancels the two cube Gibbs denominators.  In the
uncapped replica comparison of \cite[Lemma~3.2]{FUA}, the normalized
contribution of pairs with
\(|\rho|=|n^{-1}\sigma\cdot\sigma'|>C r_n\) is at most
\(C e^{-c n r_n^3}\).  Since \(|Q_{n,m}|^2\le n^{2/3}\), inserting
this square in the discarded part costs at most
\(C n^{2/3}e^{-c n r_n^3}\).

On the retained overlaps, the same comparison can be used with this
square inserted.  Here are the required uniform details.  Conditional
on \(\rho\), a Haar-rotated cube pair and a spherical pair have the
same law of an orthonormal sum-and-difference frame \((q,q')\).
Their spectral vectors can be written
\[
 Y(\rho)=\sqrt{\frac n2}
       \big(\sqrt{1+\rho}\,q+\sqrt{1-\rho}\,q'\big),\qquad
 Y'(\rho)=\sqrt{\frac n2}
       \big(\sqrt{1+\rho}\,q-\sqrt{1-\rho}\,q'\big).
\]
The ratio of the discrete overlap mass and the spherical mass on a
cell of width \(2/n\) is at most \(\exp(C+C n\rho^4)\), by the
replica comparison.  This is bounded on \(|\rho|\le C r_n\), because
\(n r_n^4=o(1)\).  The tilt exponent on this frame is
\[
 \frac n2\big((1+\rho)q^{\mathsf T}\Lambda q
                    +(1-\rho){q'}^{\mathsf T}\Lambda q'\big).
\]
Its oscillation in a cell is bounded when \(\|\Lambda\|\) is bounded.
If \(\rho,\widetilde\rho\) lie in that cell, the displayed vectors
change in Euclidean norm by at most \(C n^{-1/2}\).  After multiplying
by \(n^{-1/3}\), each vector changes by at most \(C n^{-5/6}\) and
has norm \(n^{1/6}\).  Projection onto the coordinates \(a>m\)
therefore gives, uniformly in \(m\),
\[
 |Q_{n,m}(\rho)-Q_{n,m}(\widetilde\rho)|\le C n^{-2/3}.
\]
The inequality \((u+v)^2\le2u^2+2v^2\) now bounds the inserted square
by the corresponding spherical square and \(C n^{-4/3}\), up to
the preceding bounded comparison factors.

Under two independent spherical Gibbs laws, sign symmetry kills the
cross terms in that square.  The uniform coordinate moment bound in
\cite[Proposition~2.3]{FUA} gives
\[
 (\mu^{\mathrm{sp}}\otimes\mu^{\mathrm{sp}})[Q_{n,m}^2]
 =\sum_{a>m}\big(\mu^{\mathrm{sp}}[x_a^2]\big)^2
 \le C\sum_{a>m}(\ell_{a,n}^\circ)^2.
\]
Since \(\ind_{\{L_\circ\ge1/2\}}\le4L_\circ^2\), the comparison just
proved yields
\begin{equation}
 \E_U[\ind_{\{L_\circ\ge1/2\}}V_{n,m}]
 \le C\sum_{a>m}(\ell_{a,n}^\circ)^2
       +C n^{-4/3}+C n^{2/3}e^{-c n r_n^3}.
 \label{glim:tail-haar}
\end{equation}
The constants are uniform on the retained spectral set.  The squared
resolvent tails tend to zero as \(m\to\infty\) there, and \(n r_n^3\)
is a positive power of \(n\).  Also
\(\Prob_U\{L_\circ<1/2\}=o(1)\), by
\cite[Lemma~2.5]{FUA}.  Markov's inequality, followed by exhaustion
of the tight spectral sets, proves for every \(\epsilon>0\)
\begin{equation}
 \lim_{m\to\infty}\limsup_n
       \Prob\{V_{n,m}>\epsilon\}=0.
 \label{glim:tail-probability}
\end{equation}
The same estimates make \(V_{n,m}\) tight for every fixed \(m\).
Include these countably many quantities in the representation and
extract coordinatewise limits \(V_m\).  The identity
\[
 V_{n,m}=\sum_{a,b>m}\mu[x_ax_b]^2
\]
shows that \(V_{n,m}\), and hence \(V_m\), decreases with \(m\).
Equation~\eqref{glim:tail-probability} forces \(V_m\downarrow0\)
almost surely.  Since \(V_{n,m}\to V_m\) for each fixed \(m\), this is
\eqref{glim:static-tail}.

The limiting version follows directly from the same sign symmetry
and moments.  With \(\ell_a^\circ=(g_a+b^\circ)^{-1}\), for \(M>m\),
\begin{equation}
 \left\|\sum_{m<a\le M}x_a x_a'\right\|_{L^2(\Pi_g\otimes\Pi_g)}^2
 =\sum_{m<a\le M}\Pi_g[x_a^2]^2
 \le C\sum_{a>m}(\ell_a^\circ)^2\longrightarrow0.
 \label{glim:limit-static-tail}
\end{equation}
Thus these interval sums have an \(L^2(\Pi_g\otimes\Pi_g)\) tail
whose norm tends to zero.

\paragraph{Correlation limits and relaxation.}
Orthogonality of the eigenframe gives the exact finite identities
\[
 A_{n,J}(t)=\sum_{a=1}^n\mu[x_aP_{tT_n}x_a],\qquad
 B_{n,J}(s,t)=\sum_{a=1}^n\mu[w_{n,s}x_aP_{tT_n}x_a].
\]
The preceding two-copy reduction and \eqref{glim:static-tail} give
uniform finite correlation tails on every compact positive time set.
For the limiting series, the corresponding two-time density is
\(h_s(x)k_t^g(x,x')\) relative to \(\Pi_g\otimes\Pi_g\);
for the stationary series it is \(k_t^g(x,x')\).
Equations~\eqref{glim:density-family} and~\eqref{glim:limit-kernel}
bound these densities uniformly on compact positive time sets.
The same Cauchy--Schwarz argument with
\eqref{glim:limit-static-tail} therefore gives the limiting uniform
tail control.

It remains to pass each fixed summand.  Static convergence of the mixed
pairings \(x_a\phi\) and the second moment \(x_a^2\) makes the finite
coordinate \(x_a\) strongly
convergent to the limiting \(x_a\).  If \(s_n\to s>0\), then
\(w_{n,s_n}x_a\) is weakly convergent to \(h_sx_a\).  To see this,
test first against a cylinder times a smooth cutoff of \(x_a\) at
height \(R\), and use the convergence after
\eqref{glim:time-equicontinuity}, extended from \(\mathscr C\) by the
common-support approximation above.  For a bounded cylinder \(\phi\),
the omitted part of the pairing is at most
\[
 \mathcal K(\delta)\|\phi\|_\infty
       \mu[|x_a|\ind_{\{|x_a|>R\}}]
 \le \frac{\mathcal K(\delta)\|\phi\|_\infty}{R}\mu[x_a^2]
 \quad(s_n\ge\delta).
\]
The fixed second moments are bounded and converge, so this error
vanishes uniformly as \(R\to\infty\).  Also
\(\|w_{n,s_n}x_a\|_2\le \mathcal K(\delta)\|x_a\|_2\).
The weak--strong pairing rule and
\eqref{glim:semigroup-strong} now give, whenever \(s_n\to s>0\) and
\(t_n\to t>0\),
\[
 \mu[w_{n,s_n}x_aP_{t_nT_n}x_a]
       \longrightarrow \Pi_g[h_sx_a\Sg_t x_a].
\]
The same argument without \(w_{n,s_n}\) gives the fixed stationary
summand.  Each limiting summand is continuous.  For \(t\ge\delta>0\), the kernel
bound gives
\[
 \|\Sg_t x_a\|_\infty\le \mathcal K(\delta)\Pi_g[|x_a|],\qquad
 \|x_a\Sg_t x_a\|_2\le \mathcal K(\delta)\|x_a\|_2^2.
\]
Thus \(L^2\) continuity of \(h_s\) gives continuity in \(s\),
uniformly for \(t\) in positive compact intervals.  Strong
continuity of \(\Sg_t x_a\), paired with \(h_sx_a\in L^2(\Pi_g)\),
gives continuity in \(t\).  The uniform tail bound proves the
asserted continuity and local uniform convergence of the full series.
It also proves the two locally uniform finite limits in
\eqref{glim:functional-subsequence}: failure of uniform convergence
on a compact set would provide a sequence of time points with a
convergent subsequence, contradicting the fixed-summand convergence,
the uniform tails, and continuity of the limiting sum.

Finally, the kernel of \(H_g\) consists exactly of the constants,
by \cite[Proposition~5.1]{FUA}.  The spectral theorem and
\(h_s=\Sg_{s-s_0}h_{s_0}\) give \(h_s\to\Pi_g[h_{s_0}]=1\) in
\(L^2(\Pi_g)\).  For each fixed \(a\), the last displayed bound
implies, uniformly for \(t\ge\delta\),
\[
 |\Pi_g[(h_s-1)x_a\Sg_t x_a]|
 \le \mathcal K(\delta)\|x_a\|_2^2\|h_s-1\|_2\longrightarrow0.
\]
The correlation tails are uniform also for \(s\ge s_0\), since
\(\|h_s\|_\infty\le \mathcal K(s_0)\).  Passing first through a finite sum and
then letting its length tend to infinity proves
\eqref{glim:relaxation}.
\end{proof}

The preceding proposition does not yet identify the family as a
function of \(g\).  The following estimate retains the information
from the finite uniform initialization that will make this
identification possible.

\begin{proposition}[An escape estimate from the uniform start]
\label{prop:uniform-seed-escape}
There is a measurable admissible choice \(b=b(g)\) such that, almost
surely,
\begin{equation}
 r_a:=g_a+b>0,\qquad D_g(b)>0,\qquad
 c_g(1+a)^{2/3}\le r_a\le C_g(1+a)^{2/3}\quad(a\ge1)
 \label{glim:rates}
\end{equation}
for finite positive \(c_g,C_g\).  Define
\begin{equation}
 M_g(x)=\frac12\sum_{a\ge1}\frac{x_a^2}{r_a}.
 \label{glim:weighted-size}
\end{equation}
Then \(0<M_g<\infty\) \(\Pi_g\)-almost surely.  Choose a Borel
full-measure set of edge realizations on which these assertions about
\(b\), \(r_a\), and \(M_g\) hold, with \(b\) Borel there and
\(r_a=g_a+b\).  On the product of this edge set with \(\R^\N\),
\(M_g\) is an extended nonnegative Borel function of \((g,x)\), since
its jointly Borel partial sums increase to it.  In every escape event,
\(x_i^2/M_g(x)\) denotes the quotient on the subset where
\(0<M_g(x)<\infty\), extended by zero to all remaining pairs in
\(\R^\N\times\R^\N\).  This version is jointly Borel and agrees with
the usual quotient \(\Pi_g\)-almost surely for every \(g\) in the
chosen edge set.

For every represented family in
Proposition~\ref{prop:gaussian-subsequence}, every fixed
\(i\ge1\), and every \(\epsilon>0\),
\begin{equation}
 \lim_{s\downarrow0}
 \widehat{\E}\!\left[
    \Pi_g\!\left[h_s\ind_{\{x_i^2/M_g>\epsilon\}}\right]\right]=0.
 \label{glim:seed-escape}
\end{equation}
Here the probability averages the represented environment as well as
the draw from \(h_s\Pi_g\).
\end{proposition}

\begin{proof}
We will show that \(x_i^2>s^{13/20}\) has vanishing averaged probability
under the limiting time-\(s\) law, while \(M_g<s^{14/25}\) also has
vanishing averaged probability as \(s\downarrow0\).  Away from these
two events, \(x_i^2/M_g\) is at most \(s^{9/100}\).
We obtain the numerator bound from
the finite uniform-start dynamics, then bound the denominator directly
under the conditioned edge measure.

\paragraph{The total mass in a block of finite modes.}
Fix \(\varepsilon>0\), and fix the auxiliary log-moment parameters
in Theorem~\ref{thm:warming} once and for all.  Its entropy bound and
Proposition~\ref{prop:edge-estimates} provide a quenched disorder event
\(\mathcal G_{n,\varepsilon}\) of limiting
probability at least \(1-\varepsilon\) on which their constants and
thresholds are fixed, the entropy bound holds simultaneously for
every positive time, and the uncapped cube projection tails hold at
every allowed scale.  Separately retain a spectrum-only event
\(\mathcal S_{n,\varepsilon}\) of
limiting probability at least \(1-\varepsilon\) for the band estimates
in Lemma~\ref{edge:spectral} and the tight leading-edge
bound \(|g_{1,n}|\le C_\varepsilon\).  The latter follows from the
GOE edge convergence in Section~\ref{sec:background}.  The constants in the
following calculations may depend on \(\varepsilon\).

For \(0<s<1\), put \(m=\lceil s^{-5/4}\rceil\).  For \(n\ge m\), set
\begin{equation}
 D_m(v)=\sum_{a\le m}x_a(v)^2,\qquad
 I=s^{47/100},\qquad R=s^{-53/100}.
 \label{glim:rotation-scales}
\end{equation}
Here \(x_a(v)\) is evaluated along the finite dynamics at physical
time \(vT_n\), and \(I=sR\).  Choose \(k_p=np^3=C' m\), with the
fixed \(C'\) sufficiently large.  Take \(s\) small enough that
\(C'm\) exceeds the fixed threshold, and then take \(n\) sufficiently
large that \(m\le n\) and \(p\le r_n\).  The spectral band estimates imply
that \(\dim\cH_p\ge m\) on \(\mathcal S_{n,\varepsilon}\), so \(P_p\)
contains the first \(m\) coordinates and
\(g_{m,n}\le T_np^2=(C'm)^{2/3}\).  Together with
\(|g_{1,n}|\le C_\varepsilon\), this gives on that event
\begin{equation}
 \max_{a\le m}|g_{a,n}|\le C m^{2/3}.
 \label{glim:finite-head-spectrum}
\end{equation}
At a fixed \(s\), the value
of \(k_p\) is fixed as \(n\to\infty\); this is why the extension to
the lowest scales in Proposition~\ref{prop:edge-estimates} is used here.

On \(\mathcal G_{n,\varepsilon}\cap\mathcal S_{n,\varepsilon}\) put
\(Y_p=\|P_p\sigma\|^2/M_p^2\).  Its
projection tail is
\(\mu\{Y_p>u\}\le C\exp(-a k_pu^3)\) for every \(u\ge u_0\), with
fixed \(u_0\).  Integration of this tail gives
\[
 \log\mu\exp\!\left(\tfrac a2 k_pY_p^3\right)\le C k_p.
\]
The entropy inequality for any probability density \(w\) relative
to \(\mu\), followed by Hölder's inequality, therefore gives
\[
 \mu[wY_p]\le
 C\left(1+\big(\Ent_\mu(w)/k_p\big)^{1/3}\right).
\]
Since \(D_m\le n^{-2/3}M_p^2Y_p=k_p^{1/3}Y_p\), the entropy part of
Theorem~\ref{thm:warming} yields under the direct dynamical law at
fixed disorder
\begin{equation}
 \E_{\nu_n}^J[D_m(v)]
 \le C\big(m^{1/3}+v^{-77/150}\big),\qquad 0<v\le s\le1.
 \label{glim:head-mass}
\end{equation}
This use of entropy is entirely relative to the ordinary Gibbs law
\(\mu\); the path in this expectation starts from \(\nu_n\).

Let \(Q(v)=\int_0^vD_m(w)\,dw\) and define the event
\[
 G=\{Q(s)<I,\ D_m(s)\le R\}.
\]
Integrating \eqref{glim:head-mass} and applying Markov's inequality
at the integral and the endpoint gives, on
\(\mathcal G_{n,\varepsilon}\cap\mathcal S_{n,\varepsilon}\),
\begin{equation}
 \Prob_{\nu_n}^J(G^c)
 \le C\big(s^{17/150}+s^{1/60}\big).
 \label{glim:cutoff-loss}
\end{equation}
Indeed \(m^{1/3}\le C s^{-5/12}\), and both comparisons use the
same two positive powers:
\[
 1-\frac5{12}-\frac{47}{100}
   =\frac{53}{100}-\frac5{12}=\frac{17}{150},\qquad
 1-\frac{77}{150}-\frac{47}{100}
   =\frac{53}{100}-\frac{77}{150}=\frac1{60}.
\]

\paragraph{A stopped rotation at fixed spectrum.}
For this step fix only a spectrum satisfying
\eqref{glim:finite-head-spectrum}.  Let \(\E_\Lambda\) average over
the full conditional Haar law of \(U\) and over the uniform-start
heat-bath path.  In particular, this expectation does not condition
on \(\mathcal G_{n,\varepsilon}\).

It is useful to specify the path likelihood.  In normalized time
\(v\), take independent site clocks of rate \(T_n\), a uniform
initial spin, and an independent fair mark
\(\eta\in\{-1,1\}\) at every ring, including rings at which the mark
equals the old spin.  This is a reference law; the marked path is
then determined without \(U\).  If the ring is at site \(k\), its
heat-bath probability is
\[
 p_{U,k}(\eta\mid\sigma)=
 \frac{e^{\eta b_k(\sigma)}}{2\cosh b_k(\sigma)},\qquad
 b_k(\sigma)=\sum_{\ell\ne k}W_{k\ell}\sigma_\ell.
\]
The likelihood of the dynamical path through \(s\) relative to the
reference law is
\begin{equation}
 \mathcal L_U=
 \prod_{\ell:v_\ell\le s}
 2p_{U,k_\ell}\big(\eta_\ell\mid\sigma(v_\ell-)\big).
 \label{glim:path-likelihood}
\end{equation}
There is no initial or clock factor depending on \(U\).
This likelihood fixes the marked spin outcomes while the matrix
varies; it permits differentiating the smooth heat-bath probabilities.
The absence of an initial score is the feature of the uniform start
used in this rotation.

For \(i,j\le m\), \(i\ne j\), right-rotate these two columns with
\(\dot u_i=u_j\) and \(\dot u_j=-u_i\).  Then
\[
 \dot x_i=x_j,\qquad \dot x_j=-x_i,\qquad
 \dot W=(\lambda_i-\lambda_j)
          (u_i u_j^{\mathsf T}+u_j u_i^{\mathsf T}).
\]
Writing \(\Delta_{ij}=\lambda_i-\lambda_j\), diagonal deletion in
the physical local field gives
\[
 \dot b_k=\Delta_{ij}\left[
  n^{1/3}(u_{j,k}x_i+u_{i,k}x_j)
                    -2u_{i,k}u_{j,k}\sigma_k\right].
\]
Orthogonality and
\(\sum_k u_{i,k}^2u_{j,k}^2\le1\), together with
\(|\Delta_{ij}|=T_n^{-1}|g_{i,n}-g_{j,n}|
\le C n^{-2/3}m^{2/3}\), imply
\begin{equation}
 \sum_k\dot b_k^2
 \le C n^{-4/3}m^{4/3}
       \big[n^{2/3}(x_i^2+x_j^2)+1\big].
 \label{glim:field-rotation}
\end{equation}
The final \(1\) is the contribution from deleting the diagonal.

In the filtration revealing \(U\) initially and then the marked
rings, the logarithmic derivative of \(\mathcal L_U\) through \(v\) is
\[
 S_{ij}(v)=\sum_{\ell:v_\ell\le v}
  \big(\eta_\ell-\tanh b_{k_\ell}(\sigma(v_\ell-))\big)
                 \dot b_{k_\ell}(\sigma(v_\ell-)).
\]
Under the dynamical law it is a square-integrable martingale.  In
fact each marked increment has conditional mean zero and conditional
variance \(\operatorname{sech}^2(b_k)\dot b_k^2\); its predictable
bracket is
\[
 \langle S_{ij}\rangle_v
 =T_n\int_0^v\sum_k\operatorname{sech}^2(b_k(\sigma(w-)))
                         \dot b_k(\sigma(w-))^2\,dw.
\]
Define
\[
 \tau=s\wedge\inf\{v:Q(v)\ge I\},\qquad
 I_a=\E_\Lambda\int_0^\tau x_a(v)^2\,dv\quad(a\le m).
\]
The function \(Q\) is continuous, so \(Q(\tau)\le I\).  Stopped
martingale isometry and \eqref{glim:field-rotation} give
\begin{equation}
 \E_\Lambda[S_{ij}(\tau)^2]
 \le C m^{4/3}\big(I_i+I_j+s n^{-2/3}\big).
 \label{glim:rotation-score}
\end{equation}

For a fixed marked reference history, \(D_m,Q,\tau\), and \(G\) are
exactly invariant under this rotation.  Apply Haar invariance first
to
\[
 F_{ij}=\int_0^\tau x_i(v)x_j(v)\,dv.
\]
Differentiating its integral with likelihood
\eqref{glim:path-likelihood}, and then projecting the terminal score
onto the stopped filtration, gives
\begin{equation}
 0=\E_\Lambda[\dot F_{ij}+F_{ij}S_{ij}(s)]
   =\E_\Lambda[\dot F_{ij}+F_{ij}S_{ij}(\tau)].
 \label{glim:haar-identity}
\end{equation}
Here is a justification at the stopping boundary.  The stopping
rule is constant on each rotation orbit of a fixed marked history,
so \(\dot F_{ij}\) differentiates only \(x_i x_j\).
At each fixed \(n\) and spectrum, the likelihood is bounded by
\(2^N\), where \(N\) is the number of rings, and its derivative is
bounded by this quantity times a constant (depending on \(n\) and
the spectrum) times \(N\).  Poisson exponential moments justify
differentiation of the full-time integral.  The random variable
\(F_{ij}\) is measurable at \(\tau\), and optional projection of
the square-integrable martingale replaces \(S_{ij}(s)\) by
\(S_{ij}(\tau)\).  This proves \eqref{glim:haar-identity} without
differentiating a law with a moving stopping time.

Since \(\dot x_i=x_j\) and \(\dot x_j=-x_i\), this identity reads
\[
 I_i-I_j=\E_\Lambda\!\left[
 S_{ij}(\tau)\int_0^\tau x_i(v)x_j(v)\,dv\right].
\]
Sum over \(j\le m\), \(j\ne i\).  The pathwise bound
\[
 \sum_{\substack{j\le m\\j\ne i}}\left(\int_0^\tau x_i x_j\,dv\right)^2
 \le \left(\int_0^\tau x_i^2\,dv\right)
       \left(\int_0^\tau D_m\,dv\right)
 \le I\int_0^\tau x_i^2\,dv
\]
and Cauchy--Schwarz on the product of probability and index spaces,
together with \(\sum_{j\le m} I_j\le I\) and
\eqref{glim:rotation-score}, yield
\begin{equation}
 m I_i\le I+
 \left\{C m^{4/3}
       (m I_i+I+m s n^{-2/3})\,I I_i\right\}^{1/2}.
 \label{glim:integrated-rotation}
\end{equation}
This estimate does not require independence between the scores.
To solve it, put
\[
 y=\frac{m I_i}{I},\qquad
 \alpha=C I m^{1/3},\qquad
 \kappa=\frac{m s n^{-2/3}}{I}.
\]
It gives \(y\le1+\sqrt{\alpha y(y+1+\kappa)}\).
Here \(I m^{1/3}=O(s^{4/75})\), while \(\kappa\to0\) as
\(n\to\infty\) at each fixed \(s\).  For sufficiently small \(s\)
and then sufficiently large \(n\), \(\kappa\le1\) and
\(\sqrt\alpha\le1/2\).  Using
\(\sqrt{y(y+2)}\le y+1\) proves
\begin{equation}
 I_i\le 3I/m\qquad(i\le m).
 \label{glim:individual-integral}
\end{equation}

Repeat \eqref{glim:haar-identity} with
\(F_{ij}=\ind_G x_i(s)x_j(s)\).  This random variable is measurable
at \(\tau\): it is zero if \(\tau<s\), and otherwise is measurable
at \(s=\tau\).  The indicator \(G\), including its boundary, is
constant on each rotation orbit, so its differentiation contributes
nothing.  Write \(Z_a=\E_\Lambda[\ind_G x_a(s)^2]\) for \(a\le m\).  Now
\(\sum_{a\le m} Z_a\le R\) and
\[
 \sum_{\substack{j\le m\\j\ne i}}\E_\Lambda[\ind_G x_i(s)^2x_j(s)^2]\le R Z_i.
\]
Use \eqref{glim:rotation-score} and
\eqref{glim:individual-integral} in the summed identity.  Taking \(n\)
large enough that \(m s n^{-2/3}\le I\), we obtain
\begin{equation}
 m Z_i\le R+\big(C m^{4/3} I R Z_i\big)^{1/2},
 \qquad Z_i\le C R/m.
 \label{glim:terminal-rotation}
\end{equation}
For the second inequality, \(z=mZ_i/R\) satisfies
\(z\le1+\sqrt{\alpha z}\), and hence \(z\le2+\alpha\).
As \(R/m\le C s^{18/25}\), Markov's inequality gives, at each of the
retained spectra,
\begin{equation}
 \Prob_\Lambda\{x_i(s)^2>s^{13/20},\,G\}
 \le s^{-13/20}Z_i\le C s^{7/100}.
 \label{glim:one-mode}
\end{equation}

We now combine the two probability laws.  Let \(\Prob_n\) include
the Gaussian disorder, the independent diagonal and eigenvector
signs, and the uniform-start dynamics.  Write \(\Prob_{\mathrm{dis}}\)
and \(\E_{\mathrm{dis}}\) for the completed disorder law and
\(\E_{\mathrm{eig}}\) for the eigenvalue law.  Then
\[
\begin{aligned}
 \Prob_n\{x_i(s)^2>s^{13/20}\}
 \le{}&
 \Prob_{\mathrm{dis}}\big((\mathcal G_{n,\varepsilon}
                         \cap\mathcal S_{n,\varepsilon})^c\big)\\
 &+\E_{\mathrm{dis}}\!\left[
   \ind_{\mathcal G_{n,\varepsilon}\cap\mathcal S_{n,\varepsilon}}
                   \Prob_{\nu_n}^J(G^c)\right]\\
 &+\E_{\mathrm{eig}}\!\left[
   \ind_{\mathcal S_{n,\varepsilon}}
      \Prob_\Lambda\{x_i(s)^2>s^{13/20},G\}\right].
\end{aligned}
\]
The second term uses \eqref{glim:cutoff-loss}; the last uses
\eqref{glim:one-mode} with its full conditional Haar average.
The complements of
\(\mathcal G_{n,\varepsilon}\) and \(\mathcal S_{n,\varepsilon}\)
cost at most \(2\varepsilon\) in the limit, so for
each fixed \(i\) and all sufficiently small \(s\),
\begin{equation}
 \limsup_n\Prob_n\{x_i(s)^2>s^{13/20}\}
 \le2\varepsilon+
 C_\varepsilon\big(s^{17/150}+s^{1/60}+s^{7/100}\big).
 \label{glim:finite-numerator}
\end{equation}
This separation is essential: no frame-dependent quenched event
was imposed inside the Haar identity.  Sending \(s\downarrow0\)
and then \(\varepsilon\downarrow0\) proves the finite numerator
estimate.

For a fixed positive \(s\), the cylinder convergence in
Proposition~\ref{prop:gaussian-subsequence} transfers this estimate
to \(\rho_s(dx)=h_s(x)\Pi_g(dx)\).  Indeed each coordinate marginal
of \(\Pi_g\) has a density by
\cite[Proposition~2.3]{FUA}; the same is true of \(\rho_s\).
The boundary \(x_i^2=s^{13/20}\) thus has zero limiting probability.
Weak convergence of this coordinate marginal and bounded convergence
over the represented environment yield
\begin{equation}
 \lim_{s\downarrow0}
       \widehat{\E}\big[\rho_s\{x_i^2>s^{13/20}\}\big]=0.
 \label{glim:limit-numerator}
\end{equation}
Only finite-coordinate convergence was used in this step.

\paragraph{The weighted denominator.}
Corollary~\ref{edge:limit-growth} gives the two-sided bound in
\eqref{glim:rates} for every admissible \(b\), with measurable constants
when \(b\) is measurable.  It remains to choose \(b\) with \(D_g(b)>0\).
Start with a measurable admissible \(b_0\), for example
\(b_0=1+(-g_1)_+\).  The resolvent-difference identity in
\cite[Proposition~2.3]{FUA} gives, for \(b\ge b_0\),
\[
 D_g(b)-D_g(b_0)=
 \sum_a\left(\frac1{g_a+b_0}-\frac1{g_a+b}\right).
\]
The right side increases to \(+\infty\) as \(b\to\infty\), because
\(\sum_a(g_a+b_0)^{-1}=\infty\).  Choose the first integer
increment \(b=b_0+j\) for which \(D_g(b)>0\).  This is measurable
and admissible, so Corollary~\ref{edge:limit-growth} gives
\eqref{glim:rates}.

Put \(\ell_a=r_a^{-1}\).  The uniform coordinate moment bound in
\cite[Proposition~2.3]{FUA} gives
\[
 \Pi_g[M_g]\le C_g\sum_a\ell_a^2<\infty.
\]
The first coordinate has no atom at zero, so \(M_g>0\) almost surely.
This proves that \eqref{glim:weighted-size} is well defined under
\(\Pi_g\) and under every \(\rho_s\).

Let \(N=\lceil s^{-8/5}\rceil\) and consider the coordinates
\(N\le a\le2N\).  Under the independent Gaussian law of the
conditioning construction, write \(Z_a=\sqrt{\ell_a}\xi_a\) with
independent standard normals \(\xi_a\).  Equation~\eqref{glim:rates}
implies
\[
 \frac12\sum_{a=N}^{2N}\ell_a Z_a^2
 \ge c_gN^{-4/3}\sum_{a=N}^{2N}\xi_a^2.
\]
The chi-square lower-tail bound shows that this is at least
\(c_gN^{-1/3}\), except on an event of probability at most
\(C_g e^{-c_gN}\).

The same estimate survives the norm conditioning with a fixed cost.
Leave a sufficiently large fixed head of \(m_0\) coordinates outside
this block.  Let \(q\) be the density of
\(\sum_a(Z_a^2-\ell_a)\), and \(q_{m_0}^{\rm head}\) the density of its
first \(m_0\) terms.  The tail disintegration in
\cite[Proposition~2.3]{FUA} states that
\begin{equation}
 \frac{d(\Pi_g)_{>m_0}}{d\bigotimes_{a>m_0}N(0,\ell_a)}
 =
 \frac{q_{m_0}^{\rm head}\!\left(
 D_g(b)-\sum_{a>m_0}(x_a^2-\ell_a)\right)}{q(D_g(b))}.
 \label{glim:conditioned-tail-density}
\end{equation}
The head density is bounded by Fourier inversion once \(m_0\) is
large enough, and \(q(D_g(b))>0\).  Hence the ratio in
\eqref{glim:conditioned-tail-density} is bounded by a finite
\(g\)-dependent constant.  For all sufficiently small \(s\), the
block lies outside the head and
\[
 c_g N^{-1/3}\ge s^{14/25},
 \qquad\text{since}\qquad
 \frac{14}{25}-\frac8{15}=\frac2{75}>0.
\]
It follows that
\begin{equation}
 \Pi_g\{M_g<s^{14/25}\}\le C_g e^{-c_g s^{-8/5}}.
 \label{glim:denominator-static}
\end{equation}
Using \eqref{glim:density-family}, for almost every represented
environment we obtain
\[
 \rho_s\{M_g<s^{14/25}\}
 \le C_g C\exp\!\left(C s^{-77/50}-c_gs^{-8/5}\right)
 \longrightarrow0,
\]
because \(8/5>77/50\).  These probabilities are bounded by one, so
the convergence also holds after \(\widehat{\E}\).

On the complement of the two exceptional events in
\eqref{glim:limit-numerator} and \eqref{glim:denominator-static},
\[
 \frac{x_i^2}{M_g(x)}
 \le s^{13/20-14/25}=s^{9/100}.
\]
This tends to zero and proves \eqref{glim:seed-escape}.
\end{proof}

For each represented limiting law, the averaged limit in
\eqref{glim:seed-escape} yields a deterministic rational sequence
\(s_k\downarrow0\) along which, for almost every represented
environment, the displayed conditional probabilities tend to zero
for every fixed \(i\) and \(\epsilon>0\).
Section~\ref{sec:entrance-selection} proves this extraction and uses
the resulting criterion.

\section{Selection of the entrance law}\label{sec:entrance-selection}

Proposition~\ref{prop:gaussian-subsequence} gives positive-time density
families for the edge semigroup, but an extracted family might still depend
on randomness beyond the edge points.  The uniform-start estimate of
Proposition~\ref{prop:uniform-seed-escape} is averaged over the represented
environment.  We first prove fixed-\(g\) uniqueness from a sequential
escape criterion and derive the zero-coordinate boundary from that
criterion.  We then extract a qualifying sequence for almost every
represented environment and obtain a family measurable in \(g\).

Fix \(g\) in a common Borel full-measure subset of the edge set chosen
in Proposition~\ref{prop:uniform-seed-escape} on which the conclusions
about the conditioned law and the closed gradient form in Items 1 and 3
of Proposition~\ref{prop:stationary-inputs} hold with their stated
source versions.  The rate, finiteness, and positivity conclusions of
Proposition~\ref{prop:uniform-seed-escape} hold on this set as well.
In this section write
\[
 r_i=g_i+b(g),\qquad D=D_g(b(g)),\qquad \Pi=\Pi_g,\qquad H=H_g.
\]
Thus \(D>0\), and there are constants \(0<c_g<C_g<\infty\) such that
\begin{equation}\label{entr:rates}
 c_g(1+i)^{2/3}\le r_i\le C_g(1+i)^{2/3}\qquad(i\ge1).
\end{equation}
Constants with a subscript \(g\) below may be enlarged or decreased
and depend only on this fixed realization, unless other dependencies
are displayed.
Recall the finite, strictly positive random variable
\begin{equation}\label{entr:mass}
 M(x)=\sum_{i\ge1}\frac{x_i^2}{2r_i},\qquad x\sim\Pi.
\end{equation}
Its finiteness follows from the uniform coordinate moment bound
\(\Pi[x_i^2]\le C_g/r_i\) and \(\sum_i r_i^{-2}<\infty\); its positivity
follows from the absolute continuity of a one-coordinate marginal.
These properties of \(\Pi\) are part of the conditioned edge law
\cite[Proposition~2.3]{FUA}.

\begin{theorem}[Entrance selection]\label{thm:entrance-selection}
For every \(g\) as above, there is at most one family
\((h_s)_{s>0}\) of probability densities with respect to \(\Pi_g\) with
the following properties:
\begin{enumerate}
 \item \(h_s\ge0\), \(\Pi_g[h_s]=1\), and
 \(\sup_{s_0\le s\le s_1}\|h_s\|_{L^\infty(\Pi_g)}<\infty\) whenever
 \(0<s_0<s_1<\infty\);
 \item \(h_{s+t}=S_t^g h_s\) in \(L^2(\Pi_g)\) for \(s,t>0\), where
 \(S_t^g=\exp(-t c_*H_g)\);
 \item there is a sequence \(s_k\downarrow0\) such that, for every
 \(i\ge1\) and \(\varepsilon>0\),
 \begin{equation}\label{entr:escape}
 \Pi_g\left[h_{s_k}\,
       \ind_{\{x_i^2/M(x)>\varepsilon\}}\right]\longrightarrow0.
 \end{equation}
\end{enumerate}
Any such family satisfies
\begin{equation}\label{entr:zero-boundary}
 h_s\Pi_g\ \weakto\ \delta_{\boldsymbol 0}
 \quad\text{on \(\R^{\N}\) with its product topology, as \(s\downarrow0\)}.
\end{equation}
For almost every \(g\) such a family exists.
It can be chosen measurably in \(g\), as a family of \(L^2(\Pi_g)\)
density classes, and with jointly measurable density representatives in
\((g,s,x)\).  Every represented Gaussian limit in
Proposition~\ref{prop:gaussian-subsequence} is this family almost surely.
\end{theorem}

The hypothesis \eqref{entr:escape} describes information retained from
the finite uniform initialization.  Different candidate families may
use different sequences in Item 3.  The zero-coordinate limit
\eqref{entr:zero-boundary} is a consequence of that hypothesis and
holds as \(s\) decreases to zero through all positive times.

We prove uniqueness in the time variable
\begin{equation}\label{entr:unit-time}
 u=c_*s,\qquad \widetilde h_u=h_{u/c_*},\qquad
 \widetilde S_u=e^{-uH}.
\end{equation}
Then \(\widetilde h_{u+v}=\widetilde S_v\widetilde h_u\).
The sequence \(u_k=c_*s_k\) has exactly the escape property
\eqref{entr:escape}.  All elapsed times in the next three subsections
refer to \(\widetilde S\).  We return to \(s=u/c_*\) at the end.

\subsection{Angular coordinates and a common diffusion factor}

We need a stochastic realization only after a positive time, where the
initial law has a bounded density.  Put \(m=12\) and define
\[
 R(x)=\left(\sum_{j\le m}x_j^2\right)^{1/2}.
\]
On \(\{R>0\}\), set
\[
 \theta_i(x)=\frac{x_i}{R(x)},\qquad
 \overline r(x)=\sum_{j\le m}r_j\theta_j(x)^2 .
\]
The functions \(\theta_i\) are homogeneous of degree zero: if a finite
set \(J\) contains the head \(\{1,\ldots,m\}\) and \(i\), then
\(\sum_{j\in J}x_j\partial_j\theta_i=0\).
In the finite-coordinate density of \(\Pi\), the factor imposing the
square constraint has a radial derivative.  Its contribution to the
generator therefore vanishes on \(\theta_i\).  The next lemma justifies
this cancellation for the associated diffusion, including avoidance
of \(R=0\).

\begin{lemma}[Angular realization]\label{entr:angular-realization}
Let \(w\) be a bounded probability density with respect to \(\Pi\).
There is a continuous process \(X(t)=(X_i(t))_{i\ge1}\) on
\(\R^\N\), with law \((\widetilde S_t w)\Pi\) at time \(t\), and
independent standard Brownian motions \((W_i)_{i\ge1}\) in its filtration,
with the following property.  Almost surely \(R(X(t))>0\) for all
\(t\ge0\).  With \(R(t)=R(X(t))\),
\(\Theta_i(t)=\theta_i(X(t))\), and
\(\overline r(t)=\overline r(X(t))\), one has
\begin{align}
 d\Theta_i(t)
 &=
 \frac{\sqrt2}{R(t)}
 \left(dW_i(t)-\Theta_i(t)\sum_{j\le m}\Theta_j(t)\,dW_j(t)\right)
 \notag\\
 &\quad+
 \Theta_i(t)\left(
  \frac{3-m-2\ind_{\{i\le m\}}}{R(t)^2}-r_i+\overline r(t)
 \right)dt .                                      \label{entr:angular-sde}
\end{align}
All identities hold simultaneously for \(i\ge1\).  The Brownian
increments after zero are independent of the initial sigma field.
\end{lemma}

\begin{proof}
The closed form of \(H\) is
\(\cE_g(F,G)=\Pi[\nabla F\cdot\nabla G]\).  Its domain is the closure
of the smooth compact cylinder functions, contains the constants and all
coordinates, and satisfies \(\nabla x_i=e_i\)
\cite[Proposition~5.1]{FUA}.
Smooth contractions on the core extend by closure, so this is a
Dirichlet form.  The same approximation, with smooth clipping to preserve
a common bound, gives the chain and bounded product rules.  In particular
the gradient vanishes almost everywhere on a set on which the function is
constant: apply the chain rule to smooth truncations in successively
smaller neighborhoods of the constant, and use closedness.  These rules
show strong locality.

Here are the topological hypotheses for the diffusion theorem.  The
state space \(\R^\N\) is Polish in the product topology.  The smooth
compact cylinder core consists of continuous functions and contains
a countable family separating its points.  To obtain a compact nest,
for each coordinate choose a smooth
function \(v_j\), equal to one outside a sufficiently large coordinate
interval and zero on a smaller interval, with
\(\|v_j\|_{\cE_g,1}\le\epsilon2^{-j}\).  Such a choice is possible:
coordinate tails have vanishing \(\Pi\)-mass, while a transition of
increasing width makes the gradient norm vanish.  The sum
\(\sum_jv_j\) converges in the form norm and is at least one on the
complement of a compact product box.  The definition of \(1\)-capacity
therefore bounds that complement's capacity by \(\epsilon^2\).
Taking \(\epsilon\downarrow0\) and enlarging successive boxes gives a
compact nest.  The form is thus quasi-regular.  Since \(1\) has zero
energy, its semigroup fixes \(1\); the associated process has infinite
lifetime under \(\Pi\), and hence under \(w\Pi\).

The diffusion theorem and locality criterion for quasi-regular
symmetric Dirichlet forms now give a continuous associated process
\cite[Theorems~IV.3.5 and V.1.11]{MaRockner}.  Its Fukushima
decomposition is supplied by \cite[Theorem~VI.2.5]{MaRockner}.
In particular, the coordinate martingales \(M^i\) are continuous and
have brackets
\[
 \langle M^i,M^j\rangle_t
   =2\int_0^t \nabla x_i\cdot\nabla x_j(X(v))\,dv
   =2\delta_{ij}t.
\]
The bracket identity is the energy-measure identity for the form,
given by \cite[Theorem~5.2.3 and Equation~(5.2.34)]{FOT} after the
quasi-regular regularization of \cite[Chapter~VI]{MaRockner}.
To apply its bounded formula, first take a bounded smooth clip
\(u_N=c_N(x_i)\).  For every bounded form-domain test \(f\), the chain
and product rules give
\[
 2\cE_g(u_Nf,u_N)-\cE_g(u_N^2,f)
       =2\Pi[f|c_N'(x_i)|^2].
\]
Choose the clips with \(|c_N'|\le1\) and \(c_N(x_i)\to x_i\) in form
norm.  Passing to the limit in the energy measures and polarizing gives
the displayed coordinate brackets.
The multidimensional Lévy characterization, for every finite set of
indices, gives \(M^i=\sqrt2W_i\) with the asserted countable family of
Brownian motions.  Brownian motion here is relative to the filtration of
the process, so its future increments are independent of the initial
sigma field.  The usual form exceptional set has \(\Pi\)-measure zero;
it can be discarded for \(w\Pi\).  These facts also give continuous
coordinate paths.

We verify the generator used for \(\theta_i\).  The finite-coordinate
density formula for \(\Pi\), on any finite set \(J\), is
\begin{equation}\label{entr:cylinder-density}
 \frac{q_{J^c}\!\left(D-\sum_{j\in J}(x_j^2-r_j^{-1})\right)}
      {q(D)}
 \prod_{j\in J}\left(\frac{r_j}{2\pi}\right)^{1/2}
                     e^{-r_jx_j^2/2}.
\end{equation}
Here \(q\) is the density of the full centered Gaussian square sum,
and \(q_{J^c}\) is the density after the coordinates in \(J\) are
deleted
\cite[Proposition~2.3]{FUA}.
Both densities exist, \(q(D)>0\), and \(q_{J^c}\) has bounded
derivatives of any prescribed finite order.  For the last assertion,
leave sufficiently many fixed indices outside \(J\); the product of
their characteristic-function moduli is bounded by
\(\prod (1+4t^2r_j^{-2})^{-1/4}\), which is integrable even after
multiplication by the prescribed power of \(|t|\).  Fourier inversion
then gives the assertion.  In particular \eqref{entr:cylinder-density}
is bounded by a constant times the product Gaussian density.

The set \(R=0\) is polar.  Indeed a smooth cutoff equal to one on
\(\{R\le\epsilon\}\), zero on \(\{R\ge2\epsilon\}\), and with derivative
\(O(\epsilon^{-1})\), has squared form norm
\(O(\epsilon^m+\epsilon^{m-2})\), by \eqref{entr:cylinder-density}.
Letting \(\epsilon\downarrow0\) proves zero capacity.  A diffusion
starting outside its exceptional set never hits a polar set, which
proves the stated pathwise positivity of \(R\).

For a cylinder test whose index set \(J\) contains
\(\{1,\ldots,m,i\}\), integrate by parts in
\eqref{entr:cylinder-density}.  The derivative of the factor \(q_{J^c}\)
is radial in the \(J\) coordinates.  Its pairing with the gradient of
\(\theta_i\) vanishes, because
\(\sum_{j\in J}x_j\partial_j\theta_i=0\).  Consequently the
nonpositive generator on this function is
\[
 \cL\theta_i=\Delta\theta_i-\sum_{j\in J}r_jx_j\partial_j\theta_i.
\]
This integration by parts is legitimate across \(R=0\).  More
explicitly, cut off at \(R=\epsilon\) and at large finite-coordinate
radius.  The function and its first two derivatives near zero are
bounded by
\(C_i(1+|x_i|)(1+R^{-3})\).  The products with derivatives of the
scale-\(\epsilon\) cutoff satisfy the same bound on \(R\asymp\epsilon\).
Their squared integrals near zero vanish for \(m=12>6\).
Terms with a derivative on \(q_{J^c}\) are bounded by the same Gaussian
integrals with polynomial factors, since its first derivative is
bounded.  Gaussian tails handle the outer cutoff.  Thus the cutoffs
converge in the form norm, proving \(\theta_i\) belongs to the form
domain, and the weak generator pairing passes to the limit.
The displayed \(\cL\theta_i\) is in \(L^2(\Pi)\) by the same bounds.
Extending the pairing from the core shows
\(\theta_i\in\operatorname{Dom}(H)\) and \(H\theta_i=-\cL\theta_i\).

A direct differentiation gives, with the head vector extended by zeros,
\[
 \partial_j\theta_i=\frac{\delta_{ij}-\theta_i\theta_j\ind_{\{j\le m\}}}{R},
 \qquad
 \cL\theta_i
 =\theta_i\left(\frac{3-m-2\ind_{\{i\le m\}}}{R^2}
                         -r_i+\overline r\right).
\]
The Fukushima drift is the time integral of this \(L^2\) generator.
Its martingale has brackets with the coordinate martingales equal to
\(2\int_0^t\partial_j\theta_i(X(v))\,dv\), and self-bracket
\(2\int_0^t|\nabla\theta_i(X(v))|^2\,dv\).
It therefore equals
\(\sqrt2\sum_j\int_0^t\partial_j\theta_i(X(v))\,dW_j(v)\):
the difference has zero bracket.  Only the head and index \(i\)
occur in this finite sum.  Integrability follows, for instance, from
\[
 \E\int_0^T|\nabla\theta_i(X(v))|^2\,dv
 \le T\|w\|_\infty\cE_g(\theta_i)<\infty,
\]
because \(\|\widetilde S_vw\|_\infty\le\|w\|_\infty\).
This proves \eqref{entr:angular-sde}.  Countability permits a common
null set for all indices.
\end{proof}

We now cancel the common head projection in the angular martingale and
the common angular drift with a single stochastic integrating factor.
On the process of Lemma~\ref{entr:angular-realization}, define the
positive scalar \(\mathcal U(t)\) by
\begin{equation}\label{entr:scalar-rescaling}
 \mathcal U(0)=\frac{R(0)}{\sqrt{M(X(0))}},\qquad
 d\mathcal U
 =\mathcal U\left\{
     \frac{\sqrt2}{R}\sum_{j\le m}\Theta_j\,dW_j
       +\left(\frac{m-1}{R^2}-\overline r\right)dt\right\}.
\end{equation}
This is the stochastic exponential with the displayed drift and initial
value.  On every compact time interval \(R\) has a strictly positive
minimum, so all its coefficients are locally bounded along the path;
the exponential is finite and strictly positive.  Put
\begin{equation}\label{entr:YQ}
 Q=\frac{\mathcal U}{R},\qquad f=Q^2,\qquad
 Y_i=\mathcal U\Theta_i=QX_i,\qquad y_i=Y_i(0).
\end{equation}
In the product rule for \(\mathcal U\Theta_i\), the quadratic
covariation contributes
\(2\mathcal U\Theta_i(\ind_{\{i\le m\}}-1)R^{-2}dt\).
Adding this to the two drifts in
\eqref{entr:angular-sde}--\eqref{entr:scalar-rescaling} cancels all
terms containing \(R^{-2}\) or \(\overline r\).  Their martingale
terms cancel the head projection.  Thus
\begin{equation}\label{entr:common-sde}
 dY_i=-r_iY_i\,dt+\sqrt2Q\,dW_i,\qquad
 Y_i(t)=e^{-r_it}y_i+\sqrt2\int_0^t e^{-r_i(t-v)}Q(v)\,dW_i(v).
\end{equation}
Both \(Q\) and \(f\) are positive continuous processes, including at
elapsed time zero.  The normalization at zero gives
\begin{equation}\label{entr:seed-weights}
 p_i:=\frac{y_i^2}{2r_i}\ge0,\qquad
 \sum_i p_i=1,\qquad
 y_i^2=\frac{X_i(0)^2}{M(X(0))}.
\end{equation}
In particular \eqref{entr:escape} says that these weights leave every
finite set in probability when the starting time decreases along its
specified sequence.

\subsection{The scalar constraint equation}

The coordinates in \eqref{entr:common-sde} are driven by the same
factor \(Q\).  The quadratic constraint on \(X\) determines that
factor through a scalar Volterra equation.  To identify the equation
before taking an infinite limit, view \(f\) and the time kernels causally,
by extending them by zero to negative elapsed times, and fix \(N\).  Set
\[
 E_i(t)=e^{-2r_it}\quad(t\ge0),\qquad
 L_N(t)=\sum_{i\le N}r_i^{-1}E_i(t),\qquad
 R_N(x)=D-\sum_{i\le N}(x_i^2-r_i^{-1}).
\]
Itô's formula in \eqref{entr:common-sde} gives
\[
 Y_i(t)^2=E_i(t)y_i^2+2(E_i*f)(t)
       +2\sqrt2\int_0^t E_i(t-v)Q(v)Y_i(v)\,dW_i(v).
\]
Since \(\sum_{i\le N}Y_i^2=f\sum_{i\le N}X_i^2\) and
\[
 \sum_{i\le N}X_i^2
   =D+\sum_{i\le N}r_i^{-1}-R_N(X),
\]
rearrangement gives, as an identity of causal distributions,
\begin{equation}\label{entr:finite-constraint}
 Df+(L_N*f)'=
 \sum_{i\le N}E_i y_i^2
 +2\sqrt2\sum_{i\le N} E_i*(QY_i\,dW_i)+fR_N(X).
\end{equation}
Here the derivative includes the mass at elapsed time zero:
\begin{equation}\label{entr:initial-delta}
 L_N'=\left(\sum_{i\le N}r_i^{-1}\right)\delta_0
                       -2\sum_{i\le N}E_i .
\end{equation}
Thus the coefficient of \(f\) contributed by the delta mass is retained
in \eqref{entr:finite-constraint}.  This identity is exact for finite
\(N\).  The remaining passage must remove \(fR_N(X)\) and replace
\(L_N\) by
\[
 L(t)=\sum_i r_i^{-1}E_i(t),\qquad t>0.
\]
The corresponding limiting operator on \(f\) is convolution with
\(D\delta_0+L'\).  The next lemma constructs its positive inverse
kernel and the bounds needed for the stochastic terms.

\begin{lemma}[Resolvent kernels]\label{entr:kernels}
For \(L\) defined above, let
\[
 \widehat K(z)=\frac{1}{D+z\widehat L(z)}\quad(z>0),
\]
where a hat denotes the Laplace transform and every kernel is zero on
negative times.  The function \(L\) is locally integrable, and
\(\widehat K\) is the Laplace transform of a nonnegative, decreasing
function \(K\), smooth on \((0,\infty)\).  For \(0<t\le1\),
\begin{equation}\label{entr:K-bounds}
 c_g t^{-1/2}\le K(t)\le C_g t^{-1/2},
 \qquad |K'(t)|\le C_g t^{-3/2}.
\end{equation}
Set \(K_i=K*E_i\).  Then
\begin{align}
 K_i(t)&\le C_g\frac{\sqrt t}{1+r_it},&
 |K_i'(t)|&\le C_g\frac{t^{-1/2}}{1+r_it},             \label{entr:Ki-bounds}\\
 2r_iK_i(t)&\longrightarrow K(t)                       \label{entr:Ki-limit}
\end{align}
uniformly on each compact subinterval of \((0,1]\) as \(i\to\infty\).
For \(t>0\), \(h\ge0\), \(t+h\le1\), and \(0\le\xi<1/4\),
\begin{align}
 \sup_i r_iK_i(t)^2&\le C_g,&
 \sup_i r_i|K_i(t+h)-K_i(t)|^2
   &\le C_g\min\{1,h^2/t^2\},                          \label{entr:Ki-sup}\\
 \sum_i(1+r_i)^{2\xi}K_i(t)^2
   &\le C_{g,\xi}t^{-1/2-2\xi},                       \label{entr:Ki-sum}\\
 \sum_i|K_i(t+h)-K_i(t)|^2
   &\le C_g t^{-1/2}\min\{1,h^2/t^2\}.                \label{entr:Ki-increment}
\end{align}
As causal distributions,
\begin{equation}\label{entr:kernel-inverse}
 K*(D\delta_0+L')=\delta_0.
\end{equation}
\end{lemma}

\begin{proof}
The local integrability follows directly from
\[
 \int_0^T L(t)\,dt
 =\frac12\sum_i r_i^{-2}(1-e^{-2r_iT})<\infty.
\]
For \(L_N=\sum_{i\le N}r_i^{-1}E_i\), put
\[
 F_N(z)=D+z\widehat L_N(z)
       =D+\sum_{i\le N}\frac{z}{r_i(z+2r_i)}.
\]
The rational function \(F_N\) is strictly increasing between its
negative poles, since
\(F_N'(z)=\sum_{i\le N}2(z+2r_i)^{-2}>0\).
Its zeros interlace those poles and are strictly negative, using \(D>0\).
Its reciprocal is therefore a positive constant plus a sum of positive
multiples of \((z+a)^{-1}\), \(a>0\): the residue at a zero is
\(1/F_N'(-a)>0\).  Repeated rates can be grouped before this argument.

Write these sums as Stieltjes transforms of positive measures.  Their
measures weighted by \((1+a)^{-1}\), with the constant as mass at the
point at infinity, have total mass \(F_N(1)^{-1}\le D^{-1}\).
Weak compactness on \([0,\infty]\) and monotone convergence of
\(F_N(z)\) give
\[
 \frac1{D+z\widehat L(z)}
   =c+\int_{[0,\infty)}\frac{\nu(da)}{z+a},
 \qquad \int\frac{\nu(da)}{1+a}<\infty.
\]
By \eqref{entr:rates}, splitting the sum at \(r_i=z\) gives
\begin{equation}\label{entr:laplace-scale}
 D+z\widehat L(z)\asymp_g \sqrt z\qquad(z\ge1).
\end{equation}
For sufficiently large \(z\), the part with \(r_i\le z\) is
comparable to \(\sum_{r_i\le z}r_i^{-1}\asymp\sqrt z\), and the
other part to \(z\sum_{r_i>z}r_i^{-2}\asymp\sqrt z\).
Positivity and continuity adjust the constants on the remaining
compact interval \(z\ge1\).
Hence the transform tends to zero at infinity and \(c=0\).
It is bounded as \(z\downarrow0\), so \(\nu\) has no atom at zero.
The function
\[
 K(t)=\int_{(0,\infty)}e^{-at}\nu(da)
\]
has the required Laplace transform by Tonelli's theorem.  The weighted
integrability of \(\nu\) implies smoothness for \(t>0\), and the
representation makes \(K\) positive and decreasing.

The upper bound in \eqref{entr:K-bounds} follows from
\[
 e^{-1}tK(t)\le\int_0^t e^{-v/t}K(v)\,dv
 \le\widehat K(1/t)\le C_g\sqrt t.
\]
For the lower bound choose a fixed small \(\eta>0\).  When \(t\le\eta\),
\eqref{entr:laplace-scale} bounds
\(\widehat K(\eta/t)\) below by \(c_g\sqrt{t/\eta}\), whereas the
upper bound just proved bounds its integral over \([0,t]\) by
\(C_g\sqrt t\).  Choose \(\eta\) so that the first constant is larger
than twice the second.  The integral over \([t,\infty)\) is at most
\(\eta^{-1}t e^{-\eta}K(t)\), proving the lower bound for \(t\le\eta\).
Positivity and monotonicity adjust the constant on \([\eta,1]\).
Finally
\[
 |K'(t)|=\int a e^{-at}\nu(da)
 \le \frac{C}{t}K(t/2)\le C_g t^{-3/2}.
\]

For \(r_it\le1\), integration of \(K(v)\le C_gv^{-1/2}\)
gives the first bound in \eqref{entr:Ki-bounds}.  For \(r_it>1\),
write \(K_i(t)=\int_0^t K(t-v)e^{-2r_iv}\,dv\) and split at \(t/2\).
The first half is bounded by \(C_g/(r_i\sqrt t)\); the second by
\(C_g\sqrt t\,e^{-r_it}\), which has the same bound.  For the
derivative use \(K_i'=K-2r_iK_i\).  When \(r_it\le1\) the preceding
bounds suffice.  When \(r_it>1\), write
\[
 K_i'(t)=K(t)e^{-2r_it}
  -2r_i\int_0^t e^{-2r_iv}\bigl(K(t-v)-K(t)\bigr)\,dv.
\]
On \(v\le t/2\) the derivative bound on \(K\) makes the integrand
difference at most \(C_g v t^{-3/2}\); on \(v>t/2\) the exponential
factor pays the integral of \(K\).  This gives
\(|K_i'(t)|\le C_g/(r_it^{3/2})\).  It also proves
\eqref{entr:Ki-limit}, uniformly away from zero.

The remaining estimates follow by summing
\eqref{entr:Ki-bounds}.  For example \eqref{entr:rates} gives
\[
 \sum_i\frac{(1+r_i)^{2\xi}}{(1+r_it)^2}
 \le C_{g,\xi}t^{-3/2-2\xi}\quad(0\le\xi<1/4).
\]
The restriction on \(\xi\) is exactly integrability at infinity after
the change of variable \(v=r_it\): the comparison integral there is
\(\int_0^\infty v^{1/2+2\xi}(1+v)^{-2}\,dv\).
This proves \eqref{entr:Ki-sum}.  If \(h<t\), integrate the bound on
\(K_i'\) between \(t\) and \(t+h\), and then use the same sum with
\(\xi=0\); this gives \eqref{entr:Ki-increment} and
\eqref{entr:Ki-sup}.  If \(h\ge t\), bound the two values separately.
Finally the locally integrable function
\[
 D\int_0^tK(v)\,dv+(K*L)(t)
\]
has Laplace transform \(1/z\), and hence equals \(1\) for almost
every \(t>0\).  Its causal distributional derivative gives
\eqref{entr:kernel-inverse}.
\end{proof}

For the normalized seed weights \(p_i=y_i^2/(2r_i)\), define, for
\(0<t\le1\), the candidate seed response
\[
 A_y(t)=\sum_iK_i(t)y_i^2=\sum_i p_i\,2r_iK_i(t).
\]
The bound \(2r_iK_i(t)\le C_g t^{-1/2}\), together with
\(\sum_i p_i=1\), proves that this series converges uniformly on
each compact subinterval of \((0,1]\).  Its sum \(A_y\) is continuous
there and satisfies \(0\le A_y(t)\le C_g t^{-1/2}\).
The high-index limit \(2r_iK_i\to K\) is the ingredient used below
to prove \(A_y\to K\) in probability, uniformly on compact subintervals
of \((0,1]\), when the weights leave every fixed finite set.
The bound \(K(t)\asymp_g t^{-1/2}\) determines the scale used in
the later stopped comparison of \(f=Q^2\).  The next lemma defines
the infinite stochastic convolution and controls it uniformly over
normalized seeds under the corresponding upper bound
\(|q(t)|\le C_0t^{-1/4}\).

\begin{lemma}[Stochastic kernel bounds]\label{entr:stochastic-kernels}
Let \((W_i)\) be independent Brownian motions in a filtration, let
\((y_i)\) be measurable at time zero with
\(\sum_i y_i^2/(2r_i)=1\), and let \(q\) be progressively measurable
with \(|q(v)|\le C_0v^{-1/4}\) for \(0<v\le1\).  Define
\begin{align}
 Z_i^q(t)&=\sqrt2\int_0^t e^{-r_i(t-v)}q(v)\,dW_i(v),\notag\\
 \mathsf G_q^y(t)&=2\sqrt2\sum_i\int_0^t
                K_i(t-v)q(v)e^{-r_iv}y_i\,dW_i(v),\label{entr:stochastic-parts}\\
 \mathsf N_q(t)&=2\sqrt2\sum_i\int_0^t
                K_i(t-v)q(v)Z_i^q(v)\,dW_i(v).\notag
\end{align}
The two series converge in probability in \(C([0,1])\), with value
zero at zero.  For every finite \(j\ge2\), their partial sums and their
limits satisfy
\begin{align}
 \|Z_i^q(t)\|_{L^j}^2
     &\le C_{g,C_0,j}\frac{\sqrt t}{1+r_it},           \label{entr:Z-moment}\\
 \|\mathsf G_q^y(t)\|_{L^j}+\|\mathsf N_q(t)\|_{L^j}
     &\le C_{g,C_0,j}t^{1/4},                         \label{entr:stoch-moment}\\
 \|\mathsf G_q^y(t+h)-\mathsf G_q^y(t)\|_{L^j}
  +\|\mathsf N_q(t+h)-\mathsf N_q(t)\|_{L^j}
     &\le C_{g,C_0,j}h^{1/4}\quad(0\le t<t+h\le1).
                                                               \label{entr:stoch-increment}
\end{align}
The supremum on \([0,1]\) of either limit has every finite moment,
bounded by a constant depending only on \(g,C_0\), and that order.
\end{lemma}

\begin{proof}
The bracket estimate and the Burkholder--Davis--Gundy inequality give
\[
 \|Z_i^q(t)\|_{L^j}^2
 \le C_{C_0,j}\int_0^t e^{-2r_i(t-v)}v^{-1/2}\,dv
 \le C_{C_0,j}\frac{\sqrt t}{1+r_it}.
\]
For the last inequality split at \(t/2\); on the second half use
\(v^{-1/2}\le C t^{-1/2}\), and on the first use
\(e^{-2r_i(t-v)}\le e^{-r_it}\).

Write \(p_i=y_i^2/(2r_i)\).  The bracket of the first series, at a
fixed time \(t\), is bounded pathwise by
\[
 C_{C_0}\int_0^t v^{-1/2}
       \sum_i p_i r_iK_i(t-v)^2\,dv
 \le C_{g,C_0}\sqrt t,
\]
by \eqref{entr:Ki-sup} and \(\sum_i p_i=1\).
This proves its bound in \eqref{entr:stoch-moment}.  In its increment,
the old portion of the bracket is at most
\[
 C_{g,C_0}\int_0^t v^{-1/2}
              \min\{1,h^2/(t-v)^2\}\,dv
 \le C_{g,C_0}\sqrt h ,
\]
and the new portion is bounded by
\(C\int_t^{t+h}v^{-1/2}\,dv\le C\sqrt h\).
For the displayed integral, if \(t\le2h\) discard the minimum; if
\(t>2h\), split at \(t-h\) and at \(t/2\).  These bounds prove the
first increment estimate.

For the second series, the martingale moment inequality followed by
Minkowski's inequality bounds its \(L^j\) norm squared by
\[
 C_j\int_0^t\sum_i K_i(t-v)^2
                      \|q(v)Z_i^q(v)\|_{L^j}^2\,dv
 \le C_{g,C_0,j}\int_0^t (t-v)^{-1/2}\,dv
 \le C_{g,C_0,j}\sqrt t.
\]
We used \eqref{entr:Z-moment}, the deterministic bound on \(q\), and
\eqref{entr:Ki-sum} with \(\xi=0\).
For the increment use \eqref{entr:Ki-increment} on the old portion:
its squared norm is bounded by
\[
 C_{g,C_0,j}\int_0^t (t-v)^{-1/2}
                \min\{1,h^2/(t-v)^2\}\,dv
 \le C_{g,C_0,j}\sqrt h .
\]
The new portion has the same bound by integrating
\((t+h-v)^{-1/2}\) from \(t\) to \(t+h\).

All these estimates hold for partial sums.  At each fixed time their
tails tend to zero in \(L^j\): for the first series use
\(\sum_{i>N}p_i r_iK_i(t-v)^2\to0\) and its constant bound inside
the bracket; for the second use the summable bound just displayed.
Dominated convergence applies in both cases.  Taking \(j>4\),
\eqref{entr:stoch-increment} gives uniformly small moduli of
continuity, by a dyadic chaining estimate (equivalently the
Kolmogorov continuity criterion).  Fixed-time convergence on a finite
grid then proves convergence in probability in \(C([0,1])\).
The same chaining has a convergent series
\(\sum_{k\ge0}2^{k/j}2^{-k/4}\), and bounds the \(L^j\) norm of the
supremum for \(j>4\).  Taking a higher order gives the asserted bound
for every finite order.
\end{proof}

\begin{lemma}[The renormalized Volterra equation]\label{entr:volterra}
Let the processes be as in \eqref{entr:YQ}--\eqref{entr:common-sde},
and let \(A_y\) be the candidate seed response defined above.
Almost surely, for \(0<t\le1\),
\begin{equation}\label{entr:volterra-equation}
 f(t)=A_y(t)+2\sqrt2\sum_i\int_0^t
               K_i(t-v)Q(v)Y_i(v)\,dW_i(v).
\end{equation}
The stochastic series is defined by localization and has a continuous
version on every compact subinterval of \((0,1]\).
\end{lemma}

\begin{proof}
We now justify the infinite passage in \eqref{entr:finite-constraint}.
The constraint says \(R_N(x)=\sum_{i>N}(x_i^2-r_i^{-1})\)
\(\Pi\)-almost surely.  The conditioned edge law describes the tail,
after any sufficiently large fixed head, by a bounded density relative
to its product Gaussian law
\cite[Proposition~2.3]{FUA}.
Consequently, for \(N\) beyond that head,
\begin{equation}\label{entr:tail-residual}
 \Pi[|R_N|]\le C_g\left(\sum_{i>N}r_i^{-2}\right)^{1/2}
       \longrightarrow0.
\end{equation}
Indeed the Gaussian tail is centered and has variance
\(2\sum_{i>N}r_i^{-2}\).

The density of \(X(t)\) is bounded by \(\|w\|_\infty\) for all
elapsed \(t\ge0\).  The constraint therefore holds for almost every
elapsed time almost surely, and \eqref{entr:tail-residual} gives
\[
 \E\int_0^T|R_N(X(v))|\,dv
       \le T\|w\|_\infty\Pi[|R_N|]\longrightarrow0 .
\]
The continuous process \(f\) is bounded on \([0,T]\) almost surely.
First restrict to a bound for that supremum, and then let the bound
increase.  It follows that
\(fR_N(X)\to0\) in probability in \(L^1(0,T)\).
Also
\begin{equation}\label{entr:L-tail}
 \|L-L_N\|_{L^1(0,T)}
 \le\frac12\sum_{i>N}r_i^{-2}\longrightarrow0 .
\end{equation}
Convolving \eqref{entr:finite-constraint} with \(K\), its left side
therefore tends to \(f\) in local distributions by
\eqref{entr:kernel-inverse}.  One can see this directly by writing
the difference as the derivative of
\(K*(L_N-L)*f\); its undifferentiated factor tends to zero in
\(L^1_{\mathrm{loc}}\) by \eqref{entr:L-tail}.  The convolved residual
also tends to zero in probability in \(L^1_{\mathrm{loc}}\), since
\(K\) is locally integrable.

The convolved deterministic terms are
\(\sum_{i\le N}K_i(t)y_i^2\), which converge uniformly on compact
subintervals of \((0,1]\) to \(A_y(t)\) by its definition and bound
above.

For a finite stochastic sum, stochastic Fubini gives
\[
 K*\bigl(E_i*(QY_i\,dW_i)\bigr)(t)
       =\int_0^t K_i(t-v)Q(v)Y_i(v)\,dW_i(v).
\]
It can first be used with \(QY_i\) stopped at a finite bound.
For such a bound a sufficient stochastic Fubini integral is at most
\[
 C\int_0^t K(t-s)
       \left(\int_0^s E_i(s-v)^2\,dv\right)^{1/2}ds
 \le C\int_0^t K(t-s)\sqrt s\,ds<\infty.
\]
Removing that bound by localization gives the displayed identity as
local distributions.

To specify the limit without a formal exchange of infinite stochastic
sums, for \(C_0>0\) set
\[
 q_{C_0}(v)=Q(v)\wedge(C_0v^{-1/4}),\qquad
 \zeta_{C_0}=\inf\{v>0:Q(v)>C_0v^{-1/4}\}.
\]
By Lemma~\ref{entr:stochastic-kernels}, the partial sums of
\(\mathsf G_{q_{C_0}}^y+\mathsf N_{q_{C_0}}\) converge in probability
in \(C([0,1])\) on the whole probability space.  Before
\(\zeta_{C_0}\), \eqref{entr:common-sde} gives
\(Y_i(v)=e^{-r_iv}y_i+Z_i^{q_{C_0}}(v)\), so these clipped partial
sums agree there with the convolved finite stochastic sums.

Fix a deterministic \(C_0>0\) and a rational \(a\in(0,1]\).
On the whole probability space, the clipped partial sums converge in
probability in \(C([0,a])\), and the convolved residual
\(K*(fR_N(X))\) converges in probability to zero in \(L^1(0,a)\)
as shown above.  We may therefore choose a deterministic subsequence
\(N_j\uparrow\infty\) along which both convergences hold almost surely.
Only then restrict to the event \(\{a<\zeta_{C_0}\}\).  On this event,
passage in \eqref{entr:finite-constraint} proves
\eqref{entr:volterra-equation} as a distribution on \((0,a)\).
Take \(C_0\) through the positive integers and \(a\) through the
rationals in \((0,1]\), and intersect the corresponding whole-space
probability-one convergence events.  The continuous versions then give
equality for \(0<t<\min\{\zeta_{C_0},1\}\), and by continuity at a
positive hitting time \(\zeta_{C_0}\le1\) and at \(t=1\) when
\(\zeta_{C_0}>1\).
The positive continuous \(Q\) is finite at zero, and
\(\sup_{v\le1}v^{1/4}Q(v)<\infty\) almost surely.  A sufficiently
large integer \(C_0\) therefore covers the whole interval.  This proves the
claimed identity and its precise interpretation.
\end{proof}

We next isolate the part of the equation that remembers the initial
state.  It becomes independent of that state when the weights
\eqref{entr:seed-weights} escape.

\begin{lemma}[Escaping seeds]\label{entr:escaping-seeds}
Suppose that, along a sequence of initial states possibly random,
\(\sum_i p_i=1\) and \(p_i\to0\) in probability for every fixed \(i\).
Then
\begin{equation}\label{entr:A-limit}
 A_y\longrightarrow K
 \quad\text{in probability, uniformly on compact subintervals of \((0,1]\)}.
\end{equation}
For each fixed \(t>0\), the seed stochastic term
\(\mathsf G_q^y(t)\) tends to zero in \(L^2\), uniformly over all
adapted \(q\) satisfying \(|q(v)|\le C_0v^{-1/4}\).
\end{lemma}

\begin{proof}
By \eqref{entr:Ki-limit},
\(2r_iK_i\to K\) uniformly on a fixed compact positive interval,
and these functions are uniformly bounded there.  In the convex
combination \(A_y=\sum_i p_i\,2r_iK_i\), the contribution of a
fixed finite set tends to zero in probability and the contribution
of its complement is uniformly close to \(K\).  This proves
\eqref{entr:A-limit}.

For the stochastic assertion, its squared norm is at most
\[
 C_{C_0}\int_0^t v^{-1/2}\,
       \E\!\left[\sum_i p_i r_i K_i(t-v)^2\right]\,dv .
\]
For every \(v<t\), \eqref{entr:Ki-bounds} implies
\(r_iK_i(t-v)^2\to0\), whereas \eqref{entr:Ki-sup} bounds it
uniformly.  Escape of \(p_i\) makes the expression in brackets tend
to zero in probability, and its uniform bound gives convergence of
its expectation.  Dominated convergence against \(v^{-1/2}\)
proves the assertion.
\end{proof}

\subsection{Coupling entrances}

We now prove the deterministic uniqueness and boundary assertions of
Theorem~\ref{thm:entrance-selection}.  Work in unit time
\eqref{entr:unit-time}.  Let two proposed families be given, with
starting sequences \(u_{\ell,k}\downarrow0\), \(\ell=1,2\).
Start the diffusion of Lemma~\ref{entr:angular-realization} from
\(\widetilde h^{\,\ell}_{u_{\ell,k}}\Pi\), and construct
\((Q_\ell,f_\ell,Y_{\ell,i},y_{\ell,i})\) as above.  Their seed weights
escape in probability by \eqref{entr:escape} and
\eqref{entr:seed-weights}.  The constants below depend on \(g\), and
are uniform in \(k\) and in the two families.

We may use the same Brownian motions for the two diffusions in elapsed
time even though strong existence has not been asserted.  Here is the
coupling argument.  Each joint law of the Brownian path \(W\) and all
its other continuous path variables is a probability on a standard
Borel path space.  Disintegrate it over \(W\), and, given one Brownian
path, sample the two other path variables independently from these
conditional laws.  Brownian motion remains Brownian in the joint
filtration.  To see this, if \(\Phi_\ell\) is a bounded observable of the
past of process \(\ell\) through time \(t\), independence of future
Brownian increments in that process implies
\[
 \E[\Phi_\ell\mid W\text{ on }[0,1]]
   =\E[\Phi_\ell\mid W\text{ on }[0,t]]
 \quad\text{almost surely}.
\]
Test this equality first against products of functions of the Brownian
past and future increments, and then use a monotone class argument.
For a product \(\Phi_1\Phi_2\), conditional independence in the coupling
makes its conditional expectation the product of two functions of the
Brownian past.  A second monotone class argument proves independence
of future increments from the joint past.  Completion and
right-continuous augmentation preserve this property.  Finally,
stochastic integrals are limits in probability of adapted simple
sums.  Their laws in each marginal are unchanged by the coupling,
and they are the same integrals in this joint Brownian filtration.
Thus all the preceding stochastic identities hold on the coupling.

\paragraph{Stops with uniform failure bounds.}
Choose \(C_0\) so large that
\[
 C_0^2>1+\sup_{0<t\le1}\sup_y\sqrt t\,A_y(t),
\]
where the supremum is over the normalized seeds in
\eqref{entr:seed-weights}; it is finite by the bound on \(A_y\)
following Lemma~\ref{entr:kernels}.  For each solution use the clipped process
\(q_\ell^{\,0}(v)=Q_\ell(v)\wedge(C_0v^{-1/4})\).
Let \(k_0=\inf_{0<t\le1/2}\sqrt t\,K(t)>0\), and choose
\(c\in(0,k_0/4)\).  Lemma~\ref{entr:escaping-seeds} permits a
deterministic sequence \(\eta_k\downarrow0\), as slow as necessary,
such that with probability tending to one both solutions satisfy
\begin{equation}\label{entr:seed-lower}
 \sqrt t\,A_{y_\ell}(t)\ge2c
             \qquad(\eta_k\le t\le1/2).
\end{equation}
This follows by a diagonal choice from uniform convergence on each
fixed compact interval.

Stop each solution when \(\sqrt t\,f_\ell(t)\) first reaches the
upper threshold \(C_0^2\); it starts below that threshold since
\(f_\ell\) is finite at zero.  Starting at \(\eta_k\), also stop
when \(\sqrt t\,f_\ell(t)\le c\).  Before the upper stop,
Lemma~\ref{entr:volterra} reads
\[
 f_\ell=A_{y_\ell}
       +\mathsf G_{q_\ell^{\,0}}^{y_\ell}
       +\mathsf N_{q_\ell^{\,0}}.
\]
It holds also at a positive first hit by continuity.  On the event
\eqref{entr:seed-lower}, either threshold hit by time \(\delta\le1/2\)
requires the supremum of the stochastic term on \([0,1]\) to be
at least a fixed multiple of \(\delta^{-1/2}\).
Lemma~\ref{entr:stochastic-kernels} and Markov's inequality therefore
bound its probability by \(C_j\delta^{j/2}\) for any fixed
sufficiently large \(j\), plus a term tending to zero with \(k\).

Fix \(0<\xi<1/4\).  Add a stop for each solution at the infimum,
over \(i\ge1\), of the first times
\[
 |Z_i^{q_\ell^{\,0}}(t)|\ge(1+r_i)^\xi.
\]
This costs at most a positive power of \(\delta\) by time \(\delta\).
Indeed if \(N_i(t)=\sqrt2\int_0^t q_\ell^{\,0}(v)\,dW_i(v)\),
integration by parts gives
\[
 Z_i^{q_\ell^{\,0}}(t)
 =N_i(t)-r_i\int_0^t e^{-r_i(t-v)}N_i(v)\,dv,
 \qquad
 \sup_{t\le\delta}|Z_i^{q_\ell^{\,0}}(t)|
       \le2\sup_{t\le\delta}|N_i(t)|.
\]
The last supremum has \(L^j\) norm at most \(C_j\delta^{1/4}\),
uniformly in \(i\).  Choose \(j>3/(2\xi)\), so
\(\sum_i(1+r_i)^{-j\xi}<\infty\) by \eqref{entr:rates}, and sum
Markov's inequality over \(i\).  The resulting bound is
\(C_j\delta^{j/4}\).  A countable infimum of these stopping times is
a stopping time in the right-continuous filtration.  If an infimum is
not attained, an individual stop occurs by \(2\delta\) whenever the
infimum is at most \(\delta>0\); this only changes the constant.

Let \(\tau_k\) be the minimum of all these stops for both solutions,
capped at \(1/2\).  There are constants \(b>0,C<\infty\), independent
of \(k\), such that, for sufficiently small fixed \(\delta>0\),
\begin{equation}\label{entr:stop-probability}
 \limsup_{k\to\infty}\Prob(\tau_k\le\delta)\le C\delta^b.
\end{equation}
Here and below the probability is over the initial draws and the
coupled diffusions, with \(g\) fixed.

\paragraph{Contraction of the common factors.}
Use the common stopped integrands
\[
 q_\ell(v)=Q_\ell(v)\ind_{\{v<\tau_k\}},\qquad
 d_k(v)=\sqrt v\,\E|q_1(v)-q_2(v)|^2.
\]
They satisfy \(|q_\ell(v)|\le C_0v^{-1/4}\), so \(0\le d_k(v)\le C\).
Until \(\tau_k\), the stopped versions of
\eqref{entr:volterra-equation} agree with \(f_\ell\), and
\(|Z_i^{q_\ell}(v)|\le(1+r_i)^\xi\).
For \(v\ge\eta_k\), the lower threshold gives
\begin{equation}\label{entr:sqrt-difference}
 d_k(v)\le C v\,\E\left[
       \ind_{\{v<\tau_k\}}|f_1(v)-f_2(v)|^2\right].
\end{equation}
This is the identity
\((\sqrt{f_1}-\sqrt{f_2})^2=(f_1-f_2)^2/
(\sqrt{f_1}+\sqrt{f_2})^2\), using
\(f_\ell(v)\ge c v^{-1/2}\) before the stop.

Let
\[
 F_\ell(v)=A_{y_\ell}(v)
       +\mathsf G_{q_\ell}^{y_\ell}(v)+\mathsf N_{q_\ell}(v).
\]
It agrees with \(f_\ell\) for \(v<\tau_k\).  At a fixed \(v>0\),
Lemma~\ref{entr:escaping-seeds} and its uniform bounds give
\begin{equation}\label{entr:seed-vanishes}
 \E|A_{y_1}(v)-A_{y_2}(v)|^2
  +\E|\mathsf G_{q_1}^{y_1}(v)|^2
  +\E|\mathsf G_{q_2}^{y_2}(v)|^2\longrightarrow0 .
\end{equation}
For the deterministic term, convergence in probability becomes
\(L^2\) convergence because \(|A_y(v)|\le C_gv^{-1/2}\).

It remains to compare \(\mathsf N_{q_1}\) and \(\mathsf N_{q_2}\).
In its integrand write
\[
 q_1Z_i^{q_1}-q_2Z_i^{q_2}
   =(q_1-q_2)Z_i^{q_1}+q_2(Z_i^{q_1}-Z_i^{q_2}).
\]
The first term is nonzero only before the common stop, where
\(|Z_i^{q_1}|\le(1+r_i)^\xi\).  For the second, the common
Brownian motions and Itô isometry give
\[
 \E|Z_i^{q_1}(w)-Z_i^{q_2}(w)|^2
 =2\int_0^w e^{-2r_i(w-a)}a^{-1/2}d_k(a)\,da .
\]
Using \eqref{entr:Ki-sum} with \(\xi\) for the first term and
with zero for the second yields
\begin{align}
 \E|\mathsf N_{q_1}(v)-\mathsf N_{q_2}(v)|^2
 &\le C\int_0^v (v-w)^{-1/2-2\xi}w^{-1/2}d_k(w)\,dw\notag\\
 &\quad+C\int_0^v (v-w)^{-1/2}w^{-1/2}
                 \int_0^w a^{-1/2}d_k(a)\,da\,dw .
                                                        \label{entr:contraction-integrals}
\end{align}
All sums and integrals here are justified by Tonelli's theorem and
the displayed finite upper bounds.  In the second line we used the
deterministic bound \(q_2(w)^2\le C_0^2w^{-1/2}\) and discarded
the decaying exponential after Itô isometry.  No independence of
these integrands is required.

For \(v>0\) define \(d_\infty(v)=\limsup_k d_k(v)\).
For each fixed \(v\) we eventually have \(v\ge\eta_k\).
Combine \eqref{entr:sqrt-difference}--\eqref{entr:contraction-integrals}.
The bounds \(d_k\le C\) allow the upper limit to pass under the
integrals by dominated convergence applied to upper envelopes.  The
kernels are integrable because \(\xi<1/4\).  If
\(D_\delta=\sup_{0<v\le\delta}d_\infty(v)\), the two beta integrals
then give
\[
 d_\infty(v)
 \le C\bigl(v^{1-2\xi}+v^{3/2}\bigr)D_\delta
 \qquad(0<v\le\delta).
\]
For the first power use
\(\int_0^v(v-w)^{-1/2-2\xi}w^{-1/2}\,dw=C_\xi v^{-2\xi}\);
for the second use
\(\int_0^w a^{-1/2}\,da=2\sqrt w\).
Choose \(\delta>0\) small enough that
\(C(\delta^{1-2\xi}+\delta^{3/2})<1\).  Taking the supremum proves
\begin{equation}\label{entr:q-convergence}
 d_k(v)\longrightarrow0\qquad(0<v\le\delta).
\end{equation}

At a fixed such \(\delta\), dominated convergence in
\eqref{entr:contraction-integrals} and
\eqref{entr:seed-vanishes} now shows
\[
 \E\!\left[\ind_{\{\delta<\tau_k\}}
                   |f_1(\delta)-f_2(\delta)|^2\right]\longrightarrow0.
\]
For each fixed \(i\), the same conclusion holds for
\(|Y_{1,i}(\delta)-Y_{2,i}(\delta)|^2\).  Indeed
\(y_{\ell,i}\to0\) in probability by escape and
\(|y_{\ell,i}|^2\le2r_i\), while the squared norm of the difference
of the two \(Z_i\) terms is
\(2\int_0^\delta e^{-2r_i(\delta-v)}v^{-1/2}d_k(v)\,dv\to0\)
by \eqref{entr:q-convergence}.  On \(\{\delta<\tau_k\}\) the
denominators satisfy \(f_\ell(\delta)\ge c\delta^{-1/2}\), and
\[
 |Y_{\ell,i}(\delta)|
 \le\sqrt{2r_i}+(1+r_i)^\xi .
\]
Since \(X_{\ell,i}=Y_{\ell,i}/\sqrt{f_\ell}\), it follows that for
every finite \(j\) and every \(\varepsilon>0\),
\begin{equation}\label{entr:coordinate-coupling}
 \limsup_k\Prob\left(\max_{i\le j}
       |X_{1,i}(\delta)-X_{2,i}(\delta)|>\varepsilon\right)
       \le C\delta^b.
\end{equation}

\paragraph{Equality of the families.}
The law at elapsed time \(\delta\) of process \(\ell\) is
\(\widetilde h^{\,\ell}_{u_{\ell,k}+\delta}\Pi\).  It converges
in total variation to \(\widetilde h^{\,\ell}_\delta\Pi\):
the semigroup relation writes its density as
\(\widetilde S_{u_{\ell,k}}\widetilde h^{\,\ell}_\delta\), and
strong continuity in \(L^2(\Pi)\) implies \(L^1\) convergence.
The couplings of their two elapsed-time marginals are tight on
\(\R^\N\times\R^\N\), so take a weakly convergent subsequence.
Its limit is a coupling of the two prescribed marginals at
\(\delta\).  For the open event in
\eqref{entr:coordinate-coupling}, the portmanteau theorem gives the
same upper bound for the limiting coupling.  Let
\(\varepsilon\downarrow0\) and then \(j\uparrow\infty\).
The limiting coupling assigns probability at least \(1-C\delta^b\)
to equality of all coordinates.  Hence
\[
 \|\widetilde h^{\,1}_\delta\Pi
          -\widetilde h^{\,2}_\delta\Pi\|_{\TV}\le C\delta^b.
\]
The Markov semigroup contracts total variation.  For any fixed
\(v>\delta\) the same bound holds at \(v\), and sending
\(\delta\downarrow0\) proves equality of the two marginals at \(v\).
They are equal at every positive time, and hence their density classes
are equal.  This proves uniqueness.

\paragraph{The boundary at zero.}
Use one proposed family and the preceding stops.  On
\(\{\delta<\tau_k\}\),
\[
 |X_i(\delta)|
 =\frac{|Y_i(\delta)|}{\sqrt{f(\delta)}}
 \le \frac{\sqrt{2r_i}+(1+r_i)^\xi}{\sqrt c}\,\delta^{1/4}.
\]
Let the starting index \(k\to\infty\).  The elapsed-time law
converges in total variation to the prescribed marginal at \(\delta\),
and \eqref{entr:stop-probability} bounds the excluded probability by
\(C\delta^b\).  Thus each coordinate tends to zero in probability
under \(\widetilde h_\delta\Pi\) as \(\delta\downarrow0\).
For example the bounded product metric
\(\sum_i2^{-i}\min\{1,|x_i|\}\) tends to zero in probability:
truncate its sum and use the finitely many coordinate conclusions.
This proves \eqref{entr:zero-boundary} in unit time, and then in
the theorem's critical time by \eqref{entr:unit-time}.

\subsection{Measurability and Gaussian identification}

Return now to the theorem's critical time \(s=u/c_*\), so that the
evolution is by \(S_t^g=e^{-t c_*H_g}\).  It remains to obtain the family from \(g\)
itself and to check existence.
We give the measure-theoretic passage, since the escape estimate in
Proposition~\ref{prop:uniform-seed-escape} initially averages over the
represented environment as well as \(x\).

On any represented Gaussian limiting space write
\(\mu_s=h_s\Pi_g\).  At rational \(s>0\) these probability measures
are measurably encoded by their integrals against a countable
convergence-determining collection of bounded continuous cylinders.
Equivalently
\[
 \boldsymbol\mu=(\mu_s)_{s\in\mathbb Q_{>0}}
      \in\mathcal P(\R^\N)^{\mathbb Q_{>0}}
\]
is a random element of a standard Borel space.  The asserted
measurability follows from that encoding: a countable separating
collection embeds \(\mathcal P(\R^\N)\) injectively and measurably
in a countable product, and its inverse on the image is measurable.
By the convention in Proposition~\ref{prop:uniform-seed-escape}, the
escape quotient is jointly Borel on \(\R^\N\times\R^\N\).
For each fixed \(i\) and \(\varepsilon>0\), its threshold event is a
Borel subset of this parameter/state space, so the evaluation of that
event under \(\mu_s\) is a measurable function of \((g,\mu_s)\).

The annealed escape estimate permits a deterministic rational sequence
\(s_k\downarrow0\) for which, almost surely in the represented data,
\eqref{entr:escape} holds for every \(i,\varepsilon\).
To see this choose \(s_k<\min\{2^{-k},s_{k-1}/2\}\) rational,
successively, so that
\[
 \E\sum_{i\le k}
  \mu_{s_k}\{x_i^2/M_g(x)>1/k\}\le2^{-k}.
\]
Proposition~\ref{prop:uniform-seed-escape} allows this choice.
Tonelli's theorem shows that the sum of the displayed random
quantities over \(k\) is finite almost surely.  For fixed
\(i,\varepsilon>0\), eventually \(i\le k\) and \(1/k<\varepsilon\);
the corresponding probability in \eqref{entr:escape} consequently
tends to zero.  The positive-time semigroup and boundedness properties
hold on the same full-measure set by
Proposition~\ref{prop:gaussian-subsequence}.

The finite-coordinate formula \eqref{entr:cylinder-density} makes
\(g\mapsto\Pi_g\) a measurable probability kernel.  The edge
semigroup is also measurable in \(g\) in its countable varying-\(L^2\)
encoding.  To see the precise claim needed here, take a countable
smooth cylinder core.  Its Gram entries \(\Pi_g[\phi_i\phi_j]\) and
form entries \(\Pi_g[\nabla\phi_i\cdot\nabla\phi_j]\) are measurable.
For \(\lambda>0\) the variational formula
\[
 \langle\phi,(\lambda+H_g)^{-1}\phi\rangle_{\Pi_g}
 =\sup_v\{2\Pi_g[\phi v]-\lambda\Pi_g[v^2]-\cE_g(v)\},
\]
where \(v\) ranges over the rational span of that core, proves
measurability of the resolvent pairings; polarization gives mixed
pairings.  Choose the countable core sequence linearly independent
as cylinder functions, with dense real span.  The positive
finite-coordinate densities make every finite Gram matrix positive
definite.  The coefficients of the orthogonal projection of an encoded
\(L^2\) vector on a finite core span are its cylinder pairings
multiplied by the inverse Gram matrix; they are measurable.
These projections converge in \(L^2\), proving the corresponding
measurability for arbitrary encoded vectors.  Continuous functional
calculus then gives measurability of \(S_t^g\) in the same encoding,
jointly for \(t>0\); one can approximate \(e^{-tc_*H_g}\) by continuous
functions of \((1+H_g)^{-1}\), and use strong continuity in \(t\).
This is also the countable-core construction in the functional
identification proof for the stationary edge law
\cite[proof of Proposition~9.2]{FUA}.

The normalization, boundedness, evolution, and escape conditions are
properties of the countable code together with \(g\).
Normalization and bounded domination by \(\Pi_g\) can be checked on
a countable cylinder algebra, with an integer bound on each rational
compact time interval.  Evolution can be checked against the same
tests at rational times; it then extends to real times by strong
\(L^2\) continuity.  The preceding core construction makes the semigroup
pairings measurable in \(g\).  Thus these conditions, together with
the displayed countable escape tests, define a measurable full set
of codes on which the deterministic theorem applies.

Disintegrate the law of \((g,\boldsymbol\mu)\) over \(g\).
For almost every \(g\), two independent conditional samples of
\(\boldsymbol\mu\) both extend to the positive-time families just
considered and both pass \eqref{entr:escape}.  Fixed-\(g\) uniqueness
therefore makes these samples equal.  A probability measure on a
standard Borel space whose two independent draws agree almost surely
is a point mass (apply the assertion to a countable separating
collection of Borel sets).  Thus the conditional law of
\(\boldsymbol\mu\) given \(g\) is a point mass almost surely.
Its point depends measurably on \(g\): the map
\(z\mapsto\delta_z\) is a measurable injection into the space of
probability measures and has a measurable inverse on its image.
This constructs measurable \(\mu_s^g\) for rational \(s\), and
establishes existence.

We spell out the extension to density classes and real times.  A measurable
Radon--Nikodym theorem applied to the kernels
\(\mu_q^g\ll\Pi_g\) gives measurable representatives \(h_q^g(x)\)
for rational \(q\).

For real \(s>0\), choose any rational \(q<s\) and set
\begin{equation}\label{entr:real-extension}
 h_s^g=S_{s-q}^g h_q^g .
\end{equation}
The rational evolution relation makes this independent of \(q\).
Choosing \(q=q(s)\) by a fixed measurable rational rule and using
the semigroup measurability established above shows that the density
classes in \eqref{entr:real-extension} are measurable in \((g,s)\).
The corresponding measures form a measurable kernel; the measurable
Radon--Nikodym theorem once more gives jointly measurable
representatives \(h_s^g(x)\).  They are the \(L^2\)-continuous
extension furnished by the semigroup, with the stated normalization
and local boundedness.

For a second extracted Gaussian law, condition its code over the same
GOE law of \(g\).  Its conditional samples pass the criterion on some
sequence by the same argument, so fixed-\(g\) uniqueness makes its
point mass the one already constructed, outside a common null set.
Thus every subsequential Gaussian limit is determined by \(g\).
The locally uniform correlation convergence and relaxation in
Proposition~\ref{prop:gaussian-subsequence} now hold for this
\(h^g\), and the subsequence criterion proves the asserted joint
Gaussian convergence.  Finally, undoing \eqref{entr:unit-time}
returns the evolution \(h_{s+t}^g=e^{-t c_*H_g}h_s^g\).
The only time substitution in this section was the temporary variable
\(u=c_*s\); the coefficient in \(S_t^g\) remains \(c_*\).
\qed

\section{Smooth truncation of clock histories}\label{sec:history-cutoffs}

The replacement of the coupling variables requires bounds for clock paths
whose duration is of order $T=n^{2/3}$.  We obtain them by capping transition
densities on every scale of a binary time tree.  This section constructs the
caps, proves that they discard negligible Gaussian mass, and records the
conditional path laws to which the score estimates will apply.

All times in Sections~\ref{sec:history-cutoffs}--\ref{sec:clock-scores} are
in site-clock units.  A stationary query at lag $t$ is one path of length
$tT$ started from $\mu$.  A quench query $(s,t)$ is one path started from
$\nu_n$, with consecutive legs of lengths $sT$ and $tT$.  Fix a compact set
of query times.  Every complete leg then has length in $[cT,CT]$, where
$0<c<C<\infty$ are fixed.  We use Taylor orders at most
$K=24$.  The constant $\beta>0$ will be chosen sufficiently small after
$K$.  The notation $a\lesssim_\beta b$ means
$a\le C_0 n^{C_1\beta}b$, with $C_1$ depending at most on $K$ and the fixed
query set, but not on $\beta$.  Fixed logarithmic factors can be included in
this convention.  Moment orders and the precision of probability estimates
may be increased without changing $C_1$; their constants and the lower bound
on $n$ may change.  For a vector indexed by
$\mathscr E_n=\{\{i,j\}:1\le i<j\le n\}$, put
\[
 |v|_p=\left(\frac1{\binom n2}\sum_{e\in\mathscr E_n}|v_e|^p\right)^{1/p},
 \qquad |v|_\infty=\max_{e\in\mathscr E_n}|v_e|.
\]

\subsection{A static support valid for every interaction matrix}

Write $\theta_{ij}=J_{ij}/\sqrt n$.  In entry replacement, an order-$k$
derivative in $\theta_e$ is multiplied by $n^{-k/2}$; summing over the
order-$n^2$ entries gives the factor $n^{2-k/2}$.  The Gaussian and
Rademacher laws match through order three.  The averaged fourth derivative
of the truncated trace tests must therefore tend to zero, while higher
derivatives may grow within the remaining negative power of $n$.  The
estimates below are organized around this fourth-order requirement;
$K=24$ supplies a global remainder bound.

The Gibbs probability $\mu$ is strictly positive for every finite matrix
$\theta$.  For independent spins
$X^1,\ldots,X^4$ define
\[
 \rho_2(X^a,X^b)=\frac1n\sum_iX_i^aX_i^b,
 \qquad
 \rho_4(X^1,X^2,X^3,X^4)=\frac1n\sum_iX_i^1X_i^2X_i^3X_i^4.
\]
Let $\mathscr D_n$ be the dyadic numbers from $1$ through a fixed multiple
of $CT$, including the first dyadic number beyond that multiple.  Set
\begin{equation}\label{hist:static-scales}
 r_2(d)=n^{2\beta}d^{-1/2},\quad r_4(d)=n^{10\beta}d^{-3/4},\quad
 a_2(d)=n^{3\beta}nd^{-3/2},\quad a_4(d)=n^{8\beta}nd^{-3/2}.
\end{equation}
For $r_2(d)<1$ let $A_2(d)=\{|\rho_2|>r_2(d)\}$.  For $r_4(d)<1$ let
\[
 A_4(d)=\left\{\max_{a<b}|\rho_2(X^a,X^b)|\le r_2(d),\quad
                    |\rho_4(X^1,X^2,X^3,X^4)|>r_4(d)\right\}.
\]
The displayed four-replica probability below is an unconditional
$\mu^{\otimes4}$ probability.

\begin{lemma}[Static support]\label{hist:static-support}
There is a smooth function $\chi_n(\theta)\in[0,1]$ with the following
properties.  On its support, and on the support of any of its derivatives
of order at most $K$, the matrix satisfies
\begin{equation}\label{hist:static-conditions}
 \begin{gathered}
 \mu^{\otimes2}(A_2(d))\le e^{-a_2(d)},\qquad
 \mu^{\otimes4}(A_4(d))\le e^{-a_4(d)},\qquad d\in\mathscr D_n,\\
 n^{-6}\sum_{i<j}\theta_{ij}^2\le2.
 \end{gathered}
\end{equation}
An assertion involving an omitted event is void.  For some $b=b(\beta)>0$,
every derivative of $\chi_n$ of positive order at most $K$ in one coefficient
$\theta_e$, and every differentiated factor in its product construction,
is $O(n^{-b})$, uniformly in $e$ and $\theta$.  For Gaussian couplings,
$\chi_n=1$ with probability tending to one.
\end{lemma}

\begin{proof}
Choose a fixed smooth increasing function $f$ that is zero on $(-\infty,1]$
and one on $[2,\infty)$.  For each nonempty $A_j(d)$ use the factor
$f(-\log\mu^{\otimes j}(A_j(d))/a_j(d))$, and omit empty events.  Multiply
these factors and a smooth function of $n^{-6}\sum_e\theta_e^2$ that is one
below $1$ and zero above $2$.  This gives \eqref{hist:static-conditions} and
a plateau equal to one when each probability there is at most $e^{-2a_j(d)}$
and the size statistic is at most one.

To check derivatives, the $k$th derivative of the logarithm of a nonempty
replica probability in $\theta_e$ is the difference of two cumulants of
$\sum_{a=1}^jX_i^aX_j^a$, under the restricted and unrestricted replica
probabilities.  The variable is bounded by $j\le4$, so every fixed-order
cumulant is bounded by a constant depending only on its order.  Since
$a_2(d)\ge c_0 n^{3\beta}$ and $a_4(d)\ge c_0 n^{8\beta}$ on
$\mathscr D_n$, every derivative of positive order of a replica factor gains at
least $n^{-3\beta}$.  Derivatives of the size factor gain at least $n^{-3}$:
where a derivative is nonzero, $|\theta_e|\le\sqrt2 n^3$.  There are
$O(\log n)$ factors.  The product rule now proves the stated bound, for
example with any sufficiently small $b<\min(3\beta,3)$.

We prove the Gaussian assertion with the diagonal and eigenvector sign
randomization fixed as in Section~\ref{sec:edge-estimates}.  On spectral
events of limiting probability as close to one as desired,
Lemma~\ref{edge:spherical-tails} gives, for the uncapped zero-field spherical
law and a sufficiently small fixed $\rho_0>0$,
\[
 \Prob_{\rm sph}^{\otimes2}\{ |\rho_2|>w\}
       \le C e^{-c_0nw^3},\qquad C_0p\le w\le\rho_0,
\]
where one can take any permitted reference scale with $np^3$ above the
tight threshold.  The smallest pair threshold $r_2(d)$ in
\eqref{hist:static-scales} is $\gg n^{-1/3}$, so such a $p$ is available
below all the nonvacuous shrinking pair thresholds.  In a band
$z<|\rho_2|\le2z$, truncated at $\rho_0$, we apply the weighted
two-replica overlap-cell consequence of Lemma~\ref{edge:gram-cell}.
On the retained event $\|W\|\le3$, the logarithm of the positive zero-field
angular weight has overlap derivative $O(n)$, uniformly over the
orthonormal two-frame.  It therefore oscillates by at most a fixed constant
on each cell of width $2/n$.  Summing the weighted cell inequality over the
band enlarges its spherical endpoints by at most $1/n$ and costs at most
$\exp(Cnz^4+C\log n)$ after increasing $C$.  For the pair thresholds $w$
at issue, $z\ge w\gg n^{-1/3}$, so the enlargement is absorbed by
decreasing the tail constant $c_0$.  Shrink $\rho_0$ so that this
cost is bounded by $\exp(c_0nz^3/2+C\log n)$.  Summing the dyadic bands from
$w$ to $\rho_0$ therefore bounds the Haar mean of the restricted cube
partition function, divided by the squared spherical partition function,
by $\exp(-c_1nw^3)$ for all the thresholds in question.  Here
$nw^3\ge c n^{6\beta}$ absorbs $O(\log n)$.  The likelihood ratio
$L_\circ=Z_{\rm cube}/Z_{\rm sp}$ is at least $1/2$ with probability
tending to one on these restrictions, by Equation~\eqref{edge:haar-bias}.  Markov's
inequality consequently proves the pair plateau, because
\[
 \frac{nr_2(d)^3}{a_2(d)}=n^{3\beta}\longrightarrow\infty.
\]
For overlaps bounded away from zero, the fixed-overlap estimate in
\cite[Lemma~8.2]{FUA} gives an $e^{-c n}$ tail with Gaussian probability tending to
one.  This pays $2a_2(d)=2n^{1-3\beta}r_2(d)^3=o(n)$ for any such
nonvacuous threshold.  %

Here is the corresponding four-replica calculation.  If $r_4(d)<1$, then
$r_2(d)\le n^{-14\beta/3}=o(1)$ and
$r_2(d)^2/r_4(d)\le n^{-6\beta}$.  For one row of four independent uniform
signs, let $Y\in\{-1,1\}^6$ be its six pair products and let $H$ be its
four-product.  For $n$ independent rows write
$\bar Y=n^{-1}\sum_iY_i$ and $\bar H=n^{-1}\sum_iH_i$.  Put
$\Lambda(v)=\log\E e^{v\cdot Y}$, and let $\Prob_v$ be the row probability
with density $e^{v\cdot Y-\Lambda(v)}$, with expectation $\E_v$.  For
several rows the same notation denotes its product probability.  Around the origin,
\[
 \Lambda(v)=\tfrac12|v|^2+O(|v|^3),\qquad
 \E_v H=O(|v|^2).
\]
Indeed the six coordinates of $Y$ are centered and have identity covariance, while
$\E H=\E(HY)=0$; Taylor's theorem on this finite space gives both uniform
remainders.  If the empirical pair vector is the lattice value $u$ with
$|u|\le\sqrt6 r_2(d)$, put $E_d=\{|\bar H|>r_4(d)\}$.  The change of
measure gives
\[
 \Prob_{\rm cube}^{\otimes4}\{\bar Y=u,\ E_d\}
 =e^{-n|u|^2+n\Lambda(u)}
       \Prob_u\{\bar Y=u,\ E_d\}
 \le e^{-n|u|^2+n\Lambda(u)}\Prob_u(E_d).
\]
Under $\Prob_u$ the rows remain independent.  Since
$|\E_uH|\le r_4(d)/2$ for large $n$, Hoeffding's inequality
\cite{Hoeffding1963} bounds the last probability by
$2e^{-nr_4(d)^2/8}$.  Together with the expansion of $\Lambda$, this gives
\begin{equation}\label{hist:cube-quartet}
 \Prob_{\rm cube}^{\otimes4}\{\bar Y=u,\ |\bar H|>r_4(d)\}
 \le 2\exp\{-\tfrac n2|u|^2+Cn|u|^3-cn r_4(d)^2\}.
\end{equation}

For completeness, the reverse Gram cost is also explicit.  Put
$G_{aa}=1$ and $G_{ab}=u_{ab}$.  The Gram density of four independent
uniform points of the sphere of radius $\sqrt n$ is a polynomial
normalization times $\det(G)^{(n-5)/2}$ on positive correlation matrices.
For $|u|=o(1)$,
$\log\det G=-|u|^2+O(|u|^3)$.  A cell of width $c/n$ about $G$ thus has
probability at least
\[
 \exp\{-\tfrac n2|u|^2-Cn|u|^3-C\log n\}.
\]
At a fixed spectrum with bounded norm, let $F\in\R^{n\times4}$ be a uniform
orthonormal frame and put
\[
 \mathcal A_\lambda(G)=Z_{\rm sp}^{-4}\E_F
      \exp\{(n/2)\Tr(F^{\mathsf T}\diag(\lambda)F G)\}.
\]
The exponent has Lipschitz constant $O(n)$ in $G$, so the integrand changes
by at most a factor $e^{O(1)}$ across the cell.  The integral of
$\mathcal A_\lambda$ against the spherical Gram probability equals one.
Restricting it to the cell therefore bounds the Haar mean at Gram $G$ by
\[
 \exp\{\tfrac n2|u|^2+Cn|u|^3+C\log n\}.
\]
Multiply this by \eqref{hist:cube-quartet} and sum the at most $(n+1)^6$
possible Grams.  The remaining error $Cnr_2(d)^3+C\log n$ is negligible
relative to $nr_4(d)^2$, since
\[
 nr_2(d)^3=n^{6\beta}nd^{-3/2},\qquad
 nr_4(d)^2=n^{20\beta}nd^{-3/2}.
\]
On $L_\circ\ge1/2$, division by the cube partition function changes the
four-replica numerator by at most $16$.  Markov's inequality proves its
plateau $e^{-2a_4(d)}$, since $nr_4(d)^2/a_4(d)=n^{12\beta}$.  A union over
$O(\log n)$ scales is harmless.  Finally, the Gaussian mean of
$n^{-6}\sum_e\theta_e^2$ is $O(n^{-5})$.  The size plateau follows from
Markov's inequality.  Enlarging the spectral restrictions completes the
proof that $\chi_n=1$ with probability tending to one.
\end{proof}

Only this last verification used Gaussian frame averages.  In the remainder
of this section and in the score estimates, $\theta$ is any deterministic
matrix satisfying \eqref{hist:static-conditions}.

\subsection{The tree kernels and their survival probabilities}

For each leg $I=[a,b]$, repeatedly bisect time intervals until their lengths
belong to $[1,2]$.  (Changing the harmless endpoint convention when a length
is exactly $2$ leaves all statements below unchanged.)  A node and its two
children are denoted by $I$ and $I^-,I^+$, in chronological order.  Define
\begin{equation}\label{hist:node-level}
 N_I=N(|I|),\qquad N(d)=n^\beta(T/d)^{3/2}=n^{1+\beta}d^{-3/2}.
\end{equation}
All these levels lie between $c_0n^\beta$ and $C_0n^{1+\beta}$.

\begin{lemma}[A smooth cap]\label{hist:cutoff-hazards}
For every level in \eqref{hist:node-level} there is a smooth
$W_N:\R\to(0,1)$ such that
\begin{equation}\label{hist:cap-and-plateau}
 e^lW_N(l)\le e^{CN},\qquad
 1-W_N(l)\le e^{-(\log n)^2/2}\quad(l\le N/2).
\end{equation}
For $1\le k\le K$ its derivative in $l$ satisfies globally
\begin{equation}\label{hist:global-ratios}
 \frac{|W_N^{(k)}|}{W_N}+\frac{|W_N^{(k)}|}{1-W_N}\le C_K,
 \qquad
 \frac{|W_N^{(k)}|}{1-W_N}\le\frac{C_K(\log n)^{2K}}N.
\end{equation}
For each fixed $L$, wherever $W_N(l)\ge n^{-L}$ one also has
\begin{equation}\label{hist:survival-hazard}
 \frac{|W_N^{(k)}(l)|}{W_N(l)}
 \le \frac{C_{K,L}(\log n)^{2K}}N\bigl(1-W_N(l)\bigr).
\end{equation}
Moreover, uniformly over these levels and for sufficiently large $n$,
\begin{equation}\label{hist:log-cap-growth}
 |\log W_N(l)|+|\log(1-W_N(l))|
       \le C\bigl(1+N+|l|+(\log n)^2\bigr).
\end{equation}
\end{lemma}

\begin{proof}
Let $\psi$ be smooth and nondecreasing, zero on $(-\infty,2]$, with
$0\le\psi'\le1$ and $\psi'=1$ on $[3,\infty)$.  Put
\[
 a_N(z)=(\log n)^2(z-1)+3N\psi(z),\qquad
 W_N(l)=\bigl(1+e^{a_N(l/N)}\bigr)^{-1}.
\]
For $l\le3N$ the first inequality in \eqref{hist:cap-and-plateau} is
immediate.  For $z=l/N\ge3$, $\psi(z)\ge z-3$, so
$l-a_N(z)\le3N$; the same inequality follows from $W_N\le e^{-a_N}$.
If $l\le N/2$, then $a_N(l/N)\le-(\log n)^2/2$, proving the plateau.

All $l$-derivatives of $a_N(l/N)$ of positive order at most $K$ are bounded:
the first is $(\log n)^2/N+3\psi'(l/N)$ and the others are
$3N^{1-k}\psi^{(k)}(l/N)$.  Derivatives of the logistic function have a
factor $W_N(1-W_N)$ times a bounded polynomial in $W_N$, which proves the
first part of \eqref{hist:global-ratios}.  If $W_N\ge n^{-L}$, then
$a_N\le L\log n+O(1)<(\log n)^2$; hence $l/N<2$ for large $n$.  In this
region $a_N$ is affine, and its slope in $l$ is $(\log n)^2/N$.
The same derivative formula proves \eqref{hist:survival-hazard}.  It also
proves the second part of \eqref{hist:global-ratios} for $l/N\le2$.  In the
remaining region $W_N\le e^{-(\log n)^2}$, so the first global bound with
its logistic factor gives the second bound there as well, since $N$ is
polynomial in $n$.  Finally the two logarithms are bounded by
$C(1+|a_N(l/N)|)$.  Since $\psi$ is Lipschitz and
$(\log n)^2/N\le1$ for large $n$, this proves \eqref{hist:log-cap-growth}.
\end{proof}

Let $p_d(x,y)=P_d(x,y)/\mu(y)$ be the true transition density.  Define
the pre-density $\bar k_I$, the survival factor $w_I$, and the post-density
$k_I$, all relative to the normalized probability $\mu$, by
\begin{equation}\label{hist:kernel-recursion}
 \begin{aligned}
 \bar k_I(x,y)&=p_{|I|}(x,y) &&\text{if $I$ is a leaf},\\
 \bar k_I(x,y)&=\sum_z\mu(z)k_{I^-}(x,z)k_{I^+}(z,y)
                                      &&\text{otherwise},\\
 w_I(x,y)&=W_{N_I}(\log\bar k_I(x,y)),\qquad
 k_I(x,y)=\bar k_I(x,y)w_I(x,y).
 \end{aligned}
\end{equation}
All these numbers are positive.  In particular, the logarithms are defined.
The transition $k_I(x,y)\mu(y)$ is a subprobability.  At a fixed leg we
may equivalently sample an ordinary marked heat-bath path and retain it with
the product of its $w_I(X_a,X_b)$ over all nodes.  Apply checks at their
right endpoints, ordering equal right endpoints from the smallest interval
to the largest.  The equivalence follows inductively from the Markov
property and \eqref{hist:kernel-recursion}.

\begin{lemma}[Deterministic kernel consequences]\label{hist:kernel-caps}
For every node $I$ of length $d$,
\begin{equation}\label{hist:kernel-bounds}
 0<k_I(x,y)\le\bar k_I(x,y)\le p_d(x,y),\qquad k_I(x,y)\le e^{CN(d)}.
\end{equation}
Both $k_I$ and $\bar k_I$ are symmetric in $(x,y)$.  If an interval of
length $h$ lies inside a checked tree, with $h$ larger than a fixed constant,
the weighted transition across it has density at most $e^{C'N(h)}$ after
all weights other than one contained complete node are discarded.  The same
assertion holds when the top check of the tree is omitted or is a failure
factor $1-w_I$.
\end{lemma}

\begin{proof}
Positivity and domination by $p_d$ follow by induction and the semigroup
property of $P_d$.  The cap is \eqref{hist:cap-and-plateau}.  At a split the
two children have equal lengths and use the same rule, so their kernels are
identical.  The convolution of this symmetric kernel with itself is
symmetric; multiplication by the endpoint function $W_N(\log\bar k_I)$
preserves symmetry.  This starts from detailed balance of $p_d$.

Every subinterval of length $h$ larger than a fixed constant contains a
tree node $J$ of length at least $c_1h$.  One may require $J$ to be a proper
descendant at the cost of reducing $c_1$.  Keep the checks in $J$ and discard
the other weights, all of which lie in $[0,1]$.  The remaining transition
is an ordinary kernel, followed by $k_J$, followed by an ordinary kernel.
Since ordinary kernels preserve $\mu$ and constants, its density is at most
$e^{CN(|J|)}\le e^{C'N(h)}$.  A discarded failure or top factor also lies
in $[0,1]$.
\end{proof}

\begin{lemma}[Gaussian survival]\label{hist:gaussian-survival}
For every $A<\infty$ and $\varepsilon>0$, there are Gaussian disorder
events of limiting lower probability at least $1-\varepsilon$ on which
$\chi_n=1$ and the mass lost by the checks in any complete query is at most
$n^{-A}$.  This bound is uniform over its query times in the fixed compact
set and over its initial probability on the cube.
\end{lemma}

\begin{proof}
The log-moment assertion of Theorem~\ref{thm:warming} says that, for each fixed
integer $p\ge1$ and sufficiently small fixed $\xi>0$, on disorder restrictions
of arbitrarily high limiting probability,
\[
 \sup_x\left(1+\sum_yP_d(x,y)\log_+^p p_d(x,y)\right)^{1/p}
 \le C_{p,\xi}\left(1+(T/d)^{1/(2/3-\xi)}\right),\quad 1\le d\le CT.
\]
Choose $\xi$ so that $\xi/(2/3-\xi)<\beta/2$.  Division by $N(d)$ shows
that the right side is at most $C_{p,\xi}n^{-\beta/2}N(d)$ uniformly in
this range: the largest possible extra factor is
$T^{1/(2/3-\xi)-3/2}=n^{\xi/(2/3-\xi)}$.  Markov's inequality therefore
gives
\[
 \sup_xP_d(x,\{y:\log p_d(x,y)>N(d)/2\})
       \le C_{p,\xi}n^{-p\beta/2}.
\]
By \eqref{hist:kernel-bounds}, $\bar k_I\le p_{|I|}$.  At each check, its
failure probability under an ordinary path is thus bounded by the last
display plus $e^{-(\log n)^2/2}$.  The probability of a first failure is
at most the sum of these probabilities, because all preceding weights are
at most one.  There are $O(T)$ nodes.  Choose $p$ large after $A$ and use
Lemma~\ref{hist:static-support} for the plateau.  All estimates were uniform
in the starting state and in $d$, which proves the asserted uniformity.
\end{proof}

\subsection{Tested histories and their conditional interiors}

We record precisely which changes of path measure will be used.  An
ordinary marked path includes the time, site, and new sign at every ring,
including a resampling that keeps the current sign.  Its law on a time
interval $I$ from $x$ is denoted by $\mathsf P_x^I$.

\begin{definition}[Tested path law]\label{hist:tested-law}
A checked history is one of the following subprobabilities, with initial
probability either $\mu$ or the uniform cube probability:
\begin{enumerate}
\item[(i)] an entire query path with every survival factor in its legs;
\item[(ii)] the prefix through a specified first failing check $F$, with all
preceding survival factors and the factor $1-w_F$ at that check.
\end{enumerate}
The check order is the one following \eqref{hist:kernel-recursion}; the
choice of $F$ and all tree shapes is fixed independently of the path and
of $e$.  In case (ii), the maximal checked intervals of the prefix form a
forest: the last is $F$, whose descendants survive, and every preceding
root and its descendants survive.  Designate the initial state, all leg
boundaries, and all boundaries of these maximal roots as its external
endpoints, retaining only boundaries present in a stopped prefix.  Multiply
the history by a nonnegative function $\Phi$ of these
endpoints with $\|\Phi\|_\infty\le n^{M_+}$, where $M_+$ is fixed and
$\Phi$ is the same for every $e\in\mathscr E_n$.  When the resulting mass
$m_\Phi$ is at least $n^{-M_-}$ for a fixed $M_-$, its normalization is a
tested path law, denoted by $\mathsf Q_\Phi$.

Products of a fixed number of such histories use independent ordinary
paths and a separate external test on each history.  Signed tests are
decomposed into positive and negative parts separately.  A statement that
holds with arbitrarily high inverse-polynomial probability means that for
every fixed $A$ its exception has $\mathsf Q_\Phi$ probability at most
$n^{-A}$ for all sufficiently large $n$, uniformly in the fixed choices
above and in disorder satisfying \eqref{hist:static-conditions}.
\end{definition}

There are at most polynomially many nodes.  In the last, possibly partial,
leg of a first-failure history, its preceding maximal roots are the earlier
siblings encountered in the binary subdivision leading to $F$.  In their
chronological order their lengths are nonincreasing and the last such
length is at least $|F|$.  Complete earlier legs have lengths comparable
to $T$.  These facts also show that earlier maximal roots have lengths at
least a fixed multiple of later ones, with the multiple depending only on
the query set.  The density of $\mathsf Q_\Phi$ relative to the ordinary
forward law with the same initial probability is at most
$n^{M_++M_-}$; future endpoints are not revealed in that forward law's
natural filtration.

For a node $I=[a,b]$, write
$V_I=\prod_{J\subsetneq I}w_J(X_{\partial J})$, with the product over its
proper descendants.  The normalized interior law with boundary $(x,y)$ is
\begin{equation}\label{hist:bridge-definition}
 \mathsf C_I^{x,y}(d\omega)
 =\frac{V_I(\omega)\ind_{\{X_b=y\}}}{\mu(y)\bar k_I(x,y)}
             \mathsf P_x^I(d\omega).
\end{equation}
The denominator equals the numerator's integral by the kernel recursion.
This definition excludes the top check.  A top survival, failure, or omitted
check is constant after its boundary is fixed.

\begin{lemma}[Conditional interiors and prefix bounds]\label{hist:conditional-interior}
Let $I$ be a fixed node of a checked root in a tested path law.  Given the marked
history outside the interior of $I$ and its two boundary states, the
conditional law inside $I$ is $\mathsf C_I^{X_a,X_b}$.  In particular,
resampling this interior independently from that law preserves the tested
marginal.  For an internal node its midpoint $z$ under $\mathsf C_I^{x,y}$
has probability
\begin{equation}\label{hist:midpoint-probability}
 \pi_I^{x,y}(z)=
   \frac{\mu(z)k_{I^-}(x,z)k_{I^+}(z,y)}{\bar k_I(x,y)},
\end{equation}
and conditional on $z$ the two interiors are independent with laws
$\mathsf C_{I^-}^{x,z}$ and $\mathsf C_{I^+}^{z,y}$.

For every $A$, there is $L=L(A,M_+,M_-)<\infty$ such that
$\bar k_I(X_a,X_b)\ge n^{-L}$ outside an endpoint event of probability
$n^{-A}$ under any such tested law.  If $d=|I|$ exceeds a fixed constant,
then for these boundaries the density of its first two thirds under
$\mathsf C_I^{x,y}$ relative to the ordinary unconditioned prefix from $x$
is at most
\begin{equation}\label{hist:prefix-density}
 n^L e^{C'N(d)}.
\end{equation}
The reversed last two thirds have the same bound relative to the ordinary
prefix from $y$ in reversed time.
\end{lemma}

\begin{proof}
The tree is laminar: every check whose interior intersects the interior of
$I$ is a descendant or an ancestor of $I$.  Disjoint siblings may share an
endpoint, but their factors are fixed by the exterior and boundary data.
The ancestor factors, the top factor, and the external test are fixed by
the same data.  The remaining factors are exactly
$V_I$.  The Markov property gives \eqref{hist:bridge-definition} as the
conditional law, and then gives \eqref{hist:midpoint-probability} and the
conditional independence.  This proves the resampling assertion as well.
It applies to tree nodes; an arbitrary subinterval can have crossing checks.

At a fixed start $x$,
\[
 \sum_y\mu(y)\bar k_I(x,y)\ind_{\{\bar k_I(x,y)<\delta\}}\le\delta.
\]
Integrate the node first in the tested history and discard all other weights.
The bound on $\Phi$ and the lower bound on its mass show
$\mathsf Q_\Phi\{\bar k_I(X_a,X_b)<\delta\}\le n^{M_++M_-}\delta$.
This proves the endpoint assertion by choosing $\delta=n^{-L}$.

Let $\mathcal F_t^\to$ contain just the marked prefix of the node from $a$
to $t$.  For $t-a\le2d/3$, its conditional density is exactly
\[
 \frac{d(\mathsf C_I^{x,y}|_{\mathcal F_t^\to})}
      {d(\mathsf P_x^I|_{\mathcal F_t^\to})}
 =\frac{\mathsf E_x[V_I\ind_{\{X_b=y\}}\mid\mathcal F_t^\to]}
        {\mu(y)\bar k_I(x,y)}.
\]
The interval remaining after $t$ contains a proper node of length at least
$c_1d$.  Keep its checks and discard all others in the numerator.  By
Lemma~\ref{hist:kernel-caps}, that numerator is at most $\mu(y)e^{C'N(d)}$.
This proves \eqref{hist:prefix-density}.  Detailed balance for marked paths
includes self-resampling rings, and the proper-descendant weights are
invariant under reversal because equal-length children use identical
symmetric kernels.  Reversal therefore sends $\mathsf C_I^{x,y}$ to
$\mathsf C_I^{y,x}$ on the reflected node, proving the last assertion.
No reversal of the entire stopped forest is needed.
\end{proof}

\begin{lemma}[Overlaps along a checked history]\label{hist:history-overlaps}
With arbitrarily high inverse-polynomial probability under a tested path
law, simultaneously for times in any one of its checked root intervals,
\begin{equation}\label{hist:path-overlaps}
 |\rho_2(X(t_1),X(t_2))|\lesssim_\beta(1+m)^{-1/2},\qquad
 |\rho_4(X(t_1),X(t_2),X(t_3),X(t_4))|
                    \lesssim_\beta(1+m)^{-3/4}.
\end{equation}
In each assertion the times are increasingly ordered, and $m$ is the
minimum of their consecutive gaps and of the gap from the root's start to
$t_1$.  The assertion includes a failing root.  It also holds for the
full history obtained by each of a fixed finite collection of interior
resamplings in Lemma~\ref{hist:conditional-interior}, separately within
each such history.  All times in an overlap are taken from that one history.
\end{lemma}

\begin{proof}
First use deterministic times with $m$ larger than a fixed constant.
Inside each gap choose a contained proper checked node of length comparable
to $m$.  They are disjoint.  Dropping all other weights and applying
Lemma~\ref{hist:kernel-caps} shows that the joint law of the observed states,
given any initial state, is bounded by their product $\mu$ law times
$\exp(C_2N(m))$.  This estimate counts the retained weights without
conditioning on survival.  Apply \eqref{hist:static-conditions} at a
dyadic scale comparable to $m$, first to all relevant pairs, and then to
the quartet on the event that its pairs obey the same threshold.  The
ratios $a_2(m)/N(m)=n^{2\beta}$ and $a_4(m)/N(m)=n^{7\beta}$ tend to
infinity.  Their exponents are at least a positive power of $n$, so they
pay the cap and any polynomial test density with any prescribed
inverse-polynomial precision.  A vacuous threshold or bounded $m$ gives
the bound directly from $|\rho_j|\le1$.

Use a time mesh of sufficiently small inverse-polynomial spacing.  The
probability that any mesh cell has two rings is arbitrarily small after
paying the polynomial test density: under ordinary paths the total clock
rate is $n$, and the union bound is $O(n^2T)$ times the spacing.  A single
ring changes either overlap by at most $2/n$ per affected time.  This is
smaller than the displayed thresholds, whose smallest quartet value has
order $n^{-1/2+10\beta}$.  A union over the polynomially many grid tuples
proves the simultaneous assertion.  A resampled interior has the same
tested marginal by Lemma~\ref{hist:conditional-interior}; a union proves the
last claim for any fixed number of copies.
\end{proof}

\subsection{Local differentiability and normalizations}

We finish by recording a coarse derivative bound, used to control exceptional
histories.  For $e=\{i,j\}$ write $q_e(x)=x_ix_j$ and perturb the interaction
$\theta_e$ by $\vartheta$, putting $u=\tanh\vartheta$.  The new stationary
probability is exactly
\begin{equation}\label{hist:stationary-perturbation}
 \mu_u(x)=\mu(x)\frac{1+u q_e(x)}{1+u\mu q_e},\qquad |u|<1.
\end{equation}
Let $t_i(x)=\mu[X_i\mid X_{-i}=x_{-i}]$.  At a ring of site $i$ with new
sign $\zeta$, the likelihood ratio of the perturbed mark is
\begin{equation}\label{hist:ring-likelihood}
 M_{e,i,\zeta}(u)=\frac{1+u\zeta x_j}{1+u x_jt_i(x)};
\end{equation}
interchange $i,j$ at the other site, and use one at all remaining sites.
The site rates stay equal to one.  Thus the full path likelihood
$L_{I,e}(u)$ is the product of \eqref{hist:ring-likelihood} at its rings.
The formula includes a mark that retains the current sign.

Superscript $u$ on $\bar k_I,k_I,w_I$ denotes the construction
\eqref{hist:kernel-recursion} for the perturbed dynamics and its normalized
stationary probability $\mu_u$.  Throughout, $[u^l]$ means the $l$th Taylor
coefficient at zero (the derivative divided by $l!$).  For the pre-density,
the exact change of measure is
\begin{equation}\label{hist:pre-ratio}
 \frac{\bar k_I^u(x,y)}{\bar k_I(x,y)}
 = (1+u\mu q_e)\,
 \mathsf C_I^{x,y}\left[
    \frac{L_{I,e}(u)}{1+u q_e(y)}
      \prod_{J\subsetneq I}\frac{w_J^u(X_{\partial J})}{w_J(X_{\partial J})}
                         \right].
\end{equation}
The product is over all proper descendants; it contains only survival
factors.  In particular $1+u\mu q_e$ is required by normalization of the
perturbed stationary probability.

\begin{lemma}[Coarse coefficients and moments]\label{hist:crude-ratios}
On the size support in \eqref{hist:static-conditions}, for all nodes,
boundaries and $e$, and $1\le l\le K$,
\[
 \left|[u^l]\frac{\bar k_I^u}{\bar k_I}\right|
 +\left|[u^l]\frac{k_I^u}{k_I}\right|
 +\left|[u^l]\frac{w_I^u}{w_I}\right|
 +\left|[u^l]\frac{1-w_I^u}{1-w_I}\right|\le n^{C_K}.
\]
The exponent $C_K$ is independent of $\beta$ in a fixed sufficiently small
range and of any later number of conditional copies.  Absolute fixed-order
moments of clock counts and raw path coefficients are polynomial under
tested laws and their resampled marginals.  Conditional on an arbitrary
boundary pair of one node, the same moments under $\mathsf C_I^{x,y}$ are
bounded by a polynomial (with exponent allowed to depend on the moment
order).  Finally, conditional on arbitrary designated external endpoints,
the absolute normalized coefficient through total order $K$ of an entire
checked or first-failure history is bounded by $n^{C_K}$, after enlarging
$C_K$.  Thus its unnormalized coefficient with any nonnegative external
endpoint test held fixed under the interaction perturbation is bounded by
$n^{C_K}$ times the test's base mass, however
small that mass is.
\end{lemma}

\begin{proof}
The size support gives $\sum_e|\theta_e|\le Cn^4$, hence
$\min_x\mu(x)\ge e^{-Cn^4}$.  By Cauchy--Schwarz, each absolute one-site
field is at most $Cn^{7/2}$.
On a leaf, prescribe one ring of each site in successive subintervals of
length $1/n$, with its sign set to the desired endpoint, and no other rings.
The clock event has probability $e^{-n|I|}n^{-n}$, and each prescribed mark
has probability at least $\tfrac12e^{-Cn^{7/2}}$.  Thus every leaf transition
probability is at least $e^{-Cn^5}$, uniformly in its endpoints.  Its
logarithmic density has absolute value at most $Cn^5$.  The last assertion of
Lemma~\ref{hist:cutoff-hazards} and the kernel recursion propagate this property
up the tree: composing two positive kernels costs a fixed multiple of the
previous logarithmic bound plus a polynomial, and applying $W_N$ has the
same property.  The depth is $O(\log n)$, so all these absolute logarithms
remain polynomial, with a fixed exponent.  In particular
$\mu(y)\bar k_I(x,y)\ge e^{-n^{C_0}}$ for a fixed $C_0$.

The total number of ordinary rings on a node is Poisson with parameter
$n|I|\le CnT$.  Its tail above any sufficiently large polynomial is smaller
than $e^{-2n^{C_0}}$ times any inverse power of $n$.  Dividing by the last
lower bound and using $V_I\le1$ proves polynomial conditional count moments
for every boundary.  Each fixed Taylor coefficient of the ring product is
a polynomial of fixed degree in its number of rings, since each single-ring
coefficient in \eqref{hist:ring-likelihood} is bounded.  This proves the
claimed raw moments, both conditionally and under the polynomially tilted
tested laws.  The copy statement follows from their preserved marginals.

At a leaf, \eqref{hist:pre-ratio} has an empty product; the preceding count
bound and the bounded coefficients of $(1+u\mu q_e)/(1+u q_e(y))$ give
the coarse bound for its pre-density.  Suppose normalized coefficients of
order $l\le K$ at one tree level are bounded by $D^l$.  For a composition,
division by its base value writes its ratio as a conditional average of the
product of the two child ratios and $\mu_u(z)/\mu(z)$.  Its order-$l$
coefficient is therefore at most $(C_K(D+1))^l$.  Taylor's chain rule for
$W_N(\log\bar k_I^u)$ uses the global ratios
\eqref{hist:global-ratios}; coefficients of the logarithm are polynomials
in nonconstant pre-ratio coefficients, each product having total order $l$.
The same bound applies to either survival or failure ratios and to the
post-density.  Starting with a polynomial $D$ at the leaves and iterating
$D\mapsto C_K(D+1)$ for $O(\log n)$ levels proves these coefficient bounds.  Formula
\eqref{hist:pre-ratio} itself follows by inserting the path likelihood and
then dividing by $\mu_u(y)/\mu(y)$ in the definition of the pre-density.

To check the last forest assertion, fix its designated external endpoints.
The interiors of its maximal roots are the corresponding conditional laws
\eqref{hist:bridge-definition}, and their top survival or failure factors
are fixed.  Expand a coefficient of total order at most $K$ by the product
rule.  A differentiated node contributes one of the normalized ratios just
bounded; the raw ring factors have the conditional polynomial moments
proved above.  There are polynomially many nodes and at most $K$ positive
orders, so the number and magnitude of the resulting terms are bounded by
$n^{C_K}$.  Stationary initialization, if present, contributes only the
bounded coefficients in \eqref{hist:stationary-perturbation}.  Integrating
this conditional bound against the nonnegative endpoint test and its base
endpoint measure proves the asserted factor of the base mass.
\end{proof}

Here and below, estimates on a set of arbitrarily small inverse-polynomial
probability can be used inside conditional averages.  The precise elementary
rule is as follows.  If $\mathsf Q(E^c)\le n^{-A'}$ and
$\|Y\|_{L^p(\mathsf Q)}\le n^C$ for some $p>1$, then for any boundary
$\sigma$-field $\mathcal B$,
\[
 \mathsf Q\left\{\mathsf E_{\mathsf Q}
       [|Y|\ind_{E^c}\mid\mathcal B]>n^{-A}\right\}
 \le n^{A+C-(1-1/p)A'}.
\]
Choose $A'$ after $A$ to make this as small as required.  The coarse bounds
of Lemma~\ref{hist:crude-ratios} supply these moments for all coefficients used
below; a fixed number of copies is handled by H\"older's inequality.
For a tested mass below a chosen inverse-polynomial precision, use instead
the final forest assertion of Lemma~\ref{hist:crude-ratios} before
normalization: its unnormalized coefficient is at most $n^{C_K}$ times
that mass.  Polynomially many deterministic nodes,
grid points, or choices of any fixed number of copies are handled by a union
bound.  These conventions never give a uniform assertion for every boundary
pair when the sharper estimate was proved only for typical boundaries.

\section{Projecting the clock scores}\label{sec:clock-scores}

We now prove the derivative estimates for the stopped histories.  The
ordinary likelihood coefficients of a path of length $D$ have size
$D^{l/2}$ at order $l$.  We gain in the contribution of a raw coefficient
to a tested expectation by compensating its last ring before taking absolute
values.  For the sharper estimates, we first average an earlier raw block
conditional on its endpoints while later insertions are still raw; those
endpoints are then in the past of the later block in the chosen orientation.
A drift obtained after conditioning on a future endpoint is not thereby
predictable for a smaller filtration.

Throughout this section the interaction matrix is fixed on
\eqref{hist:static-conditions}, and all weights and tested laws are at this
base matrix.  They are common to every $e\in\mathscr E_n$.  Only the ring
likelihood $L_{I,e}(u)$ varies when a clock coefficient is formed.  Write
\[
 \ell_{l,e}(I)=[u^l]L_{I,e}(u),\qquad \ell_{0,e}(I)=1.
\]
Superscripts $\to$ and $\leftarrow$ specify forward and reversed time
orientation.  The reversed object records the reversed marks, including
the old sign as the reversed new sign.
Whenever a bound written $\lesssim_\beta$ absorbs a factor $n^\eta$ from
a path estimate, choose $\eta$ beforehand as a sufficiently small fixed
multiple of $\beta$.  Only a number of such losses depending on $K$ occurs
in this section.

\subsection{Changes of path law and their compensators}

We will use two kinds of oriented probability law.  The first is a tested
path law in its natural forward filtration.  This filtration contains the
initial state and the marked past; future external endpoints are still
unrevealed.  The second is $\mathsf C_I^{x,y}$ on its first two thirds in
the forward orientation, or on its last two thirds in the reversed
orientation, for fixed $x,y$ satisfying the endpoint bound in
Lemma~\ref{hist:conditional-interior}.  Its reference probability is the
ordinary \emph{unconditioned} prefix from its initial endpoint.  The other
endpoint is a fixed parameter, and no additional future path information
is revealed.  An independent history, or a conditional copy sampled using
only already revealed data and independently of the subsequent path, can
be included among the information available at the start.  These will be
called the admissible oriented laws.
When such auxiliary data are adjoined, use the same conditional sampling
kernel under the reference and changed laws, independently of future rings
given the information already present.  This leaves both the reference
intensities and the likelihood ratio unchanged.

A bridge expectation in an identity is taken at its fixed boundary pair.
Every high-probability assertion about its pathwise bounds, however, is
with respect to the tested law of the boundary followed by that conditional
interior.  It is not a uniform conditional assertion for all pairs with
$\bar k_I(x,y)\ge n^{-L}$: the overlap and count events are inherited under
this joint law.  This convention applies to every admissible bridge portion
below.

In either orientation the reference predictable intensity of the mark
$(i,\zeta)$, $\zeta\in\{-1,1\}$, is
\begin{equation}\label{score:reference-rates}
 r_{i,\zeta}(t)=\frac{1+\zeta t_i(X(t-))}{2}.
\end{equation}
Their sum at a site is one.  Under an admissible law $\mathsf Q$, write its
predictable intensities as $r_{i,\zeta}b_{i,\zeta}$.  Such predictable
ratios exist on histories of positive $\mathsf Q$ probability by the
intensity change formula \cite[Theorem~4.1]{Jacod1975}.  For example,
conditional waiting and mark densities under a probability absolutely
continuous with respect to the ordinary path law give them by dividing
the conditional waiting density by the conditional no-ring survival.
Localization covers zeros of that survival.  For a time interval $H$ in
the oriented portion put
\begin{equation}\label{score:hellinger-cost}
 \mathcal J_H=\int_H\sum_{i,\zeta}r_{i,\zeta}(t)
                    (\sqrt{b_{i,\zeta}(t)}-1)^2\,dt,
 \qquad
 \mathcal J_{H,i}=\int_H\sum_\zeta r_{i,\zeta}
                    (\sqrt{b_{i,\zeta}}-1)^2\,dt.
\end{equation}
The cost is additive on disjoint intervals.  The probabilities and hence
the ratios $b_{i,\zeta}$ do not depend on the entry $e$ being differentiated.

\begin{lemma}[Likelihood and intensity bounds]\label{score:hellinger}
Let $\mathcal F_t$ be the natural filtration of a nonexplosive marked path,
augmented by an initial $\sigma$-field, and let $\mathsf Q\ll\mathsf P$ be a change of law up to a
finite stopping time $\tau$, with density $Z$ there and with intensities $rb$ relative to
the reference intensities $r$ conditional on that initial information.
For the cost in \eqref{score:hellinger-cost},
\begin{equation}\label{score:hellinger-inequality}
 \mathsf Q\{\mathcal J_{[0,\tau]}>A,\ \log Z\le B\}
          \le e^{(B-A)/2}.
\end{equation}
Consequently, for every fixed $\eta>0$, with arbitrarily high
inverse-polynomial probability,
\begin{equation}\label{score:action-bounds}
 \begin{cases}
 \mathcal J_{\rm path}\le n^\eta,&\text{under a tested forward law},\\
 \mathcal J_{\rm portion}\le C N(|I|)+n^\eta,
       &\text{under either oriented portion of }\mathsf C_I^{x,y},\quad |I|\ge d_0.
 \end{cases}
\end{equation}
Here $d_0$ is a fixed constant large enough for
Equation~\eqref{hist:prefix-density}.
The second assertion is over the tested boundary law followed by the
conditional interior.  In all these cases the integrated new intensity
at any one site on each specified unit-length block is at most $n^\eta$
with the same probability convention.  These integrated intensities and
$\mathcal J$ have polynomial fixed-order moments, also after averaging
over the boundaries of a conditional interior.
\end{lemma}

\begin{proof}
Include a possible initial density $z_0=d\mathsf Q|_{\mathcal F_0}/
d\mathsf P|_{\mathcal F_0}$.  With bounded integrated rates, the marked
point-process likelihood formula \cite[Theorem~5.1]{Jacod1975} gives
\[
 \sqrt Z\,e^{\mathcal J_{[0,\tau]}/2}
 =\sqrt{z_0}\prod_{s\le\tau\,\mathrm{ring}}\sqrt{b_{i_s,\zeta_s}(s)}
   \exp\left\{-\int_0^\tau\sum_{i,\zeta}r_{i,\zeta}(\sqrt{b_{i,\zeta}}-1)\,dt
       \right\}.
\]
The product allows zero factors; a mark with $b=0$ has no mass under
$\mathsf Q$.  The product-exponential following $\sqrt{z_0}$ is the
stochastic exponential for intensities $r\sqrt b$ relative to $r$.
Its expectation conditional on the initial
data is at most one, as follows by integrating the successive conditional
waiting and mark densities, also when a mark factor is zero; this is the
nonnegative likelihood supermartingale of \cite[Proposition~4.3]{Jacod1975}.
Cauchy--Schwarz
gives $\mathsf E_{\mathsf P}\sqrt{z_0}\le1$.  Thus the right side has
$\mathsf P$ expectation at most one.  For general ratios, stop only when
the integrated reference, changed, and square-root rates reach a finite
level.  Let $Y=\sqrt Z e^{\mathcal J/2}$ on $\{Z>0\}$, with extended
values, and set $Y=0$ on $\{Z=0\}$.  The localized likelihood expressions
converge to $Y$ on $\{Z>0\}$; assigning zero on $\{Z=0\}$ can only decrease
their lower limit.  Fatou's lemma gives $\mathsf E_{\mathsf P}Y\le1$.
No positivity restriction is imposed on $b$ in this localization.
Since $\mathsf Q\{Z=0\}=0$, this expectation bound also makes
$\mathcal J<\infty$ $\mathsf Q$-almost surely.  On the event in
\eqref{score:hellinger-inequality},
$\sqrt Z e^{-\mathcal J/2}\le e^{(B-A)/2}$.  Multiplying this by $Y$
and integrating under $\mathsf P$ proves the inequality.

For a tested forward law, $\log Z\le(M_++M_-)\log n$ by
Definition~\ref{hist:tested-law}.  For a conditional portion,
\eqref{hist:prefix-density} gives $\log Z\le CN(|I|)+L\log n$ outside an
arbitrarily small boundary event.  Since $N(|I|)\ge c n^\beta$, applying
\eqref{score:hellinger-inequality} with these bounds proves
\eqref{score:action-bounds} after increasing $C$.

Let $C_{i,H}$ be the number of rings at site $i$ in a unit block $H$, and
$\Lambda_{i,H}=\int_H\sum_\zeta r_{i,\zeta}b_{i,\zeta}\,dt$.  The
exponential count supermartingale \cite[Proposition~4.3]{Jacod1975} gives
\[
 \mathsf E_{\mathsf Q}\exp\{-C_{i,H}+(1-e^{-1})\Lambda_{i,H}\}\le1,
\]
with the conditional version when the boundary is fixed.  Ordinary unit
counts have every fixed moment bounded; paying a polynomial test density
and using the preserved copy marginals shows $C_{i,H}\le n^{\eta/2}$ with
arbitrarily high inverse-polynomial probability, simultaneously on any
polynomial collection of blocks.  The last exponential bound then gives
$\Lambda_{i,H}\le n^\eta$ with the same precision.  Its tail inequality
also bounds every moment of $\Lambda_{i,H}$ by a constant times a higher
moment of $1+C_{i,H}$.  The conditional assertions here are first averaged
over the boundary law.  Finally,
$(\sqrt b-1)^2\le b+1$ bounds $\mathcal J$ by the sum of the new and ordinary
integrated intensities.  The moment assertions follow, using the polynomial
count moments from Lemma~\ref{hist:crude-ratios}.
\end{proof}

For later use we make the unprojected size and reversal precise.  At a ring
of site $i\in e=\{i,j\}$, the order-$h\ge1$ coefficient of
\eqref{hist:ring-likelihood} is
\begin{equation}\label{score:ring-coefficient}
 \gamma_{h,e}(t,i,\zeta)
       =X_j(t-)^h(\zeta-t_i(X(t-)))(-t_i(X(t-)))^{h-1}.
\end{equation}
It is bounded by $2$, and
$\sum_\zeta r_{i,\zeta}\gamma_{h,e}(t,i,\zeta)=0$.

\begin{lemma}[Raw coefficient size and reversal]\label{score:raw-size}
For every fixed $1\le p<\infty$ and $1\le l\le K$, under ordinary forward
dynamics on an interval of length $d\ge1$,
\begin{equation}\label{score:martingale-moment}
 \left\|\sup_{t\in I}|\ell_{l,e}([\inf I,t])|\right\|_{L^p}
       \le C_{p,K}d^{l/2}.
\end{equation}
For every marked path the exact likelihood reversal is
\begin{equation}\label{score:reversal}
 \frac{L_{I,e}^{\to}(u)}{1+u q_e(X_{\sup I})}
   =\frac{L_{I,e}^{\leftarrow}(u)}{1+u q_e(X_{\inf I})}.
\end{equation}
It follows that, for every fixed $\eta>0$, with arbitrarily high
inverse-polynomial probability under tested laws and their resampled
marginals, simultaneously over entries and specified tree intervals,
\begin{equation}\label{score:raw-high-probability}
 |\ell_{l,e}^{\to}(I)|+|\ell_{l,e}^{\leftarrow}(I)|
       \le n^\eta |I|^{l/2},\qquad 1\le l\le K.
\end{equation}
On a unit-length interval this also holds for partial products, with an
arbitrarily small power allowance, when a bounded fixed-order product of
earlier coefficients is included.
\end{lemma}

\begin{proof}
The order-$l$ prefix coefficient is a martingale starting from zero.  At a
ring its jump is
$\sum_{h=1}^l\gamma_{h,e}\ell_{l-h,e}([\inf I,t))$; the prefix factors
are predictable and every $\gamma_h$ is centered.  The martingale
square-function inequality \cite[Theorem~1.1]{BurkholderDavisGundy1972},
applied at successive ring times, and H\"older's inequality bound its $L^p$
supremum by a constant times
\[
 \sum_{h=1}^l\|C_{e,I}^{1/2}\|_{L^{2p}}
       \left\|\sup_t|\ell_{l-h,e}([\inf I,t])|\right\|_{L^{2p}},
\]
where $C_{e,I}$ is the number of rings at the two sites.  It is Poisson
with mean $2d$, whose fixed moments are $O(d^p)$ for $d\ge1$.  Induction
in $l$, with the same assertion at the higher moment order $2p$, proves
\eqref{score:martingale-moment}.

Detailed balance for marked paths states that the forward path measure
multiplied by its initial stationary mass equals the reversed measure
multiplied by its final stationary mass.  Apply this once at the base
matrix and once using \eqref{hist:stationary-perturbation}, and divide.
This gives \eqref{score:reversal}, including at self-resampling rings.
The Taylor coefficients of the two endpoint factors in that identity are
bounded through order $K$.  Thus the reversed bounds follow from the forward
ones.  Choose $p$ sufficiently large in \eqref{score:martingale-moment} to
pay the polynomial test density, the number of entries and intervals, and
the desired inverse-polynomial precision.  This proves
\eqref{score:raw-high-probability}; resampling preserves the marginal.
On a unit block a fixed-order coefficient is a fixed-degree polynomial in
the two-site count, which gives the last assertion in the same way.
\end{proof}

\subsection{The gain from the last ring}

The compensator identity will always be used before taking absolute values.
If $A_e$ is known before an oriented interval $H$ and all displayed
products are integrable, then the last-ring
expansion and ordinary centering give the exact identity
\begin{equation}\label{score:compensation-identity}
 \begin{split}
 \mathsf E_{\mathsf Q}[A_e\ell_{l,e}(H)]
 =\mathsf E_{\mathsf Q}\sum_{h=1}^l\int_H\sum_{i\in e,\zeta}
     A_e\ell_{l-h,e}([\inf H,t))\gamma_{h,e}(t,i,\zeta)
           r_{i,\zeta}(t)(b_{i,\zeta}(t)-1)\,dt .
 \end{split}
\end{equation}
In reversed time the interval notation is interpreted in that orientation.
The first step is predictable compensation at rates $rb$
\cite[Theorem~2.1]{Jacod1975}; subtracting $r$ uses the
centering of \eqref{score:ring-coefficient}.  All prefix factors are
predictable.  A multiplier depending on an unrevealed future path cannot
be inserted in this identity; a common external test of that future is
already part of the probability $\mathsf Q$ and hence of its intensities.

\begin{lemma}[One compensated insertion]\label{score:one-insertion}
Let $H$ have length $d\ge1$ within an admissible oriented portion of a
checked root.  Suppose $v_i(t)$ is predictable, $|v_i(t)|\le1$, and
does not depend on the index $j$ of the other endpoint or on the new mark
$\zeta$.  Define
\[
 D_{\{i,j\}}(H)=\int_H X_j(t-)v_i(t)
       \sum_\zeta r_{i,\zeta}(b_{i,\zeta}-1)(\zeta-t_i)\,dt
       +\int_H X_i(t-)v_j(t)
       \sum_\zeta r_{j,\zeta}(b_{j,\zeta}-1)(\zeta-t_j)\,dt .
\]
With arbitrarily high inverse-polynomial probability, using the joint
boundary-and-interior convention for a bridge portion,
\begin{equation}\label{score:one-insertion-bounds}
 |D(H)|_2\lesssim_\beta d^{1/4}\sqrt{\mathcal J_H/n},\qquad
 |D(H)|_4\lesssim_\beta d^{5/16}n^{-1/4}\sqrt{\mathcal J_H}.
\end{equation}
For an insertion of order $h>1$, or with bounded predictable weights that
may depend on $e$, the same bounds with $d^{1/2}$ in place of $d^{1/4}$ and
$d^{5/16}$ hold without using overlap estimates.  In particular this
version applies on unit blocks to previously accumulated bounded
coefficients, with their power allowance included.
\end{lemma}

\begin{proof}
For each site write $\alpha_i$ for the signed measure in the first integral
with $X_j$ omitted.  Split it according to $b\le2$ and $b>2$.  On the first
set $|b-1|\le C|\sqrt b-1|$, so Cauchy--Schwarz and $\sum_\zeta r_{i,\zeta}=1$
bound the square integral of its absolute density by $C\mathcal J_{H,i}$.
On the second set $|b-1|\le C(\sqrt b-1)^2$, so its total variation is at
most $C\mathcal J_{H,i}$.  On each unit block this variation is also bounded
by a constant times the new and ordinary integrated site intensities, hence
by $n^\eta$ on the event of Lemma~\ref{score:hellinger}.  If $a_{ik}$ is the sum
of the two variations on the $k$th unit block, these facts imply
\begin{equation}\label{score:block-masses}
 \sum_k a_{ik}^2\le Cn^\eta\mathcal J_{H,i}.
\end{equation}
In particular the large ratios have been treated by total mass, not by a
pointwise square-integrability assertion.

First suppose all times in $H$ have age at least a fixed multiple of $d$
from the \emph{forward} start of its checked root.  Squaring the contribution
of a fixed site $i$ and averaging its other endpoint gives the kernel
$n^{-1}\sum_jX_j(t)X_j(s)=\rho_2(X(t),X(s))$.  By
\eqref{hist:path-overlaps}, on unit block indices this is bounded by
$n^{O(\beta)}(1+|k-k'|)^{-1/2}$.  Adjacent blocks use the bound one.  The
$\ell^1$ sum of this kernel over one index is $O(d^{1/2})$, so
\[
 \frac1n\sum_j\left|\int_H X_j\,d\alpha_i\right|^2
   \le n^{O(\beta)}d^{1/2}\sum_k a_{ik}^2.
\]
Omitting $j=i$ adds at most $n^{-1}(\sum_k a_{ik})^2$, which is absorbed
by the right side since $d\le CT$.  Summing over the contributing site and
dividing by $\binom n2$ gives, using \eqref{score:block-masses},
\[
 |D(H)|_2^2\le \frac{n^{O(\beta)}d^{1/2}}n
                    \sum_i\sum_k a_{ik}^2
       \le n^{O(\beta)}d^{1/2}\mathcal J_H/n.
\]

For the fourth power use the kernel
$n^{-1}\sum_j\prod_{h=1}^4X_j(t_h)=\rho_4(X(t_1),\ldots,X(t_4))$.
The age assumption and \eqref{hist:path-overlaps} bound its unit-block
kernel by a constant times
\[
 n^{O(\beta)}\sum_{1\le h<h'\le4}(1+|k_h-k_{h'}|)^{-3/4}.
\]
For each chosen pair the convolution is at most
$Cd^{1/4}\sum_k a_{ik}^2$; the other two $\ell^1$ sums contribute at most
$d\sum_k a_{ik}^2$.  Thus the fourth-power integral is at most
$n^{O(\beta)}d^{5/4}(\sum_k a_{ik}^2)^2$.  The omitted-diagonal term is
$n^{-1}(\sum a_{ik})^4$ and is again absorbed, since $d^{3/4}/n\le Cn^{-1/2}$.
Now sum over $i$ and divide by $\binom n2$.  Since
$\sum_i\mathcal J_{H,i}^2\le\mathcal J_H^2$, this gives
\[
 |D(H)|_4^4\le\frac{n^{O(\beta)}d^{5/4}}n
                   \sum_i\left(\sum_k a_{ik}^2\right)^2
       \le n^{O(\beta)}d^{5/4}\mathcal J_H^2/n.
\]
The fourth root is the second bound.

For a general $H$, partition it into unit blocks near the forward root
start and pieces whose lengths grow geometrically and are at most a fixed
multiple of their ages.  Apply the preceding estimates on each piece.
For $\gamma=1/4$ or $5/16$, Cauchy--Schwarz gives
$\sum_jd_j^\gamma\sqrt{\mathcal J_{H_j}}
 \le C d^\gamma\sqrt{\mathcal J_H}$, because $\sum_jd_j^{2\gamma}=O(d^{2\gamma})$.
This partition uses the forward age even when the compensation orientation
is reversed; it is a deterministic division of the integrals and changes
no filtration.  Finally, replacing both overlap kernels by one gives
the exponent $1/2$ for either norm.  That proof uses only a common upper
bound on the absolute block masses, so bounded $e$-dependent predictable
weights and higher insertion orders are permitted in this last assertion.
\end{proof}

\subsection{Conditional coefficients and the midpoint identity}

For a node $I=[a,b]$ with fixed boundaries, keep its descendant weights
fixed and define the conditional coefficient series by
\begin{equation}\label{score:R-definition}
 R_e(I;u)=\mathsf C_I^{x,y}\left[
                   \frac{L_{I,e}^{\to}(u)}{1+u q_e(y)}\right],
 \qquad R_{l,e}(I)=[u^l]R_e(I;u).
\end{equation}
In particular $R_0=1$.  Reversal \eqref{score:reversal} and the reversal
of the bridge give the same series from the reversed path with endpoint
$x$.  Write $c_{l,e}^I=\mathsf C_I^{x,y}[\ell_{l,e}^{\to}(I)]$ for the
bridge average of the raw coefficient.  Exactly,
\begin{equation}\label{score:raw-conditional}
 c_l^I=R_l(I)+q_e(y)R_{l-1}(I),\qquad l\ge1.
\end{equation}
For a reversed raw coefficient its end in this formula is the original
start.  These definitions keep all descendant weights fixed.  Their
derivatives are handled separately through \eqref{hist:pre-ratio}.

At the midpoint $v=(a+b)/2$, set $z=X_v$.  With $x,y$ fixed, $z$ has the
base midpoint probability \eqref{hist:midpoint-probability}, and the two
interiors are independent conditional on $z$.  Reversal
\eqref{score:reversal} gives the pathwise factorization
\[
 \frac{L_{I,e}^{\to}(u)}{1+u q_e(y)}
 =(1+u q_e(z))
   \frac{L_{I^-,e}^{\to}(u)}{1+u q_e(z)}
   \frac{L_{I^+,e}^{\leftarrow}(u)}{1+u q_e(z)}.
\]
Conditioning the two interiors on $z$ therefore gives the exact identity
\begin{equation}\label{score:midpoint-series}
 R_e(I;u)=\sum_z\pi_I^{x,y}(z)(1+u q_e(z))
                     R_e(I^-;u)R_e(I^+;u).
\end{equation}
Here the right factor is written in its reversed orientation, and the
midpoint probability is the fixed base one.  The pathwise factorization
shows why the endpoint normalizations leave exactly the single factor
$1+u q_e(z)$.

For $l\ge2$, taking coefficients and grouping the two raw terms gives
\begin{equation}\label{score:midpoint-coefficient}
 \begin{split}
 R_l(I)=\mathsf E_{\pi_I^{x,y}}\bigg[
    c_l^{I^-}+c_l^{I^+,\leftarrow}
    +\sum_{\substack{a+b=l\\a,b\ge1}}R_a(I^-)R_b(I^+)
    +q_e(z)\sum_{\substack{a+b=l-1\\a,b\ge1}}R_a(I^-)R_b(I^+)
                         \bigg].
 \end{split}
\end{equation}
No second or higher midpoint mass coefficient occurs.  This cancellation
is the reason for retaining the endpoint denominator in
\eqref{score:R-definition}.

The two linear terms and the product terms in this recursion will be
estimated differently.  For the linear terms, the conditional interior
identity gives the vector equalities
\begin{equation}\label{score:raw-before-midpoint-norm}
 \begin{aligned}
 \mathsf E_{\pi_I^{x,y}}c_l^{I^-}
   &=\mathsf C_I^{x,y}[\ell_l^{\to}(I^-)],\\
 \mathsf E_{\pi_I^{x,y}}c_l^{I^+,\leftarrow}
   &=\mathsf C_I^{x,y}[\ell_l^{\leftarrow}(I^+)].
 \end{aligned}
\end{equation}
These identities integrate the midpoint back out before compensation.
The right-hand sides use the parent oriented laws with $x,y$ fixed and
$z$ neither conditioned on nor revealed in advance; the natural path
reveals its midpoint state when it reaches that time.  We take vector norms
only after the compensation identity.  The product terms instead use the
conditional child $R$ bounds under the tested boundary law followed by the
conditional interior.  In particular, the estimate of a linear term does
not replace the norm of its midpoint average by an average of raw child norms.

Before estimating this recursion we give the raw estimate that is also
used for a final block of a tested forward path.

For a node $J$ in an oriented subtree, write $J_{\mathrm{early}}$ and
$J_{\mathrm{late}}$ for the children encountered first and second in that
orientation.  They are $(J^-,J^+)$ in forward time and $(J^+,J^-)$ in
reversed time; $J^-,J^+$ continue to denote the chronological children in
the kernel and midpoint identities.  In the raw-representation arguments
below we suppress the arrows on $\ell$ and $c$ when the orientation is fixed.
The state $X_{\mathrm{end}}$ in
$c_j=R_j+q_e(X_{\mathrm{end}})R_{j-1}$ is then the boundary state at the
end of the relevant interval in that orientation.

\begin{lemma}[Representations after compensating the last score]\label{score:raw-projection}
Fix an admissible oriented law $\mathsf Q$, its oriented filtration, and
any permitted auxiliary-data augmentation.  Fix a dyadic subtree $H$ of
length $d$ in that oriented portion and an order $1\le l\le K$.  There is
a finite family of time-integral vectors $V_\alpha^{(l)}(H)$, chosen from
this fixed setup without reference to a subsequent choice of $A$, such that
for every vector $A_e$ measurable at the oriented start of $H$ for which
the displayed products are integrable, coordinatewise,
\begin{equation}\label{score:raw-representation}
 \mathsf E_{\mathsf Q}[A_e\ell_{l,e}(H)]
       =\mathsf E_{\mathsf Q}\left[A_e\sum_\alpha V_{\alpha,e}^{(l)}(H)\right].
\end{equation}
The vectors and $A$ may depend on the same initial information.  With arbitrarily high
inverse-polynomial probability, using the joint boundary-and-interior
convention for a bridge portion, their sum of norms obeys
\begin{equation}\label{score:raw-projection-bounds}
 \begin{split}
 \sum_\alpha|V_\alpha^{(l)}(H)|_2
       &\lesssim_\beta d^{(l-1)/2+1/4}\sqrt{\mathcal J_H/n},\\
 \sum_\alpha|V_\alpha^{(l)}(H)|_4
       &\lesssim_\beta d^{(l-1)/2+5/16}n^{-1/4}\sqrt{\mathcal J_H}.
 \end{split}
\end{equation}
Thus earlier vector factors can be included by H\"older's inequality after
the identity.  This is a representation inside expectation, not a
pointwise assertion about the raw coefficient $\ell_l(H)$.
\end{lemma}

\begin{proof}
For $l=1$, \eqref{score:compensation-identity} is the representation, and
Lemma~\ref{score:one-insertion} gives both bounds.  For $l\ge2$, expand a raw
coefficient by the smallest dyadic node containing all its positive-order
ring insertions.  If all insertions lie in a leaf, keep that leaf term.
Otherwise, at that unique node $J$ the term is
$\ell_{j,e}(J_{\mathrm{early}})\ell_{l-j,e}(J_{\mathrm{late}})$ for some
$1\le j<l$.  This is an exact partition of the ring expansion, including
rings carrying Taylor order
larger than one.

On a leaf, expand according to the last ring and its assigned order as in
\eqref{score:compensation-identity}.  The remaining prefix coefficients
are predictable and are bounded by a power allowance on this unit block,
by Lemma~\ref{score:raw-size}.  The unit-block conclusion of
Lemma~\ref{score:one-insertion} applies.  At a split of length $y$, the earlier
coefficient is known before the later child and is at most
$n^\eta y^{j/2}$ by \eqref{score:raw-high-probability}.  Apply the inductive
representation to the later child, then include this earlier factor by its
supremum norm.  At any one scale the later children are disjoint and number
at most $Cd/y$.  Consequently
\[
 \sum_{J:\,|J|\asymp y}\sqrt{\mathcal J_{J_{\mathrm{late}}}}
      \le C\sqrt{(d/y)\mathcal J_H}.
\]
For the $L^2$ or $L^4$ bound the power before this summation is
$y^{(l-1)/2+\gamma}$, with $\gamma=1/4$ or $5/16$.  Relative to the claimed
right side, this scale thus has the factor
\[
 (y/d)^{(l-2)/2+\gamma}.
\]
It is geometrically summable for $l\ge2$.  The leaf terms cost at most
$C\sqrt{d\mathcal J_H}$ with the respective factor $n^{-1/2}$ or
$n^{-1/4}$, and are also bounded by the claimed expression.  All
compensations used the expectation identity with the earlier coefficient
still predictable.  After absolute values, each representation vector is
bounded by a fixed polynomial in raw clock counts and prefix coefficients,
and in
$\int_H\sum_{i,\zeta}r_{i,\zeta}|b_{i,\zeta}-1|
 \le\int_H\sum_{i,\zeta}r_{i,\zeta}b_{i,\zeta}+n|H|$.
For the sharper representations below, an earlier endpoint projection adds
a conditional coefficient with the uniform coarse bound of
Lemma~\ref{hist:crude-ratios}.  That lemma, the intensity moments in
Lemma~\ref{score:hellinger}, and H\"older's inequality therefore give
polynomial fixed moments for the representation norms, including a fixed
number of copies.  The conditional exceptional-set rule then applies to
each of the polynomially many terms and proves the stated probability
formulation.
\end{proof}

\begin{proposition}[Conditional clock scores and raw representations]\label{prop:clock-scores}
\textit{(i) Conditional coefficients.}
For every deterministic matrix on \eqref{hist:static-conditions}, every
tested path law, and every node $I$ of length $D$ in its checked roots, the
vectors in \eqref{score:R-definition} satisfy, with arbitrarily high
inverse-polynomial probability over their boundary states,
\begin{equation}\label{score:R-basic}
 |R_l(I)|_\infty\lesssim_\beta D^{l/2},\qquad
 |R_l(I)|_2\lesssim_\beta D^{l/2-1},\qquad
 |R_l(I)|_4\lesssim_\beta D^{l/2-1/2},\qquad 1\le l\le K.
\end{equation}
Also,
\begin{equation}\label{score:R-l1}
 |R_3(I)|_1\lesssim_\beta D^{-1/4},\qquad
 |R_l(I)|_1\lesssim_\beta D^{(l-4)/2},\quad 4\le l\le K.
\end{equation}
If $D\ge Tn^{-0.005}$, the $L^4$ bound for $R_1$ has the additional factor
$n^{-0.01}$, the $L^2$ bounds for $R_2,R_3$ have the additional factor
$n^{-0.001}$, and the $L^1$ bound for $R_4$ has the additional factor
$n^{-0.0007}$.  The conclusions hold simultaneously for any polynomial
collection of specified nodes and any fixed number of conditional copies,
with the same exceptional-set convention.

\textit{(ii) Raw representations.}
Fix an admissible oriented law, its oriented filtration, and any permitted
auxiliary-data augmentation, as in Lemma~\ref{score:raw-projection}.  Then
fix a dyadic subtree $H$ of length $d$ in that portion and an order
$1\le l\le K$.  The following representations satisfy
\eqref{score:raw-representation} for every earlier vector $A$ in that
lemma's class, with the representation chosen from this fixed setup before
$A$.  If $l\ge2$ and $d\ge Tn^{-0.02}$, there is a representation for which
\begin{equation}\label{score:raw-improvement}
 \sum_\alpha|V_\alpha^{(l)}(H)|_2
       \lesssim_\beta n^{-0.004}
                   d^{(l-1)/2+1/4}\sqrt{\mathcal J_H/n}.
\end{equation}
For every $1\le l\le K$, there is a representation for which
\begin{equation}\label{score:raw-l1-projection}
 \sum_\alpha|V_\alpha^{(l)}(H)|_1
       \lesssim_\beta \max\{\sqrt d,d^{l/2-5/4}\}\sqrt{\mathcal J_H/n}.
\end{equation}
For $l=4$ and $d\ge Tn^{-0.006}$, the latter estimate has the additional
factor $n^{-0.001}$.  The two families here may differ from each other and
from the family giving the two basic bounds in Lemma~\ref{score:raw-projection}.
These pathwise norm estimates have arbitrarily high inverse-polynomial
probability with that lemma's scope: for a bridge portion the probability
is under the tested boundary law followed by the conditional interior.
They are not uniform conditional assertions at every fixed boundary pair.
\end{proposition}

\begin{proof}
All conditional estimates in this proof are first established under the
tested boundary law followed by its interior sampling.  We then use the
conditional exceptional-set rule following Lemma~\ref{hist:crude-ratios}.
The polynomial coefficient and representation moments proved in
Lemmas~\ref{hist:crude-ratios} and~\ref{score:raw-projection} permit any
required stronger working precision.
This yields the stated high-probability bounds on conditional expectations,
without asserting them for every possible boundary.

\paragraph{Order one and its stronger $L^4$ bound.}
The infinity bound follows at once by applying
\eqref{score:raw-high-probability} inside
\eqref{score:R-definition} and expanding the bounded endpoint denominator.
For $R_1$, split instead at a variable $v$ in the middle third of $I$.
The pathwise reversal identity gives
\[
 R_1(I)=\mathsf C_I^{x,y}\bigl[
       \ell_{1}^{\to}([a,v])+\ell_{1}^{\leftarrow}([v,b])-q_e(X_v)\bigr].
\]
Average this identity uniformly over that middle third.  Compensating the
two ring terms in the parent forward and reversed portions uses bounded
deterministic time weights, common for the other-site index.  By
\eqref{score:action-bounds} and \eqref{score:one-insertion-bounds} their
$L^2$ cost is at most
\[
 n^{O(\beta)}D^{1/4}\sqrt{N(D)/n}
       =n^{O(\beta)}D^{-1/2}.
\]
For the boundary average $Q_e=(3/D)\int_{a+D/3}^{a+2D/3}q_e(X_v)\,dv$,
the exact edge sum
$\sum_{i<j}q_{ij}(X_v)q_{ij}(X_w)
 =\tfrac12(n^2\rho_2(X_v,X_w)^2-n)$ gives
\[
 |Q|_2^2\le\frac{C}{D^2}\int_{a+D/3}^{a+2D/3}\int_{a+D/3}^{a+2D/3}
                       \rho_2(X_v,X_w)^2\,dv\,dw
       \lesssim_\beta D^{-1}.
\]
The last step uses the square of \eqref{hist:path-overlaps}, whose kernel
is $(1+|v-w|)^{-1}$ on this middle third, and absorbs $\log(2+D)$.
This proves the $L^2$ bound for $R_1$.  Interpolation with its infinity
bound gives the ordinary $L^4$ estimate.

We will retain a stronger intermediate $L^4$ improvement.  Its ring terms
can be bounded directly by the second estimate in
\eqref{score:one-insertion-bounds}, giving
\[
 n^{O(\beta)}D^{5/16}n^{-1/4}\sqrt{N(D)}
       =n^{1/4+O(\beta)}D^{-7/16}.
\]
If $D\ge Tn^{-0.03}$ this is at most $n^{-0.02+O(\beta)}$.
The boundary term has infinity norm at most one and $L^2$ norm
$\lesssim_\beta D^{-1/2}$, so its $L^4$ norm is even smaller.  Hence
\begin{equation}\label{score:R1-intermediate}
 |R_1(I)|_4\lesssim_\beta n^{-0.02},\qquad D\ge Tn^{-0.03}.
\end{equation}
For bounded $D$, where the conditional prefix estimate was not invoked,
Lemma~\ref{score:raw-size} gives all the ordinary bounds after a constant change.

\paragraph{Joint induction for the basic bounds and the sharper raw $L^2$ representation.}
We next prove the $L^2$ and $L^4$ bounds by induction on $l$, together with
\eqref{score:raw-improvement}.  Assume \eqref{score:R-basic} in orders below
$l\ge2$.  To prove the raw improvement when $d\ge Tn^{-0.02}$, keep the
basic representation of Lemma~\ref{score:raw-projection} on split scales
$y<dn^{-0.02}$.  The scale factor in that lemma's proof is at most
$(y/d)^{1/4}$, so these scales gain
$n^{-0.005}$.  Their leaf remainder has a still larger gain.

On a larger split, project the earlier raw coefficient \emph{before}
compensating anything in the later child.  In the expectation at this
stage, the remaining raw insertions are outside the earlier child.  By
Lemma~\ref{hist:conditional-interior}, for any such exterior-measurable
integrand $H_e$ the exact identity is
\begin{equation}\label{score:earlier-projection}
 \mathsf E_{\mathsf Q}[\ell_{j,e}(J_{\mathrm{early}})H_e]
       =\mathsf E_{\mathsf Q}
          [c_{j,e}^{J_{\mathrm{early}}}(X_{\partial J_{\mathrm{early}}})H_e].
\end{equation}
For a prefix law this follows by first conditioning on the full exterior
and then taking the prefix marginal.
The new factor is known at the start of $J_{\mathrm{late}}$.  Its $L^4$ norm is at most
$n^{O(\beta)}y^{j/2-1/2}$ by the induction hypothesis and
\eqref{score:raw-conditional}; for $j=1$ the term $q_eR_0$ is simply
bounded by one.  Apply the basic $L^4$ representation to the later
child and pair these factors in $L^2$.  Compared with the basic split
$L^2$ estimate the gain is
\[
 n^{1/4}y^{-7/16}.
\]
Indeed the new power of $y$ is
$j/2-1/2+(l-j-1)/2+5/16=l/2-11/16$, while the old one is $l/2-1/4$.
Here $y\ge dn^{-0.02}\ge Tn^{-0.04}$, so the last ratio is at most
$n^{-0.024+o(1)}$.  Summing these scales and the smaller scales proves
\eqref{score:raw-improvement}, with ample room for the stated exponent.
The identity \eqref{score:earlier-projection} was applied while the later
coefficient was still raw.  It is not applied to a drift from an earlier
compensation, which could retain the earlier interior in its intensities.

Consider now \eqref{score:midpoint-coefficient}.  For each linear raw term,
use \eqref{score:raw-before-midpoint-norm} and apply
Lemma~\ref{score:raw-projection} on the corresponding half in the
\emph{parent} oriented filtration with just $x,y$ fixed.  After the
compensation identity, use the triangle inequality and the pathwise norm
bounds.  The action bound for
that parent portion is $\mathcal J\lesssim_\beta N(D)$, so its $L^2$ cost
is
\[
 n^{O(\beta)}D^{(l-1)/2+1/4}\sqrt{N(D)/n}
       =n^{O(\beta)}D^{l/2-1}.
\]
For the product terms in \eqref{score:midpoint-coefficient}, the induction
hypothesis and H\"older give
$|R_aR_b|_2\le |R_a|_4|R_b|_4\lesssim_\beta D^{(a+b)/2-1}$.
The terms with $q_e$ have the still smaller exponent $(l-1)/2-1$.
Using the conditional exceptional-set rule proves the $L^2$ bound at order
$l$.  Interpolating it with the infinity bound proves its $L^4$ bound.

For clarity we track the low-order margins needed below.  When
$D\ge Tn^{-0.019}$, a half has length at least $Tn^{-0.02}$ for large $n$.
Use \eqref{score:raw-improvement} on the two raw order-two terms.  Its
product term uses \eqref{score:R1-intermediate} on both children.  Thus
\begin{equation}\label{score:R2-intermediate}
 |R_2(I)|_2\lesssim_\beta n^{-0.004},\qquad
 |R_2(I)|_4\lesssim_\beta n^{-0.002}D^{1/2},
             \qquad D\ge Tn^{-0.019},
\end{equation}
where the second bound follows by interpolation.  Similarly, if
$D\ge Tn^{-0.018}$, the raw order-three terms use
\eqref{score:raw-improvement}; their product terms use
\eqref{score:R1-intermediate} and \eqref{score:R2-intermediate} on the
children.  The terms with $q_e$ have an additional factor $D^{-1/2}$
relative to the target bound.  This gives
\begin{equation}\label{score:R3-intermediate}
 |R_3(I)|_2\lesssim_\beta n^{-0.0015}D^{1/2},
                          \qquad D\ge Tn^{-0.018}.
\end{equation}
These deductions use the induction hypothesis at strictly smaller orders,
so the proof of \eqref{score:R-basic} and of these improvements is complete.

\paragraph{The raw $L^1$ representation.}
With the basic conditional bounds established, we prove
\eqref{score:raw-l1-projection} by induction on $l$.  In every split of the
smallest-node decomposition, first apply \eqref{score:earlier-projection}.
In the term $R_j$ of $c_j=R_j+q_eR_{j-1}$, use its $L^2$ bound
$y^{j/2-1}$ and the basic later-child $L^2$ representation.  The power
before the summation at scale $y$ is $y^{l/2-5/4}$.  The term $q_eR_{j-1}$
has an extra $y^{-1/2}$ when $j\ge2$ and is bounded the same way.  Summing
the square roots of the actions at this scale contributes $\sqrt{d/y}$.
Thus the sum over scales is bounded by
\[
 C\sqrt d\sum_{\substack{y\le d\\y\ \mathrm{dyadic}}}
       y^{l/2-7/4}\sqrt{\mathcal J_H/n}
 \le C\max\{\sqrt d,d^{l/2-5/4}\}\sqrt{\mathcal J_H/n},
\]
with a logarithm allowed at an equality of exponents.

When $j=1$, the remaining term in $c_1$ is $q_e$ alone.  It is known before
the later child.  Use \eqref{score:raw-l1-projection} inductively for the
later order $l-1$; order one follows from the $L^2$ estimate in
Lemma~\ref{score:one-insertion}.  Its scale sum is at most
\[
 C\sqrt d\max\{1,d^{l/2-9/4}\}\sqrt{\mathcal J_H/n},
\]
which is no larger than the claimed expression.  Leaf terms cost
$C\sqrt{d\mathcal J_H/n}$ as in Lemma~\ref{score:raw-projection}.  This proves
the ordinary bound at every order.  Throughout, the earlier conditional
coefficient is left inside expectation until the later compensation is
performed, where it is predictable.

At order four the leading scale contribution, divided by its value at
$y=d$, is $(y/d)^{1/4}$.  If $y<dn^{-0.01}$ this gains $n^{-0.0025}$.
If $y\ge dn^{-0.01}$ and $d\ge Tn^{-0.006}$, then
$y\ge Tn^{-0.016}$.  For an earlier order $j=2$ or $3$, its $R_j$ term has
the improved $L^2$ bound \eqref{score:R2-intermediate} or
\eqref{score:R3-intermediate} on that child.  Its $q_eR_{j-1}$ term has
the extra $y^{-1/2}$ just noted.  For $j=1$ with its $R_1$ term, the later
order is three, and \eqref{score:raw-improvement} applies there.  Constant
halves are absorbed by the slack in these thresholds.  Finally the leaf
and $q_e$-alone branches cost only $O(\sqrt d)$ in place of $d^{3/4}$.
Each case gains at least $n^{-0.001+O(\beta)}$.  This proves the additional
factor in \eqref{score:raw-l1-projection}.

\paragraph{Conditional $L^1$ bounds.}
Apply the $L^1$ representation to each raw term in
\eqref{score:midpoint-coefficient}, using the order of conditioning in
\eqref{score:raw-before-midpoint-norm}.  Since
\begin{equation}\label{score:node-action-scale}
 \sqrt{N(D)/n}=n^{\beta/2}D^{-3/4},
\end{equation}
the two raw order-three terms are $\lesssim_\beta D^{-1/4}$, and raw order
$l\ge4$ is $\lesssim_\beta D^{(l-4)/2}$.  For the product terms use the
two $L^2$ estimates in \eqref{score:R-basic}.  At order three they cost at
most $D^{-1/2}$ (and the term with $q_e$ at most $D^{-1}$); at order
$l\ge4$ they cost at most $D^{(l-4)/2}$, with an extra $D^{-1/2}$ in the
terms containing $q_e$.  This proves \eqref{score:R-l1}.

If $D\ge Tn^{-0.005}$, a half satisfies the threshold for the improved
order-four representation.  In its product terms, $R_1R_3$ uses the
improved $L^2$ bound for $R_3$, $R_2R_2$ uses that for $R_2$, and the terms
with $q_e$ already have the extra $D^{-1/2}$.  The claimed factor
$n^{-0.0007}$ follows.  The stated weaker late factors for $R_1,R_2,R_3$
follow immediately from \eqref{score:R1-intermediate}--\eqref{score:R3-intermediate}.
The conditional exceptional-set rule and a union over the specified nodes
and fixed number of copies complete the proof.
\end{proof}

The representations in Lemma~\ref{score:raw-projection} and
Proposition~\ref{prop:clock-scores}(ii) also apply to a last differentiated
root of a tested forward path.  Each is used for its respective norm bound.
In that use \eqref{score:action-bounds} gives
$\sqrt{\mathcal J_H}\le n^\eta$ after adjusting $\eta$.  Earlier root
coefficients can first be conditioned on their endpoints by
\eqref{score:earlier-projection}; these endpoints are then in the past of
the last root.  This is the form used in the replacement argument.

\section{Replacement and functional universality}\label{sec:replacement}

We now transfer the Gaussian limit to Rademacher couplings.  The comparison
uses the spin autocorrelations themselves.  Its principal input is the
averaged clock-score estimate of Proposition~\ref{prop:clock-scores}; the remaining
point is to differentiate every survival and failure probability in the
history construction.  Conditional copies of a checked path provide these
derivatives while preserving the probability laws to which the score
estimates apply.

\begin{proposition}[Transfer to Rademacher disorder]\label{prop:universality}
For the Rademacher coupling law,
\[
 \Law_{\mathrm{Rad}}\bigl(A_{n,J},B_{n,J}\bigr)
 \weakto \Law_g\bigl(A_g,B_g\bigr)
\]
in
$C_{\mathrm{loc}}((0,\infty))\times
C_{\mathrm{loc}}((0,\infty)^2)$, where $(A_g,B_g)$ is the Gaussian limit
obtained by combining Proposition~\ref{prop:gaussian-subsequence} with the
identification of the entrance family in Theorem~\ref{thm:entrance-selection}.
In particular the entrance family and the semigroup in this limit are the same $h_s^g$ and
$S_t^g=\exp(-t c_*H_g)$ as in the Gaussian case, and the finite time scale is
$T=T_n=n^{2/3}$.
\end{proposition}

Throughout this section a length $D$ of a path interval is in site-clock
units.  All query legs have lengths in $[cT,CT]$, where $0<c<C<\infty$ are
fixed by a compact set of query times.  We put $K=24$.  The small parameter
$\beta>0$ and the convention $\lesssim_\beta$ are those of the preceding two
sections.  Constants in an exponent $O(\beta)$ will depend on $K,c,C$, but
not on the number of conditional-copy generations introduced below.

\subsection{Derivatives of tested histories}

We first state the derivative estimate in the form needed for smooth tests
of several traces.  Let $e=\{i,j\}$, $q_e(x)=x_i x_j$,
$\bar q_e=\mu q_e$, and $\bar q=(\bar q_e)_e$.  Add $\vartheta$ to the
interaction coefficient $J_{ij}/\sqrt n$ and set $u=\tanh\vartheta$.
For a checked full history or prefix $\alpha$ of
Definition~\ref{hist:tested-law}, define the unnormalized checked integral
\begin{equation}\label{repl:tested-integral}
 Z_{\alpha,e}(u)=
 \mathbb E_{\alpha,e,u}\left[
   v_\alpha
   \prod_{I\in\mathscr S_\alpha}w_{I,e}(u)
   \prod_{I\in\mathscr F_\alpha}(1-w_{I,e}(u))\right].
\end{equation}
Here $\mathbb E_{\alpha,e,u}$ is ordinary dynamics on the prescribed full
history or prefix, with uniform or perturbed stationary initialization as
specified by $\alpha$.  The set $\mathscr S_\alpha$ consists of its included
survival checks, and $\mathscr F_\alpha$ is empty or consists of its specified
first failure.  The nonnegative function $v_\alpha$ depends only on the
designated initial, leg, and forest endpoints; it is independent of $u$ and
common to every $e$.  Write $z_\alpha=Z_{\alpha,e}(0)$, which does not depend
on $e$.  When $z_\alpha$ satisfies the inverse-polynomial lower bound in
Definition~\ref{hist:tested-law}, normalizing this base integral gives the
tested path law.  A product of such integrals always uses independent histories;
the external tests factor separately across them.

For a vector indexed by $e$, the norms $|\cdot|_p$ use uniform probability
on the $\binom n2$ edges, as in Proposition~\ref{prop:clock-scores}.  Coefficients
$[u^j]$ below are at $u=0$.  The distinction between coefficients and
derivatives costs constants depending only on $K$.

\begin{lemma}[Averaged derivatives of histories]\label{repl:history-derivatives}
There are constants $a>0$ and $C_K<\infty$ with the following property for
all sufficiently small fixed $\beta>0$.  Fix $A,B<\infty$, at most $K$
histories of the form \eqref{repl:tested-integral}, and integers
$j_\alpha\ge0$ with $1\le k=\sum_\alpha j_\alpha\le K$.  Suppose
$0\le v_\alpha\le n^B$.  On the deterministic static support of
Lemma~\ref{hist:static-support}, put
\[
 \mathscr D_e=\prod_\alpha [u^{j_\alpha}]Z_{\alpha,e}(u).
\]
Then, uniformly in the selected histories,
\begin{align}
 |\mathscr D|_1
 &\le C\left(\prod_\alpha z_\alpha\right)
 \begin{cases}
 n^{-a+C_K\beta},&1\le k\le4,\\
 n^{C_K\beta}T^{(k-4)/2},&5\le k\le K,
 \end{cases}
 +n^{-A},                                                     \label{repl:averaged-history}\\
 |\mathscr D|_\infty
 &\le C n^{C_K\beta}T^{k/2}\prod_\alpha z_\alpha+n^{-A}.
                                                               \label{repl:sup-history}
\end{align}
The constant $C$ and the threshold for $n$ may depend on $A,B$ and on the
fixed rectangles.  The exponent constant $C_K$ is independent of these
precision parameters and of $\beta$.  The same estimates hold for specified
derivatives in $\vartheta$ in place of the coefficients in $u$.  Signed tests
are covered by applying the statement to their positive and negative parts.
\end{lemma}

\begin{proof}
The additive precision in the lemma is useful for very small tested
masses.  To prove it, first suppose each $z_\alpha\ge n^{-M}$ for a fixed
$M$.  Normalize each integral by $z_\alpha$ and use the tested probability
law.  The estimates below then have the factor $\prod_\alpha z_\alpha$.
All exceptional contributions may be made smaller than $n^{-A'}$ for an
arbitrary fixed $A'$ by the precision convention in
Definition~\ref{hist:tested-law} and Proposition~\ref{prop:clock-scores}.  If some $z_\alpha<n^{-M}$,
Lemma~\ref{hist:crude-ratios} and the polynomial path moments bound its
unnormalized coefficients by $n^{C(K,B)}z_\alpha$.  Choosing $M$ after $A$
discards all such terms at cost $n^{-A}$.  We shall use this reduction
throughout the proof.

\paragraph{Coefficients with fixed survival and failure weights.}
We start with the terms in which no $w_I$ is differentiated.  Recall from
Proposition~\ref{prop:clock-scores} that, for a node of length $D$,
\begin{equation}\label{repl:score-table}
\begin{aligned}
 |R_l(D)|_\infty&\lesssim_\beta D^{l/2},&
 |R_l(D)|_2&\lesssim_\beta D^{l/2-1},&
 |R_l(D)|_4&\lesssim_\beta D^{l/2-1/2},\\
 |R_3(D)|_1&\lesssim_\beta D^{-1/4},&
 |R_l(D)|_1&\lesssim_\beta D^{(l-4)/2}\quad(l\ge4).
\end{aligned}
\end{equation}
When $D\ge Tn^{-0.005}$, the $L^4$ bound for $R_1$, the $L^2$ bounds for
$R_2,R_3$, and the $L^1$ bound for $R_4$ have additional factors
$n^{-0.01},n^{-0.001},n^{-0.0007}$, respectively.  The notation in
\eqref{repl:score-table} records the length, with the node and its fixed
endpoints understood.

For a last raw coefficient of order $l$ in a forward tested history, the
compensation representations of
Lemma~\ref{score:raw-projection} and Proposition~\ref{prop:clock-scores}(ii), with
$\mathcal J_I\lesssim_\beta1$ by \eqref{score:action-bounds}, give the
following costs on an interval of length
$d$:
\begin{equation}\label{repl:forward-costs}
\begin{array}{c|c}
\text{edge norm}&\text{cost after the compensation identity}\\ \hline
L^2& n^{-1/2}d^{l/2-1/4}\\
L^4& n^{-1/4}d^{l/2-3/16}\\
L^1& n^{-1/2}\max(d^{1/2},d^{l/2-5/4}).
\end{array}
\end{equation}
Each line allows a factor $n^{O(\beta)}$.  These are estimates for the
integrands in the representation after taking the compensator, with earlier
measurable factors kept inside the expectation.  They do not estimate the
expected absolute value of an unprojected score.  The representations used
for different lines or their improvements may be different; the table does
not assert simultaneous bounds on one family of vectors.  The $L^2$ cost gains
$n^{-0.004}$ for $l\ge2$ and $d\ge Tn^{-0.02}$, and the $L^1$ cost at order
four gains $n^{-0.001}$ for $d\ge Tn^{-0.006}$.

Stationary initialization contributes
\[
 \frac{1+u q_e(X_0)}{1+u \bar q_e}.
\]
We need a gain when these are the only differentiated factors.  If $p$ is
a nonnegative stationary initial density with $\mu p=1$ and
$\|p\|_\infty\le n^{C_0}$, then the exact replica identity is
\begin{equation}\label{repl:initial-overlap}
 \frac1{\binom n2}\sum_e \mu[pq_e]^2
 =\frac{n\mu^{\otimes2}[p(X)p(X')\rho_2(X,X')^2]-1}{n-1}
 \le C n^{4\beta}T^{-1}+n^{-A'}
\end{equation}
for every fixed $A'$.  Indeed choose the first dyadic $d_\circ\ge cT$, so
$d_\circ\le2cT$, $r_2(d_\circ)\le c^{-1/2}r_2(T)$, and
$a_2(d_\circ)\ge c_1 n^{3\beta}$.  This scale is among the query scales
in Lemma~\ref{hist:static-support}.  Split at $|\rho_2|=r_2(d_\circ)$;
the complementary probability is
$\exp(-c n^{3\beta})$, which pays $\|p\|_\infty^2$.  This is the weighted
replica identity used in the stationary comparison
\cite[Section~8]{FUA}. %
In particular
\begin{equation}\label{repl:stationary-factor}
 |\bar q|_2\lesssim_\beta T^{-1/2},\qquad
 |\bar q|_4\lesssim_\beta T^{-1/4},\qquad |\bar q|_\infty\le1 .
\end{equation}
The normalized initial marginal of a tested stationary history is such a
$p$: its density is at most $\|v_\alpha\|_\infty/z_\alpha$.  The coefficient
of positive order $j$ in the initialization factor is
$(-\bar q_e)^j+q_e(X_0)(-\bar q_e)^{j-1}$.  A numerator supplies at most one
$q_e(X_0)$ on each history.  Thus if all positive orders come from
initialization, at least one factor is either $\bar q$ or the integrated vector
$\mu[pq_e]$; \eqref{repl:initial-overlap} and boundedness of the other
factors give a fixed power saving.

Suppose now a raw clock likelihood has a positive order.  Each path is a
forest with $O(\log n)$ roots, or consists of the one or two complete legs.
Chronologically earlier roots have lengths at least a constant times later
ones.  Independent paths can be interleaved preserving this property:
merge their lists in decreasing order of the greatest remaining root length,
which is comparable to the current root length on each list.
Project every differentiated root preceding the last such root onto its
endpoints.  The conditional coefficient is
\[
 c_l^I=R_l(I)+q_e(X_{\sup I})R_{l-1}(I),
 \qquad R_0=1,
\]
by \eqref{score:raw-conditional}.  This projection is within the expectation,
using Lemma~\ref{hist:conditional-interior}; the endpoint factors of preceding
roots are measurable before the last root.  After expanding their two
monomials, all initialization coefficients can therefore be bounded by one
in this calculation.

If these preceding roots supply no positive-order $R$ factor, the last
root, of order $l\le k$ and length $d\le CT$, uses the $L^1$ line of
\eqref{repl:forward-costs}.  Its first term is $O(n^{-1/6})$.
Since $T^{3/4}=\sqrt n$, its second term is at most
$CT^{(l-4)/2}$ when $l\ge3$; when $l\le2$ it is at most $T^{-3/4}$,
because $d\ge1$.  For $l<4$ there is a fixed power saving.  Put
\[
 \varepsilon_0=10^{-4}.
\]
For $l=4$ the bound saves $n^{-3\varepsilon_0/4}$ if
$d\le Tn^{-\varepsilon_0}$, while the late $L^1$ improvement saves
$n^{-0.001}$ on the remaining range.  This proves the required bounds for
these terms.

Otherwise, let $p$ be the total order in the preceding positive-order $R$
factors and $D$ their greatest length.  Use the $L^2$ bound on one factor
of length $D$, the supremum bound on the other preceding factors, and the
$L^2$ line of \eqref{repl:forward-costs} for the last order $l$ and length
$d$.  The result is
\begin{equation}\label{repl:unmarked-product}
 n^{O(\beta)}D^{p/2-1}d^{l/2-1/4}T^{-3/4},
 \qquad d\le C D,\qquad p+l\le k.
\end{equation}
For $p+l<4$ this has a fixed power saving.  For $p+l>4$ it is bounded by
$n^{O(\beta)}T^{(p+l-4)/2}$.  To check both assertions when $p=1$, write
the length factor as $(d/D)^{1/2}d^{l/2-3/4}T^{-3/4}$; when $l=1$ it is
even smaller, and otherwise its power of $d$ is positive.

Only $p+l=k=4$ is a borderline case.  Write $D=Tx$, $d=Ty$ with
$y\le Cx$.  Apart from $n^{O(\beta)}$, the three possible factors in
\eqref{repl:unmarked-product} are
\[
 \begin{array}{c|c}
 (p,l)&\text{factor}\\ \hline
 (1,3)&(y/x)^{1/2}y^{3/4}\\
 (2,2)&y^{3/4}\\
 (3,1)&x^{1/2}y^{1/4}.
 \end{array}
\]
If $D$ or $d$ is at most $Tn^{-\varepsilon_0}$, these bounds save at least
$n^{-\varepsilon_0/4}$.  Otherwise every contributing earlier root also
has length at least a constant times $Tn^{-\varepsilon_0}$.  If $l\ge2$,
the late $L^2$ improvement in \eqref{repl:forward-costs} applies.  If
$l=1$ and a preceding order is two or three, choose that factor for the
$L^2$ bound and use its $n^{-0.001}$ improvement.  Using its length in
place of the greatest length costs at most $n^{\varepsilon_0/2}$, since
all those lengths lie between a constant times $Tn^{-\varepsilon_0}$
and $CT$.  Finally, if the preceding orders are $(1,1,1)$, apply
$L^4$ to all four factors.  The three $R_1$ factors have their late
improvements, and the last factor is at most
$n^{O(\beta)}T^{5/16}n^{-1/4}=n^{-1/24+O(\beta)}$.
These are fixed power savings in all cases.  The $O(\log n)^K$ choices of
roots are absorbed into $n^{O(\beta)}$.  This proves
\eqref{repl:averaged-history} for terms with fixed cutoff weights.

\paragraph{The exact expansion at a differentiated cutoff.}
For a node $I=[a,b]$, use the pre-density $\bar k_I$ and conditional
interior law $\mathsf C_I^{x,y}$ from Lemma~\ref{hist:conditional-interior}.
For the perturbation at $e$, put
\[
 H_{I,e}(u)=\frac{\bar k_{I,e}(u)}{\bar k_I},\quad H_{I,e}(0)=1,
 \qquad
 V_{I,e}^+(u)=\frac{w_{I,e}(u)}{w_I},\quad
 V_{I,e}^-(u)=\frac{1-w_{I,e}(u)}{1-w_I}.
\]
All denominators are positive.  Write $W_N^+=W_N$ and $W_N^-=1-W_N$.
For $j\ge1$, the scalar chain rule gives
\begin{equation}\label{repl:cutoff-chain}
\begin{split}
 [u^j]V_{I,e}^{\sigma}(u)
 &=\sum_{\ell=1}^j
 \frac{(W_{N_I}^{\sigma})^{(\ell)}(\log\bar k_I)}
      {\ell!\,W_{N_I}^{\sigma}(\log\bar k_I)}
 \sum_{\substack{a_1+\cdots+a_\ell=j\\a_h\ge1}}
       \prod_{h=1}^{\ell} b_{I,e,a_h},\\
 b_{I,e,a}
 &:=[u^a]\log H_{I,e}(u)\\
 &=\sum_{v=1}^a\frac{(-1)^{v+1}}v
       \sum_{\substack{d_1+\cdots+d_v=a\\d_h\ge1}}
       \prod_{h=1}^{v}[u^{d_h}]H_{I,e}(u),
 \qquad \sigma\in\{+,-\}.
\end{split}
\end{equation}
Thus a positive-order cutoff coefficient supplies one positive-order
derivative ratio of $W_N^\sigma$ and a product of positive-order pre-density
coefficients whose orders sum to its assigned order.  We call this occurrence
of a derivative ratio a \emph{cutoff mark} at $I$.

The pre-density ratio has the exact representation
\begin{equation}\label{repl:pre-density-copy}
 H_{I,e}(u)
 =(1+u \bar q_e)\,
 \mathsf C_I^{x,y}\left[
    \frac{L_{I,e}(u)}{1+u q_e(y)}
    \prod_{\substack{J\subsetneq I\\J\text{ a checked node}}}
           V_{J,e}^+(u)\right].
\end{equation}
Indeed, the perturbed stationary probability is
$\mu_u(y)=\mu(y)(1+u q_e(y))/(1+u \bar q_e)$.  Dividing the perturbed
pre-truncation transition mass by its base value gives the conditional
expectation of $L_{I,e}(u)\prod_{J\subsetneq I}V_{J,e}^+(u)$.
The further factor $\mu(y)/\mu_u(y)$ gives
\eqref{repl:pre-density-copy}; see also Equation~\eqref{hist:pre-ratio}.
It is essential here that $\bar k_I$ is a density with respect to the
normalized stationary probability.

Each product of coefficients of $H_{I,e}$ in \eqref{repl:cutoff-chain}
is represented using independent interiors with law $\mathsf C_I^{x,y}$
conditional on the same endpoints.  Recursively do the same for every
positive-order descendant cutoff coefficient in
\eqref{repl:pre-density-copy}.  Given all already sampled paths and the
endpoints of a new copy, its interior is sampled independently with this
law.  Ancestor and exterior factors depend only on data outside the copied
interior.  By Lemma~\ref{hist:conditional-interior}, replacing that interior in
the full tested history preserves its law.  This remains true for a copy
inside a copy by iteration.  In particular the marginal estimates for each
copy are those for a base tested path, averaged over its endpoints; no
pointwise assertion for arbitrary endpoints is being made.

There is a useful order invariant in this expansion.  Assign to each
coefficient slot its positive Taylor order.  Formula
\eqref{repl:cutoff-chain} splits an order into positive orders with the
same sum, and expanding a slot by \eqref{repl:pre-density-copy} has the
same property.  The total order in all unexpanded slots, raw likelihood
slots, and endpoint or $\bar q_e$ factors is therefore at most $k$ at every
stage.  Cutoff marks record transferred orders; they do not add to this sum.
There may be a long chain of cutoff marks carrying a single order.

A second invariant gives a projected factor when expansion terminates.
For a positive-order pre-density slot of order $r$ whose descendants receive
no positive cutoff order, keep the endpoint denominator together with the
conditional likelihood.  Its coefficient is exactly
\begin{equation}\label{repl:terminal-mark}
 [u^r](1+u \bar q_e)R_e(I;u)
       =R_r(I)+\bar q_e R_{r-1}(I).
\end{equation}
Each summand contains a positive-order $R_l(I)$, or the order-one factor
$\bar q_e$ (when $r=1$ the second summand is $\bar q_e$ itself).
Induction from terminal slots shows that every terminal branch descending
from a cutoff mark contributes a positive-order conditional coefficient or an
$\bar q_e$ factor.  The endpoint denominator in
\eqref{repl:terminal-mark} is never estimated separately.  By
\eqref{repl:stationary-factor}, $\bar q_e$ can be treated as an $R_1$ factor
of length $T$: it satisfies all the corresponding bounds in
\eqref{repl:score-table}, including the late $L^4$ improvement.

\paragraph{Summation over cutoff marks and termination of the expansion.}
Let $\lambda_I=1-w_I$ for a survival cutoff mark and $\lambda_I=1$ for a
failure cutoff mark.  For any fixed $\eta>0$,
Lemma~\ref{hist:cutoff-hazards} gives, outside an arbitrarily small
inverse-polynomial contribution,
\begin{equation}\label{repl:mark-gain}
 \left|
 \frac{(W_{N_I}^{\sigma})^{(\ell)}(\log\bar k_I)}
      {\ell!\,W_{N_I}^{\sigma}(\log\bar k_I)}
 \right|
 \le C n^\eta\frac{\lambda_I}{N_I},
 \qquad
 N_I=n^\beta(T/|I|)^{3/2}.
\end{equation}
For survival this estimate is used on $w_I\ge n^{-L}$.
The discarded event $w_I<n^{-L}$ is paid from its included base survival
factor, with $L$ chosen arbitrarily large.  For failure the estimate holds
without this restriction.  The global derivative-ratio bounds of
Lemma~\ref{hist:cutoff-hazards} control exceptional contributions.

Here is a precise way to sum the choices of nodes in
\eqref{repl:mark-gain}.  On a normalized tested history with mass
$z_\alpha\ge n^{-M}$ and test bounded by $n^B$, put
$H=\sum_{I\in\mathscr S_\alpha}(1-w_I)$.  Since
$\prod_I w_I\le e^{-H}$,
\[
 \mathbb P_{\mathrm{test}}\{H>h\}\le n^{M+B}e^{-h}.
\]
One possible failure adds at most one to this hazard sum.  The same
statement holds for the marginal of every conditional copy, by its
law-preserving construction.  To make these statements simultaneous, for
each deterministic prefix of at most $P+K$ choices of nodes carrying cutoff
marks, sample the possible next copies conditional on that prefix.  All continuations
of a prefix share its already sampled paths.  There are $n^{O_K(P)}$
prefixes for fixed $P$.  The arbitrary precision convention permits the
hazard and moment bounds to hold for all of them simultaneously.  This union
uses the joint copy law and the marginal estimates for each indexed path.

All expectation and compensator identities used for this summation are
established in their stated filtrations before restriction to the
simultaneous event.  At a fixed sampled prefix, let $\mathcal I\ne\varnothing$
be the eligible next nodes on one shared path.  For each $I$, let $B_I\ge0$
bound the absolute value of the resulting continuation integrand or its edge norm on the common
sampled tree, including all later factors and gains but excluding the factor
$C n^\eta\lambda_I/N_I$ used to bound the current mark ratio.  The bound $B_I$
may depend on later sampled data.  On this common realization,
\[
 \sum_{I\in\mathcal I}\frac{\lambda_I}{N_I}B_I
 \le \left(\sum_{I\in\mathcal I}\lambda_I\right)
          \sup_{I\in\mathcal I}\frac{B_I}{N_I}.
\]
On the simultaneous event, the parenthesis is at most $n^\eta+1$.
Combining this with the current $C n^\eta$ from \eqref{repl:mark-gain}
costs at most $C n^{2\eta}$ while retaining $N_I^{-1}$ with the continuation
for the same choice.  Apply this bound separately to eligible slots and paths,
paying the fixed combinatorial multiplicity.  Iteration gives a supremum over
indexed full branches.  On the simultaneous event, the score bounds together
with the length estimates below control this supremum before integration.  The hazards are
scalars independent of $e$, so the same argument applies to edge-norm
estimates after projection.

The indicator of the simultaneous event is used only to bound these
resulting integrands.  The complementary integrals are
controlled by the polynomial moments in Lemmas~\ref{hist:crude-ratios}
and \ref{score:hellinger} and in the proof of
Lemma~\ref{score:raw-projection}, choosing their precision after $P$.
Thus no future-dependent hazard event is used
as a predictable factor.

It suffices to stop after a fixed number $P$ of cutoff marks.  If this
threshold is reached, leave the remaining pre-density slots unexpanded.  By
Lemma~\ref{hist:crude-ratios} and the fixed-order path moment bounds, their
product, together with all other positive-order slots, is at most
$n^{C_{\mathrm{cr}}}$ in the normalized estimate, for a constant
$C_{\mathrm{cr}}=C_{\mathrm{cr}}(K)$ independent of $P,\beta$.
Indeed the sum of their positive orders is at most $K$; the derivative
ratios already extracted at cutoff marks have uniform global bounds.  Every
node has length at most $CT$, so $N_I^{-1}\le Cn^{-\beta}$.  The contribution
stopped at $P$ cutoff marks, with an overshoot of at most $K$, is therefore
at most
\begin{equation}\label{repl:depth-remainder}
 C_{K,P}n^{C_{\mathrm{cr}}-P\beta+2(P+K)\eta}
\end{equation}
apart from exceptional precision.

The parameter choices can now be made in order.  First $K=24$ is fixed.
After choosing a sufficiently small $\beta$ for the score estimates, take
$P>(C_{\mathrm{cr}}+3)/\beta$, and then
$\eta\le\beta/(4(P+K))$.  For this fixed $P$, all scalar multiplicities
from \eqref{repl:cutoff-chain} and powers of $\log n$ from partitions cost
at most $n^{\beta/2}$ for large $n$.  Together with
$2(P+K)\eta\le\beta/2$, \eqref{repl:depth-remainder} is at most $n^{-2}$.
Finally choose the exceptional-set and small-mass precisions.  In a fully
expanded term at most $K$ positive-order factors use the score estimates.
Their total exponent loss is thus $C_K\beta$ independent of $P$.

\paragraph{The edge norms in fully expanded terms with cutoff marks.}
It remains to estimate terms which terminate before depth $P$.  We first
form a partition from the locations of their cutoff marks, before deciding
which endpoints to condition on.  Split each root into dyadic pieces only
at nodes which properly contain a descendant carrying a cutoff mark.  Make
the same partition on a copy using its own cutoff marks.

Let $m$ be the least length of a node carrying a cutoff mark, capped at $T$.
In each fully expanded monomial, select one factor $R_l(I)$ with $l\ge1$, or
one factor $\bar q_e$, supplied by a terminal $H_{I,e}$ slot through
\eqref{repl:terminal-mark}.  Make this choice from the indexed monomial,
hence the same way for every $e$.  The slot was created by a cutoff mark at
$I$, so the selected factor has length $|I|\ge m$ in the $R_l(I)$ case and
assigned length $T\ge m$ in the $\bar q_e$ case.  Every other terminal factor
remains in the product.  The gain from a shortest cutoff mark will be retained
separately below.  All partition pieces have length at
least a constant times $m$, except possibly untouched original forest
roots.  Indeed, a split made to isolate a descendant $M$ carrying a cutoff
mark creates two children, each of length at least $|M|\ge m$; further
splits satisfy the same statement with their own such descendants.  A copy
starts at a node carrying a cutoff mark, so it has no untouched smaller
root.  This also gives $O_P(1+\log n)$ pieces.  The untouched roots smaller
than a sufficiently small constant times $m$ can occur only in a stopped
forest, and occur later than all cutoff marks on that path: the lengths of
its original roots decrease chronologically.  Choose that constant below
the constants for complete legs, so no complete leg is included in this
exception.

Projecting the raw order-one coefficient on a root of length $D_1$ would
replace it by $c_1=R_1(D_1)+q_{\mathrm{end}}$.  If paired with $R_3(D_3)$,
the $R_3(D_3)R_1(D_1)$ summand would have $L^2$-times-$L^2$ cost
$n^{O(\beta)}(D_3/D_1)^{1/2}$ by \eqref{repl:score-table}, which can grow
when $D_1$ is very short.  We therefore reserve one allocation before
projecting piece coefficients.
An untouched original root $I$ of length $z<cm$ may carry a raw coefficient
of order one while the sole other factor of positive order is a conditional
coefficient of order three.
This is the small-root $(3,1)$ allocation.  Its terminal state
is not revealed in advance, and its complete raw coefficient
$\ell_{1,e}(I)=[u]L_{I,e}(u)$ is retained until forward compensation.
The total order is four, so no other positive-order factor is present.
The small root carries no cutoff mark and cannot supply the selected terminal
factor; hence in this allocation the selected factor is the $R_3$ factor.

For every other allocation, condition on the relevant endpoints of nodes
carrying cutoff marks and on the partition endpoints.  The fresh piece
interiors then have their conditional laws $\mathsf C_J$ and factor after
the endpoint variables and their shared identifications are fixed: lower
weights stay inside their pieces, and every other included weight depends
only on the conditioned endpoints.  Shared exterior variables occur once;
this does not assert independence of the completed copied histories.
Whenever a raw coefficient on a piece is projected, it becomes
$c_l=R_l+q_{\mathrm{end}}R_{l-1}$.  Nonterminal denominators in
\eqref{repl:pre-density-copy} and initialization may contribute endpoint
powers of absolute value one.  At a terminal slot use
\eqref{repl:terminal-mark} instead.  The resulting terms are products of
$R$ factors, $\bar q$ factors counted as $R_1(T)$, and bounded endpoint
factors.  The sum of their positive orders is at most $k$, and the selected
terminal factor remains in this product.  All these conditional identities
are integrated under the joint copy law when the sharp estimates are used;
they give no uniform sharp bound at an arbitrary fixed endpoint graph.

For total $R$ order at most four, \eqref{repl:score-table} gives the
following complete list.  Unused H\"older exponents are filled with the
constant function on the edge probability space.
\begin{equation}\label{repl:low-patterns}
\begin{array}{c|c|c}
\text{orders of the }R\text{ factors}&\text{norms used}
   &\text{bound apart from }n^{O(\beta)}\\ \hline
1^r,\ 1\le r\le4& L^4\text{ for each}&1\\
2,1^r,\ 0\le r\le2&L^2,L^4,\ldots,L^4&1\\
2,2&L^2,L^2&1\\
3&L^1&D_3^{-1/4}\\
4&L^1&1\\
3,1&L^2,L^2&(D_3/D_1)^{1/2}.
\end{array}
\end{equation}
Thus only $(3,1)$ can cost more than a constant.  If its order-one length
$D_1$ is at least a constant times $m$, its cost is at most
$C(T/m)^{1/2}$.  Retain the gain from a shortest cutoff mark,
\begin{equation}\label{repl:minimum-mark-gain}
 \frac1{N(m)}=n^{-\beta}(m/T)^{3/2};
\end{equation}
when the shortest node carrying a cutoff mark has length greater than $T$,
the same formula holds up to a constant.  The remaining gains are at most a constant times
$n^{-\beta}$ each and pay their hazard summations.
If $m\le Tn^{-\varepsilon_0}$,
\eqref{repl:minimum-mark-gain} leaves at most $(m/T)$ in the $(3,1)$
case and $(m/T)^{3/2}$ in all other cases, after the total $n^\beta$
allowance for hazards and logarithms.  This is a fixed power saving.

For the reserved small-root allocation, retaining $\ell_{1,e}(I)$ keeps its
$R_1$ and $q_{\mathrm{end}}$ conditional terms together before a triangle
inequality.  The selected order-three factor is earlier on the same path or
belongs to an independent original history.  The endpoints of all nodes carrying
cutoff marks on the same path precede $I$.  Their conditional copies can be
sampled before $I$, independently of its future given those already observed
endpoints.  Factors on other original histories can also be exposed first,
because those histories and their external tests factor separately.  This
auxiliary sampling is fixed as part of the admissible law before choosing
the earlier vector factor.

After projecting the earlier pieces, we return to the joint tested
expectation with the terminal state of $I$ unrevealed in advance, and
compensate its raw order-one coefficient in the forward filtration.
Under the forward tested law on this root, the external terminal and
possible failure tests are included in its tilted intensities.
The $L^2$ compensation in \eqref{repl:forward-costs} is therefore valid
with these earlier factors retained and gives
$n^{O(\beta)}z^{1/4}/\sqrt n$.  Paired with the order-three $L^2$ bound,
the result is
\begin{equation}\label{repl:small-root}
 n^{O(\beta)}\frac{D_3^{1/2}z^{1/4}}{\sqrt n}
 =n^{O(\beta)}(D_3/T)^{1/2}(z/T)^{1/4}
 \le n^{O(\beta)}.
\end{equation}
This calculation is made before projecting the small root onto its
endpoint, and the compensation identity precedes restriction to the
hazard event.  It proves the same early-$m$ saving using
\eqref{repl:minimum-mark-gain}.

For completeness, suppose $m\ge Tn^{-\varepsilon_0}$.  The selected terminal
factor has length at least $m$ and is in the late range.  Except for $(3,1)$,
the norms in \eqref{repl:low-patterns} give a fixed improvement at that
factor: a selected $R_1$ uses its late $L^4$ bound, a selected $R_2$ its late
$L^2$ bound, and a selected $R_4$ its late $L^1$ bound.
If the selected factor is the singleton $R_3$, then
$D_3^{-1/4}\le T^{-1/4}n^{\varepsilon_0/4}$.
For the regular $(3,1)$ case, put $\delta_0=0.005$ and use the
$R_3$ late $L^2$ gain $n^{-0.001}$.  Since $D_1\ge cm$,
\[
 \begin{cases}
 n^{-0.001}(D_3/D_1)^{1/2}
       \le C n^{-0.001+\varepsilon_0/2},
       &D_3\ge Tn^{-\delta_0},\\
 (D_3/D_1)^{1/2}
       \le C n^{-(\delta_0-\varepsilon_0)/2},
       &D_3<Tn^{-\delta_0}.
 \end{cases}
\]
The two savings are $n^{-0.00095}$ and $n^{-0.00245}$.
In \eqref{repl:small-root}, the corresponding cases gain $n^{-0.001}$
and $n^{-\delta_0/2}$, respectively.  These margins, and the early
$m$ margin, remain positive after an exponent $C_K\beta$ is included.
Together with the fixed-weight calculation, they prove the low-order
part of \eqref{repl:averaged-history}; one may take, for example,
$a=10^{-7}$.

For $k>4$ the product estimates need no late gains or lower length bound.
The following allocation of H\"older norms makes the power count explicit.
A factor of order $l\ge4$ uses its $L^1$ bound and saves two powers of
$T$ relative to the supremum bounds with all lengths replaced by $CT$.
If the largest order
is three, pair it with an order two or three in $L^2$, or with two order
ones using $L^2,L^4,L^4$; again the saving is two powers of $T$.
A singleton order three is bounded by one.  An order three paired with
exactly one order one costs at most $T^{1/2}$, which is allowed because
$k\ge5$.  If all orders are at most two, two order-two factors use
$L^2,L^2$, one order two and two order ones use $L^2,L^4,L^4$, and four
order ones use $L^4$ four times.  Each of these saves two powers of $T$;
when fewer factors are present their product is bounded by one.
All unused factors use their supremum bounds.  Thus if the total $R$
order is $h\le k$, the cost is at most
$n^{O(\beta)}T^{(k-4)/2}$.  This proves the high-order part.

Finally the same expansion with the supremum bounds in
Proposition~\ref{prop:clock-scores}, the raw fixed-order path bounds, and
$\sum l\le k$ gives \eqref{repl:sup-history}.  The depth remainder is
already smaller, and the exceptional and small-mass reductions at the
start give the stated additive precision.  Since the Taylor coefficients
of $\tanh\vartheta$ at zero are fixed constants, a derivative of order $k$
in $\vartheta$ is a fixed linear combination of coefficients of orders at
most $k$ in $u$.  Applying this to each factor of a product leaves total
$u$ order at most the assigned total $\vartheta$ order.  The lower-order
estimates fit the asserted bounds for that total order.
\end{proof}

\subsection{Localized traces and entrywise comparison}

Let $\mathcal Q_n$ be a finite set of queries in the fixed compact
rectangles, with $|\mathcal Q_n|\le n^{d_0}$ for a fixed $d_0$.  It may
contain both stationary and quench queries.  For a query $q$, let
$\zeta_q$ be its endpoint observable
$n^{-2/3}\sum_i X_iY_i$, so $|\zeta_q|\le n^{1/3}$.  Define
\[
 \widetilde Y_q=
 \mathbb E_q\left[\zeta_q\prod_{I\text{ in }q}w_I\right],
 \qquad
 \ell_q=1-\mathbb E_q\left[\prod_{I\text{ in }q}w_I\right].
\]
The expectations have the actual initialization of the query.  In
particular the quench history starts from the uniform cube law.
The surviving trace $\widetilde Y_q$ is not renormalized.
Order checks chronologically as in Definition~\ref{hist:tested-law}.  Telescoping
the product gives the exact first-failure decomposition
\begin{equation}\label{repl:first-failure}
 \ell_q=\sum_{I\text{ in }q}\ell_{q,I},\qquad
 \ell_{q,I}:=\mathbb E_q\left[
       \prod_{J\text{ preceding }I}w_J\,(1-w_I)\right]\ge0 .
\end{equation}
The summand is the mass of the corresponding checked first-failure prefix;
its unneeded future integrates to one.

Let $\Gamma(J)=\chi_n((J_e/\sqrt n)_e)$ be the static cutoff of
Lemma~\ref{hist:static-support}.
Choose a smooth $\chi:\mathbb R\to[0,1]$ equal to one on $(-\infty,1]$
and zero on $[2,\infty)$, and put
\begin{equation}\label{repl:outer-cutoff}
 G(J)=\Gamma(J)\chi\left(n^{10}\sum_{q\in\mathcal Q_n}\ell_q(J)\right).
\end{equation}
The Gaussian static estimate and Lemma~\ref{hist:gaussian-survival}, used at a
precision chosen after $d_0$, imply
\begin{equation}\label{repl:gaussian-cutoff}
 \mathbb E_{\mathrm{G}}G\longrightarrow1.
\end{equation}

We compare bounded smooth functions $F=F_n$ on
$\mathbb R^{\mathcal Q_n}$ satisfying, for fixed constants independent of
$n$,
\begin{equation}\label{repl:test-class}
 \|F\|_\infty\le C_F,\qquad
 \sup_y\sum_{q_1,\ldots,q_j\in\mathcal Q_n}
 \left|\partial_{q_1}\cdots\partial_{q_j}F(y)\right|
 \le C_F(\log n)^{C_F},\quad 1\le j\le K .
\end{equation}
This class includes $F=1$.  Write
$\Phi(J)=G(J)F((\widetilde Y_q(J))_{q\in\mathcal Q_n})$ and define
\[
 \partial_e^j\Phi(J)
 =\left.\frac{d^j}{d\vartheta^j}
       \Phi(J_{-e},J_e+\sqrt n\,\vartheta)\right|_{\vartheta=0}.
\]
Thus $\partial_e$ differentiates the interaction coefficient, rather than
the unscaled entry $J_e$.

\begin{lemma}[Derivatives of the localized observable]\label{repl:observable-derivatives}
The parameter $\beta>0$ can be chosen so that there are
$a_0>0$ and $0<\delta<0.01$ with, for every disorder matrix $J$,
\begin{align}
 \frac1{\binom n2}\sum_e|\partial_e^4\Phi(J)|
     &\le n^{-a_0},                                           \label{repl:fourth-derivative}\\
 \frac1{\binom n2}\sum_e|\partial_e^k\Phi(J)|
     &\le n^\delta T^{(k-4)/2},\quad 5\le k<K,               \label{repl:higher-derivative}\\
 \max_e|\partial_e^j\Phi(J)|
     &\le n^\delta T^{j/2},\quad 1\le j\le K.                \label{repl:global-derivative}
\end{align}
The constants may depend on the fixed exponents in
\eqref{repl:test-class} and on $d_0,c,C$.
\end{lemma}

\begin{proof}
On the support of a term containing a differentiated history, the static
support conditions hold and
$\sum_q\ell_q\le2n^{-10}$.  This follows from the supports of $\Gamma$
and its derivatives, and of $\chi$ and its derivatives.  Split each
$\zeta_q$ into its positive and negative parts.  The base masses of these two
checked integrals are at most $n^{O(\beta)}$.  Indeed the post-check kernel
caps at the ends of the one or two legs dominate the joint law of the two observation spins
by
$\exp(C'N(cT))\mu^{\otimes2}$.  For a stationary query the first spin
already has law $\mu$; for a uniform-start query the first leg cap holds
from each initial spin.  Choose the first dyadic $d_*\ge cT$.  Then
$d_*\le2cT$, $r_2(d_*)\le r_2(cT)$, and $a_2(d_*)\ge c_1n^{3\beta}$.
The pair cutoff therefore gives
\[
 \mathbb E_q\left[|\zeta_q|\prod_Iw_I\right]
 \le n^{1/3}r_2(cT)
     +n^{1/3}\exp(C'N(cT)-c n^{3\beta})
 \le C n^{2\beta}+n^{-A'}
\]
for arbitrary fixed $A'$.  We used $N(cT)=O(n^\beta)$.

Apply the ordinary chain rule to $\chi(n^{10}\sum\ell_q)$ and to
$F(\widetilde Y)$.  A term has at most $K$ positive-order history
coefficients, with total order at most the differentiating order; repeated
histories are represented by independent copies.  The coefficients from
derivatives of $F$ are evaluated at the base $\widetilde Y(J)$, hence are
common across $e$, and their total absolute sums are bounded by
\eqref{repl:test-class}.  For loss factors, expand
\eqref{repl:first-failure}.  The mass factors in
Lemma~\ref{repl:history-derivatives}, summed over all choices of queries and
first-failure nodes and multiplied by their $n^{10}$ factors, give at most
\[
 \left(n^{10}\sum_{q,I}\ell_{q,I}\right)^r\le2^r
\]
when there are $r$ loss factors.  Trace mass factors cost at most
$n^{O_K(\beta)}$.  The additive precisions in that lemma are chosen after
the polynomial number of nodes and queries and after the factors
$n^{10K}$, so their total is negligible.

If a static cutoff is also differentiated, its pointwise derivative bound
from Lemma~\ref{hist:static-support} multiplies these estimates.  If only the
static cutoff is differentiated, that bound is $O(n^{-b})$ for some
$b=b(\beta)>0$, and its remaining outer factor is bounded.  Consequently
the fourth averaged derivative is at most
$n^{-a+C'_K\beta}(\log n)^{C'}+O(n^{-b})$, higher averaged derivatives
are at most $n^{C'_K\beta}(\log n)^{C'}T^{(k-4)/2}$, and the supremum
bound is at most $n^{C'_K\beta}(\log n)^{C'}T^{j/2}$.
These estimates hold everywhere: outside the indicated supports the
corresponding product-rule terms vanish.  Choose $\beta$ after $K$ so
that $C'_K\beta<a/2$ and $C'_K\beta<0.005$, while satisfying the preceding
sections' smallness requirements.  Absorb logarithms by a still smaller
fixed power.  This gives \eqref{repl:fourth-derivative}--%
\eqref{repl:global-derivative} for some $a_0>0$ and $\delta<0.01$.
\end{proof}

\begin{lemma}[Entrywise replacement]\label{repl:comparison}
For every sequence of query sets and tests satisfying
\eqref{repl:test-class},
\[
 \mathbb E_{\mathrm{Rad}}\!\left[G F(\widetilde Y)\right]
 -\mathbb E_{\mathrm{G}}\!\left[G F(\widetilde Y)\right]\longrightarrow0 .
\]
\end{lemma}

\begin{proof}
We use the Taylor replacement principle \cite{Chatterjee2006}, in the finite
mixture form of the stationary comparison \cite[Section~8]{FUA}.
Put
$\nu_\lambda=(1-\lambda)\mathsf N(0,1)
 +\lambda(\delta_{-1}+\delta_1)/2$ for $0\le\lambda\le1$.
Differentiating its finite product measure gives exactly
\begin{equation}\label{repl:mixture-derivative}
 \frac d{d\lambda}\mathbb E_{\nu_\lambda^{\otimes\binom n2}}\Phi
 =\sum_e\mathbb E_{\lambda,-e}
    \left[\mathbb E_\xi\Phi(J_{-e},\xi)
                 -\mathbb E_Z\Phi(J_{-e},Z)\right],
\end{equation}
where $\xi$ is a symmetric sign, $Z$ is standard normal, and
$\mathbb E_{\lambda,-e}$ averages the other entries.

Independently of $J_{-e},\xi,Z$, sample $V$ from $\nu_\lambda$ and set
$J^0=(J_{-e},V)$.  Taylor-expand at this completed matrix:
\begin{equation}\label{repl:centered-taylor}
 \Phi(J_{-e},x)=
 \sum_{k=0}^{K-1}
   \frac{\partial_e^k\Phi(J^0)}{k!\,n^{k/2}}(x-V)^k
       +\mathcal R_e(x,V),
 \qquad
 |\mathcal R_e(x,V)|
 \le\frac{|x-V|^K}{K!\,n^{K/2}}\,n^\delta T^{K/2}.
\end{equation}
The remainder bound is global by \eqref{repl:global-derivative}, including
every point of the segment.  Conditional on $J_{-e},V$, the orders
$k\le3$ cancel exactly between the two proposals.  In fact
\begin{equation}\label{repl:centered-moments}
 \Delta_k(V):=\mathbb E_\xi(\xi-V)^k-\mathbb E_Z(Z-V)^k
 =\sum_{j=4}^k\binom kj(-V)^{k-j}
      \bigl(\mathbb E\xi^j-\mathbb EZ^j\bigr),
\end{equation}
where the sum is empty for $k\le3$.  Independence of the completed entry
from both proposals is what permits this conditional cancellation.

For $4\le k<K$, on $|V|\le\log n$ the absolute value of
\eqref{repl:centered-moments} is at most $C_K(\log n)^K$.  The complementary
mixture tail, even with any fixed moment of $V$, is smaller than every
inverse power of $n$.  Its contribution is negligible by the global
polynomial bounds \eqref{repl:global-derivative}.  When $e$ is uniform,
$J^0$ has the full iid mixture law independently of $e$.  Therefore
\eqref{repl:mixture-derivative}--\eqref{repl:centered-taylor} yield,
uniformly in $\lambda$,
\begin{align}
 \left|\frac d{d\lambda}\mathbb E_{\nu_\lambda^{\otimes\binom n2}}\Phi\right|
 &\le C_K(\log n)^K
   \sum_{k=4}^{K-1}n^{2-k/2}
     \mathbb E_\lambda\!\left[
       \frac1{\binom n2}\sum_e|\partial_e^k\Phi(J)|\right]\notag\\
 &\hspace{2cm}+C_K n^{2-K/2+\delta}T^{K/2}+o(1).             \label{repl:master-comparison}
\end{align}
The remainder uses the bounded $K$th moments of $\xi,Z,V$.
By \eqref{repl:fourth-derivative}, the $k=4$ term tends to zero.
For $k\ge5$, \eqref{repl:higher-derivative} and $T=n^{2/3}$ give
\[
 n^{2-k/2+\delta}T^{(k-4)/2}
      =n^{-(k-4)/6+\delta}=o(1).
\]
The remainder is $n^{2-K/6+\delta}=n^{-2+\delta}$ because $K=24$.
Thus \eqref{repl:master-comparison} tends to zero uniformly in $\lambda$.
Integration over $[0,1]$ proves the lemma.
\end{proof}

\subsection{Removing the cutoff and transferring tightness}

Take $F=1$ in Lemma~\ref{repl:comparison}.  In view of
\eqref{repl:gaussian-cutoff},
\begin{equation}\label{repl:both-cutoffs}
 \mathbb E_{\mathrm{Rad}}(1-G)\longrightarrow0,
 \qquad
 \mathbb E_{\mathrm{G}}(1-G)\longrightarrow0.
\end{equation}
This step transfers the high probability of the cutoff to Rademacher
disorder using derivative estimates that held for every matrix on their
deterministic supports.

Let $Y_q$ be the actual trace.  On $G>0$,
\begin{equation}\label{repl:remove-loss}
 |Y_q-\widetilde Y_q|
 \le n^{1/3}\ell_q
 \le2n^{1/3-10},\qquad q\in\mathcal Q_n.
\end{equation}
The first-derivative bound in \eqref{repl:test-class} makes $F$ Lipschitz
for the coordinate supremum norm with constant $C_F(\log n)^{C_F}$.
It follows from \eqref{repl:both-cutoffs}--\eqref{repl:remove-loss} that,
for either entry law,
\[
 \mathbb E F(Y)-\mathbb E[G F(\widetilde Y)]\longrightarrow0.
\]
Hence Lemma~\ref{repl:comparison} applies to the actual traces at finite query
sets and at all polynomial grids with tests in
\eqref{repl:test-class}.  At fixed finite sets this proves the equality
of limiting joint finite-dimensional laws with the Gaussian pair.

We give the tightness transfer because a comparison only at fixed times
would not suffice for the asserted topology.  The finite heat-bath
generator $\mathcal L_J=\sum_i(P_i-I)$ has
$\|\mathcal L_J\|_{\infty\to\infty}\le2n$, and its semigroup is a
contraction in supremum norm.  Differentiating the one or two semigroups
in a trace therefore gives, for every $J$ and on the fixed rectangles,
\begin{equation}\label{repl:time-lipschitz}
 |\partial_t A_{n,J}(t)|\le2n^2,\qquad
 |\partial_s B_{n,J}(s,t)|+|\partial_t B_{n,J}(s,t)|\le4n^2.
\end{equation}
Indeed each derivative in $s$ or $t$ inserts $T\mathcal L_J$ into its
corresponding semigroup, and the undifferentiated trace is bounded by $n^{1/3}$.

Choose grids of mesh at most $h_n=n^{-4}$ on the two rectangles.  They and
all pairs of their points have polynomial size, and
\eqref{repl:time-lipschitz} makes their interpolation error $O(n^{-2})$.
For grid coordinates $y$, the signed forms used to test a supremum are
$\pm y_q$, and those for a modulus of continuity are
$\pm(y_q-y_{q'})$ for pairs in the same rectangle at distance at most
a specified radius.  Include the zero form.  Each such linear form has
coefficient $\ell^1$ norm at most two.  For a collection of $M_n$ forms
$a_r(y)$, with $M_n$ polynomial, put
\[
 \operatorname{smax}_{\lambda_n}(a(y))
 =\lambda_n^{-1}\log\sum_{r=1}^{M_n}e^{\lambda_n a_r(y)},
 \qquad \lambda_n=(\log n)^2.
\]
It lies between $\max_r a_r$ and
$\max_r a_r+(\log M_n)/\lambda_n$, so its error is $o(1)$.
Its derivatives obey
\begin{equation}\label{repl:softmax-derivatives}
 \sum_{q_1,\ldots,q_j}
 \left|\partial_{q_1}\cdots\partial_{q_j}
      \operatorname{smax}_{\lambda_n}(a(y))\right|
 \le C_j\lambda_n^{j-1},\qquad 1\le j\le K.
\end{equation}
For example, the derivative of order $j$ is $\lambda_n^{j-1}$ times
the joint cumulant of the coefficient coordinates of a random form chosen
with its normalized exponential weight.  The cumulant partition formula,
followed by the coefficient $\ell^1$ bound two in each factor, proves
\eqref{repl:softmax-derivatives}.  Composing with a fixed bounded smooth
function whose derivatives through order $K$ are bounded now satisfies
\eqref{repl:test-class} by the ordinary chain rule.

Let $\omega_n(\rho)$ be the greater of the two moduli of continuity on
the fixed rectangles, with the supremum metric on each parameter domain.
Use the preceding smoothed maximum over grid pairs at distance at most
$2\rho$, followed by a smooth increasing function which is zero below
$\epsilon/3$ and one above $2\epsilon/3$.  For fixed $\rho,\epsilon>0$,
the interpolation and smoothing errors are smaller than $\epsilon/12$
for large $n$.  Comparison of the actual grid tests consequently gives
\begin{equation}\label{repl:modulus-transfer}
 \limsup_{n\to\infty}\mathbb P_{\mathrm{Rad}}\{\omega_n(\rho)>\epsilon\}
 \le
 \limsup_{n\to\infty}\mathbb P_{\mathrm{G}}\{\omega_n(2\rho)>\epsilon/4\}.
\end{equation}
Here a pair of original points at distance at most $\rho$ is moved to
nearest grid points at distance at most $2\rho$.  The same construction
with forms $\pm y_q$ transfers tightness of the supremum norm.  Gaussian
tightness, proved in Proposition~\ref{prop:gaussian-subsequence} and Theorem~\ref{thm:entrance-selection},
makes the right side of \eqref{repl:modulus-transfer} tend to zero as
$\rho\downarrow0$ and makes the supremum tails vanish as their threshold
tends to infinity.  The Arzel\`a--Ascoli tightness criterion therefore
gives tightness of the Rademacher restrictions to these rectangles.

Apply this argument on the exhaustion
$[j^{-1},j]$ and $[j^{-1},j]^2$, $j\ge2$.  The usual diagonal form of
the same criterion gives tightness in the product of the two
$C_{\mathrm{loc}}$ spaces.  Every subsequential limit has the Gaussian
finite-dimensional laws established above; evaluations on a countable
dense set determine a law on continuous functions.  The subsequential
limit is therefore $\Law_g(A_g,B_g)$, proving
Proposition~\ref{prop:universality}.

Every query in the comparison used the original times $sT,tT$ and the
original uniform initialization for $B_{n,J}$.  The replacement changes
only the coupling law.  Thus the coefficient $c_*$ and the entrance
$h_s^g$ obtained from the Gaussian construction pass unchanged to the
Rademacher limit.

\section{Completion of the proof}\label{sec:completion}

\begin{proof}[Proof of Theorem~\ref{thm:main}]
In Gaussian disorder, Proposition~\ref{prop:edge-estimates} supplies the
cap estimates used in Theorem~\ref{thm:warming}. Apply that theorem to
the initial law \(\nu_n\). Proposition~\ref{prop:gaussian-subsequence}
then supplies, on a
further represented subsequence of every sequence, a normalized
positive-time density family for \(S_t^g=e^{-tc_*H_g}\), together
with locally uniform limits of both finite autocorrelations.
Proposition~\ref{prop:uniform-seed-escape} gives the averaged uniform-start
escape estimate. The sequence extraction and fixed-\(g\) uniqueness in
Theorem~\ref{thm:entrance-selection} identify every represented family
with the same measurable family \(h^g\) for almost every \(g\).
The subsequence criterion gives the Gaussian joint convergence.

The density family is bounded at each positive time, hence belongs
to \(L^2(\Pi_g)\); it is nonnegative, has mass one, and satisfies the
semigroup relation. Its weak product-topology boundary is
\(\delta_0\) by Theorem~\ref{thm:entrance-selection}.
Proposition~\ref{prop:gaussian-subsequence} proves the locally uniform
convergence and continuity of the series defining \(B_g\).
Each finite partial sum is continuous in \((s,t)\), and its evaluations
at rational pairs are measurable in \(g\) by the density and semigroup
constructions in Theorem~\ref{thm:entrance-selection}. It is therefore
Borel as a \(C_{\mathrm{loc}}((0,\infty)^2)\)-valued map, and so is its
locally uniform limit.

Proposition~\ref{prop:gaussian-subsequence} gives
\(\|h_s^g-1\|_{L^2(\Pi_g)}\to0\), and hence the required
\(L^1\) convergence. It also gives locally uniform convergence
\(B_g(s,\cdot)\to A_g(\cdot)\) for almost every \(g\) on which the
selected family is defined.
This implies the stated convergence in probability over \(g\).

Proposition~\ref{prop:universality} transfers the joint function law
to Rademacher disorder. Its comparison uses the uniform initial law
for each quench trace, so the limiting family is the one already
selected in the Gaussian argument. The added Gaussian diagonal and
the auxiliary eigenvector signs do not affect either finite
autocorrelation. For both laws the finite dynamics use rate-one site
clocks and are observed in the time unit \(T_n=n^{2/3}\). The temporary
substitution \(u=c_*s\) was reversed in the entrance proof, so the
limiting semigroup retains exactly the coefficient \(c_*\) in
\eqref{eq:edge-semigroup},
completing all the assertions.
\end{proof}

\begingroup
\raggedright
\bibliographystyle{plain}
\bibliography{references}
\endgroup
\end{document}
